\documentclass[11pt,a4paper]{article}
\usepackage[a4paper,left=25mm,right=25mm,top=24mm,bottom=24mm]{geometry}
\usepackage{amsmath,amssymb,amsthm,mathtools,mathrsfs}
\usepackage{microtype,booktabs,longtable,array,enumitem,xurl,xcolor}
\usepackage[hidelinks,hypertexnames=false,bookmarksnumbered]{hyperref}
\setlength{\emergencystretch}{4em}
\setlength{\parskip}{2pt}
\setcounter{tocdepth}{1}
\numberwithin{equation}{section}
\newtheorem{theorem}{Theorem}[section]
\newtheorem{lemma}[theorem]{Lemma}
\newtheorem{proposition}[theorem]{Proposition}
\newtheorem{corollary}[theorem]{Corollary}
\theoremstyle{definition}
\newtheorem{definition}[theorem]{Definition}
\newtheorem{assumption}[theorem]{Assumption}
\theoremstyle{remark}
\newtheorem{remark}[theorem]{Remark}
\newcommand{\E}{\mathbb E}
\newcommand{\Pp}{\mathbb P}
\newcommand{\Prob}{\mathbb P}
\newcommand{\R}{\mathbb R}
\newcommand{\N}{\mathbb N}
\newcommand{\Z}{\mathbb Z}
\newcommand{\Q}{\mathbb Q}
\newcommand{\ii}{\mathrm i}
\newcommand{\e}{\mathrm e}
\newcommand{\dd}{\,\mathrm d}
\newcommand{\CE}{C_{\mathrm E}}
\newcommand{\PE}{P_{\mathrm E}}
\newcommand{\Dn}{\Delta_n}
\newcommand{\D}{\Delta_n}
\newcommand{\phiZ}{\phi_0}
\newcommand{\phiZero}{\phi_0}
\newcommand{\phio}{\phi(0)}
\newcommand{\eps}{\varepsilon}
\newcommand{\pos}[1]{(#1)_+}
\newcommand{\cR}{\mathcal R}
\newcommand{\cP}{\mathcal P}
\newcommand{\one}{\mathbf 1}
\newcommand{\Sym}{\operatorname{Sym}}
\newcommand{\Lam}{\frac{59}{250}}
\newcommand{\rhoMax}{\frac{15}{8}}
\newcommand{\CertificateCells}{111}
\newcommand{\CertificateCoefficients}{6,549}
\DeclareMathOperator{\Var}{Var}
\DeclareMathOperator{\Cov}{Cov}
\DeclareMathOperator{\supp}{supp}
\DeclareMathOperator{\dist}{dist}
\DeclareMathOperator{\sgn}{sgn}
\title{Moment defects and sharp iid Berry--Esseen bounds\\
\large Structure, lattice asymptotics, and effective estimates}
\author{Kacper Rodziewicz \and Bartosz Ko{\l}odziejek}
\date{Working draft --- 12 October 2026}
\begin{document}
\maketitle
\begin{center}\small
This is a working draft. A revised and corrected version is intended
for submission to arXiv.
\end{center}
\begin{abstract}
The known sharp third-moment inequality supplies an exact defect
that identifies the standardized two-point family and controls
transport and characteristic functions. We develop the
corresponding Fourier reduction and give the qualitative
eventual upper-bound argument. A penalized extremal problem yields
eventual two-point optimality and a uniform linear deficit, with sharp
cubic Wasserstein stability. A fifth-order lattice saddle gives the
full oscillatory \(n^{-2}\) correction for the optimal constants.
A separate effective argument proves \(C_n<C_{\mathrm E}\) for every
\(n\ge31\), by finite-index certificates and a first-failure argument.
Exact polynomial certificates give bounds for two and three summands.
For four summands, a common-measure Fourier comparison and rational
certificates give the unrestricted bound \(C_4\le0.4097<C_{\mathrm E}\);
a sharper estimate \(C_{4,\le3}<0.409\) holds on at most three points.
The opening argument explains an elementary historical bound \(C\le0.45\)
and the passage from a symmetric reference to the matching two-point law.
Its finite arithmetic is recorded in an appendix.
ChatGPT and Codex assisted the development of the proofs and computational
certificates and the preparation of this manuscript, using the
GPT-5.6 Sol, GPT-6.1 Sol and GPT-6 Astra models.
\end{abstract}
\clearpage
\section*{Introduction}
The common normalization is mean zero, variance one and finite
third absolute moment. Put
\[
 R_n(P)=\frac{\sqrt n\,\Delta_n(P)}{\rho(P)},\qquad
 C_n=\sup_P R_n(P),\qquad
 C_{\mathrm E}=\frac{3+\sqrt{10}}{6\sqrt{2\pi}}.
\]
In Part I we use
\[
 \beta(P)=\rho(P),\qquad f(\gamma)=\rho_2(\gamma),\qquad
 B_n=C_Q^*(n),\qquad d(P)=\rho(P)-\rho_2(\gamma(P)).
\]
The normalized defect \(q(P)=d(P)/\rho(P)\) is distinct from the
Bernoulli probability \(q=1-p\). Lattice spans and conditional-noise
variances are introduced locally.

The principal qualitative conclusion is that, for all sufficiently
large \(n\), every standardized law satisfies
\[
 R_n(P)+10^{-4}\frac{d(P)}{\rho(P)}\le C_Q^*(n)=C_n<C_{\mathrm E}.
\]
Theorem~\ref{stab:thm:penalty} proves this statement. The asymptotic
refinement, including its arithmetic oscillation, is given in
Corollary~\ref{main:full-second-order}. The effective theorem
\(\,C_n<C_{\mathrm E}\) for \(n\ge31\) is
Theorem~\ref{cut31:global:main}; it does not assert two-point
optimality at every index from 31 onward.

Part~\ref{part:structural} develops the defect, Fourier stability,
coherent response and two-point saddle, followed by the quantitative
and second-order refinements. Part~\ref{part:effective} proves the
effective estimates by analytic reductions and complete arithmetic
certificates. The unrestricted fourth-index bound,
Theorem~\ref{n4all:thm:main}, localizes a hypothetical extremizer and
retains one defect measure throughout the Fourier comparison.
Literature inputs are stated where used. The unrestricted indices
\(5,\ldots,30\) remain open in this work.

He--Cheng \cite{HC2026v3} prove the eventual upper bound
\(C_n\le\CE\), with the explicit threshold
\(2\lceil\exp(10^{17})\rceil\). The equality
\(C_n=C_Q^*(n)\), the stability statement and the lattice correction
are separate structural assertions; an effective strict upper bound
does not by itself imply them.

The manuscript and reproducibility material are maintained in the
\href{https://github.com/rodziewiczk/berry-esseen-045-certificate/tree/main/moment-defects}
{\texttt{moment-defects} directory of the accompanying repository}.
The original \(0.45\) certificate is preserved separately there.

\clearpage
\begingroup
\normalsize
\raggedright
\setlength{\parskip}{0pt}
\setlength{\abovedisplayskip}{4pt}
\setlength{\belowdisplayskip}{4pt}
\setlength{\abovedisplayshortskip}{3pt}
\setlength{\belowdisplayshortskip}{3pt}
\section*{Principal results}
\label{main:results-at-a-glance}
Let \(P\) have mean zero, variance one and finite
\(\rho=\E|X|^3\); write \(\gamma=\E X^3\),
\(R_n(P)=\sqrt n\,\Delta_n(P)/\rho\), \(C_n=\sup_P R_n(P)\),
and \(C_Q^*(n)=\sup_{Q\ {\rm two\mbox{-}point}}R_n(Q)\).
Put \(\CE=(3+\sqrt{10})/(6\sqrt{2\pi})\).
\begin{enumerate}[leftmargin=*,label=\arabic*.,itemsep=8pt,topsep=8pt,parsep=0pt]
\item \textbf{Exact two-point defect.}
With \(Q_\gamma\) the matching standardized two-point law,
\[
 d(P)=\rho-\frac{\gamma^2+2}{\sqrt{\gamma^2+4}}\ge0,
 \qquad d(P)=0\ \Longleftrightarrow\ P=Q_\gamma.
\]
Known sharp moment inequality \cite{Shevtsova2014};
Theorem~\ref{orig:thm:defect} records its defect form.

\item \textbf{Linear Fourier stability.}
For every \(t\in\R\),
\[
 |\varphi_P(t)-\varphi_{Q_\gamma}(t)|
 \le\frac{|t|^3d(P)}{4(1-|c_\gamma|)},\qquad
 c_\gamma=\frac{\gamma(\gamma^2+6)}{(\gamma^2+4)^{3/2}}.
\]
Theorem~\ref{orig:thm:cf}.

\item \textbf{Eventual two-point optimality.}
For all sufficiently large \(n\),
\[
 C_n=C_Q^*(n)<\CE.
\]
Every maximizing law is two-point. Theorem~\ref{stab:thm:penalty}.

\item \textbf{The full oscillatory second-order asymptotic.}
With the explicit constants and phase of Section~\ref{bin:section},
\[
 C_n=\CE-\frac{\kappa_\infty}{n}
 -\frac{\kappa_2+A_2\omega_n^2}{n^2}+o(n^{-2}).
\]
The limit set of \(n^2(C_n-\CE+\kappa_\infty/n)\) is exactly
\([-\kappa_2-A_2/4,-\kappa_2]\).
Corollary~\ref{main:full-second-order}.

\item \textbf{An effective unrestricted threshold.}
\(\ C_n<\CE\) for every integer \(n\ge31\).
Theorem~\ref{cut31:global:main}.
This threshold concerns the strict upper bound, not two-point optimality.

\item \textbf{Finite-index certificates.}
\[
 C_2\le\frac{2\sqrt2}{7},\qquad
 C_3\le\frac{59\sqrt3}{250},\qquad
 C_4\le\frac{4097}{10000}<\CE.
\]
All three bounds are unrestricted. On at most three points the sharper
bound is \(C_{4,\le3}<409/1000\).
Theorems~\ref{n2:thm:main}, \ref{n3:thm:main},
\ref{n4all:thm:main} and \ref{n4:thm:main}.

\end{enumerate}
\endgroup
\clearpage
{\small\tableofcontents}
\clearpage
\section{From a symmetric comparison to the moment defect}\label{opening:comparison}
We begin with a simpler bound. It isolates the comparison principle
without the extremizer or lattice arguments used later.
Throughout this section, \(\E X=0\), \(\E X^2=1\),
\(\rho=\E|X|^3\), \(\gamma=\E X^3\), and
\(f_n(t)=\varphi_X(t/\sqrt n)^n\).

\subsection{The Rademacher reference}
The least possible third absolute moment is one. Equality holds
only for the symmetric law on \(\{-1,1\}\): Jensen gives
\(\E|X|^3\ge(\E X^2)^{3/2}\), and equality forces \(|X|=1\).
Thus \(d_R=\rho-1\) is a defect from one fixed reference law.
Mattner--Shevtsova \cite[Theorem~1.2 and Example~1.4]{MattnerShevtsova2019}
give the distribution-free comparison
\begin{equation}\label{opening:rad}
 |\varphi_X(u)-\cos u|\le \frac{B(\rho)|u|^3}{6},\qquad
 B(\rho)^2=\frac{\rho^2+\rho\sqrt{\rho^2+8}}2-2.
\end{equation}
This uses the sharp third-moment geometry of Shevtsova
\cite[Theorem~6]{Shevtsova2014}. In particular,
\[
 B(\rho)=\sqrt{\frac83(\rho-1)}+o\bigl(\sqrt{\rho-1}\bigr)
 \qquad(\rho\downarrow1).
\]
The small parameter in this comparison is therefore a square root
of the symmetric defect.

Put \(u=t/\sqrt n\), \(c=|\cos u|\),
\(\delta=B(\rho)|u|^3/6\), and \(a=\min(1,c+\delta)\).
Factoring a difference of powers gives
\begin{equation}\label{opening:powers}
 |f_n(t)-\cos^n u|
 \le D_{n,\rho}(t):=
 \min\left\{2,n\delta,\delta\sum_{j=0}^{n-1}a^jc^{n-1-j}\right\}.
\end{equation}
Consequently the two quantities needed in smoothing have majorants
\[
 \begin{split}
 |f_n(t)-e^{-t^2/2}|&\le
 \min\{2,|\cos^n u-e^{-t^2/2}|+D_{n,\rho}(t)\},\\
 |f_n(t)|&\le \min\{1,c^n+D_{n,\rho}(t)\}.
 \end{split}
\]
For the second quantity one also takes the minimum with the
Prawitz--Shevtsova modulus bound described in
Appendix~\ref{app:elementary045}. Each bound applies to the original
law; the Rademacher law supplies a tractable reference, not a
restriction on the marginal.

\subsection{Smoothing and the elementary bound}
Let
\[
 K(s)=\frac{1-s}{2}+\frac{\ii}{2}
 \left((1-s)\cot(\pi s)+\frac1\pi\right),\qquad
 L(s)=K(s)-\frac{\ii}{2\pi s}.
\]
If \(R(t)\ge|f_n(t)-e^{-t^2/2}|\) and \(M(t)\ge|f_n(t)|\),
Prawitz smoothing \cite{Prawitz1972} gives, for \(T>0\), \(0<t_0<1\),
\begin{align}\label{opening:prawitz}
 \Delta_n\le{}&\int_0^{t_0}2|K(s)|R(Ts)\dd s
 +\int_{t_0}^1 2|K(s)|M(Ts)\dd s\notag\\
 &+\int_0^{t_0}2|L(s)|e^{-T^2s^2/2}\dd s
 +\frac1{2\pi}E_1(T^2t_0^2/2),
\end{align}
where \(E_1(x)=\int_x^\infty e^{-y}\dd y/y\).
The reference comparison controls the low frequencies. The modulus
estimate controls the remaining frequencies. These are distinct
roles and both are needed.

For completeness, the finite reduction is short. The one-summand
estimate of Bentkus--Kir\v{s}a \cite{BentkusKirsa1989} and Shevtsova's
bounds \cite{Shevtsova2013,Shevtsova2012} imply
\[
 R_1(P)\le .3704,\qquad
 R_n(P)\le .3322(1+.429/\rho),\qquad
 R_n(P)\le\CE+\frac{3\rho}{\sqrt n}.
\]
Put \(\rho_0=712569/589000\). The second bound is at most \(9/20\)
for \(\rho\ge\rho_0\); the third is strictly below \(9/20\) for
\(n\ge8124\) and \(\rho\le\rho_0\). It remains to bound
\eqref{opening:prawitz} on
\[
 2\le n\le8123,\qquad 1\le\rho\le\rho_0.
\]
The closed-interval certificate in Appendix~\ref{app:elementary045}
does precisely this and proves \(C\le9/20\).
The reduction is analytic; its last finite inequality is
computer-assisted. The supplementary arithmetic removes endpoint
singularities before taking upper sums and encloses differences of
nearby positive quantities without subtractive cancellation.

\subsection{Why the matching two-point law is the right reference}
An asymmetric two-point law has \(\rho>1\), so its Rademacher error
allowance in \eqref{opening:rad} does not vanish. The reference
should instead match the signed third moment.
Shevtsova's sharp moment inequality \cite[Theorem~6]{Shevtsova2014}
is equivalently
\begin{equation}\label{opening:sharp-moment}
 \rho\ge\rho_2(\gamma):=
 \frac{\gamma^2+2}{\sqrt{\gamma^2+4}}.
\end{equation}
Equality identifies the standardized two-point law \(Q_\gamma\).
Thus the appropriate defect is
\[
 d(P)=\rho(P)-\rho_2(\gamma(P)),\qquad
 d(P)=0\ \Longleftrightarrow\ P=Q_\gamma.
\]
Theorem~\ref{orig:thm:defect} gives its explicit nonnegative density.
Theorem~\ref{orig:thm:cf} then proves the linear comparison
\begin{equation}\label{opening:matched-fourier}
 |\varphi_P(t)-\varphi_{Q_\gamma}(t)|
 \le\frac{|t|^3d(P)}{4(1-|c_\gamma|)},\qquad
 c_\gamma=\frac{\gamma(\gamma^2+6)}{(\gamma^2+4)^{3/2}}.
\end{equation}
The allowance now vanishes on the entire two-point family and is
linear in the nonbinary excess. This is the structural improvement:
replace a fixed symmetric reference by the binary law matching the
first three moments. It is not a pointwise domination of distribution
functions, nor a claim that \eqref{opening:matched-fourier} is smaller
than \eqref{opening:rad} for every law and frequency.

The same matching also fixes the smoothing scale. If
\(m_1=\E|X|\) and \(s=\sqrt{\gamma^2+4}\), then
\[
 \rho_2(\gamma)+\E|Q_\gamma|=s,
 \qquad 0\le\rho+m_1-s\le K_1(\gamma)d(P)
\]
by Proposition~\ref{orig:prop:m1}. Here \(s\) is the binary lattice
span. Retaining \(m_1\), as in Shevtsova's modulus estimate
\cite{Shevtsova2012}, therefore aligns the Fourier cutoff with the
same reference. This comparison, its coherent defect measure and
the signed lattice smoothing form the starting point of Part~I.

\clearpage
\part{Structure, stability, and lattice asymptotics}\label{part:structural}
\section{Normalization and main results}

Let \(\Phi\) and \(\phi=\Phi'\) denote respectively the standard normal
cdf and density, and for any real random variable \(Z\) write
\[
                 \varphi_Z(t):=\E e^{\ii tZ}.
\]
Let $Y,Y_1,Y_2,\ldots$ be iid with
\[
  \E Y=0,\qquad \E Y^2=1,\qquad \rho=\E|Y|^3<\infty,
\]
and put
\[
  S_n=\frac{Y_1+\cdots+Y_n}{\sqrt n},\qquad
  \Dn(Y)=\sup_{x\in\R}|\Pp\{S_n\le x\}-\Phi(x)|,
\]
\[
 C_n=\sup_Y\frac{\sqrt n\,\Dn(Y)}{\rho},\qquad
 \CE=\frac{\sqrt{10}+3}{6\sqrt{2\pi}}.
\]
The supremum in $C_n$ is over centered variance-one laws with finite third
absolute moment.
For later local arguments, let $X$ be centered and standardized and put
\[
 \gamma=\E X^3,\quad \rho=\E|X|^3,\quad
 s=\sqrt{\gamma^2+4},\quad
 p={s-\gamma\over2s},\quad q=1-p.
\]
Then $\gamma=(q-p)/\sqrt{pq}$ and $s=(pq)^{-1/2}$.  If
$x=s(\eta+\theta)/\sqrt n$, $np+\eta\in\mathbb Z$, and
$|\theta|\le1/2$, write
\[
 \mathcal C_n^+(X;\eta,\theta)
 := {\sqrt n\over\E|X|^3}\left\{
 \Pp\left({X_1+\cdots+X_n\over\sqrt n}\le x\right)-\Phi(x)\right\},
\]
\[
 \mathcal C_n^-(X;\eta,\theta)
 := {\sqrt n\over\E|X|^3}\left\{
 \Phi(x)-\Pp\left({X_1+\cdots+X_n\over\sqrt n}<x\right)\right\}.
\]
Thus
\[
 \mathcal C_n(X):={\sqrt n\,\Delta_n(X)\over\rho}
 =\sup_{\eta,\theta}\max\{\mathcal C_n^+,\mathcal C_n^-\}.
\]

The following sharp moment inequality is due to Shevtsova
\cite[Theorem~6]{Shevtsova2014}; see also
Mattner--Shevtsova \cite[Example~1.3]{MattnerShevtsova2019}.
Their bound $|\gamma|\le B(\rho)$ is equivalent to
$\rho\ge\rho_2(\gamma)$, since
\[
 B(\rho)=\sqrt{\frac{\rho^2+\rho\sqrt{\rho^2+8}}2-2}
\]
is the inverse of $\rho_2$ restricted to $[0,\infty)$.
For $\gamma\ne0$, the equality case is also contained in the osculatory
interpolation of \cite[Lemma~2.4(b)]{MattnerShevtsova2019}; for
$\gamma=0$, it follows from Jensen's equality case.
We include a self-contained factorized proof, recording the
nonnegative defect density needed below.

\begin{theorem}[exact two-point defect]\label{orig:thm:defect}
Let $\gamma=\E Y^3$ and
\[
 s=\sqrt{\gamma^2+4},\qquad
 a=\frac{s-\gamma}{2},\qquad b=\frac{s+\gamma}{2}.
\]
Let $Z_\gamma$ have values $-a,b$ with probabilities $b/s,a/s$.
Then
\[
 d(Y):=\rho-\rho_2(\gamma)\ge0,\qquad
 \rho_2(\gamma)=\frac{\gamma^2+2}{\sqrt{\gamma^2+4}},
\]
and equality holds if and only if $Y\stackrel d=Z_\gamma$.
\end{theorem}

For any centered variance-one random variable \(Z\), we use the functional
notation
\[
 \gamma(Z):=\E Z^3,\qquad \rho(Z):=\E|Z|^3,\qquad
 d(Z):=\rho(Z)-\rho_2(\gamma(Z)),                       \tag{N0}
\]
and write
\[
 s_\gamma(\gamma):=\sqrt{\gamma^2+4},\qquad
 p(Z):={s_\gamma(\gamma(Z))-\gamma(Z)\over
              2s_\gamma(\gamma(Z))}.
                                                               \tag{N1}
\]
When the random variable is fixed, the corresponding unsubscripted symbols
\(\gamma,\rho,d,s,p\) retain the shorter meanings introduced above.

\begin{theorem}[linear Fourier stability]\label{orig:thm:cf}
With $c_\gamma=\gamma(\gamma^2+6)/(\gamma^2+4)^{3/2}$ and
$\chi_\gamma(t)=\E\e^{\ii tZ_\gamma}$, one has, for all $t\in\R$,
\[
 \left|\E\e^{\ii tY}-\chi_\gamma(t)\right|
 \le \frac{|t|^3}{4(1-|c_\gamma|)}d(Y).
\]
\end{theorem}

We henceforth write $Q_\gamma$ for the standardized two-point law
$Z_\gamma$ in Theorem~\ref{orig:thm:defect}.  Thus, with
$a=(\sqrt{\gamma^2+4}-\gamma)/2$ and
$b=(\sqrt{\gamma^2+4}+\gamma)/2$,
\[
 \Pp\{Q_\gamma=-a\}={b\over a+b},\qquad
 \Pp\{Q_\gamma=b\}={a\over a+b}.                       \tag{Q0}
\]
This notation always denotes the law (or a random variable having that law),
never an auxiliary polynomial.


Write
\[
 \overline{\mathbb R}:=[-\infty,+\infty]
\]
for the two-ended order compactification of $\mathbb R$.

\begin{corollary}[coherent defect representation]\label{orig:cor:rankfree}
Let $X$ be centered with variance one and third moment $\gamma$.  There is
a unique $p\in(0,1)$ such that, with $q=1-p$,
\[
 \gamma=\gamma_p:=\frac{q-p}{\sqrt{pq}}.
\]
Put $s=(pq)^{-1/2}$ and $Y=p+X/s$.  Then
\[
 \E Y=\E Y^2=\E Y^3=p.
\]
There are an explicit nonnegative defect density $W_p$, an explicit
continuous response curve $r_{p,t}:\overline\R\to\mathbb C$, and a
probability measure $\nu$ such that, with
$b_p(t)=q+p\e^{2\pi\ii t}$,
\[
 \boxed{\quad \varphi_Y(2\pi t)=b_p(t)+d\,r_\nu(t),\qquad
 r_\nu(t)=\int r_{p,t}(y)\dd\nu(y),\quad}    \tag{RF}
\]
where
\[
 d=\E|X|^3-\rho_2(\gamma_p),\qquad
 \rho_2(\gamma_p)=\frac{p^2+q^2}{\sqrt{pq}},
\]
The probability measure \(\nu\) is independent of \(t\).
The centered Bernoulli transform is
\[
 c_p(t)=e^{-2\pi\ii pt}b_p(t).
\]
When \(d=0\), the law of \(Y\) is Bernoulli and the perturbation vanishes.
\end{corollary}

For clarity we fix here, before its first use, the signed smoothing notation.
For a centered variance-one summand $X$, let
\[
 F_{n,X}(x)=\Pp\{(X_1+\cdots+X_n)/\sqrt n\le x\},\qquad
 F_{n,X}(x-)=\lim_{u\uparrow x}F_{n,X}(u),
\]
and put
\[
 D_X^+(x)=F_{n,X}(x)-\Phi(x),\qquad
 D_X^-(x)=\Phi(x)-F_{n,X}(x-).
\]
For \(\sigma\in\{+,-\}\), \(D_X^\sigma(x)\) denotes the corresponding
one of these two signed errors.  The same sign convention is used for
\(\mathcal C_n^\sigma\) and for the one-sided operators below.
Thus
\[
 \Delta_n(X)=\sup_x\max\{D_X^+(x),D_X^-(x)\}.
\]
The Prawitz kernel used throughout is
\[
 K(t)=\frac{1-t}{2}+\frac{\ii}{2}
 \left\{(1-t)\cot(\pi t)+\frac1\pi\right\},\qquad0<t<1,
\]
and
\[
 K(t)={\ii\over2\pi t}+{1\over2}+O(t),\qquad
 \overline{K(t)}=-{\ii\over2\pi t}+{1\over2}+O(t)
 \quad(t\downarrow0).
\]
The corresponding one-sided complex improper integral need not exist.  For
$\kappa\in\{K,\overline K\}$, $0<b\le1$, and
$g(t)=g(0)+O(t)$ with $g(0)\in\mathbb R$, define instead
\[
 \mathfrak I_{\kappa,b}(g):=
 2\lim_{\varepsilon\downarrow0}\operatorname {Re}
       \int_\varepsilon^b\kappa(t)g(t)\dd t,\qquad
 \mathfrak I_\kappa:=\mathfrak I_{\kappa,1}.              \tag{P0}
\]
Thus every later expression of the form
$2\operatorname {Re}\int_0^b\kappa(t)g(t)\dd t$ is understood as the
limit in (P0).

\begin{lemma}[Prawitz--Beurling half-line kernel]\label{orig:lem:BPcardinal}
Define, in the real-part improper sense of \emph{(P0)},
\[
 M(u):={1\over2}+2\Re\int_0^1K(t)e^{2\pi\ii tu}\,\dd t,
 \qquad e_B(u):=M(u)-1_{(-\infty,0]}(u).                 \tag{P1}
\]
Then
\[
 M(u)\ge1_{(-\infty,0]}(u),\qquad
 1-M(-u)\le1_{(-\infty,0)}(u),\qquad e_B(u)\ge0         \tag{P2}
\]
for every $u\in\mathbb R$.  At the integer nodes the identities are exact:
\[
             M(j)=1_{\{j\le0\}},\qquad
             1-M(-j)=1_{\{j<0\}}\quad(j\in\mathbb Z).  \tag{P3}
\]
If $U$ is integrable and $\psi_U(t)=\E e^{2\pi\ii tU}$, then all
expectations below are finite and
\[
 \E M(U)={1\over2}+\mathfrak I_K(\psi_U),\qquad
 \E\{1-M(-U)\}={1\over2}-\mathfrak I_{\overline K}(\psi_U). \tag{P4}
\]
Thus taking expectations in (P2) involves no formal interchange with the
$t^{-1}$ part of the kernel.
\end{lemma}

\begin{proof}
The one-sided smoothing inequality of \cite{Prawitz1972}, in the normalization
whose Fourier kernel is the displayed $K$, applied to the point mass at
$u$, gives $1_{(-\infty,0]}(u)\le M(u)$.  Thus the only imported assertion
here is exactly Prawitz's standard half-line inequality; the second
inequality in (P2) follows from the first at $-u$.  We verify the cardinal
and integrability facts used later.  The independent partial-fraction
calculation (4.52a)--(4.52b) below also checks the sign of the resulting
Beurling error away from the integer nodes.

For $j=0$, the real part of $K$ gives $M(0)=1$.  If
$j\in\mathbb Z\setminus\{0\}$, the elementary cosine term and the bare sine
term integrate to zero, and hence
\[
 M(j)={1\over2}-\int_0^1(1-t)\cot(\pi t)\sin(2\pi jt)\,\dd t.       \tag{P5}
\]
The product following $(1-t)$ is invariant under $t\mapsto1-t$.
Moreover, for $m\ge1$,
\[
 \cot(\pi t)\sin(2\pi mt)
 =1+\cos(2\pi mt)+2\sum_{\ell=1}^{m-1}\cos(2\pi\ell t),          \tag{P6}
\]
and the formula for $-m$ is its negative.  Thus the integral in (P5) is
$\tfrac12\operatorname {sgn}j$, proving (P3).

Finally $2\int_0^1\Re K(t)\,\dd t=1/2$, so for each real $u$
\[
 M(u)=1+2\Re\int_0^1K(t)\{e^{2\pi\ii tu}-1\}\,\dd t.             \tag{P7}
\]
This integral is absolutely convergent near zero because
$K(t)=\ii/(2\pi t)+O(1)$ and
$|e^{2\pi\ii tu}-1|\le2\pi t|u|$.  For integrable $U$, the same bound with
$\E|U|$ permits Fubini in (P7).  It gives the first identity in (P4), and
conjugation gives the second.  The same estimate supplies an integrable
affine bound for $M(U)$.
\end{proof}

We also fix the associated order operators here, before any asymptotic use.
If $X$ is real-valued and integrable, $n\ge1$, $h>0$, and $x\in\mathbb R$,
put
\[
 g_{X,n,h,x}(t)=e^{-2\pi\ii x\sqrt n\,t/h}
                 \varphi_X(2\pi t/h)^n,
\]
\[
 \mathcal U^+_{X,n,h}(x):={1\over2}+\mathfrak I_K(g_{X,n,h,x}),
 \qquad
 \mathcal U^-_{X,n,h}(x):={1\over2}-
                 \mathfrak I_{\overline K}(g_{X,n,h,x}).          \tag{P8}
\]
For $Z=(X_1+\cdots+X_n-x\sqrt n)/h$, the lemma gives
\[
 \mathcal U^+_{X,n,h}(x)=\E M(Z),\qquad
 \mathcal U^-_{X,n,h}(x)=\E\{1-M(-Z)\},                         \tag{P9}
\]
and therefore
\[
 \begin{aligned}
 F_{n,X}(x)&=\mathcal U^+_{X,n,h}(x)-\E e_B(Z),\\
 F_{n,X}(x-)&=\mathcal U^-_{X,n,h}(x)+\E e_B(-Z).
 \end{aligned}                                                    \tag{P10}
\]
Thus $\mathcal U^+$ is an upper operator for the closed half-line and
$\mathcal U^-$ a lower operator for the open half-line.  In particular,
\[
 D_X^+(x)\le\mathcal U^+_{X,n,h}(x)-\Phi(x),\qquad
 D_X^-(x)\le\Phi(x)-\mathcal U^-_{X,n,h}(x).                     \tag{P11}
\]


\begin{theorem}[two-point asymptotics]\label{orig:thm:twopoint}
If $C_Q^*(n)$ is the supremum defining $C_n$ restricted to standardized
two-point laws, then
\[
 C_Q^*(n)=\CE-\frac{\kappa_\infty}{n}+o(n^{-1}),
 \qquad
 \kappa_\infty=\frac{867+287\sqrt{10}}
 {25920\sqrt{2\pi}}.
\]
\end{theorem}

\begin{theorem}[qualitative eventual upper bound]\label{orig:thm:eventual}
There exists an integer $N_0$ such that
\[
 C_n\le\CE\qquad(n\ge N_0).
\]
\end{theorem}


\section{The exact moment defect}\label{orig:sec:defect}

The elementary identities
\[
 ab=1,\qquad a+b=s,\qquad b-a=\gamma
\]
show that $Z_\gamma$ has moments $0,1,\gamma$ of orders one, two and
three.  Its third absolute moment is
\[
 \rho_2(\gamma)=\frac{b}{s}a^3+\frac{a}{s}b^3
 =s-\frac2s=\frac{\gamma^2+2}{s}.
\]

\begin{proof}[Proof of Theorem~\ref{orig:thm:defect}]
For $\gamma\ne0$, the density $w_\gamma$ below is the normalized
osculatory error of Mattner--Shevtsova
\cite[Lemma~2.4(b)]{MattnerShevtsova2019}, applied to $x^3$ at the
nodes $-a,b$. Its formula extends continuously to $\gamma=0$.
We record its explicit factorization.
Define
\[
 c_\gamma=\frac{\gamma(\gamma^2+6)}{s^3}
\]
and
\[
 w_\gamma(x)=|x|^3-c_\gamma x^3-\frac{12}{s^3}x^2
              +\frac{6\gamma}{s^3}x+\frac4{s^3}.
\]
Notice that the denominator occurring below is strictly positive, since
\[
 1-c_\gamma^2
 =\frac{(\gamma^2+4)^3-\gamma^2(\gamma^2+6)^2}
        {(\gamma^2+4)^3}
 =\frac{4(3\gamma^2+16)}{(\gamma^2+4)^3}>0.
\]
Using $ab=1$, $a+b=s$ and $b-a=\gamma$, direct multiplication gives
\begin{align}
 w_\gamma(x)
 &=\frac{2a^2}{s^3}(x-b)^2\{2+(a+3b)x\},&&x\ge0,\label{orig:eq:wplus}\\
 w_\gamma(-t)
 &=\frac{2b^2}{s^3}(t-a)^2\{2+(3a+b)t\},&&t\ge0.\label{orig:eq:wminus}
\end{align}
Both right-hand sides are nonnegative, and they vanish only at $b$ and
$a$, respectively.  Taking expectations and using the first three moments,
\[
 \E w_\gamma(Y)=\rho-c_\gamma\gamma-\frac8{s^3}.
\]
The identity
\[
 c_\gamma\gamma+\frac8{s^3}
 =\frac{\gamma^4+6\gamma^2+8}{s^3}
 =\frac{(\gamma^2+2)(\gamma^2+4)}{s^3}
 =\rho_2(\gamma)
\]
proves $d(Y)=\E w_\gamma(Y)\ge0$.  Equality forces $Y$ to be supported
on $\{-a,b\}$.  If $r=\Pp\{Y=b\}$, the equation $\E Y=0$ gives
$r=a/(a+b)=a/s$; hence $Y\stackrel d=Z_\gamma$.  The converse follows
from the displayed factorization.
\end{proof}

\begin{corollary}[quadratic concentration]\label{orig:cor:transport}
Put $a_*=\tfrac12(s-|\gamma|)$.  Then
\[
 \E\operatorname{dist}\bigl(Y,\{-a,b\}\bigr)^2
 \le\frac{s^3}{4a_*^2}d(Y).
\]
\end{corollary}

\begin{proof}
In \eqref{orig:eq:wplus} and \eqref{orig:eq:wminus} the final linear factor is at
least $2$, while $\min(a,b)=a_*$.  Therefore
\[
 w_\gamma(x)\ge\frac{4a_*^2}{s^3}
 \operatorname{dist}(x,\{-a,b\})^2.
\]
Integrate and use $\E w_\gamma(Y)=d(Y)$.
\end{proof}

\begin{proposition}[stability of the smoothing scale]\label{orig:prop:m1}
With
\[
 b_*=\frac{s+|\gamma|}{2},\qquad
 K_1(\gamma)=\frac{s^2}{a_*(a_*+3b_*)},
\]
one has
\[
 0\le \rho+\E|Y|-s\le K_1(\gamma)d(Y).
\]
\end{proposition}

\begin{proof}
By reflection assume $\gamma\ge0$ and define
\[
 v_\gamma(x)=|x|^3+|x|-\frac\gamma s x^3
              -\frac4s x^2+\frac\gamma s x.
\]
Direct factorization gives
\[
 v_\gamma(x)=\frac{2a}{s}x(x-b)^2\quad(x\ge0),
 \qquad
 v_\gamma(-t)=\frac{2b}{s}t(t-a)^2\quad(t\ge0).
\]
Thus $v_\gamma\ge0$ and
$\E v_\gamma(Y)=\rho+\E|Y|-s$.  Dividing these factorizations by
\eqref{orig:eq:wplus}--\eqref{orig:eq:wminus}, and using monotonicity of the
resulting fractional-linear functions, yields
\[
 \frac{v_\gamma(x)}{w_\gamma(x)}
 \le\frac{s^2}{a(a+3b)},\qquad
 \frac{v_\gamma(-t)}{w_\gamma(-t)}
 \le\frac{s^2}{b(3a+b)}.
\]
The first constant is larger when $b\ge a$.  Take expectations and then
reflect to cover $\gamma<0$.
\end{proof}

\section{Proof of the characteristic-function stability theorem}

Let $H_{\gamma,t}$ be the cubic Hermite interpolant of $x\mapsto \e^{\ii tx}$
at $-a,b$, matching values and first derivatives.  Since $Y$ and
$Z_\gamma$ have identical moments through order three,
\begin{equation}\label{orig:eq:momentcancel}
 \E\e^{\ii tY}-\chi_\gamma(t)
 =\E\{\e^{\ii tY}-H_{\gamma,t}(Y)\}.
\end{equation}
By reflection we may suppose $\gamma,t\ge0$.  Put
\[
 p=\frac as,\qquad q=\frac bs=1-p,
 \qquad z=ts,
\]
and let $P_z$ be the cubic Hermite interpolant of $e^{\ii zy}$ at
$0,1$.  Write $E_z(y)=e^{\ii zy}-P_z(y)$.

For $x\ge0$, $y=(x+a)/s$ and a direct change of variables gives
\[
 \frac{w_\gamma(x)}{1-c_\gamma}
 =s^3(y-1)^2(y-\alpha_p),
 \qquad \alpha_p=\frac{p}{3-2p}.
\]
The function $p\mapsto\alpha_p$ is increasing.  Hence it suffices to
prove
\begin{align}
 |E_z(y)|&\le\frac{|z|^3}{4}(y-1)^2
 \frac{2y(1-y)}{3-2y},&&0<y<\tfrac12,\label{orig:eq:central1}\\
 |E_z(y)|&\le\frac{|z|^3}{4}(y-1)^2(y-\tfrac14),&&y\ge\tfrac12.
 \label{orig:eq:central2}
\end{align}
Indeed, when $p\le y\le1/2$,
\[
 y-\alpha_p\ge y-\alpha_y=\frac{2y(1-y)}{3-2y},
\]
and when $y\ge1/2$, one has
$y-\alpha_p\ge y-\alpha_{1/2}=y-1/4$.
On the negative half-axis the identity
\[
 E_z(1-y)=e^{\ii z}\overline{E_z(y)}
\]
follows from uniqueness of Hermite interpolation.  To make the reduction
explicit, write $t=-x\ge0$ and
$y'=(t+b)/s=1-y\ge q$, where $q=1-p\ge1/2$.  Direct substitution in
\eqref{orig:eq:wminus} gives
\[
 \frac{w_\gamma(-t)}{1-c_\gamma}
 =s^3(y'-1)^2 A_q(y'-\beta_q),\qquad
 A_q=\frac{q^2(3-2q)}{p^2(1+2q)},\qquad
 \beta_q=\frac{q}{3-2q}.
\]
Here $A_q\ge1$, so the difference
$A_q(y'-\beta_q)-(y'-1/4)$ is increasing in $y'$.  After multiplication by
the positive denominator $4p^2(1+2q)$, its value at $y'=q$ is
\[
 p(16q^3-6q^2-3q+1)
 =p(2q-1)(8q^2+q-1)\ge0,
\]
because $q\ge1/2$.  Thus
$A_q(y'-\beta_q)\ge y'-1/4$, and
\eqref{orig:eq:central1}--\eqref{orig:eq:central2} cover both half-axes.

The Hermite--Genocchi formula gives
\[
 E_z(y)=y^2(y-1)^2 f_z[0,0,y,1,1],
 \qquad f_z(u)=e^{\ii zu}.
\]
Let $M_y$ denote the probability B-spline with knots $0,0,y,1,1$.  The
normalized formula, including the simplex-volume factor, is
\[
 f_z[0,0,y,1,1]
 ={1\over4!}\int_0^1 f_z^{(4)}(u)M_y(u)\dd u.             \tag{HG}
\]
For $0<y<1$ its density is
\[
 M_y(u)=
 \begin{cases}
 \displaystyle {4u^2\over y^2}\{3y-(1+2y)u\},
       &0\le u\le y,\\[2mm]
 \displaystyle {4(1-u)^2\over(1-y)^2}\{(3-2y)u-y\},
       &y\le u\le1,
 \end{cases}
\]
with the limiting interpretation at $y=0,1$.  Direct differentiation shows
that $M_y$ is unimodal, vanishes at the endpoints, and has maximum at most
$2$; hence its total variation is at most $4$.  Stieltjes integration by
parts in (HG) yields
\begin{equation}\label{orig:eq:hermitecubic}
 |E_z(y)|\le\frac{|z|^3}{6}y^2(y-1)^2,\qquad0\le y\le1.
\end{equation}
For $0<y\le1/2$, \eqref{orig:eq:central1} follows from
$2y^2-6y+3\ge0$.  For $1/2\le y\le1$, \eqref{orig:eq:central2} follows from
$8y^2-12y+3\le0$.  It remains to treat $y\ge1$.

Put $v=z(y-1)$, $w=zy=v+z$, and
\[
 A(r)=\int_0^1e^{\ii\theta r}\dd\theta.
\]
Define the normalized error
\[
 F=\frac{e^{-\ii z}E_z(y)}{z^2(y-1)^2},
\]
with its continuous value at $z=0$ or $y=1$.  Cancelling the double zero at
$y=1$ gives the exact identity
\[
 F=\ii w\int_0^1u(1-u)e^{-\ii zu}\{1-A(wu)\}\dd u.
\]
Consequently, by Cauchy--Schwarz,
\begin{equation}\label{orig:eq:FJ2}
 |F|\le w\sqrt{J_2(w)/6},\qquad
 J_2(w)=\int_0^1u(1-u)|1-A(wu)|^2\dd u.
\end{equation}

\begin{lemma}\label{orig:lem:J2}
For every $w\ge0$,
\[
 J_2(w)\le\frac{243}{1250}=6\left(\frac9{50}\right)^2.
\]
\end{lemma}

\begin{proof}
For $r\ge0$,
\[
 |1-A(r)|^2=1-2\frac{\sin r}{r}
                 +2\frac{1-\cos r}{r^2}.
\]
On $0\le r\le10$ insert the alternating Taylor bounds
\begin{align*}
 \frac{\sin r}{r}&\ge
 \sum_{k=0}^{8}\frac{(-1)^kr^{2k}}{(2k+1)!}-\frac{r^{18}}{19!},\\
 \frac{1-\cos r}{r^2}&\le
 \sum_{k=1}^{8}\frac{(-1)^{k+1}r^{2k-2}}{(2k)!}
 +\frac{r^{16}}{18!}.
\end{align*}
These inequalities are not an appeal to a numerical alternating-series
test: the first is the global lower sine polynomial through degree $19$
and the second is the global upper polynomial for $1-\cos r$ through degree
$18$.  Repeated differentiation reduces both remainders to
$1-\cos r\ge0$, proving the stated signs for every $r\ge0$.
Let $P(r)$ be the resulting upper polynomial and
$Q(w)=\int_0^1u(1-u)P(wu)\dd u$.  Each monomial is integrated by
\[
 \int_0^1u^{j+1}(1-u)\dd u=\frac1{(j+2)(j+3)}.
\]
After $x=w^2/100$, the degree-nine polynomial
$243/1250-Q(10\sqrt x)$ has the positive Bernstein coefficients printed
in Appendix~\ref{orig:app:bernstein}, separately on $[0,1/2]$ and $[1/2,1]$.
This proves the estimate for $w\le10$.

For $w\ge10$, direct integration gives
\[
 J_2(w)=\frac16+\frac{2\operatorname{Cin}(w)}{w^2}
 -\frac4{w^2}+\frac{4\sin w}{w^3},
 \qquad
 \operatorname{Cin}(w):=\int_0^w{1-\cos t\over t}\,\dd t.
                                                               \tag{J1}\label{orig:eq:J1}
\]
We bound the tail without invoking a special-function identity.  The global
lower Taylor polynomial of degree $22$ for cosine gives, after termwise
integration on $[0,10]$,
\[
 \operatorname{Cin}(10)\le
 C_{22}:=\sum_{k=1}^{11}{(-1)^{k+1}10^{2k}\over 2k(2k)!}
 ={2702150949482225\over905566992594264}.
\]
The direction of the Taylor bound follows, as above, by repeated
differentiation down to $1-\cos t\ge0$.  Moreover, integration by parts gives
for every $w\ge10$
\[
 \left|\int_{10}^w{\cos t\over t}\,\dd t\right|
 \le {1\over w}+{1\over10}+\int_{10}^w{\dd t\over t^2}
 ={1\over5}.
\]
Consequently, with
\[
 D:=C_{22}+\frac15
 ={14416321740005389\over4527834962971320},
\]
we have
\[
 \operatorname{Cin}(w)\le D+\log(w/10)
\]
and hence, by \eqref{orig:eq:J1},
\[
 J_2(w)\le K(w):=\frac16+
 {2D+2\log(w/10)-4\over w^2}+{4\over w^3}.
\]
Since $D>5/2$,
\[
 K'(w)={10-4D-4\log(w/10)\over w^3}-{12\over w^4}<0
 \qquad(w\ge10).
\]
Finally, exact rational arithmetic yields
\[
 K(10)={43998176831418013\over226391748148566000},
 \qquad
 {243\over1250}-K(10)
 ={61895043316087\over1131958740742830000}>0.
\]
\end{proof}

Lemma~\ref{orig:lem:J2} and \eqref{orig:eq:FJ2} give $|F|\le9w/50$.  For $y\ge1$,
\[
 \frac{z(y-1/4)}4=\frac w4\left(1-\frac1{4y}\right)
 \ge\frac{3w}{16}>\frac{9w}{50}.
\]
Thus \eqref{orig:eq:central2} holds for all $y\ge1$.  Substituting the resulting
pointwise Hermite estimate into \eqref{orig:eq:momentcancel} and using
$\E w_\gamma(Y)=d(Y)$ proves Theorem~\ref{orig:thm:cf}.

\section{The coherent defect representation}

The common defect measure retains the signed Hermite error at every
frequency. Its explicit density and moment cancellation enter the outer,
shape and low-frequency estimates below.

\begin{proof}[Proof of Corollary~\ref{orig:cor:rankfree}]
The function $p\mapsto(1-2p)/\sqrt{p(1-p)}$ is continuous and strictly
decreasing from $+\infty$ to $-\infty$, so the required $p$ exists and is
unique.  From $Y=p+\sqrt{pq}\,X$, $\E X=0$, $\E X^2=1$, and
$\E X^3=\gamma_p$, direct expansion gives
$\E Y=\E Y^2=\E Y^3=p$.

The density is
\[
W_p(y)=
\begin{cases}
 \displaystyle {2q\over p\sqrt{pq}}y^2
       \{3p-(1+2p)y\},&y\le p,\\[2mm]
 \displaystyle {2p\over q\sqrt{pq}}(1-y)^2
       \{(3-2p)y-p\},&y\ge p.
\end{cases}                                             \tag{3.2}
\]
Both bracketed factors are positive on their stated half-lines: at the
common endpoint they equal $2pq$, and they improve away from it.  Hence
$W_p\ge0$, with zero set $\{0,1\}$.  To make the moment cancellation
explicit, put $r=\sqrt{pq}$.  On both half-lines the same polynomial
identity holds:
\[
 r^3\{W_p(y)-|r^{-1}(y-p)|^3\}=\mathcal Q_p(y),
\]
where
\[
 \mathcal Q_p(y)=-p^3+3p^2y+(6p^3-12p^2+3p)y^2
              +(-4p^3+6p^2-1)y^3.
\]
Since $\E Y=\E Y^2=\E Y^3=p$,
\[
 \E \mathcal Q_p(Y)=-pq(p^2+q^2).
\]
Consequently
\[
 \E W_p(Y)=\E|s(Y-p)|^3-\rho_2(\gamma_p)=d.          \tag{3.3}
\]

Let $H_t$ be the cubic Hermite interpolant of
$y\mapsto\e^{2\pi\ii ty}$ at $0$ and $1$, matching values and first
derivatives, and define
\[
 r_{p,t}(y)=\frac{\e^{2\pi\ii ty}-H_t(y)}{W_p(y)}.  \tag{3.4}
\]
The double zeros in numerator and denominator make the singularities at
$0,1$ removable.  As $y\to+\infty$ and $y\to-\infty$, separately, the
exponential term divided by $W_p(y)$ tends to zero, while
$H_t(y)/W_p(y)$ has a finite limit, since both polynomials have degree at
most three and $W_p$ has a nonzero cubic leading coefficient on each tail.
Thus (3.4) has finite, possibly different, limits at the two ends and
extends continuously to $\overline\R$ and is bounded.

If $d>0$, set
\[
 \dd\nu(y)=d^{-1}W_p(y)\dd\Pp_Y(y).
\]
Equation (3.3) says that $\nu$ is a probability measure.  Since $H_t$ is
cubic and $Y$ has the first three moments of a Bernoulli variable $B$ of
parameter $p$,
\[
 \E H_t(Y)=\E H_t(B)=\E\e^{2\pi\ii tB}=b_p(t).
\]
Therefore
\[
 \varphi_Y(2\pi t)-b_p(t)
 =\E\{\e^{2\pi\ii tY}-H_t(Y)\}
 =d\int r_{p,t}(y)\dd\nu(y),
\]
which proves (RF).  If $d=0$, take, for example, $\nu=\delta_0$; its choice
is immaterial because it is multiplied by $d$.  Equation (3.3) and the zero
set of $W_p$ force
$Y\in\{0,1\}$ almost surely; its mean fixes the Bernoulli weights and the
same identity holds with zero right side.

\end{proof}

\section{The standardized two-point branch}\label{orig:sec:binary}

Let $0<p\le1/2$, $q=1-p$, and let $Q_p$ denote the standardized
two-point law which takes the values $\sqrt{q/p}$ and $-\sqrt{p/q}$ with
probabilities $p$ and $q$, respectively.  Equivalently,
$Q_p=Q_{\gamma_p}$ for $\gamma_p=(q-p)/\sqrt{pq}$ in the notation (Q0).
Reflection sends $p$ to $q$, preserves $\rho(p)$ and the Kolmogorov
distance, and therefore makes the restriction $p\le1/2$ harmless.
If $K\sim\operatorname{Bin}(n,p)$, then
\[
 \frac{Q_{p,1}+\cdots+Q_{p,n}}{\sqrt n}
 =\frac{K-np}{\sqrt{npq}},
 \qquad
 \rho(p)=\frac{p^2+q^2}{\sqrt{pq}}.
\]
Writing $z_k=(k-np)/\sqrt{npq}$ and
$B_{n,p}(k)=\Pp\{K\le k\}$, with $B_{n,p}(-1):=0$, constancy of the distribution function
between neighbouring atoms gives the exact identity
\[
 \Dn(Q_p)=\max_{0\le k\le n}\max\left(
 |B_{n,p}(k)-\Phi(z_k)|,
 |B_{n,p}(k-1)-\Phi(z_k)|\right).
\]

Put $s_p=\sqrt{pq}$, $D=p^2+q^2$, and
$\phi_0=(2\pi)^{-1/2}$.  We use the following uniform saddle expansion.
Its purpose is not merely to identify formal coefficients: the displayed
remainder is uniform in every parameter subsequently optimized.

\begin{lemma}[uniform central binomial saddle]\label{orig:lem:binsaddle}
Let $\epsilon=n^{-1/2}\le10^{-24}$, $p\in[2/5,11/25]$, and let
$k=np+\eta$ be an integer with $|\eta|\le1$.  Then
\begin{equation}\label{orig:eq:edge2}
 \frac{B_{n,p}(k)-\Phi(\eta\epsilon/s_p)}{\rho(p)\epsilon}
 =G(p)+\epsilon^2F(p,\eta)+\mathcal R(p,\eta,\epsilon),
\end{equation}
where
\[
 G(p)=\phi_0\frac{2-p}{3D},
\]
\[
 F(p,\eta)=-\frac{\phi_0}{540pqD}
 \{270q\eta^2+45(p^2-7p+5)\eta
          +p^3+21p^2-69p+46\},
\]
and
\[
 |\mathcal R(p,\eta,\epsilon)|
 \le10^{10}\epsilon^{11/4}.
\]
On the symmetric wider compact range $p\in[1/6,5/6]$ the same formula holds with
the bound
\[
 |\mathcal R(p,\eta,\epsilon)|
 \le10^{20}\epsilon^{11/4}.
\]
Under $K\mapsto n-K$, $p\leftrightarrow1-p$ and
$\eta\leftrightarrow-\eta$, this gives the signed expansion at the left
edge of a jump.  Its leading
coefficient is
\[
 G_-(p)=\phi_0\frac{1+p}{3D}\le\phi_0\qquad(0<p\le1/2).
\]
\end{lemma}

\begin{proof}
The exact incomplete-beta identity is
\[
 B_{n,p}(k)=I_q(n-k,k+1)
 =\frac{\int_0^q t^{n-k-1}(1-t)^k\dd t}
        {\int_0^1 t^{n-k-1}(1-t)^k\dd t}.
\]
Set $t=q+\epsilon s_p y$, $A=\sqrt{p/q}$ and $B=\sqrt{q/p}$.
The common constant cancels, and the logarithm of the remaining density is
\[
 L(y)=(nq-\eta-1)\log(1+\epsilon Ay)
      +(np+\eta)\log(1-\epsilon By).                 \tag{S1}
\]
The denominator has limits $-B/\epsilon,A/\epsilon$, while the numerator
ends at zero.  Expanding the two logarithms algebraically gives
\[
 L(y)=-\frac{y^2}{2}+\sum_{j\ge1}\epsilon^j\ell_j(y),
 \quad
 \ell_j(y)=c_jy^{j+2}+d_jy^j,                         \tag{S2}
\]
where
\[
 c_j=\frac{q(-1)^{j+3}A^{j+2}-pB^{j+2}}{j+2},\qquad
 d_j=\frac{-(\eta+1)(-1)^{j+1}A^j-\eta B^j}{j}.       \tag{S3}
\]
Define $E_0=1$ and
\[
 E_m=\frac1m\sum_{j=1}^m j\ell_jE_{m-j}.
\]
Thus $\exp(\sum_{j\ge1}\epsilon^j\ell_j)
=\sum_{m\ge0}\epsilon^mE_m$.  Equations (S2)--(S3) imply
$E_m(-y)=(-1)^mE_m(y)$.  This parity is retained throughout the quotient;
in particular, no spurious order $\epsilon^4$ term is created.

We give the complete remainder ledger.  Put
$R_\epsilon=\epsilon^{-1/12}$ and restrict first to
$|y|\le R_\epsilon$.  Here $\max(A,B)<5/4$ and
$\max(A,B)\epsilon R_\epsilon<1/2$.  The coefficient norm in $y$ obeys
\[
 \|\ell_j\|_1\le
 b_j:=\frac{(5/4)^{j+2}}{j+2}+\frac{3(5/4)^j}{j}
 \quad(1\le j\le4),
\]
and hence, by the displayed recursion,
\[
 \|E_1\|_1<4.402,\quad \|E_2\|_1<12.65,\quad
 \|E_3\|_1<29.9,\quad \|E_4\|_1<62.5.             \tag{S4}
\]
The geometric remainder of the logarithms after $\ell_4$ is at most
\[
 \left(\frac{(5/4)^7}{7}+\frac{3(5/4)^5}{5}\right)
 \frac{\epsilon^5R_\epsilon^7}{1-(5/4)\epsilon R_\epsilon}
 <6\epsilon^5R_\epsilon^7.                           \tag{S5}
\]
If $P_4=\sum_{m=0}^4\epsilon^mE_m$, coefficient majorization by
$\exp(z+z^2+z^3+z^4)$ gives
\[
 \left|\exp\left(\sum_{j=1}^4\epsilon^j\ell_j\right)-P_4\right|
 \le88(12\epsilon R_\epsilon^3)^5.
\]
Together with (S5), this yields
\[
 |e^{L(y)+y^2/2}-P_4(y)|
 \le2.2\,10^7\epsilon^{15/4}.                       \tag{S6}
\]

No tail of the beta integral is hidden in (S6).  Differentiating (S1),
using $pq\ge6/25$ and $n\ge10^{48}$, gives
\[
 L''(y)\le-0.4,\qquad L(0)=0,\qquad |L'(0)|<5\epsilon.
\]
Consequently
\[
 L(y)\le32\epsilon^2-0.2(y-y_0)^2,\qquad |y_0|<12.5\epsilon,
\]
and both omitted tails are smaller than
$8R_\epsilon^{-1}e^{-0.19R_\epsilon^2}<8\epsilon^5$.
The tails of $e^{-y^2/2}P_4(y)$ are smaller than
$10^{40}e^{-R_\epsilon^2/3}<\epsilon^5$.  On $[-1,1]$ one has
$L(y)\ge-1/2-11\epsilon$, so the denominator in the beta quotient is
larger than one.

For $0\le m\le4$ define
\[
 D_m=\int_{\mathbb R}e^{-y^2/2}E_m(y)\dd y,\qquad
 A_m=\int_{-\infty}^0e^{-y^2/2}E_m(y)\dd y,
 \qquad N_4=\sum_{m=0}^4\epsilon^mA_m.
\]
The parity noted above then gives
\[
 N_4=\tfrac12(D_0+\epsilon^2D_2+\epsilon^4D_4)
       +\epsilon A_1+\epsilon^3A_3.
\]
Here is the coefficient calculation.  Write $r=s_p=\sqrt{pq}$, so that
$A=p/r$, $B=q/r$, and put $D=p^2+q^2$.  From (S3),
\[
\begin{array}{lll}
 c_1={p-q\over3r},& d_1=-{\eta+p\over r},
 &E_1=c_1y^3+d_1y,\\[1mm]
 c_2=-{1-3r^2\over4r^2},&d_2={\eta(p-q)+p^2\over2r^2},
 &E_2=c_2y^4+d_2y^2+\tfrac12(c_1y^3+d_1y)^2,\\[1mm]
 c_3={(p-q)D\over5r^3},&
 d_3=-{\eta(1-3r^2)+p^3\over3r^3},
 &E_3=(c_3y^5+d_3y^3)+E_1(c_2y^4+d_2y^2)+\tfrac16E_1^3.
\end{array}                                                   \tag{S6a}
\]
Using
\[
 D_0=\sqrt{2\pi},\qquad
 {1\over D_0}\int_{\mathbb R}y^{2j}e^{-y^2/2}\dd y=(2j-1)!!,\qquad
 \int_{-\infty}^0y^{2j+1}e^{-y^2/2}\dd y=-2^j j!,
\]
one obtains, without any numerical integration,
\begin{align}
 A_1={}&-2c_1-d_1={1\over r}\left\{\eta+{2-p\over3}\right\},
                                                               \tag{S6b}\\
 {D_2\over D_0}={}&3c_2+d_2+{15\over2}c_1^2+3c_1d_1+{d_1^2\over2}
 =:\delta_2,                                                   \tag{S6c}\\
 A_3={}&-8c_3-2d_3-48c_1c_2-8(c_1d_2+d_1c_2)-2d_1d_2\notag\\
 &\hspace{12mm}-64c_1^3-24c_1^2d_1-4c_1d_1^2-{d_1^3\over3}.    \tag{S6d}
\end{align}
In particular, direct collection over the common denominator $540r^3$
gives the checkable identity
\begin{equation}
 A_3-A_1\delta_2+{\eta^3\over6r^3}
 =-{270q\eta^2+45(p^2-7p+5)\eta
          +p^3+21p^2-69p+46\over540r^3}.                       \tag{S6e}
\end{equation}

For clarity, we also spell out the quotient bookkeeping.  Equations
(S4)--(S6), including the two truncated tails, give for the exact
denominator and numerator of the beta quotient
\begin{align*}
 \mathcal D&=D_0+\epsilon^2D_2+\epsilon^4D_4+e_D,\\
 \mathcal N&=\tfrac12(D_0+\epsilon^2D_2+\epsilon^4D_4)
       +\epsilon A_1+\epsilon^3A_3+e_N,
 \qquad |e_D|+|e_N|<1.2\,10^8\epsilon^{15/4}.
\end{align*}
The Gaussian moments and (S4) give
\[
 |A_1|<12,\quad |D_2|<500,\quad |A_3|<10^5,\quad |D_4|<10^8.
\]
Thus, with $x=(\epsilon^2D_2+\epsilon^4D_4)/D_0$, one has
$|x|<201\epsilon^2$.  The exact identity
$(1+x)^{-1}=1-x+x^2/(1+x)$ yields
\begin{equation}
 {\mathcal N\over\mathcal D}
 =\frac12+\epsilon{A_1\over D_0}
 +\epsilon^3\left\{{A_3\over D_0}-{A_1D_2\over D_0^2}\right\}
 +O(1.2\,10^8\epsilon^{15/4})+O(2.1\,10^8\epsilon^5).          \tag{S6f}
\end{equation}
Finally,
\[
 \Phi(\eta\epsilon/r)=\frac12+\phi_0{\eta\epsilon\over r}
       -\phi_0{\eta^3\epsilon^3\over6r^3}+O(10\epsilon^5).
\]
Since $D_0^{-1}=\phi_0$ and $\rho(p)=D/r$, (S6b) gives the coefficient
$G(p)=\phi_0(2-p)/(3D)$, while (S6e) gives exactly the displayed $F$.
After division by $\rho(p)\epsilon\ge\epsilon$, (S6f) gives
$|\mathcal R|<10^{10}\epsilon^{11/4}$.

For the wider interval replace $5/4$ by $9/4$.  Then
\[
 (b_1,b_2,b_3,b_4)<(10.55,14.01,22.93,40.85)
\]
and the same recursion gives
\[
 (\|E_1\|_1,\|E_2\|_1,\|E_3\|_1,\|E_4\|_1)
 <(10.55,69.7,367,1676).
\]
The logarithmic remainder constant in (S5) is then less than $160$, and
the analogue of (S6) is
\[
 |e^{L(y)+y^2/2}-P_4(y)|
 \le9.8\,10^{13}\epsilon^{15/4}.
\]
Since $pq\ge5/36$, the same concavity argument applies (already for
$n\ge120$); the raw quotient error is below
$5\,10^{14}\epsilon^{15/4}$.  The deliberately coarse coefficient bounds
$(10^4,10^8,10^{12},10^{16})$ for $(A_1,D_2,A_3,D_4)$ leave, after
$|x|<40000001\epsilon^2$, and the same exact reciprocal identity leaves a
polynomial residual below $5\,10^{19}\epsilon^5$.  Together with the
$10\epsilon^5$ normal remainder and division by $\rho\epsilon$, this gives
the asserted
$10^{20}\epsilon^{11/4}$ bound.  Reflection
$K\mapsto n-K$, $p\leftrightarrow q$ gives the last assertion.
\end{proof}

\begin{lemma}[bounded-node consequence]\label{orig:lem:boundednode}
Put
\[
 p_*={4-\sqrt{10}\over2},\quad q_*=1-p_*,\quad
 s_*=(p_*q_*)^{-1/2},\quad
 \rho_*={p_*^2+q_*^2\over\sqrt{p_*q_*}},\quad
 \eta_*=-{20+\sqrt{10}\over72},
\]
and $a_{\rm tan,*}=s_*^3q_*\phi_0/2$.
For every fixed $M<\infty$ and every fixed interval
$U$ with $p_*\in U\Subset(2/5,11/25)$, uniformly for
$p\in U$ and $|\eta|\le M$ with $np+\eta\in\mathbb Z$,
\[
 {\sqrt n\over\rho(p)}
 \{B_{n,p}(np+\eta)-\Phi(\eta/\sqrt{npq})\}
 =G(p)+{F(p,\eta)\over n}+o_M(n^{-1}).                \tag{S7}
\]
Consequently, if $p\to p_*$, then
\[
 {\sqrt n\over\rho(p)}
 \{B_{n,p}(np+\eta)-\Phi(\eta/\sqrt{npq})\}
 \le C_Q^*(n)-{a_{\rm tan,*}\over\rho_*n}
             (\eta-\eta_*)^2+o_M(n^{-1}),             \tag{S8}
\]
\end{lemma}

\begin{proof}
In the proof of Lemma~\ref{orig:lem:binsaddle}, $\eta$ occurs only in the
coefficients $d_j$ in (S3).  Replacing $|\eta|\le1$ by $|\eta|\le M$
merely replaces the fixed coefficient majorants by constants depending on
$M$; the concavity and tail arguments are unchanged.  This proves (S7).
To compare with the optimized two-point value without using the theorem
being proved, write
\[
 \eta_0(p)=-{p^2-7p+5\over12(1-p)},\qquad
 a(p)={\phi_0\over2p^{3/2}(1-p)^{1/2}}.
\]
For the integer $k=np+\eta$, set
$p'=p+(\eta-\eta_0(p))/n$.  Then
$k=np'+\eta_0(p)$, so the right-edge score of the two-point law with
parameter $p'$ at this same integer is an admissible competitor in
$C_Q^*(n)$.  Applying (S7) to $p$ and $p'$ and using the exact quadratic
identity
\[
 F(p,\eta)=F(p,\eta_0(p))
       -{a(p)\over\rho(p)}(\eta-\eta_0(p))^2
\]
gives, as $p\to p_*$,
\[
 G(p)-G(p')=o_M(n^{-1}),\qquad
 F(p',\eta_0(p))=F(p,\eta_0(p))+o_M(1).
\]
Consequently the current score is at most that competitor minus
$(a_{\rm tan,*}/\rho_*)
(\eta-\eta_*)^2/n+o_M(n^{-1})$, which is (S8).
\end{proof}

The equation $G'(p)=0$ is $2p^2-8p+3=0$, so the unique maximizer is
\[
 p_*=\frac{4-\sqrt{10}}2,
 \qquad G(p_*)=\CE.
\]
More explicitly,
\[
 G''(p)=\frac{\phi_0}{3D^3}
 (-8p^3+48p^2-36p+4).
\]
On $[2/5,11/25]$ the negative of the displayed cubic is minimized at an
endpoint and is at least $50448/15625>3.2286$.  Since $D\le13/25$ and
$\phi_0>199/500$, it follows that $-G''>3$ throughout that interval.
The quadratic $\eta\mapsto F(p_*,\eta)$ has its vertex at
\[
 r_*=-\frac{20+\sqrt{10}}{72},
\]
and exact substitution gives
\[
 \max_{\eta\in\R}F(p_*,\eta)
 =-\frac{867+287\sqrt{10}}{25920\sqrt{2\pi}}
 =-\kappa_\infty.
\]

\begin{proof}[Proof of Theorem~\ref{orig:thm:twopoint}]
For the lower bound choose an integer $k_n$ nearest to $np_*+r_*$ and put
$p_n=(k_n-r_*)/n$.  Then $p_n=p_*+O(n^{-1})$ and
$k_n-np_n=r_*$.  Equation \eqref{orig:eq:edge2} gives
\[
 C_Q^*(n)\ge \CE-\frac{\kappa_\infty}{n}+o(n^{-1}).
\]

For completeness we spell out the global localization used for the upper
bound.  Proposition~\ref{orig:prop:schulzlocal} states the imported theorem with
all four edge terms and the normalization used here.  In particular, for
$p\in[1/6,5/6]$ it bounds the normalized distance by those four terms and
$c_0\widehat G(p)$, where $c_0<0.948$, and on
$p\in[1/3,2/3]$ it gives the strict bound by $\widehat G(p)$.  On the
half-range $p\le1/2$ one has $\widehat G(p)=G(p)$.  For $p\le1/6$ the
universal estimate
\cite{Shevtsova2013}
\[
 \frac{\sqrt n\Delta_n(Q_p)}{\rho(p)}
 \le0.3322\left(1+\frac{0.429}{\rho(p)}\right)
 \le0.405740
\]
has a fixed gap from $\CE$; reflection covers $p\ge5/6$.  The term
$c_0\widehat G(p)$ likewise has a fixed gap.  The imported strict bound by $G$,
together with continuity and the unique maximum at $p_*$, first eliminates
the compact set $[1/3,1/2]\setminus(2/5,11/25)$.

It remains to inspect the four central edges on
$[1/6,1/3]\cup[2/5,11/25]$, all of which lie in the wider parameter range
of Lemma~\ref{orig:lem:binsaddle}.  On the left edges the reflected leading
coefficient satisfies $G_-(p)\le\phi_0<\CE$; on right edges with
$p\in[1/6,1/3]$ one has $G(p)\le G(1/3)<\CE$.  These positive leading
coefficients also fix the signs of the corresponding signed expansions for
all sufficiently large $n$.  Thus every sequence attaining the supremum up
to $o(n^{-1})$ eventually consists of right central endpoints with
$p\in[2/5,11/25]$.  On that interval one has
$-G''(p)>3$.  Comparison with the explicit lower-bound sequence therefore gives
$p-p_*=O(\epsilon)$.  Write $p=p_*+u\epsilon$; the numbers $u$ remain
bounded, and the central-point theorem gives $|\eta|\le1$.  Lemma
\ref{orig:lem:binsaddle}, continuity of $F$, and Taylor's theorem now give,
uniformly on bounded $u$ and $|\eta|\le1$,
\[
 \frac{\sqrt n\Delta_n(Q_p)}{\rho(p)}
 \le \CE+\epsilon^2
 \left\{\tfrac12G''(p_*)u^2+F(p_*,\eta)\right\}+o(\epsilon^2).
\]
The first term in braces is nonpositive, while the concave quadratic
$F(p_*,\eta)$ has maximum $-\kappa_\infty$ at $\eta=r_*$.  Taking upper
limits proves
\[
 \limsup_{n\to\infty}n(C_Q^*(n)-\CE)\le-\kappa_\infty.
\]
Together with the lower bound this proves the theorem.
\end{proof}

\section{The original eventual upper bound}\label{orig:sec:eventual}

We now prove Theorem~\ref{orig:thm:eventual} with its natural qualitative
quantifiers.  Put
\[
 p_*={4-\sqrt {10}\over2},\qquad q_*=1-p_*,
 \qquad s_*=(p_*q_*)^{-1/2},
\]
\[
 r_{\rm E}=1-2p_*,\qquad \gamma_*=s_*r_{\rm E},
 \qquad \eta_*:=r_*=-{20+\sqrt {10}\over72},
\]
 and let $\rho_*:=\rho(p_*)$.  We write
 \[
                  Q_*:=Q_{p_*}=Q_{\gamma_*}.             \tag{4.0}
 \]
 For a centered variance-one law $X$ write
\[
 \mathcal C_n(X)={\sqrt n\,\Delta_n(X)\over\rho},\qquad
 \gamma=\E X^3,\qquad d=\rho-\rho_2(\gamma).
\]
After reflection we use the unique $p\le1/2$ for which
$\gamma=(1-2p)/\sqrt{p(1-p)}$.  Thus $d$ is precisely the defect in
Corollary~\ref{orig:cor:rankfree}.  All little-oh statements below are sequential:
they are uniform along every triangular array satisfying the hypotheses in
which they occur.  This is enough for the final $\exists N_0\,\forall n$
statement and avoids introducing an artificial cutoff.

We use the following terminology throughout this section.  A one-sided cdf
point $(x,\sigma)$ in row $(n,X)$ is \emph{active} if its normalized signed
score $\sqrt nD_X^\sigma(x)/\rho$ is larger than $\CE$.  A positive cell is
the set of points, as $|\theta|\le1/2$ varies, written as
\[
 x={s(\eta+\theta)\over\sqrt n},\qquad np+\eta\in\mathbb Z,
 \qquad |\theta|\le\tfrac12.                              \tag{4.0a}
\]
(At a cell boundary either adjacent representation may be chosen.)  The
cell is \emph{active} if it contains an active point; when $(\eta,\theta)$
is displayed, ``active'' refers to that particular point.  A
triangular array of points is said to lie on the \emph{localized positive
branch} when $p_n\to p_*$, $d_n\to0$, its sign is positive, and $x_n\to0$.
Thus ``active'' records the inequality being contradicted, whereas
``localized'' records precisely the four sequential conclusions just
listed; neither term hides an additional estimate.

\subsection*{Outer stability}

\begin{proposition}[outer stability]\label{orig:prop:qouter}
Let $n_j\to\infty$ and let $X_j$ be centered and standardized.  If
\[
                    \mathcal C_{n_j}(X_j)>\CE,         \tag{OS0}
\]
then, after reflection and passage to a subsequence,
\[
                    p_j\longrightarrow p_*,\qquad d_j\longrightarrow0.
                                                               \tag{OS}
\]
Moreover every sequence of cdf points whose normalized one-sided error
tends to $\CE$ lies in the positive central branch and tends to zero in
the standardized cdf coordinate.
\end{proposition}

\begin{proof}
The published structural estimate \cite{Shevtsova2013}
\[
 {\sqrt n\,\Delta_n(X)\over\rho}
 \le {1661\over5000}\left(1+{429\over1000\rho}\right)       \tag{4.1}
\]
shows, by direct comparison at $\rho=37/20$, that every row in (OS0)
has $\rho<37/20$.  Hence its moment parameters are precompact.

Put
\[
 r={\gamma\over\sqrt{4+\gamma^2}},\qquad u=\sqrt{1-r^2},
 \qquad s={2\over u},\qquad
 \varrho_2(r):=\rho_2(2r/u)={1+r^2\over u},             \tag{4.2}
\]
and set $\lambda=\rho+\E|X|$.  Proposition~\ref{orig:prop:m1} gives
\[
 s\le\lambda\le s+{2d\over(1-r)(2+r)}.                 \tag{4.3}
\]
We record why the bandwidth $\lambda$ removes every nonzero leading
Fourier cell.  If $X'$ is an independent copy and $D=X-X'$, then
\[
 z^2\le |z|^3+{ |\sin(\pi z)|\over\pi}\qquad(z\in\R). \tag{4.4}
\]
For $|z|\ge1$ this is immediate; on $[0,1]$ it follows from
$\sin(\pi z)\ge\pi z(1-z)$, and symmetry finishes the proof.  For $H>0$
apply (4.4) to $D/H$.  Independence and centering give
\[
 \E|D|^3\le2(\rho+\E|X|)=2\lambda,\qquad
 \E\sin^2(\pi D/H)={1-|\varphi_X(2\pi/H)|^2\over2}.
\]
Here the first inequality is elementary but worth recording: pointwise,
\[
 |x-y|^3\le(x-y)^2(|x|+|y|).
\]
After expansion with $(x,y)=(X,X')$, the two mixed terms containing
$\E X$ vanish, while the remaining terms are
$2\E|X|^3+2\E X^2\E|X|=2(\rho+\E|X|)$.
Cauchy--Schwarz therefore yields
\[
 H\le\lambda+{H^3\over2\pi\sqrt2}
              \sqrt{1-|\varphi_X(2\pi/H)|^2}.          \tag{4.5}
\]
With $H=\lambda/t$ this becomes
\[
 \sqrt{1-|\varphi_X(2\pi t/\lambda)|^2}
 \ge {2\pi\sqrt2\over\lambda^2}t^2(1-t).             \tag{4.6}
\]
Indeed $2\le\lambda\le37/20+1=57/20$.  Fix temporarily
$0<\tau<1/4$.  Squaring (4.6) and using $1-y\le e^{-y}$ gives constants
$q_\tau<1$ and $c_\tau>0$, independent of the row, such that
\[
 \begin{array}{ll}
 |\varphi_X(2\pi t/\lambda)|^n\le q_\tau^n,
      &\tau\le t\le1-\tau,\\[1mm]
 |\varphi_X(2\pi t/\lambda)|^n
      \le e^{-c_\tau n(1-t)^2},&1-\tau\le t\le1.
 \end{array}                                                   \tag{4.6a}
\]
This is the common Cram\'er-gap partition used below.  Since $K$ is bounded
on the middle interval, that integral is $O(q_\tau^n)$.  Since
$|K(1-h)|\le C_\tau h$, the terminal integral is bounded by
\[
 C_\tau\int_0^\tau h e^{-c_\tau nh^2}\,\dd h=O(n^{-1}).           \tag{4.6b}
\]
After multiplication by $\sqrt n$ both pieces vanish.  We subsequently
shrink $\tau$, uniformly in $\lambda\in[2,57/20]$, so that the remaining
zero cell $[0,\tau]$ lies in the small-frequency region of (4.7).

We next evaluate the zero cell uniformly for triangular arrays.  This is
the point at which uniformity in the cdf coordinate and in a possibly
nonvanishing defect is needed.  We give the replacement argument
explicitly.  Apply the two one-sided Prawitz operators with physical span
parameter $h=\lambda$, equivalently with Fourier cutoff
$T_{n,\lambda}=2\pi\sqrt n/\lambda$.  At a standardized cdf coordinate $x$
their Fourier factors are
$e^{-\ii T_{n,\lambda}tx}\varphi_X(2\pi t/\lambda)^n$; hence every estimate below
is unchanged when $x$ varies.  Suppress the row index and write
\[
 \varphi(v)=\E e^{\ii vX},\qquad
 \chi(v)=\E e^{\ii vQ_\gamma}.
\]
The moment cap above bounds, uniformly in the row, the third absolute
moments of both $X$ and $Q_\gamma$.  Symmetrization and
$1-\cos z\ge z^2/2-|z|^3/6$ therefore give constants $c,\delta>0$ such
that
\[
 \max\{|\varphi(v)|,|\chi(v)|\}\le e^{-cv^2}
 \qquad (|v|\le\delta).                                \tag{4.7}
\]
Indeed, if $Z'$ is an independent copy, then $\E(Z-Z')^2=2$ and
\[
 1-|\E e^{\ii vZ}|^2
 =\E\{1-\cos(v(Z-Z'))\}
 \ge v^2-{|v|^3\over6}\E|Z-Z'|^3.                       \tag{4.7a}
\]
The just-proved difference inequality and the common third-moment bound
make the last expectation uniformly bounded for $Z=X,Q_\gamma$.
Choose $\delta$ so that the right side of (4.7a) is at least $v^2/2$.
Then $|\E e^{\ii vZ}|^2\le1-v^2/2\le e^{-v^2/2}$, which proves (4.7),
for instance with $c=1/4$.

Use the exact divided power
\[
 \mathscr P_n(v)={\varphi(v)^n-\chi(v)^n\over
                n\{\varphi(v)-\chi(v)\}}
       =\int_0^1\{(1-a)\chi(v)+a\varphi(v)\}^{n-1}\,\dd a,          \tag{4.8}
\]
with its continuous value if the denominator vanishes.  Since both
transforms equal $1-v^2/2+O(|v|^3)$, with a row-independent constant,
the difference-of-powers identity and (4.7) imply, for
$0\le w\le V_n:=n^{1/21}$,
\[
 \left|\mathscr P_n(w/\sqrt n)-e^{-w^2/2}\right|
 \le C\left\{{w^3\over\sqrt n}+{w^2+w^4\over n}\right\}e^{-cw^2}.
                                                               \tag{4.9}
\]
On $V_n<w\le\delta\sqrt n$ the left side is instead bounded by
$2e^{-cw^2}$.  These bounds are independent of the cdf phase.  The
Prawitz substitution $t=\lambda w/(2\pi\sqrt n)$ also gives, uniformly
on the zero cell,
\[
 2|K(t)|\,\dd t\le
       \left\{{1\over\pi w}+{C\over\sqrt n}\right\}\dd w.          \tag{4.10}
\]
Consequently every error in replacing the exact divided power by its
Gaussian limit is dominated by an integrable Gaussian polynomial.  In
particular, this replacement is simultaneous in the cdf coordinate
$x$, because the omitted factor $e^{-\ii wx}$ has modulus one.

It remains to identify the contribution of $\varphi-\chi$.  In the
Bernoulli coordinate $Y=p+\sqrt{p(1-p)}X$ of
Corollary~\ref{orig:cor:rankfree}, the two cubic tails of $W_p$ have leading
coefficients
\[
 C_+(p)={2p(3-2p)\over(1-p)\sqrt{p(1-p)}},\qquad
 C_-(p)={2(1-p)(1+2p)\over p\sqrt{p(1-p)}}.             \tag{4.10a}
\]
The precompact moment range keeps $p$ away from $0$ and $1$.  Let $P_z$
be the cubic Hermite interpolant of $e^{\ii zy}$ at $0,1$ and put
$E_z(y)=e^{\ii zy}-P_z(y)$.  Moment matching and the common centering
phase give
\[
 \varphi(v)-\chi(v)=e^{-\ii spv}\E E_{sv}(Y).
\]
Hermite--Genocchi and the explicit branches of $W_p$ give, uniformly in
this compact $p$-range,
\[
 |E_z(y)|\le {|z|^4\over24}y^2(1-y)^2,\qquad
 {y^2(1-y)^2\over(1+|y|)W_p(y)}\le C.                 \tag{4.10b}
\]
For the tails write
\[
 g(v)=e^{\ii v}-1-\ii v+{v^2\over2}+{\ii v^3\over6}.
\]
The phase-uniform bound
\[
                         |g(v)|\le {|v|^3\over6}       \tag{4.10c}
\]
follows by setting $A=\cos v-1+v^2/2$,
$B=\sin v-v+v^3/6$ and observing
\[
 {\dd\over\dd v}\{v^6/36-A^2-B^2\}
       ={v^3\over3}(1-\cos v)\ge0\qquad(v\ge0);
\]
conjugation covers $v<0$.  Moreover the Hermite basis gives
\[
 |E_z(y)-g(zy)|\le C|z|^4(1+|y|^3)                    \tag{4.10d}
\]
outside a fixed compact interval.

Split the defect probability
$d^{-1}W_p(y)\,\dd\Pp_Y(y)$ at $H_n=n^{1/8}$ (with the assertion
trivial when $d=0$).  Equations (4.10b) and (4.10d), followed by
(4.9)--(4.10), make the part $|y|\le H_n$ and the interpolant correction
on $|y|>H_n$ at most $CdH_n/\sqrt n$ and $Cd/\sqrt n$, respectively,
in normalized score.  The divided-power replacement itself costs
$O(d/\sqrt n)$, by the characteristic-function stability theorem and
the Gaussian moments in (4.9).  These quantities are $o(1)$ uniformly
in the row and in $x$.  On the remaining tail, (4.10c) and the exact cubic leading
terms in (4.10a) give
\[
 {n^{3/2}\over d}
 \left|\E\{E_{sw/\sqrt n}(Y);|Y|>H_n\}\right|
 \le {s^3|w|^3\over6}
       \left[\max\{C_+(p)^{-1},C_-(p)^{-1}\}+o(1)\right]
       +C{w^4\over\sqrt n}.                              \tag{4.10e}
\]
Here the $o(1)$ is independent of the row and of $w$ on
$0\le w\le V_n$.  Beyond $V_n$ the
characteristic-function stability theorem and the second bound following
(4.9) give a vanishing Gaussian tail.  Thus dominated convergence applies
to the full zero-cell integral, without first assuming $d\to0$.  Since
$p\le1/2$, one has
\[
 {C_-(p)\over C_+(p)}-1
 ={(1-2p)(1+2p-2p^2)\over p^2(3-2p)}\ge0;
\]
hence the maximum in (4.10e) is $C_+(p)^{-1}$, and
\[
 {1\over\pi}\int_0^\infty w^2e^{-w^2/2}\,\dd w=\phi_0,
 \qquad {s^3\over6C_+(p)}={1\over3(1-r)^2(2+r)}.        \tag{4.10f}
\]

To make the replacement just proved directly reusable, its conclusion is
the following uniform inequality.  With the fixed zero-cell width $\tau$
chosen after (4.6b), there is $\varepsilon_n\to0$, independent of the row,
the cdf coordinate, and the choice $\kappa\in\{K,\overline K\}$, such that
\begin{align}
&\sqrt n\sup_{x\in\mathbb R}\left|
 2\Re\int_0^\tau\kappa(t)e^{-\ii T_{n,\lambda}tx}
 \left\{\varphi(2\pi t/\lambda)^n-
              \chi(2\pi t/\lambda)^n\right\}\dd t\right| \notag\\
&\hspace{35mm}\le
 d\left\{{\phi_0s^3\over6C_+(p)}+\varepsilon_n\right\}.          \tag{4.10h}
\end{align}
When $d=0$ both sides vanish.  Indeed the part $w\le V_n$ is exactly
the product of (4.8), (4.9), and (4.10e); (4.10) supplies the factor
$1/(\pi w)$, and (4.10f) evaluates its Gaussian majorant.  The terms
$CdH_n/\sqrt n$, $Cd/\sqrt n$, and the tail $w>V_n$ listed above are
$\varepsilon_nd$.  This proves (4.10h), rather than merely a pointwise
limit for a fixed $x$.

The complementary frequencies have already been partitioned in
(4.6a)--(4.6b): their raw integrals are $O(q_\tau^n)+O(n^{-1})$ for either
kernel.  The estimates discard only a spatial factor of modulus one, so
after score normalization they are $o(1)$ uniformly in $x$.

The matching-law zero cell follows from the same expansion, now one order
further.  Uniformly for $0\le w\le V_n$,
\[
 \chi(w/\sqrt n)^n=e^{-w^2/2}
       \left(1-{\ii\gamma w^3\over6\sqrt n}\right)
       +O\left({(w^4+w^6)e^{-cw^2}\over n}\right).
\]
For $V_n<w\le\delta\sqrt n$ both the exact power and the displayed
Gaussian polynomial have a vanishing Gaussian tail by (4.7).
In signed form the kernel expansion used here is
\[
 2K(t)\,\dd t=\left\{{\ii\over\pi w}
       +{\lambda\over2\pi\sqrt n}\right\}\dd w
       +O(w/n)\,\dd w,
\]
uniformly throughout the fixed zero cell, with the conjugate expansion for the
other one-sided operator.  Together with normal inversion, this gives,
uniformly for all real $x$, the two profiles
\[
 {\phi(x)\over\rho}
       \left\{{\lambda\over2}
          \mathbin\pm {\gamma\over6}(1-x^2)\right\}+o(1).          \tag{4.10g}
\]
Because $\lambda>|\gamma|$, the positive maximum of either profile is at
$x=0$, and the plus profile is the larger one.  Hence (P11), (4.10h), and
the vanishing complementary cells give the explicit uniform assembly
\[
 \mathcal C_n(X)\le {\phi_0\over\rho}
 \left\{{\lambda\over2}+{\gamma\over6}
                 +{s^3\over6C_+(p)}d\right\}+o(1).       \tag{4.10i}
\]
Now use (4.3), $\rho=\varrho_2(r)+d$, $s=2/u$, and $\gamma=2r/u$:
\begin{align*}
 {\lambda\over2}+{\gamma\over6}
                 +{s^3d\over6C_+(p)}
 &\le {s\over2}+{\gamma\over6}
  +\left\{{1\over(1-r)(2+r)}+{s^3\over6C_+(p)}\right\}d\\
 &= {3+r\over3u}
  +\left\{{1\over(1-r)(2+r)}
       +{1\over3(1-r)^2(2+r)}\right\}d.
\end{align*}
Passing to any subsequence on which $(r_j,d_j)\to(r,d)$ therefore proves
\[
 \limsup_j\mathcal C_{n_j}(X_j)
 \le \phi_0{N(r)+L(r)d\over\varrho_2(r)+d},            \tag{4.11}
\]
where
\[
 N(r)={3+r\over3u},\qquad
 L(r)={1\over(1-r)(2+r)}
      +{1\over3(1-r)^2(2+r)}.                          \tag{4.12}
\]

It remains to prove the strict scalar assertion behind (4.11).  For
$d=0$ it is the elementary maximization
\[
 \phi_0{3+r\over3(1+r^2)}\le\CE,                       \tag{4.13}
\]
with equality only at $r=r_{\rm E}$.  Indeed the derivative of
$(3+r)/(3(1+r^2))$ has the sign of $1-6r-r^2$, whose only nonnegative
zero is $\sqrt {10}-3=r_{\rm E}$; substitution there gives
$(3+\sqrt {10})/6$.

For fixed $r$ the quotient in (4.11) is fractional linear in $d$, hence
monotone or constant, so it suffices to inspect the two endpoints of the
allowed interval.  Enlarge the moment cap to $R=37/20$.  Since
$\varrho_2'(r)=r(3-r^2)/(1-r^2)^{3/2}>0$ and
$\varrho_2(2/3)>R$, every admissible row has $r<2/3$ and
$0\le d<R-\varrho_2(r)$.  It remains to prove, for $0\le r\le2/3$,
\[
 {N(r)+L(r)\{R-\varrho_2(r)\}\over R}<{41\over40}
 <{3+\sqrt {10}\over6}.                                \tag{4.14}
\]
The second inequality is equivalent to $\sqrt {10}>63/20$, and follows on
squaring because $10>3969/400$.  For the first put
\[
 P=r^4+6r^3-7r^2-4r+2,\qquad
 B={123r^3-249r+86\over40}.
\]
The algebraic equivalence is checkable without expansion:
\[
 {41R\over40}-N-L\{R-\varrho_2(r)\}
 ={RuB-P\over3u(1-r)^2(2+r)}.
\]
Thus the first inequality in (4.14) is exactly
\[
                         RuB-P>0.                       \tag{4.15}
\]
On $0\le r\le37/100$, $B>0$ and $u\ge1-r^2$, so the left side of
(4.15) is at least $f=R(1-r^2)B-P$.  Writing $800f=F$ gives
\[
 F=-4551r^5-800r^4+8964r^3+2418r^2-6013r+1582.
\]
Here is an exact derivative ledger.  Since
\[
 B'(r)={369r^2-249\over40}<0,\qquad
 B(37/100)={100319\over40000000}>0,
\]
the asserted sign of $B$ holds.  Furthermore
\[
 F'''=53784-19200r-273060r^2,\qquad
 F''''=-19200-546120r<0,
\]
and
\[
 F'''(37/100)={4649043\over500}>0,\quad F''(0)=4836,
\]
\[
 F'(37/100)=-{22614402711\over20000000}<-1000,\qquad
 F(37/100)={957159691693\over10^{10}}>80.
\]
Thus $F'''>0$, then $F''>0$, on the entire interval; consequently $F'$ is
increasing but remains negative, and $F$ is decreasing with a positive
right-end value.

For $r\ge37/100$ one has $P<0$.  Indeed
$P''(37/100)=2407/2500>0$, $P'''=24r+36>0$, and
$P'(2/3)=-112/27<0$; hence $P'<0$ on this interval, while
$P(37/100)=-11564039/10^8<0$.  If $B\ge0$, (4.15) is immediate.  If
$B<0$, use
\[
 u\le U:=1-{r^2\over2}-{r^4\over8},\qquad
 U^2-(1-r^2)={r^6(r^2+8)\over64},
\]
so that $RuB-P\ge h:=RUB-P$.  Since $B$ is decreasing,
\[
 B(37/100)={100319\over40000000},\qquad B(3/8)=-{91\over4096},
\]
and the preceding monotonicity of $P$ gives, on $[37/100,3/8]$,
$-P>1/10$ and $|B|<1/40$.  As $0<U\le1$, this yields
$h>1/10-37/800=43/800>1/20$ whenever $B<0$.

For the remaining derivative checks, direct multiplication gives
\[
\begin{split}
h(r)={}&{791\over400}-{6013\over800}r+{4009\over800}r^2
 +{1743\over320}r^3-{4791\over3200}r^4\\
&-{8991\over6400}r^5-{4551\over6400}r^7.
\end{split}                                           \tag{4.15a}
\]
Thus
\[
 h''''(r)=-{14373\over400}-{26973\over160}r
                  -{95571\over160}r^3<0.                \tag{4.15b}
\]
The function $h''$ is concave, so its minimum on $[3/8,1/2]$ is at an
endpoint; exact substitution gives
\[
 h''(3/8)={1892502967\over104857600}>17,\qquad
 h''(1/2)={1784437\over102400}>17.                       \tag{4.15c}
\]
At $r_0=12/25$,
\[
 h(r_0)={17639835737\over2441406250000}>{7\over1000},\qquad
 h'(r_0)=-{7154739661\over195312500000},\quad
 |h'(r_0)|<{1\over25}.                                  \tag{4.15d}
\]
Taylor's formula and $h''\ge17$ yield
\[
 h(r)\ge h(r_0)-{|h'(r_0)|^2\over34}
 >{7\over1000}-{1\over34\cdot25^2}>0.                  \tag{4.15e}
\]
Finally (4.15b) and
$h'''(1/2)=-803367/51200<0$ show that $h'$ is concave on
$[1/2,2/3]$, while
\[
 h'(1/2)={128931\over409600}>{1\over4},\qquad
 h'(2/3)={549953\over194400}>2.
\]
A concave function lies above its endpoint chord, so $h'>0$ throughout
this interval; hence $h$ is increasing there.  This
proves (4.14), and consequently the second inequality in (4.11), with
equality only at $(r,d)=(r_{\rm E},0)$.

Every score in (OS0) is larger than $\CE$.  Thus (4.11) first forces the
scores to converge to $\CE$, and its equality statement forces
$r_j\to r_{\rm E}$ and $d_j\to0$, equivalently $p_j\to p_*$.  At this
point let
\[
 P_\pm(x)={\phi(x)\over\rho_*}
 \left\{{s_*\over2}\mathbin\pm{\gamma_*\over6}(1-x^2)\right\}.
\]
The derivative calculation used after (4.10g) shows that $P_+$ has its
unique global maximum $P_+(0)=\CE$ at zero, whereas
\[
 \sup_xP_-(x)=P_-(0)=\CE-{\phi_0\gamma_*\over3\rho_*}<\CE.
\]
Consequently, for every fixed $\epsilon>0$, compactness together with the
Gaussian tails supplies $c_\epsilon>0$ such that
\[
 \sup_{|x|\ge\epsilon}P_+(x)\le\CE-c_\epsilon,
 \qquad \sup_xP_-(x)\le\CE-{\phi_0\gamma_*\over3\rho_*}.
\]
The convergence in (4.10g) and the replacement bound (4.10h) are uniform
in $x$.  Hence an above-$\CE$ point must have the plus sign and
$x_j\to0$.  This proves every assertion of
Proposition~\ref{orig:prop:qouter}.
\end{proof}

\subsection*{The certified shape input and the low-frequency defect}

The only computer-assisted premise in the eventual argument is the next
one-dimensional inequality.  It is stated separately so that its trust
boundary is not mixed with the analytic proof.

\begin{lemma}[shape--contraction]\label{orig:lem:shapecontraction}
For $0<p<1$ put $q=1-p$, and let $H_p$ be the cubic Hermite interpolant at
$0,1$ of
\[
                         g_p(y)=(y-p)e^{2\pi\ii y}.
\]
Put $E_p^{\rm sh}=g_p-H_p$ and
\[
 W_p(y)=
 \begin{cases}
 \displaystyle {2q\over p\sqrt{pq}}y^2
       \{3p-(1+2p)y\},&y\le p,\\[2mm]
 \displaystyle {2p\over q\sqrt{pq}}(1-y)^2
       \{(3-2p)y-p\},&y\ge p.
 \end{cases}                                             \tag{SC0}
\]
At the certified parameter $p=p_*$, for every real $y$,
\[
 \boxed{\quad |E_p^{\rm sh}(y)|
       \le {19\over16}W_p(y)+|e^{2\pi\ii y}-1|^2.\quad} \tag{SC}
\]
Only the compact inequality at the single parameter $p_*$ is
computer-certified: the checker uses exact outward-rounded dyadic interval
arithmetic on $187$ accepted intervals.  The tail proof and the passage from
$p_*$ to nearby $p$ given below are analytic.  No floating-point number
enters an acceptance decision.
\end{lemma}

\begin{proof}
The exact real and imaginary parts are
\[
 \Re E_p^{\rm sh}(y)=(y-p)(\cos2\pi y-1),
\]
\[
 \Im E_p^{\rm sh}(y)=(y-p)\sin2\pi y
 -2\pi\{(1-2p)y^3+(3p-1)y^2-py\}.                    \tag{SC1}
\]
The tails require no interval arithmetic.  Direct substitution in the
radical defining $p_*$, followed only by squaring positive quantities,
gives
\[
 .418<p_*<.419,\quad .581<q_*<.582,\quad
 0<1-2p_*<.164,\quad \sqrt{p_*q_*}<1/2,
 \quad 2\pi<44/7,
\]
and
\[
 {2q_*\over p_*\sqrt{p_*q_*}}>{2324\over419},\qquad
 {2p_*\over q_*\sqrt{p_*q_*}}>{836\over291}.
\]
Consequently, with $y=-x$ and $x\ge1$, (SC0)--(SC1) give
\[
 |E_p^{\rm sh}(-x)|<{11\over10}x^3+9x^2,
 \qquad W_p(-x)>10x^3+6x^2.
\]
Their difference is positive at $x=1$ and has positive derivative
thereafter.  With $y=1+x$ and $x\ge2$, the same expansion gives
\[
 |E_p^{\rm sh}(1+x)|<{11\over10}x^3+{47\over10}x^2
                       +{67\over10}x+{9\over5},
\]
\[
 W_p(1+x)>{31\over5}x^3+5x^2.
\]
The difference is positive at $x=2$ and has positive derivative
thereafter.  Thus the stronger inequality $|E_p^{\rm sh}|<W_p$ holds
outside $[-1,3]$.

For completeness we describe the exact compact certificate.  Put
\[
 c(t)={1-\cos(2\pi t)\over t^2},\quad
 s_0(t)={\sin(2\pi t)\over2\pi t},\quad
 b(t)={1-s_0(t)\over t}.
\]
On the left branch the checker divides the triangle gap
$(19/16)W_p+|e^{2\pi\ii y}-1|^2-|\Re E_p^{\rm sh}|-|\Im E_p^{\rm sh}|$
by $y^2$ and evaluates the regularized identities
\[
 {\Re E_p^{\rm sh}(y)\over y^2}=-(y-p)c(y),
\]
\[
 {\Im E_p^{\rm sh}(y)\over y^2}
 =2\pi\{s_0(y)-(1-2p)y-(3p-1)+pb(y)\},               \tag{SC2}
\]
\[
 {W_p(y)\over y^2}={2q\over p\sqrt{pq}}
       \{3p-(1+2p)y\},\qquad
 {|e^{2\pi\ii y}-1|^2\over y^2}=2c(y).
\]
The right branch is the same calculation after $t=y-1$, divided by
$t^2$.  The entire alternating series of $s_0,b,c$ remove the apparent
singularities.  The checker encloses $\sqrt {10}$ by integer square roots,
encloses $\pi$ through Machin's identity and alternating rational series,
and rounds every rational operation outward on a fixed dyadic grid.  It
subdivides only when the resulting lower endpoint of the normalized
triangle gap is not yet positive.  All $187$ terminal intervals have a
strictly positive rational lower endpoint, cover $[-1,3]$, and overlap
across the irrational branch point $y=p_*$.  Since both terms on the
right of (SC) are nonnegative and $|E|\le|\Re E|+|\Im E|$, positivity of
the triangle gap proves (SC).
The endpoints $y=0,1$ follow directly from (SC1).
\end{proof}

We next make precise the continuity step which transports this
single-parameter certificate. This continuity step is analytic.  For $p$ near $p_*$ and $y\ne0,1$ define
\[
 \mathscr R(p,y):=
 {\bigl(|E_p^{\rm sh}(y)|-|e^{2\pi\ii y}-1|^2\bigr)_+\over W_p(y)}.
                                                               \tag{SC3}
\]
Here $W_p>0$ off $\{0,1\}$.  The following endpoint expansions are
uniform when $p$ ranges over a sufficiently small closed neighbourhood of
$p_*$:
\begin{align*}
 E_p^{\rm sh}(y)
 &=\{2\pi^2p+2\pi\ii(2-3p)\}y^2+O(|y|^3),\qquad y\to0,\\
 |e^{2\pi\ii y}-1|^2&=4\pi^2y^2+O(y^4),\qquad
 W_p(y)={6(1-p)\over\sqrt{p(1-p)}}y^2+O(|y|^3),       \tag{SC4}\\
 E_p^{\rm sh}(1+t)
 &=\{-2\pi^2(1-p)+2\pi\ii(3p-1)\}t^2+O(|t|^3),\qquad t\to0,\\
 |e^{2\pi\ii(1+t)}-1|^2&=4\pi^2t^2+O(t^4),\qquad
 W_p(1+t)={6p\over\sqrt{p(1-p)}}t^2+O(|t|^3).         \tag{SC5}
\end{align*}
At $p=p_*$ the moduli of the two complex coefficients in (SC4)--(SC5)
are strictly smaller than $4\pi^2$ (squaring and using
$.418<p_*<.419$ verifies both inequalities).  They therefore remain
strictly smaller on a fixed neighbourhood of $p_*$.  Thus
$\mathscr R$ is identically zero in fixed punctured neighbourhoods of
$0$ and $1$, and extends there jointly continuously by the value zero.
On the two tails, uniformly in the same $p$-neighbourhood,
\begin{align*}
 E_p^{\rm sh}(y)&=-2\pi\ii(1-2p)y^3+O(y^2),\\
 W_p(y)&=C_+(p)y^3+O(y^2) &&(y\to+\infty),\\
 W_p(y)&=C_-(p)|y|^3+O(y^2) &&(y\to-\infty),          \tag{SC6}
\end{align*}
where $C_\pm$ are defined in (4.10a).  Hence the two remaining removable
limits (in the notation of (SC3)) are
\[
 \mathscr R(p,+\infty)={2\pi(1-2p)\over C_+(p)},\qquad
 \mathscr R(p,-\infty)={2\pi(1-2p)\over C_-(p)}.       \tag{SC7}
\]
Equations (SC1) and (SC4)--(SC7), including the agreement of the two
branches of $W_p$ at $y=p$, show that $\mathscr R$ is jointly continuous
on $U_0\times\overline{\mathbb R}$ for some compact
$U_0\Subset(0,1/2)$ containing $p_*$.  Since the second factor is the
two-point compactification, uniform continuity and (SC) give a modulus
$\omega(\delta)\downarrow0$ such that
\[
 |E_p^{\rm sh}(y)|\le
 \{19/16+\omega(|p-p_*|)\}W_p(y)+|e^{2\pi\ii y}-1|^2             \tag{SC8}
\]
for all $p\in U_0$ and $y\in\mathbb R$.

For a Bernoulli-coordinate law with the first three moments equal to $p$,
write
\[
 \zeta=\E e^{2\pi\ii Y}=r e^{\ii\alpha},\qquad
 m=\E\{(Y-p)e^{2\pi\ii Y}\}.
\]
Moment matching annihilates $H_p$ (its Bernoulli expectation is
$qg_p(0)+pg_p(1)=0$), while (3.3) gives $\E W_p(Y)=d$.  Therefore (SC8)
gives, whenever $p\to p_*$ and $d\to0$,
\[
 |m|\le\{19/16+o(1)\}d+(1-r^2)+O(d^2).                \tag{4.16}
\]
Indeed $2(1-\Re\zeta)=1-r^2+|1-\zeta|^2$, while the linear Fourier
stability theorem gives $|1-\zeta|=O(d)$ uniformly for $p\in U_0$.
Thus the $o(1)$ in (4.16) is uniform along every array with $p\to p_*$;
there is no unproved pointwise-to-uniform passage.  We shall repeatedly use
\[
 r^{\,n-1}(1-r^2)\le {2\over n}(1-r^n),              \tag{4.17}
\]
which follows after cancelling $1-r$ from
$1-r^n=(1-r)\sum_{j<n}r^j$.

At the matching physical bandwidth $s$, define the zero-frequency piece
at a cdf point $x$ by
\[
 L_{X,n}(x):={1\over2}-\Phi(x)
 +2\Re\int_0^{1/4}K(t)e^{-2\pi\ii x\sqrt n\,t/s}
        \varphi_X(2\pi t/s)^n\,\dd t,                 \tag{4.17a}
\]
and let $L_{Q,n}$ be the same expression for the matching two-point law.

For $p\in(0,1)$ and $d\ge0$ define, before its first use,
\[
 q=1-p,\qquad s(p):=\{p(1-p)\}^{-1/2},\qquad
 \gamma(p):={1-2p\over\sqrt{p(1-p)}},
\]
\[
 \begin{aligned}
 \rho_2(p)&:=\rho_2(\gamma(p))
       ={p^2+(1-p)^2\over\sqrt{p(1-p)}},
 &\qquad \rho(p,d)&:=\rho_2(p)+d,\\
 \Lambda_{\rm abs}(p,d)&:={s(p)^3\phi_0\over6\rho(p,d)}
       \max\{C_+(p)^{-1},C_-(p)^{-1}\}.
 \end{aligned}                                             \tag{4.18a}
\]
Here and below, $\rho_2(p)$ is only the explicitly declared
Bernoulli-parameter shorthand $\rho_2(\gamma(p))$; the function
$\rho_2(\gamma)$ itself retains the meaning fixed in Theorem~\ref{orig:thm:defect}.
Also put
\[
 \Lambda_{\rm abs,*}:=\Lambda_{\rm abs}(p_*,0),\qquad
 c_{\rm sh}(p,d):={s(p)^3\phi_0\over4\pi\rho(p,d)},\qquad
 c_{\rm sh,*}:=c_{\rm sh}(p_*,0).                     \tag{4.18c}
\]

\begin{lemma}[phase-free low-frequency bound]\label{orig:lem:abslow}
Fix a compact interval $U\Subset(0,1)$ containing $p_*$.  Let
$n_j\to\infty$, let $X_j$ be centered and variance one, and suppose that
their matching parameters satisfy
\[
 p_j\in U,\qquad d_j:=\E|X_j|^3-\rho_2(p_j)\longrightarrow0.       \tag{4.18b}
\]
Let \(Q_j:=Q_{p_j}\) denote the matching standardized two-point law.
Then the difference of the zero-frequency parts of the calibrated
one-sided operators satisfies
\[
 \sup_{x\in\mathbb R}{\sqrt {n_j}\over\rho(p_j,d_j)}
 |L_{X_j,n_j}(x)-L_{Q_j,n_j}(x)|
 \le\{\Lambda_{\rm abs}(p_j,d_j)+o_U(1)\}d_j.          \tag{4.18}
\]
Here $o_U(1)\to0$ along every array satisfying (4.18b), independently of
the cdf coordinate.  In particular, no assertion is made for arrays on
which $d_j$ fails to tend to zero.
\end{lemma}

\begin{proof}
The case $d_j=0$ is separate and immediate from
Theorem~\ref{orig:thm:defect}; below pass to the rows with $d_j>0$ and suppress
the index $j$.  Write $\varphi=\varphi_X$ and let $\chi_p$ be the
characteristic function of the standardized matching two-point law.  If
$E_z$ denotes the cubic Hermite error for $y\mapsto e^{\ii zy}$, the
 defect representation from Corollary~\ref{orig:cor:rankfree} is
\[
 \varphi(u)-\chi_p(u)=e^{-\ii s(p)pu}\,d
       \int {E_{s(p)u}(y)\over W_p(y)}\,\dd\nu(y),\qquad
 \dd\nu={W_p\,\dd\Pp_Y\over d}.                       \tag{4.20}
\]
Direct expansion of the four interpolation data, uniformly for $p\in U$
and small $|z|$, gives
\[
 |E_z(y)-g(zy)|\le C_U|z|^4(1+|y|^3),\qquad
 g(w)=e^{\ii w}-1-\ii w+w^2/2+\ii w^3/6.              \tag{4.21}
\]
Choose $H=n^{1/4}$.  Hermite--Genocchi and the explicit branches of
$W_p$ give
\[
 |E_z(y)|\le {|z|^4\over24}y^2(1-y)^2,\qquad
 {y^2(1-y)^2\over W_p(y)}\le C_U(1+|y|),              \tag{4.21a}
\]
while $|g(w)|\le |w|^3/6$ by (4.10c).  Consequently, for each fixed
$v\ge0$, with $z=s(p)v/\sqrt n$,
\begin{align}
 n^{3/2}\left|\int {E_z(y)\over W_p(y)}\dd\nu(y)\right|
 \le{}&{s(p)^3v^3\over6}
   \left[\max\{C_+(p)^{-1},C_-(p)^{-1}\}+O_U(H^{-1})\right]
 \notag\\
 &+C_U{v^4(H+1)\over\sqrt n}.                         \tag{4.21c}
\end{align}
Indeed, (4.21a) gives the last bound on $|y|\le H$; on $|y|>H$,
(4.21) handles the interpolant correction and the exact cubic leading
terms give
\[
 { |y|^3\over W_p(y)}\le
 \max\{C_+(p)^{-1},C_-(p)^{-1}\}\{1+O_U(H^{-1})\}.
\]

We now give an explicit frequency partition.  The exact divided power is
\[
 \mathscr P_n(u)={\varphi(u)^n-\chi_p(u)^n\over
             n\{\varphi(u)-\chi_p(u)\}}
 =\int_0^1\{(1-a)\chi_p(u)+a\varphi(u)\}^{n-1}\dd a,             \tag{4.21d}
\]
with its continuous value when the denominator vanishes.  The third
absolute moments are uniformly bounded.  From
\[
 1-|\E e^{\ii uZ}|^2=\E\{1-\cos(u(Z-Z'))\},\qquad
 1-\cos w\ge w^2/2-|w|^3/6,
\]
and $\E(Z-Z')^2=2$, there are $c>0$ and $\delta>0$, depending only on
$U$, such that
\[
 \max\{|\varphi(u)|,|\chi_p(u)|\}\le e^{-cu^2}\qquad(|u|\le\delta).
                                                               \tag{4.21e}
\]
The same estimate, with a smaller $c$, holds for every convex combination
in (4.21d).  Hence
\[
 \mathscr P_n(v/\sqrt n)\longrightarrow e^{-v^2/2},\qquad
 |\mathscr P_n(v/\sqrt n)|\le e^{-cv^2}\qquad(0\le v\le\delta\sqrt n), \tag{4.21f}
\]
where the convergence is uniform for $v$ in a fixed compact interval.

For clarity, the rest of the zero cell is not covered by the words
``Cram\'er gap.''  Its standardized frequency is
$u=2\pi t/s(p)$ and its right endpoint is $b(p)=\pi/(2s(p))$.
Decrease $\delta$ so that $0<\delta<\min_{p\in U}b(p)$.  Since
\[
 |\chi_p(u)|^2=1-2p(1-p)\{1-\cos(s(p)u)\},
\]
compactness gives
\[
 q_U:=\sup_{p\in U,\ \delta\le |u|\le b(p)}|\chi_p(u)|<1.       \tag{4.21g}
\]
Theorem~\ref{orig:thm:cf} gives
$|\varphi(u)-\chi_p(u)|\le C_Ud|u|^3$.  Thus, for all sufficiently
large rows, both endpoints and every convex combination in (4.21d) have
modulus at most $q'_U:=(1+q_U)/2<1$ throughout the interval in
(4.21g).  Its normalized operator contribution is therefore bounded by
\[
 C_Un^{3/2}d(q'_U)^{n-1}=o_U(d),                       \tag{4.21h}
\]
uniformly in the cdf phase.  This proves the required common gap without
any restriction on how slowly $d\downarrow0$.

It remains only to justify a single dominated-convergence passage on the
root-scale interval.  Put
\[
 B_n(v):=n^{3/2}\int {E_{s(p)v/\sqrt n}(y)\over W_p(y)}\dd\nu(y).
\]
Theorem~\ref{orig:thm:cf} applied in (4.20) supplies the global local-cell
majorant $|B_n(v)|\le C_Uv^3$.  Under the substitution
\[
 t={s(p)v\over2\pi\sqrt n},\qquad
 2K(t)\,\dd t=\left\{{\ii\over\pi v}
       +{s(p)\over2\pi\sqrt n}+O_U(v/n)\right\}\dd v,             \tag{4.21i}
\]
the absolute value of the complete normalized integrand, after the
spatial phase is discarded, is at most
\[
 C_Ud\{v^2+v^3/\sqrt n+v^4/n\}e^{-cv^2}.              \tag{4.21j}
\]
This is integrable on $[0,\infty)$ and independent of the cdf coordinate.
Equations (4.21c), (4.21f), and (4.21j), first on a fixed $v$-interval
and then with its Gaussian tail removed, give
\[
 {d\,s(p)^3\over6\pi\rho(p,d)}
 \max\{C_+(p)^{-1},C_-(p)^{-1}\}
 \int_0^\infty v^2e^{-v^2/2}\dd v+o_U(d).
\]
Because $\pi^{-1}\int_0^\infty v^2e^{-v^2/2}\dd v=\phi_0$, this is
the right side of (4.18).  The bounded and $O(v/n)$ kernel terms vanish
under the same majorant.  The conjugate one-sided kernel has identical
absolute bounds.  No linearization of the divided power, and hence no
unrecorded $d^2$ remainder, has been used.  This proves (4.18).
\end{proof}

We isolate the elementary arithmetic needed later.  Direct substitution of
$p_*=(4-\sqrt {10})/2$, followed only by squaring positive radicals, gives
\[
 {\CE\over\rho_*}>{39\over100},\qquad
 c_{\rm sh,*}<{51\over200},
 \qquad \Lambda_{\rm abs,*}<{17\over200},             \tag{4.22}
\]
\[
 G_{\rm span,*}:={\CE-\phi_0K_1(\gamma_*)/2\over\rho_*}
 >{9\over50}.                                          \tag{4.23}
\]
In particular the two margins which occur below are strictly positive:
\[
 {\CE\over\rho_*}-{19\over16}c_{\rm sh,*}
                    -\Lambda_{\rm abs,*}>{7\over3200},
 \quad
 G_{\rm span,*}-\Lambda_{\rm abs,*}>{19\over200}.     \tag{4.24}
\]
A direct outward-rational evaluation of the same radical formulas gives the
following exact integer enclosures (the common denominator is $10^{15}$):
\[
\begin{gathered}
10^{15}c_{\rm sh,*}\in[254150792637701,254150792637702],\\
10^{15}\Lambda_{\rm abs,*}\in[84254560311404,84254560311405],\\
10^{15}{\CE\over\rho_*}\in[393927535745323,393927535745324],\\
10^{15}{19c_{\rm sh,*}\over16}\in[301804066257270,301804066257271],\\
10^{15}G_{\rm span,*}\in[182181753528660,182181753528661],\\
10^{15}\left({\CE\over\rho_*}-{19c_{\rm sh,*}\over16}
 -\Lambda_{\rm abs,*}\right)
 \in[7868909176648,7868909176649],\\
10^{15}(G_{\rm span,*}-\Lambda_{\rm abs,*})
 \in[97927193217255,97927193217256].
\end{gathered}                                               \tag{4.24e}
\]
All endpoints in (4.24e) are integers, not rounded floating-point output;
cross multiplication proves the two rational inequalities in (4.24).

\subsection*{Uniform endpoint and tangent expansions}

For later bookkeeping we now define the bandwidth-dependent operators
explicitly.  For any real-valued law $X$, integer $n\ge1$, $h>0$, and
$x\in\mathbb R$, put
\[
 g_{X,n,h,x}(t):=
 e^{-2\pi\ii x\sqrt n\,t/h}\,
 \varphi_X(2\pi t/h)^n,\qquad 0\le t\le1.
\]
With
\[
 \epsilon_+=1,\quad\epsilon_-=-1,\qquad
 \widetilde K_+=K,\quad\widetilde K_-=\overline K,
\]
define
\[
 \boxed{\quad
 \mathcal U^\sigma_{X,n,h}(x):=
 {1\over2}+\epsilon_\sigma
   \mathfrak I_{\widetilde K_\sigma}(g_{X,n,h,x}),
 \qquad \sigma\in\{+,-\}.
 \quad}                                                        \tag{4.24u}
\]
Thus, written without the notation from (P0),
\[
 \mathcal U^+_{X,n,h}(x)
 ={1\over2}+2\Re\int_0^1K(t)g_{X,n,h,x}(t)\,\dd t,\qquad
 \mathcal U^-_{X,n,h}(x)
 ={1\over2}-2\Re\int_0^1\overline {K(t)}g_{X,n,h,x}(t)\,\dd t.
\]
All integrals meeting zero are understood in the real-part improper sense
of (P0).  For an interval $I\subset[0,1]$, denote its integral contribution
(without the common $1/2$) by
\[
 \mathcal J^\sigma_{X,n,h;I}(x):=
 2\epsilon_\sigma\Re\int_I
 \widetilde K_\sigma(t)g_{X,n,h,x}(t)\,\dd t.             \tag{4.24v}
\]

To record the half-line conventions, let
\[
 Z={X_1+\cdots+X_n-x\sqrt n\over h},\qquad
 M(u)={1\over2}+2\Re\int_0^1K(t)e^{2\pi\ii tu}\,\dd t,
\]
and let $e_B(u):=M(u)-1_{(-\infty,0]}(u)$ be the upper-majorant
excess over the closed negative half-line.  Then
\[
 \mathcal U^+_{X,n,h}(x)=\E M(Z),\qquad
 \mathcal U^-_{X,n,h}(x)=\E\{1-M(-Z)\},
\]
and hence the exact identities are
\[
 \begin{aligned}
 F_{n,X}(x)&=\mathcal U^+_{X,n,h}(x)-\E e_B(Z),\\
 F_{n,X}(x-)&=\mathcal U^-_{X,n,h}(x)+\E e_B(-Z).
 \end{aligned}                                               \tag{4.24w}
\]
In particular,
\[
 F_{n,X}(x)\le\mathcal U^+_{X,n,h}(x),\qquad
 F_{n,X}(x-)\ge\mathcal U^-_{X,n,h}(x),
\]
so
\[
 D_X^+(x)\le\mathcal U^+_{X,n,h}(x)-\Phi(x),\qquad
 D_X^-(x)\le\Phi(x)-\mathcal U^-_{X,n,h}(x).
\]
Fix $\delta:=1/4$, put
$x_{\eta,\theta}=s(\eta+\theta)/\sqrt n$ and
$k=np+\eta\in\mathbb Z$.  All pieces in the next display are in raw cdf
units.  Indeed, for the Bernoulli coordinate $Y=p+X/s$ and matching
bandwidth $h=s$,
\[
 g_{X,n,s,x_{\eta,\theta}}(t)
 =e^{-2\pi\ii(k+\theta)t}\varphi_Y(2\pi t)^n,             \tag{4.24x}
\]
which is the common phase used below:
\[
\begin{aligned}
D^{(0)}_{X,n}(\eta,\theta)&={1\over2}-\Phi(x_{\eta,\theta})
 +2\Re\int_0^\delta K(t)e^{-2\pi\ii(k+\theta)t}
                  \varphi_Y(2\pi t)^n\,\dd t,\\
D^{(1)}_{X,n}(\eta,\theta)&=2\Re\int_{1-\delta}^1K(t)
 e^{-2\pi\ii(k+\theta)t}\varphi_Y(2\pi t)^n\,\dd t,\\
M_{X,n}(\eta,\theta)&=2\Re\int_\delta^{1-\delta}K(t)
 e^{-2\pi\ii(k+\theta)t}\varphi_Y(2\pi t)^n\,\dd t.
\end{aligned}                                                   \tag{4.24a0}
\]
Set \(\Psi_Q(z):=q+pe^{\ii z}\).  The pieces $D^{(j)}_{Q,n}$ and
$M_{Q,n}$ are obtained by replacing $\varphi_Y$ by $\Psi_Q$ with the same
cuts, phase and normal increment.  Let
$\mathcal P_n^+(X;\eta,\theta)$ be the
normalized Prawitz upper operator.  The actual positive cdf score
$\mathcal C_n^+(X;\eta,\theta)$ satisfies
$\mathcal C_n^+(X;\eta,\theta)\le\mathcal P_n^+(X;\eta,\theta)$, and the
exact operator decomposition is
\begin{equation}\tag{4.24a}
\begin{split}
 \mathcal P_n^+(X;\eta,\theta)={\sqrt n\over\rho}\big[&
 D^{(0)}_{Q,n}(\eta,\theta)+D^{(1)}_{Q,n}(\eta,\theta)
 +\{D^{(0)}_{X,n}(\eta,\theta)-D^{(0)}_{Q,n}(\eta,\theta)\}\\
 &+\{D^{(1)}_{X,n}(\eta,\theta)-D^{(1)}_{Q,n}(\eta,\theta)\}
 +M_{X,n}(\eta,\theta)\big].
\end{split}
\end{equation}
For the matching law write, without changing any cdf convention,
\[
 N_{Q,n}(\eta,\theta):={1\over2}-\Phi(x_{\eta,\theta}),\qquad
 Z_{Q,n}(\eta,\theta):=D^{(0)}_{Q,n}(\eta,\theta)
                         -N_{Q,n}(\eta,\theta),        \tag{4.24b}
\]
and define the complete integer-node value
\[
 E_{Q,n}(\eta):=N_{Q,n}(\eta,0)+Z_{Q,n}(\eta,0)
       +D^{(1)}_{Q,n}(\eta,0)+M_{Q,n}(\eta,0).         \tag{4.24c}
\]
By the cardinal identity (P3),
$\sqrt nE_{Q,n}(\eta)/\rho_2(p)$ is the exact right-node Bernoulli score.
These definitions will make explicit which phase increment owns
each later term.
The matching endpoint is (4.29), the zero-cell phase increment is (4.31),
the joint saddle is (4.40), and the actual endpoint is first factorized by
(4.26); its complete remainder is (4.43).  This identity is the owner of
the later assembly, so no zero or endpoint row is silently duplicated.
For reference, the ownership ledger is
\[
\begin{array}{c|c|c}
\text{quantity}&\text{owner}&\text{used in}\\ \hline
N_{Q,n}(\eta,0)+Z_{Q,n}(\eta,0)&E_{Q,n}(\eta)&(4.40)\\
D^{(1)}_{Q,n}(\eta,0)&E_{Q,n}(\eta)&(4.40)\\
M_{Q,n}(\eta,0)&E_{Q,n}(\eta)&(4.40)\\
D^{(0)}_{Q,n}(\eta,\theta)-D^{(0)}_{Q,n}(\eta,0)&(4.31)&(4.34)\\
\Re\{\zeta^nJ_{Q,n}(\eta,\theta)-J_{Q,n}(\eta,0)\}&(4.43b)&(4.43e),(4.43g)\\
M_{Q,n}(\eta,\theta)-M_{Q,n}(\eta,0)&\text{middle gap}&\text{exponentially small}.
\end{array}                                                   \tag{4.24d}
\]

For the rest of the proof use Bernoulli coordinates
\[
 \begin{gathered}
 Y=p+{X\over s},\qquad q=1-p,\qquad s=(pq)^{-1/2},\qquad
 \rho_2(p)={p^2+q^2\over\sqrt{pq}},\\
 \Psi(z)=\E e^{\ii zY},\qquad \Psi_Q(z)=q+pe^{\ii z}
 \quad\text{(as above)}.
 \end{gathered}
\]
For an integer $k$ and a circular phase $|\theta|\le1/2$, put
\[
 \eta=k-np,\qquad x={s(\eta+\theta)\over\sqrt n}.       \tag{4.25}
\]
The terminology in this paragraph repeats (4.0a).  For a fixed row
$(n,X)$, the positive cell indexed by $k$ is the set of physical cdf
points in (4.25) as $\theta$ ranges over $[-1/2,1/2]$.  It is
\emph{active} if at least one of those points has
$\mathcal C_n^+(X;\eta,\theta)>\CE$; when a particular point is under
discussion, ``active'' means that this strict inequality holds at that
point.  Every sequence of active points on the
\emph{localized positive branch} after the reduction in
Proposition~\ref{orig:prop:qouter} satisfies, explicitly,
\[
 p\to p_*,\qquad d\to0,\qquad x\to0,\qquad
 \eta/\sqrt n\to0,                                    \tag{4.25a}
\]
and the positive, rather than reflected negative, one-sided score is being
used.  The last condition in (4.25a) follows from the preceding three and
$|\theta|\le1/2$.  Thus later uses of either phrase carry no additional
unstated localization assumption.

\begin{lemma}[endpoint factorization]\label{orig:lem:endfactor}
If $p$ remains in a fixed neighbourhood of $p_*$, then, uniformly for
$0\le h\le1/4$,
\[
 \Psi(2\pi(1-h))
 =\zeta\Psi_Q(2\pi(1-h))-2\pi\ii hm+O(h^2d),           \tag{4.26}
\]
where $\zeta=\Psi(2\pi)$ and
$m=\E\{(Y-p)e^{2\pi\ii Y}\}$.
\end{lemma}

\begin{proof}
Let $A(h):=\Psi(2\pi(1-h))$ and
$B(h):=\zeta\Psi_Q(2\pi(1-h))$.  Then
$A(0)=B(0)=\zeta$ and
$A'(0)-B'(0)=-2\pi\ii m$.  Twice differentiating the remaining error
reduces, apart from a fixed factor, to the expectation of
\[
 e^{2\pi\ii y}\{y^2e^{-2\pi\ii hy}-pe^{-2\pi\ii h}\}. \tag{4.27}
\]
Subtract its cubic Hermite interpolant at $0,1$.  Its expectation is zero
by the three moment identities.  The quotient of the Hermite remainder by
$W_p(y)$ extends through the two double zeros and stays bounded at infinity,
because numerator and denominator have cubic growth.  Compactness in
$(p,h)$ therefore bounds (4.27) by $Cd$.  Integrating twice proves (4.26).
\end{proof}

\begin{lemma}[uniform first endpoint shape mode]\label{orig:lem:endshape}
Fix $U\Subset(0,1)$.  In the notation of
Lemma~\ref{orig:lem:endfactor}, let $p\in U$, let $d\to0$, and let $n\to\infty$.
The contribution of the term
\[
 -2\pi\ii nhm\,\zeta^{\,n-1}
       \Psi_Q(2\pi(1-h))^{n-1}                         \tag{4.27a}
\]
to the normalized returned-endpoint operator is, uniformly in the integer
node and in the full circular phase $|\theta|\le1/2$, at most
\[
 \{c_{\rm sh}(p,d)+o_U(1)\}r^{\,n-1}|m|.             \tag{4.27b}
\]
The $o_U(1)$ is independent of the rates of $p,d$ and of the node and
phase.  After one derivative in the cdf phase the same term is bounded by
$C_Ur^{\,n-1}|m|$.  Neither assertion uses a fourth moment or a previously
localized $\eta$-window.
\end{lemma}

\begin{proof}
The spatial and circular phase factors have modulus one.  The exact
Bernoulli identity
\[
 |\Psi_Q(2\pi(1-h))|^2
 =1-2p(1-p)\{1-\cos(2\pi h)\}                         \tag{4.27c}
\]
gives, uniformly for $p\in U$, a Gaussian bound $e^{-c_U(n-1)h^2}$
on $0\le h\le\delta$ and a fixed Cram\'er gap on
$\delta\le h\le1/4$.  Thus the latter part of the endpoint cell is
$O_U(n^{3/2}e^{-c_Un})r^{n-1}|m|$ after normalization.

On the root-scale part first restrict to $0\le v=\sqrt n h\le n^{1/12}$.
Directly from the displayed
kernel definition,
\[
 2K(1-h)=h+O(h^2),                                    \tag{4.27d}
\]
uniformly on this subcell.  The absolute value of the integral generated by
(4.27a), including the factor $\sqrt n/\rho(p,d)$, is consequently
\[
 {2\pi n\sqrt n\over\rho(p,d)}r^{n-1}|m|
 \int_0^{n^{-5/12}}h^2
 e^{-2\pi^2p(1-p)(n-1)h^2}\{1+O_U(h+n h^4)\}\dd h.   \tag{4.27e}
\]
The $nh^4$ term follows by expanding the logarithm in (4.27c); it is not
asserted as a relative expansion on the rest of the endpoint cell.  After
the substitution it is dominated on this subcell by a fixed Gaussian
polynomial.  On $n^{1/12}<v\le\delta\sqrt n$, the exact bound (4.27c)
makes the integrated contribution $O_U(e^{-c_Un^{1/6}})$, and the fixed
remainder was treated above.  Dominated convergence is therefore uniform
for $p\in U$.  Since
\[
 \int_0^\infty v^2e^{-2\pi^2v^2/s(p)^2}\dd v
       ={s(p)^3\phi_0\over8\pi^2},
\]
the limiting coefficient in (4.27e) is
\[
 {2\pi\over\rho(p,d)}{s(p)^3\phi_0\over8\pi^2}
       ={s(p)^3\phi_0\over4\pi\rho(p,d)}
       =c_{\rm sh}(p,d).
\]
One phase derivative only inserts the bounded factor
$-2\pi\ii(1-h)$ before this argument, proving the last assertion.
\end{proof}

\begin{lemma}[uniform signed tangent expansion]\label{orig:lem:tangentexp}
Let $p\to p_*$ and $d\to0$.  Define
\[
 A(p)={\phi_0\over\rho_2(p)}
       \left({s\over2}+{\gamma\over6}\right),\qquad
 a(p)={s^3q\phi_0\over2},\qquad
 \eta_0(p)=-{p^2-7p+5\over12q}.                       \tag{4.28}
\]
Write $a_{\rm tan,*}:=a(p_*)$.
Then $A(p)\le\CE$, with equality only at $p_*$, and
$\eta_0(p_*)=\eta_*$.  The following assertions hold uniformly.

\begin{enumerate}[label=(\alph*)]
\item If $|\eta|=o(\sqrt n)$, then the matching two-point high endpoint is
\[
 2\int_{3/4}^1K(t)e^{-2\pi\ii(k+\theta)t}
       (q+pe^{2\pi\ii t})^n\,\dd t
 =e^{-2\pi\ii\theta}\left\{{C_0(p)\over n}
 +{\ii C_1(p,\eta+\theta)\over n^{3/2}}\right\}
 +O\left({1+\eta^2\over n^2}\right),                  \tag{4.29}
\]
with one uniform $\theta$ derivative, where
\[
 C_0(p)={s^2\over4\pi^2},\qquad
 C_1(p,z)={s^3\phi_0\over24\pi}
              \{6z+3(1-2p)+1\}.                     \tag{4.30}
\]

\item The zero-end increment, after subtracting the normal increment, is
\[
 D^{(0)}_{Q,n}(\eta,\theta)-D^{(0)}_{Q,n}(\eta,0)
 ={ -a(p)\theta^2+\{\ell(p)-2a(p)\eta\}\theta
       \over n^{3/2}}+R^{(0)}_n,                       \tag{4.31}
\]
where
\[
 \ell(p)=-{s^3\phi_0\over48}\{(1-2p)^2+12(1-2p)+7\}
         =2a(p)\eta_0(p),                              \tag{4.32}
\]
and, whenever $|\eta|\le C\sqrt{nd+1}$,
\[
 \sqrt n\,|R^{(0)}_n|
 \le C\left\{{1+\eta^4\over n^{3/2}}
              +{1+|\eta|^5\over n^2}\right\}.        \tag{4.33}
\]

\item If $d=o(n^{-1/2})$, $nd\to\infty$, and
$|\eta|\le C\sqrt{nd+1}$, the integer-node matching score together with
the zero-cell phase increment, namely
\[
 {\sqrt n\over\rho_2(p)}\{E_{Q,n}(\eta)
       +D^{(0)}_{Q,n}(\eta,\theta)-D^{(0)}_{Q,n}(\eta,0)\},
\]
is at most
\[
 A(p)-{a(p)\over\rho_2(p)n}
       \{\eta+\theta-\eta_0(p)\}^2+o(d).              \tag{4.34}
\]
The returned-endpoint phase increment (4.29) is deliberately not included
in (4.34); it is inserted exactly once in Lemma~\ref{orig:lem:pertendpoint}.
\end{enumerate}
\end{lemma}

\begin{proof}
For (a), put $t=1-v/\sqrt n$, $z=\eta+\theta$ and
$\xi=z/\sqrt n$.  After extracting the Bernoulli quadratic term the
integrand equals
\[
 e^{-2\pi\ii\theta}e^{-2\pi^2v^2/s^2}e^{2\pi\ii\xi v}
 \left[1+{\ii B(p)v^3\over\sqrt n}+{P_{p,n}(v)\over n}\right],
 \qquad B(p)={4\pi^3\over3}pq(1-2p),                  \tag{4.35}
\]
where $P_{p,n}$ is bounded by a fixed polynomial under the Gaussian.
Keep $e^{2\pi\ii\xi v}$ exact.  Its Gaussian integrals are entire in
$\xi$; Taylor's formula at zero gives an error $O(z^2/n^2)$ in the
leading $v/n$ integral and $O((1+|z|)/n^2)$ in the next integrals.
Using
\[
 \int_0^\infty ve^{-2\pi^2v^2/s^2}\,\dd v={s^2\over4\pi^2},
\]
\[
 \int_0^\infty v^2e^{-2\pi^2v^2/s^2}\,\dd v
 ={s^3\phi_0\over8\pi^2},\qquad
 \int_0^\infty v^4e^{-2\pi^2v^2/s^2}\,\dd v
 ={3s^5\phi_0\over32\pi^4}
\]
gives (4.29)--(4.30).  Splitting at $v=n^{1/12}$ and using the exact
Bernoulli Gaussian envelope on the complement proves the stated uniform
remainder and its phase derivative.

At the zero endpoint put again $z=\eta+\theta$ and
$C_z(v)=2\pi zv+B(p)v^3$.  The centered Bernoulli cumulants through order
five give
\begin{equation}\tag{4.36}
\begin{split}
 e^{-2\pi\ii(k+\theta)t}(q+pe^{2\pi\ii t})^n
 =e^{-2\pi^2v^2/s^2}\big[&1-\ii n^{-1/2}C_z(v)\\
 &+n^{-1}\{E_4v^4-C_z(v)^2/2\}\\
 &+\ii n^{-3/2}\{J_5v^5-C_z(v)E_4v^4+C_z(v)^3/6\}\\
 &+O(n^{-2}P_{12}(v,z))\big],
\end{split}
\end{equation}
where
\[
 E_4={2\pi^4\over3}pq(1-6pq),\qquad
 J_5={4\pi^5\over15}pq(1-2p)(1-12pq).
\]
Also
\[
 2K(v/\sqrt n)n^{-1/2}
 ={\ii\over\pi v}+n^{-1/2}
 -n^{-1}\left(v+{\ii\pi v\over3}\right)
 +{\ii\pi v^2\over3n^{3/2}}+O(n^{-2}P_3(v)).          \tag{4.37}
\]
Here is the uniform remainder statement implicit in these two displays.
There are a neighbourhood $U_1$ of $p_*$, constants $c,C>0$, and fixed
positive polynomials $P_j$ such that, for $p\in U_1$,
$0\le v\le n^{1/20}$, and $|z|\le\varepsilon_0\sqrt n$ for a fixed
small $\varepsilon_0>0$, the remainder inside the
brackets in (4.36), together with one $z$ derivative, is bounded by
\[
 {C(1+|z|^4)v^4P_{8}(v)\over n^2}.                    \tag{4.37a}
\]
After multiplication by (4.37), its integrable majorant is the same
bound times $(1+v^{-1})e^{-cv^2}$.  To verify (4.37a), expand
$n\log\{e^{-2\pi\ii pt}(q+pe^{2\pi\ii t})\}$ through the fifth
derivative; on $|t|\le n^{-9/20}$ its sixth derivative is uniformly
bounded for $p\in U_1$.  Taylor's integral remainder for the exponential
is then bounded by
$Cn^{-2}v^4P_8(v)(1+|z|^4)e^{-cv^2}$.  Here the spatial factor is
expanded separately, using the global integral-remainder inequality for
$e^{-\ii w}$; hence no assumption $z=O(n^{1/4})$ is hidden.
Differentiating in $z$ merely removes one factor of $z$ and inserts
$v/\sqrt n$.  On
$n^{1/20}<v\le\sqrt n/4$, the exact identity
\[
 |q+pe^{2\pi\ii v/\sqrt n}|^2
 =1-2pq\{1-\cos(2\pi v/\sqrt n)\}
\]
gives $Ce^{-cv^2}$ up to a fixed small cell and $Ce^{-cn}$ on its
complement.  The same estimates hold after one $z$ derivative.  Thus the
root-scale expansion may be integrated and differentiated term by term;
this is a quantitative $C^1$ statement, not a formal Edgeworth
calculation.

For transparency, the parity ledger for the real part of the product of
(4.36) and (4.37) is
\[
\begin{array}{c|c|c}
\text{order}&\text{real coefficient before Gaussian integration}
   &\text{$z$-dependence after the normal term}\\ \hline
1&0&0\\
n^{-1/2}&1+C_z(v)/(\pi v)&\text{linear term cancels}\\
n^{-1}&-v&\text{none}\\
n^{-3/2}&R_p(z,v)&\text{degrees }0,1,2,3,
\end{array}                                                   \tag{4.37b}
\]
because the omitted coefficients in the first three rows are purely
imaginary.  In the second row,
$2z\int_0^\infty e^{-2\pi^2v^2/s^2}\dd v=s\phi_0z$, which cancels the
$-s\phi_0z$ from $1/2-\Phi(sz/\sqrt n)$.  The order-$n^{-1}$ row is
independent of $z$ and therefore cancels when two phases in the same cell
are subtracted.  The last-row polynomial is
\[
 R_p(z,v)=E_4v^4-{C_z(v)^2\over2}
 -{J_5v^5-C_z(v)E_4v^4+C_z(v)^3/6\over\pi v}
 -{\pi vC_z(v)\over3}.                                \tag{4.38}
\]
Using
\[
 I_2={s^3\phi_0\over8\pi^2},\quad
 I_4={3s^5\phi_0\over32\pi^4},\quad
 I_6={15s^7\phi_0\over128\pi^6},\quad
 I_8={105s^9\phi_0\over512\pi^8},                    \tag{4.38a}
\]
where $I_j=\int_0^\infty v^je^{-2\pi^2v^2/s^2}\dd v$, direct collection
of the four powers of $z$ gives
\[
\begin{array}{c|c}
z^0&\displaystyle
 \beta_0(p)=\phi_0s^3\left(-{1\over32}-{7\tau\over160}
                  -{\tau^2\over96}+{\tau^3\over4320}\right)\\[1mm]
z^1&\displaystyle
 -{s^3\phi_0\over48}(\tau^2+12\tau+7)=\ell(p)\\[1mm]
z^2&\displaystyle-{s^3q\phi_0\over2}=-a(p)\\[1mm]
z^3&\displaystyle-{s^3\phi_0\over6},
\end{array}\qquad \tau=1-2p.                           \tag{4.38b}
\]
For example, the linear row before simplification is
\[
 s^3\phi_0\left[-{1\over12}+{1-6pq-2\tau\over8}
                         -{5\tau^2\over24}\right],
\]
which is the displayed value of $\ell(p)$; this also independently checks
the sign.  Finally
\[
 \Phi(sz/\sqrt n)=\tfrac12+\phi_0
 \left({sz\over\sqrt n}-{s^3z^3\over6n^{3/2}}\right)
 +O(|z|^5n^{-5/2})                                    \tag{4.39}
\]
shows that the $z^3$ row of (4.38b) is cancelled exactly by the normal
term.  The value $\beta_0(p)$ and every earlier $z$-independent row cancel
on subtracting two phases.  To state this without overloading the
two-argument cell notation, define
\[
 \mathscr D^{(0)}_{Q,n}[z]:={1\over2}-\Phi(sz/\sqrt n)
 +2\Re\int_0^{1/4}K(t)e^{-2\pi\ii zt}
       \{e^{-2\pi\ii pt}\Psi_Q(2\pi t)\}^{n}\,\dd t .
                                                               \tag{4.38b1}
\]
Thus
$D^{(0)}_{Q,n}(\eta,\theta)=\mathscr D^{(0)}_{Q,n}[\eta+\theta]$.
More precisely, the estimates just proved and
Taylor's remainder in (4.39) give a function $\mathcal E_{n,p}$ such
that
\begin{align}
 \mathscr D^{(0)}_{Q,n}[z]-\mathscr D^{(0)}_{Q,n}[0]
 &=n^{-3/2}\{\ell(p)z-a(p)z^2\}+\mathcal E_{n,p}(z),\notag\\
 |\partial_z\mathcal E_{n,p}(z)|
 &\le C\left\{{1+|z|^4\over n^2}
                   +{1+|z|^5\over n^{5/2}}\right\}              \tag{4.38c}
\end{align}
uniformly for $p\in U_1$ and $|z|\le\varepsilon_0\sqrt n$.

For completeness, the replacement of the finite root cell by the positive
half-line in the leading row is uniform at the sharper spatial scale needed
below.  Write $G_{0,p}(v)$ for the unphased leading coefficient in the
score-normalized root integral corresponding to the first two rows of
(4.37b), including the bounded Gaussian smoothing term and the uniformly
bounded factor $1/\rho$.  For two phases in one cell put
$x=s(p)\eta/\sqrt n$ and $y=s(p)(\eta+\theta)/\sqrt n$.
On $p\in U_1$ there are a fixed positive polynomial $P$ and $c,C>0$ such
that $|G_{0,p}(v)|\le CP(v)e^{-cv^2}$.  Hence, for $A>0$,
\[
 \left|\int_A^\infty
   \Re\{(e^{-\ii yv}-e^{-\ii xv})G_{0,p}(v)\}\,\dd v\right|
 \le {C|y-x|\over A^2}
       \int_0^\infty v^3P(v)e^{-cv^2}\,\dd v .             \tag{4.38c1}
\]
Indeed $|e^{-\ii yv}-e^{-\ii xv}|\le |y-x|v$, and on $v>A$ one
inserts $1\le v^2/A^2$.  The real cosine smoothing summand satisfies the
same estimate directly from
$|\cos(yv)-\cos(xv)|\le |y-x|v$.  Taking
$A=\delta\sqrt n$ and using
$|y-x|=s(p)|\theta|/\sqrt n\le Cn^{-1/2}$ makes the omitted spatial tail
$O(n^{-3/2})$.
Thus passing to the half-line in the coefficient calculation neither loses
an $n^{-1}$ term nor hides a merely $O(n^{-1})$ remainder.

Apply the mean-value
theorem between $z=\eta$ and $z=\eta+\theta$.  Since
$|\theta|\le1/2$ and the window in (b) has
$|\eta|/\sqrt n=O(\sqrt d+n^{-1/2})=o(1)$, (4.38c) proves (4.31) and
exactly the stated bound (4.33).

It remains to retain the node loss rather than maximize it away.  The
exact beta-integral representation used in Lemma~\ref{orig:lem:binsaddle}, with
$\epsilon=n^{-1/2}$ and $\xi=\epsilon\eta$ kept as a parameter, gives the
following explicit joint saddle.  On a fixed rectangle about $(p_*,0)$ put
\[
 \begin{aligned}
 P_1(p,\xi;y)&={s(p-q)\over3}y^3
       +{s^2\xi(p-q)\over2}y^2-spy,\\
 P_2(p,\xi;y)&=-{s^2(1-3pq)\over4}y^4+{s^2p^2\over2}y^2
       -{s^3\xi(1-3pq)\over3}y^3,
 \end{aligned}                                                   \tag{4.39a}
\]
and $W_\xi(y)=\exp(-y^2/2-s\xi y)$.  For $j=0,1,2$ define
\[
 \mathcal Q_0=1,\qquad \mathcal Q_1=P_1,\qquad
 \mathcal Q_2=P_2+\tfrac12P_1^2,
\]
\[
 \mathcal D_j(p,\xi)=\int_{\mathbb R}W_\xi(y)\mathcal Q_j(p,\xi;y)\dd y,
 \qquad
 \mathcal N_j(p,\xi)=\int_{-\infty}^0W_\xi(y)\mathcal Q_j(p,\xi;y)\dd y.  \tag{4.39b}
\]
Taylor expansion of the two logarithms in the beta density gives
\[
 L(y)=-y^2/2-s\xi y+\epsilon P_1(p,\xi;y)
       +\epsilon^2P_2(p,\xi;y)+\epsilon^3R_\epsilon(p,\xi;y).     \tag{4.39c}
\]
Consequently, if $R_1,R_2$ denote the two quotient coefficients, then
\begin{align}
 R_1={}&{\mathcal N_1\over\mathcal D_0}
       -{\mathcal N_0\mathcal D_1\over\mathcal D_0^2},\notag\\
 R_2={}&{\mathcal N_2\over\mathcal D_0}
       -{\mathcal N_1\mathcal D_1\over\mathcal D_0^2}
       +\mathcal N_0\left\{{\mathcal D_1^2\over\mathcal D_0^3}
                    -{\mathcal D_2\over\mathcal D_0^2}\right\}. \tag{4.39d}
\end{align}
The leading quotient is
$\mathcal N_0/\mathcal D_0=\Phi(s\xi)$.  Thus the required analytic
functions are, explicitly,
\[
                         H_0={R_1\over\rho_2(p)},\qquad
                         H_1={R_2\over\rho_2(p)}.                  \tag{4.39e}
\]
They satisfy, uniformly on that rectangle,
\[
 {\sqrt n\over\rho_2(p)}E_{Q,n}(\eta)
 =H_0(p,\xi)+\epsilon H_1(p,\xi)+O(\epsilon^2),        \tag{4.40}
\]
with the same estimate after one $\xi$ derivative.  Here is a quantitative
version of the remainder assertion.  On
$|y|\le\epsilon^{-1/10}$, both $R_\epsilon$ and
$\partial_\xi R_\epsilon$ are bounded by $C(1+|y|^M)$ for one fixed
$M$, and Taylor's integral formula gives
\begin{align}
 W_\xi(y)\big|e^{\epsilon P_1+\epsilon^2P_2+\epsilon^3R_\epsilon}
       -\{1+\epsilon\mathcal Q_1+\epsilon^2\mathcal Q_2\}\big|
 &\le C\epsilon^3(1+|y|^{3M})e^{-y^2/3},\notag\\
 W_\xi(y)\big|\partial_\xi
  [e^{\epsilon P_1+\epsilon^2P_2+\epsilon^3R_\epsilon}
       -\{1+\epsilon\mathcal Q_1+\epsilon^2\mathcal Q_2\}]\big|
 &\le C\epsilon^3(1+|y|^{3M+1})e^{-y^2/3}.            \tag{4.39f}
\end{align}
Strict concavity makes the two complementary tails, with the same one
derivative, $O(e^{-c\epsilon^{-1/5}})$.  Moreover,
\[
 \partial_\xi L(y)=\epsilon^{-1}
 \log{1-\epsilon sqy\over1+\epsilon spy};
\]
on the central interval this is polynomially bounded, and at either
endpoint its logarithmic factor is absorbed by the beta power there.
Thus the same tail estimate holds after one $\xi$ derivative.  In
particular, if $\mathcal N_\epsilon,\mathcal D_\epsilon$ denote the exact
truncated and full beta integrals, (4.39f) implies
\[
 \left\|{\mathcal N_\epsilon\over\mathcal D_\epsilon}
 -{\mathcal N_0+\epsilon\mathcal N_1+\epsilon^2\mathcal N_2
       \over
    \mathcal D_0+\epsilon\mathcal D_1+\epsilon^2\mathcal D_2}
 \right\|_{C^1(\xi)}\le C\epsilon^3.                  \tag{4.39g}
\]
The denominator is bounded away from zero on the fixed rectangle.
Expanding the quotient in (4.39g), subtracting its leading normal term,
and multiplying by $\epsilon^{-1}=\sqrt n$ proves the asserted genuine
$O_{C^1}(\epsilon^2)$ remainder in (4.40).

For an independently checkable evaluation of the coefficients, put
$z=s\xi$ and
\[
 I_m(z)=\int_{-\infty}^0y^m e^{-y^2/2-zy}\dd y.
\]
Integration by parts gives
\[
 I_0(z)=\sqrt{2\pi}e^{z^2/2}\Phi(z),\qquad
 I_1(z)=-1-zI_0(z),\qquad
 I_m(z)=(m-1)I_{m-2}(z)-zI_{m-1}(z).                  \tag{4.40a}
\]
The corresponding full-line moments are obtained by differentiating
$\sqrt{2\pi}e^{z^2/2}$.  Substitution of (4.39a) in (4.39d), followed by
(4.40a), gives
\[
 H_0(p,\xi)={\phi(s\xi)\over\rho_2(p)}
 \left\{{s\over2}+{\gamma\over6}(1-s^2\xi^2)\right\}.            \tag{4.40b}
\]
At $\xi=0$, $Q_2$ is even, so
$\mathcal N_2=\mathcal D_2/2$, while
$\mathcal D_1=0$ and $\mathcal N_0=\mathcal D_0/2$; hence $H_1(p,0)=0$.
Here and below a prime means $\partial_\xi|_{\xi=0}$.  The moment
recursion (4.40a) gives the scalar reductions
\[
 \mathcal N_1={s(2-p)\over3},\qquad
 \mathcal D_1'={s^2\over2}\mathcal D_0,\qquad
 \mathcal D_0'=0,\quad \mathcal N_0'=s,
\]
\[
 \mathcal D_2={s^2(1-pq)\over12}\mathcal D_0,\qquad
 \mathcal D_2'=0,\quad \mathcal N_2'={s^3p\over3}.             \tag{4.40c}
\]
Differentiating the second line of (4.39d) at zero therefore gives
\[
 \begin{aligned}
 R_2'(p,0)
 &=\phi_0s^3\left\{{p\over3}-{2-p\over6}-{1-pq\over12}\right\}\\
 &=-{\phi_0(p^2-7p+5)\over12(pq)^{3/2}}\\
 &=-{s^3\phi_0\over48}
             \{(1-2p)^2+12(1-2p)+7\}=\ell(p).
 \end{aligned}                                           \tag{4.40d}
\]
Thus
\[
 H_1(p,0)=0,\qquad
 \partial_\xi H_1(p,0)={\ell(p)\over\rho_2(p)}.       \tag{4.41}
\]
As a direct cross-check, (S6e) says exactly
\[
 F(p,\eta)-F(p,0)=-{a(p)\over\rho_2(p)}\eta^2
                         +{\ell(p)\over\rho_2(p)}\eta;
\]
this is also obtained by putting $\xi=\epsilon\eta$ in the
$C^1$ expansion (4.40).

Since $\phi(s\xi)=\phi_0(1-s^2\xi^2/2+O(\xi^4))$ and
$s+\gamma=2sq$, the quadratic coefficient in (4.40b) is
$-\phi_0s^2(s+\gamma)/(4\rho_2)=-a/\rho_2$.  Together with (4.41),
this gives
\[
 H_0=A(p)-{a(p)\over\rho_2(p)}\xi^2+O(\xi^4),\qquad
 H_1={\ell(p)\over\rho_2(p)}\xi+O(\xi^2).             \tag{4.42}
\]
On the scale in (c),
\[
 \xi^2=O(d+n^{-1})=O(d),\qquad n^{-1}=o(d).
\]
Thus the three node remainders in (4.42) and (4.40) are, respectively,
$O(\eta^4/n^2)=O(d^2)$,
$O(\eta^2/n^{3/2})=O(d/\sqrt n)$, and $O(n^{-1})$, all $o(d)$.
The zero-cell remainder is also $o(d)$: (4.33) and
$\eta^2=O(nd)$ give
\[
 \sqrt n|R_n^{(0)}|\le
 O(n^{-3/2})+O(\sqrt n\,d^2)
        +O(n^{-2})+O(\sqrt n\,d^{5/2})=o(d),          \tag{4.42a0}
\]
because $nd\to\infty$, $\sqrt n\,d\to0$, and $d\to0$.
Equations (4.40)--(4.41) are the expansion of the complete node
$E_{Q,n}(\eta)$, so its high-end node value is already present.  Adding
only the zero-cell increment (4.31), and using $\ell=2a\eta_0$, completes
the exact square in (4.34).  The high-end phase increment from (4.29) has
not been added here and is therefore not double counted.  Finally $A(p)$
is exactly the function in (4.13) under the change $r=1-2p$, proving its
last assertion.
\end{proof}

\begin{lemma}[coarse phase localization]\label{orig:lem:coarsephase}
Assume $p\to p_*$, $d\to0$, $\sqrt n\,d\to0$, and let an active positive
cell have $\xi=\eta/\sqrt n=o(1)$ and $|\theta|\le1/2$.  Before using the
windowed remainder (4.33), there exist fixed
$c_{\rm loc},C_{\rm loc}>0$ such that, uniformly,
\[
 \mathcal C_n^+(X;\eta,\theta)
 \le A(p)-c_{\rm loc}{\eta^2\over n}+C_{\rm loc}d+C_{\rm loc}n^{-1}
       +o(d+\eta^2/n).                                \tag{4.42a}
\]
Consequently every cell whose score exceeds $\CE$ satisfies
$\eta^2=O(nd+1)$.
\end{lemma}

\begin{proof}
We give the complete assembly because this estimate is used to obtain the
$\eta$-window and therefore cannot itself use that window.  All constants
below are uniform for $p$ in a sufficiently small fixed neighbourhood of
$p_*$, and we continue to write $\rho=\rho_2(p)+d$.  Let
\[
 \begin{aligned}
 \mathcal E_{0,n}&={\sqrt n\over\rho}
       \{D^{(0)}_{X,n}(\eta,\theta)-D^{(0)}_{Q,n}(\eta,\theta)\},\\
 \mathcal E_{M,n}&={\sqrt n\over\rho}
       \{M_{X,n}(\eta,\theta)-M_{Q,n}(\eta,\theta)\}.
 \end{aligned}                                             \tag{4.42b}
\]
Put $\zeta=re^{\ii\alpha}$ and choose
$\vartheta\in[-1/2,1/2)$ so that
$n\alpha=2\pi(j+\vartheta)$; this is possible for all sufficiently large
rows because $|1-\zeta|=O(d)$.  When $d=0$ take
$r=1$, $\alpha=\vartheta=0$.  With $u=\theta-\vartheta$, the exact
decomposition (4.24a), followed only by the endpoint factorization
(4.26), has the following ledger:
\[
 J_{Q,n}(\eta,\theta):=2\int_{3/4}^1K(t)
 e^{-2\pi\ii(k+\theta)t}\Psi_Q(2\pi t)^n\,\dd t,
 \qquad D^{(1)}_{Q,n}=\Re J_{Q,n}.
\]
\begin{equation}\tag{4.42c}
\begin{split}
 \mathcal P_n^+(X;\eta,\theta)={}&{\sqrt n\over\rho}E_{Q,n}(\eta)\\
 &+{\sqrt n\over\rho}
   \{D^{(0)}_{Q,n}(\eta,\theta)-D^{(0)}_{Q,n}(\eta,0)\}\\
 &+{\sqrt n\over\rho}\Re\{\zeta^nJ_{Q,n}(\eta,\theta)
                                      -J_{Q,n}(\eta,0)\}\\
 &+{\sqrt n\over\rho}
   \{M_{Q,n}(\eta,\theta)-M_{Q,n}(\eta,0)\}\\
 &+\mathcal E_{0,n}+\mathcal E_{M,n}
   +\mathcal S_n+\mathcal R_n .
\end{split}
\end{equation}
Here $\mathcal S_n$ is the normalized returned-endpoint integral generated
by
\[
 -2\pi\ii nhm\,\zeta^{n-1}
       \Psi_Q(2\pi(1-h))^{n-1},                       \tag{4.42d}
\]
and, writing $R_n^{\rm end}(h)$ for the power remainder defined below,
\[
 \begin{aligned}
 \mathcal S_n&={\sqrt n\over\rho}\,2\Re\int_0^{1/4}K(1-h)
 e^{-2\pi\ii(k+\theta)(1-h)}
 \{-2\pi\ii nhm\,\zeta^{n-1}
       \Psi_Q(2\pi(1-h))^{n-1}\}\,\dd h,\\
 \mathcal R_n&={\sqrt n\over\rho}\,2\Re\int_0^{1/4}K(1-h)
 e^{-2\pi\ii(k+\theta)(1-h)}R_n^{\rm end}(h)\,\dd h.
 \end{aligned}                                             \tag{4.42d1}
\]
Thus
(4.42c) is an identity, not a comparison of cdf expansions at two sides
of a lattice jump.

We first bound the four matching-law terms.  Equations (4.40)--(4.42),
with $\xi=\eta/\sqrt n$, give
\[
 {\sqrt n\over\rho_2(p)}E_{Q,n}(\eta)
 =A(p)-{a(p)\over\rho_2(p)}\xi^2
   +{\ell(p)\over\rho_2(p)}{\xi\over\sqrt n}
   +O\left(\xi^4+{\xi^2\over\sqrt n}+{1\over n}\right).
                                                               \tag{4.42e}
\]
Since $\xi=o(1)$ and $a(p)/\rho_2(p)$ stays bounded away from zero,
Young's inequality yields fixed $c_1,C_1>0$ such that
\[
 {\sqrt n\over\rho}E_{Q,n}(\eta)
 \le A(p)-c_1\xi^2+{C_1\over n}.                     \tag{4.42f}
\]
Indeed, the left side of (4.42e) is $A(p)+o(1)>0$, uniformly along the
rows under consideration, and
\[
 {\sqrt n\over\rho}E_{Q,n}(\eta)
 ={\rho_2(p)\over\rho}
       {\sqrt n\over\rho_2(p)}E_{Q,n}(\eta)
 \le {\sqrt n\over\rho_2(p)}E_{Q,n}(\eta).
\]

For the zero-cell phase increment define, at real root-scale coordinate
$x$,
\[
 Z_{0,n}(p,x):={\sqrt n\over\rho_2(p)}\left\{
 {1\over2}-\Phi(sx)+2\Re\int_0^{1/4}K(t)
       e^{-2\pi\ii\sqrt n\,x t}c_p(t)^n\,\dd t\right\}.             \tag{4.42f0}
\]
Thus the normal term is included.  Since both the normal argument and the
Fourier phase depend on $\eta+\theta$, one has the exact identity
\[
 {\sqrt n\over\rho}
 \{D^{(0)}_{Q,n}(\eta,\theta)-D^{(0)}_{Q,n}(\eta,0)\}
 ={\rho_2(p)\over\rho}
 \{Z_{0,n}(p,\xi+\theta/\sqrt n)-Z_{0,n}(p,\xi)\}.       \tag{4.42f1}
\]
On a fixed compact set of $(p,x)$, the substitution $t=v/\sqrt n$ gives
\[
 Z_{0,n}(p,x)=H_{0,0}(p,x)+n^{-1/2}H_{1,0}(p,x)
                         +O_{C^1}(n^{-1}),             \tag{4.42g}
\]
where $H_{1,0}$ has bounded $C^1$ norm and
\[
 H_{0,0}(p,x)={\phi(sx)\over\rho_2(p)}
       \left\{{s\over2}+{\gamma\over6}(1-s^2x^2)\right\}          \tag{4.42g1}
\]
is even.  Indeed, Taylor expansion one order beyond the leading profile
gives, uniformly for $v\le n^{1/10}$,
\[
 \begin{aligned}
 c_p(v/\sqrt n)^n
 &=e^{-2\pi^2pqv^2}\left\{1+n^{-1/2}\ii\beta_3(p)v^3
       +n^{-1}P_6(p,v)+O(n^{-3/2}P_9(v))\right\},\\
 2n^{-1/2}K(v/\sqrt n)
 &={\ii\over\pi v}+n^{-1/2}+n^{-1}\kappa_Kv
       +O(n^{-3/2}v^2),
 \end{aligned}                                                \tag{4.42g2}
\]
where $\beta_3(p):=-(4\pi^3/3)p(1-p)(1-2p)$,
$P_6(p,\cdot)$ has degree at most six with coefficients uniformly
bounded near $p_*$, $P_9$ is a fixed positive polynomial majorant, and
$\kappa_K$ is fixed.  The order-one singular integral cancels the normal
term, by
\[
 {1\over2}-\Phi(sx)+{1\over\pi}\int_0^\infty
       e^{-2\pi^2pqv^2}{\sin(2\pi xv)\over v}\,\dd v=0.           \tag{4.42g3}
\]
The coefficients of $n^{-1/2}$ and $n^{-1}$ in the raw zero-cell
integral give respectively $\rho_2(p)H_{0,0}$ and
$\rho_2(p)H_{1,0}$.  Gaussian integration gives (4.42g1).
After one $x$ derivative, the unnormalized remainder in (4.42g2),
including all products, is bounded by
$Cn^{-3/2}P_{12}(v)e^{-cv^2}$.  The complement of this root-scale part has
the exact Bernoulli Gaussian bound, also after that derivative.  Keeping
$e^{-2\pi\ii x v}$ exact and multiplying by $\sqrt n/\rho_2(p)$ therefore
proves the $O_{C^1}(n^{-1})$ remainder in (4.42g).  Since
$\partial_x H_{0,0}(p,0)=0$, (4.42f1) and the mean-value theorem give
\[
 {\sqrt n\over\rho}
 \{D^{(0)}_{Q,n}(\eta,\theta)-D^{(0)}_{Q,n}(\eta,0)\}
 =O\left({|\xi|\over\sqrt n}+{1\over n}\right).       \tag{4.42h}
\]
The matching middle-cell phase increment is
$O(\sqrt n q_0^n)=o(n^{-1})$, for a fixed $q_0<1$.

It remains to keep the returned endpoint, including its value and shape
modes.  The expansion (4.29), which is uniform for
$|\eta|=o(\sqrt n)$ and has one uniform $\theta$ derivative, gives
\begin{equation}\tag{4.42i}
\begin{split}
 {\sqrt n\over\rho}\Re\{\zeta^nJ_{Q,n}(\eta,\theta)
                                  -J_{Q,n}(\eta,0)\}
 ={}&{C_0(p)\over\rho\sqrt n}
       \{r^n\cos(2\pi u)-1\}\\
 &+{r^nC_1(p,\eta+\theta)\over\rho n}\sin(2\pi u)
   +O\left({1+\eta^2\over n^{3/2}}\right).
\end{split}
\end{equation}
In particular, its value part contains
\[
 -{C_0(p)\over\rho\sqrt n}(1-r^n)
 +{C_0(p)r^n\over\rho\sqrt n}\{\cos(2\pi u)-1\},     \tag{4.42j}
\]
and the second term is nonpositive.  Lemma~\ref{orig:lem:endshape}, used here
before any localization of $\eta$ or of the circular phase, and (4.16)
give
\begin{align*}
 \mathcal S_n
 &\le\{c_{\rm sh}(p,d)+o(1)\}r^{n-1}|m|\\
 &\le Cd+{2c_{\rm sh}(p,d)+o(1)\over n}(1-r^n).
                                                               \tag{4.42k}
\end{align*}
Here we used $|1-\zeta|=O(d)$, the $O(d^2)$ term was included in $Cd$,
and (4.17) was used for the displayed radial term.  Since $C_0(p)/\rho$
stays bounded away from zero while $c_{\rm sh}(p,d)$ stays bounded,
\[
 -\left\{{C_0(p)\over\rho\sqrt n}
       -{2c_{\rm sh}(p,d)+o(1)\over n}\right\}(1-r^n)\le0          \tag{4.42l}
\]
for all sufficiently large $n$.  This explicitly absorbs the radial
shape mode into the returned value mode.

For completeness, we also bound the endpoint power remainder rather than
hiding it in the word ``Taylor.''  Write
\[
 A_h=\Psi(2\pi(1-h)),\quad
 B_h=\zeta\Psi_Q(2\pi(1-h)),\quad
 l_h=-2\pi\ii hm,\quad e_h=A_h-B_h-l_h.
\]
Lemma~\ref{orig:lem:endfactor} gives $|e_h|\le Cdh^2$.  The Bernoulli gap and
$|1-\zeta|+|m|\le Cd$ imply, uniformly for $0\le\tau\le1$,
\[
 \max\{|A_h|,|B_h|,|B_h+\tau(A_h-B_h)|\}
        \le \exp\{-c h^2+C d^2\}.                       \tag{4.42m}
\]
Indeed, for some fixed $c_0>0$, $|B_h|\le e^{-c_0h^2}$ and
$|A_h-B_h|\le |l_h|+|e_h|\le Cdh$.  On the root-scale part this gives
$|B_h+\tau(A_h-B_h)|\le1-c_1h^2+Cdh$, and
$Cdh\le c_1h^2/2+Cd^2$ proves (4.42m); the fixed part of the endpoint cell
is covered by the strict Bernoulli gap.
Taylor's exact integral identity
\[
 (a+b)^n-a^n-na^{n-1}l
 =na^{n-1}e+n(n-1)b^2
       \int_0^1(1-\tau)(a+\tau b)^{n-2}\,\dd\tau,
\]
with $a=B_h$, $b=l_h+e_h$, and $e=e_h$, now gives
\[
 |R_n^{\rm end}(h)|
 \le C\{ndh^2+n^2|m|^2h^2+n^2d^2h^4\}
              e^{-cnh^2+Cnd^2}.                       \tag{4.42n}
\]
Here $R_n^{\rm end}(h)=A_h^n-B_h^n-nB_h^{n-1}l_h$, and
$\mathcal R_n$ in (4.42c) is its normalized returned-endpoint integral
with the common spatial phase and kernel.
Because $nd^2\to0$, $|K(1-h)|\le Ch$,
$\int_0^\infty h^3e^{-cnh^2}\dd h=O(n^{-2})$, and
$\int_0^\infty h^5e^{-cnh^2}\dd h=O(n^{-3})$, its normalized integral
satisfies
\[
 |\mathcal R_n|
 \le C\left\{{d\over\sqrt n}+\sqrt n|m|^2
                         +{d^2\over\sqrt n}\right\}
 =o(d+n^{-1}).                                         \tag{4.42o}
\]

Finally Lemma~\ref{orig:lem:abslow} gives
$|\mathcal E_{0,n}|\le\{\Lambda_{\rm abs}(p,d)+o(1)\}d=O(d)$.
On the middle cell the matching transform has a fixed Cram\'er gap and
Theorem~\ref{orig:thm:cf} gives a uniform $O(d)$ transform difference; hence
\[
 |\mathcal E_{M,n}|\le Cn^{3/2}d q_1^{\,n-1}=o(d)
\]
for some $q_1<1$ (and it is exactly zero when $d=0$).

Insert (4.42f), (4.42h)--(4.42o) into the exact ledger (4.42c).
Since $C_1(p,z)$ is affine in $z$, all still unsigned mixed terms are
bounded by
\[
 C{|\xi|\over\sqrt n}+{C\over n}
       +O\left({1+\eta^2\over n^{3/2}}\right).
\]
For any fixed $\varepsilon>0$, Young's inequality gives
$C|\xi|/\sqrt n\le\varepsilon\xi^2+C_\varepsilon/n$; moreover
$(1+\eta^2)n^{-3/2}=o(n^{-1}+\xi^2)$.  Choose $\varepsilon$ smaller
than half the loss in (4.42f), absorb every pure $o(n^{-1})$ term into
$C_{\rm loc}/n$, and use
$\mathcal C_n^+\le\mathcal P_n^+$.  This proves (4.42a) with fixed
$c_{\rm loc},C_{\rm loc}>0$.

Since $A(p)\le\CE$, a score larger than $\CE$ forces, after absorbing
the sequential $o(d+\xi^2)$ term,
$\xi^2\le C(d+n^{-1})$.  Equivalently,
$\eta^2=O(nd+1)$, as asserted.
\end{proof}

We now record the signed endpoint consequence of the factorization.  This
is the common input for all three defect scales below.

\begin{lemma}[perturbed endpoint assembly]\label{orig:lem:pertendpoint}
Assume $p\to p_*$, $d\to0$, and $\sqrt n\,d\to0$.  Since
$|1-\zeta|=O(d)$, one has $\zeta\ne0$ for all sufficiently large $n$.
Write $\zeta=re^{\ii\alpha}$ and choose
$\vartheta\in[-1/2,1/2)$ so that
$n\alpha=2\pi(j+\vartheta)$ for an integer $j$.  For each physical
$\theta\in[-1/2,1/2]$ put $u=\theta-\vartheta$; only the trigonometric
class of $u$ modulo $\mathbb Z$ will be used.  Define the complex matching
endpoint integral
\[
 J_{Q,n}(\eta,\theta):=2\int_{3/4}^1K(t)
 e^{-2\pi\ii(k+\theta)t}\Psi_Q(2\pi t)^n\,\dd t;
 \qquad D^{(1)}_{Q,n}=\Re J_{Q,n}.
\]
The phase-modulated analytic baseline is
\[
 \mathcal B^+_{n,p,\zeta}(\eta,\theta)
 :={\sqrt n\over\rho}\left\{
 D^{(0)}_{Q,n}(\eta,\theta)
 +\Re\bigl(\zeta^nJ_{Q,n}(\eta,\theta)\bigr)
 +M_{Q,n}(\eta,\theta)\right\}.                    \tag{4.42P}
\]
Every term here uses the same integer $k$, the same physical phase
$\theta$, and hence the same cdf point $x_{\eta,\theta}$.  This is only an
analytic comparison object; unless $\zeta=1$ it is not asserted to be the
Prawitz operator of a probability law.
Uniformly under these hypotheses,
\begin{equation}\tag{4.43}
\begin{split}
 \Psi(2\pi(1-h))^n={}&\zeta^n\Psi_Q(2\pi(1-h))^n\\
 &-2\pi\ii nhm\,\zeta^{n-1}
          \Psi_Q(2\pi(1-h))^{n-1}+\mathcal R_n(h),
\end{split}
\end{equation}
where
\[
 {\sqrt n\over\rho}\,2\int_0^{1/4}|K(1-h)|
             |\mathcal R_n(h)|\,\dd h
 \le C\{d/\sqrt n+\sqrt n\,|m|^2\}.                 \tag{4.43a}
\]
The complete value mode relative to the calibrated two-point node is
\begin{align*}
 &{\sqrt n\over\rho}\Re\{\zeta^nJ_{Q,n}(\eta,\theta)
                         -J_{Q,n}(\eta,0)\}\\
 &\quad={C_0(p)\over\rho\sqrt n}
              \{r^n\cos(2\pi u)-1\}
 +{r^nC_1(p,\eta+\theta)\over\rho n}\sin(2\pi u)
 +O\left({1+\eta^2\over n^{3/2}}\right).             \tag{4.43b}
\end{align*}
The first shape mode is at most
\[
 \{c_{\rm sh}(p,d)+o(1)\}r^{n-1}|m|,
 \qquad c_{\rm sh}(p,d)\longrightarrow c_{\rm sh,*}. \tag{4.43c}
\]
Consequently its radial part is absorbed by the value mode:
\[
 -\left\{{C_0(p)\over\rho\sqrt n}
       -{2c_{\rm sh}(p,d)+o(1)\over n}\right\}(1-r^n)\le0.
                                                                  \tag{4.43d}
\]

For every fixed $D,M<\infty$, under $nd\le D$ and $|\eta|\le M$, the
global circular-phase comparison, followed by the zero-cell increment,
gives uniformly in the admissible $\eta$ and $|\theta|\le1/2$
\[
 \mathcal B^+_{n,p,\zeta}(\eta,\theta)
 \le C_Q^*(n)-{a_{\rm tan,*}\over\rho_*n}
          (\eta+\theta-\eta_*)^2+o_{D,M}(n^{-1}).      \tag{4.43e}
\]
On the localized positive branch the phase-free low-frequency bound gives
\[
 \mathcal C_n^+(X;\eta,\theta)
 \le C_Q^*(n)-\{\CE/\rho_*-{19\over16}c_{\rm sh,*}
                   -\Lambda_{\rm abs,*}-o(1)\}d
 -{a_{\rm tan,*}\over\rho_*n}(\eta+\theta-\eta_*)^2
 +o_{D,M}(n^{-1}).
                                                                  \tag{4.43f}
\]
If instead $d=o(n^{-1/2})$, $nd\to\infty$, and
$\eta^2=O(nd+1)$, then
\begin{equation}\tag{4.43g}
 \max\{\mathcal C_n^+(X;\eta,\theta),
          \mathcal C_n^-(X;\eta,\theta)\}
 \le C_Q^*(n)-
 \left\{{C_Q^*(n)\over\rho}
          -{19\over16}c_{\rm sh,*}
          -\Lambda_{\rm abs,*}-o(1)\right\}d+o(d).
\end{equation}
In particular, the coefficient in braces tends to the first positive
margin in (4.24).
\end{lemma}

\begin{proof}
Write $A_h=\Psi(2\pi(1-h))$, $B_h=\zeta\Psi_Q(2\pi(1-h))$,
$l_h=-2\pi\ii hm$, and $e_h=h^2dR_{p,h}$, where $|R_{p,h}|\le C$ by
(4.26).  The Bernoulli gap and $|1-\zeta|+|m|\le Cd$ imply, uniformly for
$0\le t\le1$,
\[
 \max\{|A_h|,|B_h|,|B_h+t(A_h-B_h)|\}
 \le \exp\{-c_1h^2+C_1d^2\}.                         \tag{4.43a0}
\]
Near zero this follows from
$e^{-ch^2}+Cdh\le1-c'h^2+Cdh$ and completion of the square; outside a
smaller fixed cell it follows from the strict Bernoulli gap.  Taylor's
literal integral identity is
\[
 (a+b)^n-a^n-na^{n-1}l
 =na^{n-1}e+n(n-1)b^2\int_0^1(1-t)(a+tb)^{n-2}\,\dd t,              \tag{4.43a1}
\]
where $a=B_h$, $b=l_h+e_h$, and $e=e_h$.  Equations (4.43a0)--(4.43a1)
give
\[
 |\mathcal R_n(h)|\le C\{ndh^2+n^2|m|^2h^2+n^2d^2h^4\}
 e^{-cnh^2+Cnd^2}.
\]
Since $nd^2\to0$, $|K(1-h)|\le Ch$,
$\int h^3e^{-cnh^2}\,\dd h=O(n^{-2})$, and
$\int h^5e^{-cnh^2}\,\dd h=O(n^{-3})$, this proves
(4.43)--(4.43a).  Differentiating the endpoint operator in the cdf phase
only multiplies the integrand by the bounded factor $-2\pi\ii(1-h)$, so
the same bound holds for one phase derivative.  No fourth moment is used.
Multiplying (4.29) by $\zeta^n$, using
$e^{-2\pi\ii\theta}\zeta^n=r^n e^{-2\pi\ii u}$, and taking real parts gives
(4.43b).  Lemma~\ref{orig:lem:endshape}, which is independent of the present
phase localization, gives (4.43c).  Equations (4.16)--(4.17) and
$r^n\cos(2\pi u)-1\le-(1-r^n)$ give (4.43d).
Indeed, separately from the case $r=1$, where the radial factor vanishes,
$p\to p_*$ and $d\to0$ imply
\[
 {C_0(p)\over\rho\sqrt n}
 -{2c_{\rm sh}(p,d)+o(1)\over n}
 \ge {c\over\sqrt n}-{C\over n}>0                    \tag{4.43d1}
\]
for all sufficiently large $n$.  Thus no division by $1-r$, $1-r^n$, or
$d$ is hidden in the absorption.

Fix $D,M<\infty$, assume $nd\le D$ and $|\eta|\le M$, and split off the
radial part in (4.43b)
and put
\[
 a_n={C_0(p)r^n\over\rho\sqrt n},\qquad
 b_n(\theta)={r^nC_1(p,\eta+\theta)\over\rho n},\qquad
 B_n=\sup_{|\theta|\le1/2}|b_n(\theta)|.
\]
Characteristic-function stability gives
$1-r\le |1-\zeta|\le Cd\le CD/n$, so $r^n\ge c_D>0$ for all sufficiently
large $n$.  Hence $a_n\asymp_D n^{-1/2}$ and
$B_n=O_{D,M}(n^{-1})$.
For every physical $\theta$, not merely at a maximizing phase,
\[
 a_n\{\cos(2\pi u)-1\}+b_n(\theta)\sin(2\pi u)
 \le a_n\{\cos(2\pi u)-1\}+B_n|\sin(2\pi u)|
 \le\sqrt{a_n^2+B_n^2}-a_n\le {B_n^2\over2a_n}.
\]
Thus the possible positive phase residue is $O_{D,M}(n^{-3/2})$ uniformly in
$\vartheta$ and $\theta$.  Retaining the radial part, (4.43b) gives
\[
 {\sqrt n\over\rho}\Re\{\zeta^nJ_{Q,n}(\eta,\theta)
                         -J_{Q,n}(\eta,0)\}
 \le -{C_0(p)\over\rho\sqrt n}(1-r^n)+O_{D,M}(n^{-3/2}). \tag{4.43b1}
\]

To identify the reference object, note that the already defined quantity
(4.24c) is equivalently
\[
 E_{Q,n}(\eta)=D^{(0)}_{Q,n}(\eta,0)
       +\Re J_{Q,n}(\eta,0)+M_{Q,n}(\eta,0).
\]
At $\theta=0$ the cardinal excess vanishes on the integer-valued Bernoulli
sum by (P3), so
\[
 {\sqrt n E_{Q,n}(\eta)\over\rho_2(p)}
\]
is exactly the right-node Bernoulli score to which
Lemma~\ref{orig:lem:boundednode} applies.  Direct
subtraction, with no change of the physical point inside the baseline,
gives
\begin{align*}
 \mathcal B^+_{n,p,\zeta}(\eta,\theta)
 ={\sqrt n\over\rho}E_{Q,n}(\eta)
 &+{\sqrt n\over\rho}\{D^{(0)}_{Q,n}(\eta,\theta)
                         -D^{(0)}_{Q,n}(\eta,0)\}\\
 &+{\sqrt n\over\rho}\Re\{\zeta^nJ_{Q,n}(\eta,\theta)
                         -J_{Q,n}(\eta,0)\}\\
 &+{\sqrt n\over\rho}\{M_{Q,n}(\eta,\theta)
                         -M_{Q,n}(\eta,0)\}.
\end{align*}
The last line is exponentially small uniformly on the bounded node/phase
set.  Put
\[
 \lambda_n:={\rho_2(p)\over\rho},\qquad
 \tau_n:={a_{\rm tan,*}\over\rho_*n}
              (\eta+\theta-\eta_*)^2 .
\]
Lemma~\ref{orig:lem:boundednode}, (4.31), and (4.43b1) first give the naturally
normalized inequality
\[
 \mathcal B^+_{n,p,\zeta}(\eta,\theta)
 \le \lambda_n\{C_Q^*(n)-\tau_n\}
 -{C_0(p)\over\rho\sqrt n}(1-r^n)
 +\widetilde R_n,
 \qquad \widetilde R_n=o_{D,M}(n^{-1}).                  \tag{4.43e0}
\]
Here the factor $\lambda_n$ multiplies the complete tangent square; it
cannot be removed merely from $\lambda_n\le1$, because the tangent enters
with a negative sign.  Instead, $\rho=\rho_2(p)+d$ gives
\[
 1-\lambda_n={d\over\rho},
 \qquad
 \lambda_n\{C_Q^*(n)-\tau_n\}
 =\lambda_n C_Q^*(n)-\tau_n+R_n^{\rm norm},
\]
where
\[
 R_n^{\rm norm}:=(1-\lambda_n)\tau_n
 ={d\over\rho}\,{a_{\rm tan,*}\over\rho_*n}
       (\eta+\theta-\eta_*)^2 .
\]
This is uniformly $o_{D,M}(n^{-1})$.  Indeed, under the quantifiers in
(S.3),
\[
 |\eta+\theta-\eta_*|
 \le M+\tfrac12+|\eta_*|,\qquad d\le {D\over n},
 \qquad \rho\longrightarrow\rho_*>0,
\]
and hence, for all sufficiently large $n$,
\[
 0\le R_n^{\rm norm}
 \le {2a_{\rm tan,*}D\over\rho_*^2}
       \left(M+\tfrac12+|\eta_*|\right)^2 n^{-2}
 =o_{D,M}(n^{-1}).
\]
After absorbing $R_n^{\rm norm}$ into $\widetilde R_n$, use
$\ell=2a\eta_0$ to obtain algebraically
\[
 -{a_{\rm tan,*}\over\rho_*n}\{(\eta-\eta_*)^2+\theta^2
                  -2(\eta_* -\eta)\theta\}
 =-{a_{\rm tan,*}\over\rho_*n}(\eta+\theta-\eta_*)^2.
\]
Consequently,
\[
 \mathcal B^+_{n,p,\zeta}(\eta,\theta)
 \le {\rho_2(p)\over\rho}C_Q^*(n)
 -{a_{\rm tan,*}\over\rho_*n}(\eta+\theta-\eta_*)^2
 -{C_0(p)\over\rho\sqrt n}(1-r^n)+o_{D,M}(n^{-1}).   \tag{4.43e1}
\]
Only at this point do $C_Q^*(n)>0$ and $\rho_2(p)/\rho\le1$ imply the
displayed weaker form (4.43e).  In particular, the retained phase $\zeta^n$
has not been replaced by one.

It remains to compare the actual operator with this baseline.  The exact
decomposition (4.24a) counts the zero cell, returned endpoint, and middle
once each.  The zero-cell difference is at most
$\{\Lambda_{\rm abs,*}+o(1)\}d$.  The endpoint linear term is at most
\[
 \{c_{\rm sh}(p,d)+o(1)\}r^{n-1}|m|
 \le \{\tfrac{19}{16}c_{\rm sh,*}+o(1)\}d
      +\{c_{\rm sh}(p,d)+o(1)\}r^{n-1}(1-r^2)+O(d^2).
\]
The second term is at most
$\{2c_{\rm sh}(p,d)+o(1)\}(1-r^n)/n$ and is absorbed by the negative radial
term in (4.43b1).  The endpoint remainder is $o_D(n^{-1})$, and
$M_{X,n}-M_{Q,n}$ is exponentially small.  Finally, since
$\rho=\rho_2(p)+d$ and the reference right-node score tends to $\CE$,
\[
 {\rho_2(p)\over\rho}C_Q^*(n)
 =C_Q^*(n)-\{\CE/\rho_*+o(1)\}d.
\]
There is no rate conversion hidden here.  Indeed (4.16) and
$|1-\zeta|=O(d)$ give $|m|=O(d)$.  Under $nd\le D$ and $|\eta|\le M$,
the terms not displayed with a sign satisfy the explicit ledger
\[
\begin{array}{c|c}
\text{term}&\text{bound}\\ \hline
\text{endpoint power remainder (4.43a)}
 &O(d/\sqrt n+\sqrt n\,d^2)=O_D(n^{-3/2})\\
\text{$C_1$ phase residue after the sine--cosine completion}
 &O_{D,M}(n^{-3/2})\\
\text{shape-continuity and quadratic errors}&o(1)d+O(d^2)=o_D(n^{-1})\\
\text{low-frequency remainder}&o(1)d=o_D(n^{-1})\\
\text{middle-cell difference}&O(n^{3/2}dq_1^n)=o_D(n^{-1}).
\end{array}                                                   \tag{4.43f0}
\]
The second row uses the global completion preceding (4.43b1), not an
assertion that the affine $C_1$ term has a sign by itself.  On the
localized branch Lemma~\ref{orig:lem:coarsephase} supplies
$\eta^2=O(nd+1)$, so boundedness of $\eta$ in this regime is a conclusion,
not an additional hypothesis.
Combining these signed terms proves (4.43f); its coefficient is positive
by (4.24).  Notice that the loss $\CE d/\rho_*$ and the radial absorption
belong to the combined baseline comparison, not to a false standalone
inequality asserting that the actual operator is already below the
baseline by a fixed multiple of $d$.

In the growing window use the two signed rows at the same active point,
rather than extracting a tangent.  The complete upper and lower
two-point cells at the calibrated integer node are at most
$\{\rho_2(p)/\rho\}C_Q^*(n)$.  The corresponding zero-cell spatial
increments have a common error
$O((1+|\eta|)/n)=o(d)$.  Indeed $\eta^2=O(nd)$ and $nd\to\infty$ give
\[
 {1+|\eta|\over nd}=O((nd)^{-1}+(nd)^{-1/2})=o(1),
 \qquad
 {1+\eta^2\over n^{3/2}d}=o(1)+O(n^{-1/2})=o(1).      \tag{4.43g0}
\]
The signed middle-cell error is $o(n^{-1})=o(d)$ because $nd\to\infty$.
The physical returned-endpoint owner contributes
$\{19c_{\rm sh,*}/16+o(1)\}d+o(d)$ simultaneously for both signs, and
the actual-minus-matching zero/middle owner contributes
$\{\Lambda_{\rm abs,*}+o(1)\}d+o(d)$.  These are the same owners as in
(4.24d), each inserted once.  Moreover (4.16), (4.43a), and
$\sqrt n\,d\to0$ give
\[
 |m|=O(d),\qquad
 {d\over\sqrt n}+\sqrt n|m|^2
 \le d\{n^{-1/2}+O(\sqrt n\,d)\}=o(d).               \tag{4.43g1}
\]
The remaining low-frequency and exponential errors are $o(d)$.  Finally
the exact denominator identity
\[
 {\rho_2(p)\over\rho}C_Q^*(n)
 =C_Q^*(n)-{C_Q^*(n)\over\rho}d,\qquad
 {C_Q^*(n)\over\rho}\longrightarrow {\CE\over\rho_*}, \tag{4.43g2}
\]
turns the paired owner ledger into (4.43g).  Its coefficient tends to
$\CE/\rho_*-19c_{\rm sh,*}/16-\Lambda_{\rm abs,*}>0$ by (4.24).
This proof is uniform in the active point and does not require the
more detailed $A(p)$ tangent expansion.
\end{proof}

\subsection*{The active signed subcritical estimates}

This supplement makes explicit two points which are easy to lose in the
compressed proof of Lemma~\ref{orig:lem:pertendpoint}.  First, all pieces of a
row are evaluated at one and the same integer, phase, and physical cdf
point.  Second, in the growing window the Gaussian boundary/quadrature and
the signed binomial stop--loss prefix are one quantity.  Neither of them is
estimated in absolute value before they have cancelled.

\paragraph{Signed orientation and uniform quantifiers.}
The upper and lower signed estimates use the same orientation. Their
node--zero, middle, and returned contributions are treated together.
We display the upper notation,
\[
 k=np+\eta\in\mathbb Z,\qquad |\theta|\le\tfrac12,
 \qquad x_{\eta,\theta}={s(p)(\eta+\theta)\over\sqrt n}.       \tag{S.1}
\]
The lower row uses the conjugate kernel and its own lower node--zero
identity at this same physical point.  Thus no formula combines an upper
integer-node score with a lower returned endpoint, and every occurrence of
\(k\), \(\theta\), and \(x_{\eta,\theta}\) refers to the same physical
point. Reflection gives an equivalent proof, but is unnecessary for
these paired estimates.

We use the following precise meaning of uniform little-oh.  Fix
\(D,M<\infty\).  An assertion
\[
 R_n(\eta,\theta)=o_{D,M}(n^{-1})                         \tag{S.2}
\]
means that for every sequence of admissible rows with
\[
 p_n\to p_*,\qquad d_n\to0,\qquad \sqrt n\,d_n\to0,
 \qquad nd_n\le D,
\]
one has
\[
 \sup_{\substack{|\eta|\le M,\ |\theta|\le1/2,\\
                  np_n+\eta\in\mathbb Z}}
       n\,|R_n(\eta,\theta)|\longrightarrow0.            \tag{S.3}
\]
Only cells occurring in the row enter the supremum.  Similarly,
\(R_n=o_{C_\eta}(d_n)\) in the growing window means that the supremum of
\(|R_n|/d_n\) tends to zero along every admissible sequence satisfying
\[
 d_n=o(n^{-1/2}),\qquad nd_n\to\infty,\qquad
 \eta^2\le C_\eta(nd_n+1).                              \tag{S.4}
\]
Thus the quantifier over the active cell and physical phase is inside the
limit.  The coarse localization lemma is used only to supply a
finite \(M=M(D)\) in the bounded branch and a finite \(C_\eta\) in the
growing branch.  The assembly itself applies to every cell satisfying the
displayed bounds.  An \emph{active point} has exactly the meaning from
(4.0a), namely \(\mathcal C_n^+(X;\eta,\theta)>\CE\); it carries no extra
analytic assumption.

\paragraph{One exact finite-row ledger.}
Let \(J_{Q,n}\) be as in Lemma~\ref{orig:lem:pertendpoint}.  In oriented upper
notation put
\begin{equation}\tag{S.5}
\begin{aligned}
 T_n&={\sqrt n\over\rho}\left\{E_{Q,n}(\eta)
       +D^{(0)}_{Q,n}(\eta,\theta)-D^{(0)}_{Q,n}(\eta,0)\right\},\\
 V_n&={\sqrt n\over\rho}\Re\{\zeta^nJ_{Q,n}(\eta,\theta)
                                      -J_{Q,n}(\eta,0)\},\\
 H_n&={\sqrt n\over\rho}\{M_{Q,n}(\eta,\theta)
                                      -M_{Q,n}(\eta,0)\},\\
 L_n&={\sqrt n\over\rho}\{D^{(0)}_{X,n}(\eta,\theta)
                                      -D^{(0)}_{Q,n}(\eta,\theta)\},\\
 E_n&={\sqrt n\over\rho}\{M_{X,n}(\eta,\theta)
                                      -M_{Q,n}(\eta,\theta)\}.
\end{aligned}
\end{equation}
Let \(S_n\) be the normalized real part of the first-order endpoint term
\[
 -2\pi\ii nhm\,\zeta^{\,n-1}
       \Psi_Q(2\pi(1-h))^{n-1}                               \tag{S.6}
\]
integrated on \(0\le h\le1/4\) with the returned Prawitz kernel, and let
\(R_n^{\rm pow}\) be the integral of the literal remainder
\(\mathcal R_n(h)\) in (4.43).  These are definitions, not estimates.  The
exact decomposition (4.24a), with (4.24c) added and subtracted once, gives
\[
 \boxed{\quad
 \mathcal C_n^+(X;\eta,\theta)
 \le \mathcal P_n^+(X;\eta,\theta)
 =T_n+V_n+H_n+L_n+S_n+R_n^{\rm pow}+E_n .
 \quad}                                                       \tag{S.7}
\]
No term uses an independently chosen node or phase.  In particular,
\(E_{Q,n}\) owns all four upper integer-node pieces in (4.24c), and the
returned endpoint occurs only in \(V_n+S_n+R_n^{\rm pow}\).

There are constants \(C,c>0\), \(0<q_1<1\), and uniform moduli
\(\omega_{\rm lo},\omega_{\rm sh}\to0\) such that
\begin{align}
 |L_n|&\le\{\Lambda_{\rm abs,*}+\omega_{\rm lo}(n)\}d,
                                                               \tag{S.8}\\
 |E_n|&\le Cn^{3/2}d q_1^{\,n-1},                              \tag{S.9}\\
 |H_n|&\le C\sqrt n\,q_1^n,                                   \tag{S.10}\\
 |R_n^{\rm pow}|&\le
 C\{d/\sqrt n+\sqrt n\,|m|^2+d^2/\sqrt n\},                   \tag{S.11}\\
 |m|&\le C d.                                                  \tag{S.12}
\end{align}
Here (S.8) is (4.18), (S.9) is the middle-cell estimate, (S.10) is
the matching middle gap, and (S.11) follows from (4.43a).  The harmless
last summand in (S.11) merely displays all quadratic endpoint remainders.
Equation (S.12) follows from (4.16) and \(|1-\zeta|=O(d)\).

\subsubsection*{The bounded active-localized branch}

Fix \(D,M<\infty\), assume \(nd\le D\), and restrict to active cells with
\(|\eta|\le M\).  The bounded-node expansion, zero-cell increment,
returned-endpoint sine--cosine completion, and (S.10) give the single
signed inequality
\begin{align}
 T_n+V_n+H_n
 \le{}&{\rho_2(p)\over\rho}C_Q^*(n)
 -{a_{\rm tan,*}\over\rho_*n}
       (\eta+\theta-\eta_*)^2\notag\\
 &-{C_0(p)\over\rho\sqrt n}(1-r^n)
 +R_n^{\rm base},
 \qquad R_n^{\rm base}=o_{D,M}(n^{-1}).              \tag{S.13}
\end{align}
The possible positive returned-endpoint residue is bounded before it is
placed in \(R_n^{\rm base}\).  Indeed, with
\[
 a_n={C_0(p)r^n\over\rho\sqrt n},\qquad
 B_n=\sup_{|\theta|\le1/2}
       {r^n|C_1(p,\eta+\theta)|\over\rho n},
\]
one has \(a_n\asymp_Dn^{-1/2}\), \(B_n=O_{D,M}(n^{-1})\), and
\[
 a_n\{\cos(2\pi u)-1\}+B_n|\sin(2\pi u)|
 \le {B_n^2\over2a_n}=O_{D,M}(n^{-3/2}).               \tag{S.14}
\]
Thus (S.13) is uniform in both the retained phase \(\vartheta\) and the
physical phase \(\theta\).  It does not assign a sign to the affine
\(C_1\)-term by itself.

More explicitly, one may decompose
\[
 R_n^{\rm base}=R_n^{\rm node}+R_n^{0,\rm loc}
                 +R_n^{0,\rm tail}+R_n^{\rm circ}+R_n^{\rm coef}
                 +R_n^{\rm matchmid},                         \tag{S.14a}
\]
where, uniformly under the quantifiers in (S.3),
\[
\begin{aligned}
 R_n^{\rm node}&=o_M(n^{-1}),&
 |R_n^{0,\rm loc}|&\le C_Mn^{-3/2},\\
 |R_n^{0,\rm tail}|&\le C_Mn^{-3/2},&
 |R_n^{\rm circ}|&\le C_{D,M}n^{-3/2},\\
 R_n^{\rm coef}&=o_{D,M}(n^{-1}),&
 |R_n^{\rm matchmid}|&\le C\sqrt n\,q_1^n.
\end{aligned}                                                \tag{S.14b}
\]
The first row is the uniform remainder in the bounded-node lemma; the
second is (4.33) with bounded \(\eta\); the third is the finite-cell to
half-line remainder (4.38c1); the fourth contains (S.14), the remainder in
(4.43b); the fifth contains the \(o(1)/n\) coefficient replacements; and the
last is (S.10).  Thus every remainder in (S.13) has a named owner.
In particular, the fifth owner contains the normalization correction
\[
 R_n^{\rm norm}
 =\left(1-{\rho_2(p)\over\rho}\right)
   {a_{\rm tan,*}\over\rho_*n}
      (\eta+\theta-\eta_*)^2,
 \qquad
 0\le R_n^{\rm norm}\le C_{D,M}n^{-2}.
\]
Thus moving $\rho_2(p)/\rho$ off the negative tangent square is an explicit
$o_{D,M}(n^{-1})$ coefficient replacement, rather than an application of
$\rho_2(p)/\rho\le1$.

The shape row is retained until its radial part is paired with (S.13).
Equations (4.16), (4.17), and (4.43c) give
\begin{align}
 S_n\le{}&\{\tfrac{19}{16}c_{\rm sh,*}
                    +\omega_{\rm sh}(n)\}d
 +{2c_{\rm sh}(p,d)+\omega_{\rm sh}(n)\over n}(1-r^n)
 +Cd^2.                                                       \tag{S.15}
\end{align}
The radial rows in (S.13) and (S.15) obey, uniformly for all sufficiently
large \(n\),
\[
 \left[-{C_0(p)\over\rho\sqrt n}
       +{2c_{\rm sh}(p,d)+\omega_{\rm sh}(n)\over n}\right]
       (1-r^n)\le0.                                           \tag{S.16}
\]
No division by \(d\), \(1-r\), or \(1-r^n\) occurs in (S.16).

The denominator loss is the exact identity
\[
 {\rho_2(p)\over\rho}C_Q^*(n)
 =C_Q^*(n)-b_n d,\qquad b_n={C_Q^*(n)\over\rho}
       \longrightarrow{\CE\over\rho_*}.                      \tag{S.17}
\]
Substituting (S.8)--(S.17) into (S.7) gives
\begin{align}
 \mathcal C_n^+(X;\eta,\theta)
 \le{}&C_Q^*(n)
 -\left\{{\CE\over\rho_*}-{19\over16}c_{\rm sh,*}
                -\Lambda_{\rm abs,*}-\varepsilon_{n,D,M}\right\}d\notag\\
 &-{a_{\rm tan,*}\over\rho_*n}
       (\eta+\theta-\eta_*)^2+R_{n,D,M},                      \tag{S.18}
\end{align}
where \(\varepsilon_{n,D,M}\to0\) and
\begin{align}
 |R_{n,D,M}|\le{}&|R_n^{\rm base}|
 +C\{d/\sqrt n+\sqrt n\,d^2+d^2\}
 +\{\omega_{\rm lo}(n)+\omega_{\rm sh}(n)\}d\notag\\
 &+Cn^{3/2}dq_1^{\,n-1}.                                     \tag{S.19}
\end{align}
Every term on the right of (S.19) is \(o_{D,M}(n^{-1})\): since
\(d\le D/n\), the first two endpoint terms are \(O_D(n^{-3/2})\),
\(d^2=O_D(n^{-2})\), and exponential decay dominates every power of
\(n\).  This is (4.43f) with its finite-row, active-localized, and signed
quantifiers exposed.

\subsubsection*{The growing branch: the direct paired active ledger}

Assume (S.4) with $nd\to\infty$.  For each sign, let $T_n^\pm$ be the
complete matching two-point cell at the calibrated integer node plus the
signed matching zero-cell spatial increment from that node to the active
point.  Let $R_{n,\pm}^{\rm mid}$, $R_{n,\pm}^{\rm ret}$, and
$R_{n,\pm}^{\rm act}$ denote respectively the matching middle, physical
returned-endpoint, and actual-minus-matching zero/middle owners.  The exact
finite-row grouping is
\[
 \mathcal C_n^\pm(X;\eta,\theta)
 \le T_n^\pm+R_{n,\pm}^{\rm mid}
       +R_{n,\pm}^{\rm ret}+R_{n,\pm}^{\rm act}.          \tag{S.G0}
\]
The calibrated upper and lower node--zero blocks admit one signed error
$r_n^{\rm node}$ such that
\[
 T_n^\pm\le {\rho_2(p)\over\rho}C_Q^*(n)+r_n^{\rm node},
 \qquad {r_n^{\rm node}\over d}\longrightarrow0.          \tag{S.G1}
\]
The estimate is obtained before taking the maximum of the two signs.  The
integer-node terms use the exact two-point bound, while the two zero-cell
increments are bounded by $C(1+|\eta|)/n$; (4.43g0) makes the latter
$o(d)$.  No binomial prefix is separated from a Gaussian boundary term.

With the same cutoff and the same physical point, the remaining concrete
owners give errors satisfying
\begin{align}
 |R_{n,\pm}^{\rm mid}|&\le r_n^{\rm mid},&
 r_n^{\rm mid}/d&\longrightarrow0,\notag\\
 R_{n,\pm}^{\rm ret}
   &\le\{\tfrac{19}{16}c_{\rm sh,*}+o(1)\}d+r_n^{\rm ret},
       &r_n^{\rm ret}/d&\longrightarrow0,\notag\\
 R_{n,\pm}^{\rm act}
   &\le\{\Lambda_{\rm abs,*}+o(1)\}d+r_n^{\rm act},
       &r_n^{\rm act}/d&\longrightarrow0.                   \tag{S.G2}
\end{align}
Here $0\le r_n^{\rm mid}=o(n^{-1})=o(d)$, the returned owner includes the
power error (4.43a), and the actual owner is the full zero/middle
actual-minus-matching difference.  Hence the exact finite-row assembly
(S.G0) gives
\[
 \mathcal C_n^\pm(X;\eta,\theta)
 \le {\rho_2(p)\over\rho}C_Q^*(n)
 +\{\tfrac{19}{16}c_{\rm sh,*}
       +\Lambda_{\rm abs,*}+o(1)\}d+r_n,
 \qquad r_n/d\longrightarrow0.                              \tag{S.G3}
\]
Since $\rho=\rho_2(p)+d$,
\[
 {\rho_2(p)\over\rho}C_Q^*(n)
 =C_Q^*(n)-{C_Q^*(n)\over\rho}d,\qquad
 {C_Q^*(n)\over\rho}\longrightarrow{\CE\over\rho_*}.    \tag{S.G4}
\]
Taking the maximum in (S.G3) and using (S.G4) proves (4.43g).

\subsection*{The transition to one adaptive lattice}

The scale $\sqrt n\,d\asymp1$ cannot be treated by expanding around a
fixed two-point resonance: the resonance itself may move.  The next lemma
supplies the required signed compactness without approximating the
discontinuous sawtooth.

In this subsection an \emph{active near-extremal positive-branch
sequence} means a sequence of points \(x_n\) which are active in the sense
of (4.0a), are chosen on a putative extremizing sequence, and satisfy, with
\(\rho_n=\E|X_n|^3\),
\[
 {\sqrt n\over\rho_n}\{F_{n,X_n}(x_n)-\Phi(x_n)\}\longrightarrow\CE
\]
and the positive error is the competitive one.  The outer localization
proved above then gives $p_n\to p_*$, $d_n\to0$, $x_n\to0$, while the
reflected branch has a fixed gap.  We shall always state explicitly when
only the corresponding Prawitz operator, rather than the cdf error, is
being used; thus ``active'' carries no unmentioned Fourier estimate.

\begin{lemma}[spectral selection and the vanishing-power branch]
\label{orig:lem:spectral}
Let $X_n$ be centered, variance-one laws, set
$\gamma_n=\E X_n^3$, let $p_n$ be the matching Bernoulli parameter, and
put $d_n=\E|X_n|^3-\rho_2(\gamma_n)$.  Suppose
$p_n\to p_*$ and $d_n\to0$.  Put
\[
 R_n(u)=|\varphi_{X_n}(u)|^2,\qquad
 u_n^*={2\pi\over s(p_n)}.
\]
There are fixed $\delta,c>0$ such that, for all sufficiently large $n$,
$R_n$ has a unique interior maximizer $u_n$ on
$I_n=[u_n^*-\delta,u_n^*+\delta]$ and $R_n''\le-c$ there.
Choose once and for all
\[
 \max\{s_*/2,\gamma_*\}<h_0<s_*.
\]
If $|\varphi_{X_n}(u_n)|^n\to0$, then
\[
 \limsup_n\mathcal C_n(X_n)
 \le {\phi_0\over\rho_*}
       \left({h_0\over2}+{\gamma_*\over6}\right)<\CE. \tag{4.43h}
\]
The reflected one-sided score has a strictly smaller limit.
\end{lemma}

\begin{proof}
The defect identity and $d_n\to0$ concentrate the law at the two zeros of
the defect density.  The first two moments determine their limiting
weights, so $X_n$ converges to its matching two-point law.  Since the
second and third absolute moments also converge to those of that law,
$|X_n|^2$ and $|X_n|^3$ are uniformly integrable.  For an independent
copy $X_n'$ set $D_n=X_n-X_n'$.  Then
\[
 R_n''(u)=-\E\{(X_n-X_n')^2\cos(u(X_n-X_n'))\}.
\]
Uniform integrability of $D_n^2$, followed by a fixed tail truncation,
shows that both $R_n'$ and $R_n''$ converge uniformly on every fixed
compact $u$-interval.  In particular, uniformly on a fixed neighbourhood
of $u_n^*$,
\[
 -R_n''(u)\longrightarrow
 2\cos\{s_*(u-2\pi/s_*)\}.
\]
Choose $\delta>0$ and $c_0>0$ so that the last cosine is at least $4c_0$
when $|u-2\pi/s_*|\le2\delta$.  For all large $n$,
\[
 R_n''(u)\le-2c_0\quad(u\in I_n),\qquad
 R_n'(u_n^*-\delta)>c_0\delta,\qquad
 R_n'(u_n^*+\delta)<-c_0\delta.                       \tag{4.43h1}
\]
The derivative signs follow by integrating the limiting curvature from
the limiting critical point.  Thus strict concavity gives the asserted
unique $u_n\in\operatorname {int}I_n$.  Uniform convergence of $R_n'$
and uniqueness of the zero of its limit on this interval also give
$u_n\to2\pi/s_*$, and hence $R_n(u_n)\to1$.  Taylor's theorem at $u_n$
then gives, after decreasing $\delta$ if necessary,
\[
 |\varphi_{X_n}(u)|^n
 \le |\varphi_{X_n}(u_n)|^n
       e^{-c_1n(u-u_n)^2},\qquad |u-u_n|\le\delta,     \tag{4.43h2}
\]
with fixed $c_1>0$.  Indeed (4.43h1) gives
$R_n(u)\le R_n(u_n)-c_0(u-u_n)^2$, while
$R_n(u_n)\to1$; divide by $R_n(u_n)$ and use
$(1-y)^{n/2}\le e^{-ny/2}$.  This proves (4.43h2) after a uniform
decrease of $c_1$.

We make the frequency split used here explicit.  At physical bandwidth
$h_0$, put $U_0=2\pi/h_0$.  The assumptions $s_*/2<h_0<s_*$ imply
$2\pi/s_*<U_0<4\pi/s_*$.  Hence the limiting two-point modulus is one on
$[0,U_0]$ precisely at $0$ and $2\pi/s_*$.  Choose fixed disjoint
neighbourhoods
\[
 Z=[0,\delta_0],\qquad
 E_n=[u_n-\delta_1,u_n+\delta_1]\subset(0,U_0),
\]
with $\delta_0,\delta_1>0$ small enough for the Gaussian zero-frequency
bound and (4.43h2), respectively.  On the compact remainder
$M_n=[0,U_0]\setminus(Z\cup E_n)$, uniform convergence of characteristic
functions and the displayed description of the limiting equality set
give constants $q<1$ and $n_0$ such that
\[
             \sup_{u\in M_n}|\varphi_{X_n}(u)|\le q
             \qquad(n\ge n_0).                       \tag{4.43h3}
\]
Thus the normalized middle contribution is $O(\sqrt n q^n)$.  On $E_n$
the kernel and the change of physical frequency are bounded, and
(4.43h2) bounds the normalized contribution by
\[
 C\sqrt n\,|\varphi_{X_n}(u_n)|^n
       \int_\mathbb R e^{-c_1nv^2}\,\dd v=o(1).       \tag{4.43h4}
\]
This proves, with uniform constants, that only the zero cell remains.
Uniform integrability of the third moment gives the uniform Peano expansion
 \[
  \varphi_{X_n}(w)=1-{w^2\over2}-{\ii\gamma_nw^3\over6}
                          +|w|^3\eta_n(w),
 \]
 where
 \[
  \lim_{\delta\downarrow0}\limsup_n
       \sup_{|w|\le\delta}|\eta_n(w)|=0,
 \]
 together with a uniform Gaussian power envelope
 $|\varphi_{X_n}(w)|^n\le Ce^{-cnw^2}$ for small $w$.  In the Prawitz
 integral put
 \[
  t={h_0v\over2\pi\sqrt n},
  \qquad\text{equivalently}\qquad
  w={2\pi t\over h_0}={v\over\sqrt n}.
 \]
 The kernel expansion is then
 \[
  2K(t)\,\dd t=
  \left\{{\ii\over\pi v}+{h_0\over2\pi\sqrt n}
                    +O(v/n)\right\}\dd v,
 \]
 with the conjugate expansion for the reflected operator.  After
 multiplication by $\sqrt n$, every remainder is bounded, independently
 of the cdf coordinate, by a constant times
 \[
             (v^2+v^3+v^4)e^{-cv^2}.
 \]
 First pass to the limit on $0\le v\le M$ and then let $M\to\infty$.
 This two-stage domination yields, uniformly in the cdf coordinate, the
 two signed limits
\[
 {\phi(x)\over\rho_*}
 \left\{{h_0\over2}\mathbin\pm
             {\gamma_*\over6}(1-x^2)\right\}.        \tag{4.43i}
\]
For the plus sign the positive maximum is at zero.  For the minus sign it
is also at zero because $h_0>\gamma_*$, and it is smaller.  Finally
\[
 {\phi_0\over\rho_*}
 \left({h_0\over2}+{\gamma_*\over6}\right)
 <{\phi_0\over\rho_*}
 \left({s_*\over2}+{\gamma_*\over6}\right)=\CE.
\]
This proves the lemma.
\end{proof}

\begin{lemma}[non-circular adaptive majorant]\label{orig:lem:adaptiveupstream}
Let $X_n$ be centered variance-one laws with matching parameters
$p_n\to p_*$, defects $d_n\to0$, third moments $\gamma_n$, and third
absolute moments $\rho_n$.  If $h_n\to s_*$, define the exact Prawitz
operators at physical bandwidth $h_n$.  Then, uniformly in $x\in\mathbb R$,
\[
 {\sqrt n\over\rho_n}\{\mathcal U^+_{X_n,n,h_n}(x)-\Phi(x)\}
 ={\phi(x)\over\rho_n}
 \left\{{h_n\over2}+{\gamma_n\over6}(1-x^2)\right\}+o(1),          \tag{4.43j}
\]
and
\[
 {\sqrt n\over\rho_n}\{\Phi(x)-\mathcal U^-_{X_n,n,h_n}(x)\}
 ={\phi(x)\over\rho_n}
 \left\{{h_n\over2}-{\gamma_n\over6}(1-x^2)\right\}+o(1).          \tag{4.43k}
\]
Consequently the positive adaptive majorant is at most $\CE+o(1)$ and
the reflected one has a fixed limiting gap.  No residual-power conclusion
and no residual-variance estimate of order $o(n^{-1})$ is used.
\end{lemma}

\begin{proof}
Write $\varphi_n:=\varphi_{X_n}$.
Uniformly for $\gamma$ near $\gamma_*$, the exact defect density satisfies
\[
 w_\gamma(x)\ge c|x|^3\quad(|x|\ge R_0),\qquad
 \E w_{\gamma_n}(X_n)=d_n.                            \tag{4.43l}
\]
Thus $\{|X_n|^3\}$ is uniformly integrable: for all sufficiently large
$n$ its tail is at most $Cd_n$, and the finitely many earlier rows are
handled by increasing the cutoff.  Quadratic defect localization and the
moment equations also give $X_n\Rightarrow Q_*$.  Hence
\[
 \omega(r):=\sup_n\E\{|X_n|^3\min(r|X_n|,1)\}\longrightarrow0
 \quad(r\downarrow0),                                 \tag{4.43m}
\]
and
\[
 \varphi_n(w)=1-{w^2\over2}-{\ii\gamma_nw^3\over6}
                  +|w|^3\epsilon_n(w),\qquad
 |\epsilon_n(w)|\le C\omega(|w|).                    \tag{4.43n}
\]
Symmetrization and $1-\cos y\ge y^2/2-|y|^3/6$ give fixed
$c,\delta_0>0$ such that $|\varphi_n(w)|\le e^{-cw^2}$ for
$|w|\le\delta_0$.

Put
\[
 f_n(v)=\varphi_n(v/\sqrt n)^n,\qquad
 r_n(v)=f_n(v)-e^{-v^2/2}
 \left(1-{\ii\gamma_nv^3\over6\sqrt n}\right),\qquad B_n=n^{1/21}.
\]
The exact difference of powers and (4.43n) give, for $v\le B_n$,
\[
 |r_n(v)|\le Ce^{-cv^2}\left\{
 {v^3\over\sqrt n}\omega(B_n/\sqrt n)+{v^4+v^6\over n}\right\},
\]
whereas for $B_n\le v\le\delta_0\sqrt n$ the Gaussian envelope gives
$|r_n(v)|\le Ce^{-cv^2}(1+v^3/\sqrt n)$.  Therefore
\[
 \sqrt n\int_0^{\delta_0\sqrt n}{|r_n(v)|\over v}\,\dd v=o(1),
 \qquad
 \int_0^{\delta_0\sqrt n}|f_n(v)-e^{-v^2/2}|\,\dd v=o(1),          \tag{4.43o}
\]
and $\int v|f_n(v)|\,\dd v=O(1)$ on that interval.

On the zero cell put $t=h_nv/(2\pi\sqrt n)$.  The signed kernel expansion is
\[
 2K(t)\,\dd t=\left\{{\ii\over\pi v}
       +{h_n\over2\pi\sqrt n}+O(v/n)\right\}\dd v,                \tag{4.43p}
\]
with the singular sign reversed for the conjugate kernel.  Equation
(4.43o) makes the replacement by the Gaussian-skewness expression an
$o(n^{-1/2})$ error uniformly in $x$, because the omitted spatial phase
has modulus one.  Normal inversion and
\[
 {1\over\pi}\int_0^\infty e^{-v^2/2}\cos(vx)\,\dd v=\phi(x),
 \quad
 {1\over\pi}\int_0^\infty v^2e^{-v^2/2}\cos(vx)\,\dd v
 =(1-x^2)\phi(x)
\]
give (4.43j)--(4.43k) on the zero cell.

Here is the promised uniform partition of the rest of the physical
$t$-interval.  Fix $0<\delta<1/4$ so small that the zero-frequency
estimate used above holds on $I_0=[0,\delta]$ and the curvature interval
from Lemma~\ref{orig:lem:spectral} contains
$2\pi I_1/h_n$, where $I_1=[1-\delta,1]$, for all sufficiently large
$n$.  On $I_m=[\delta,1-\delta]$ one has, uniformly in $n$,
\[
 \sup_{t\in I_m}|\varphi_n(2\pi t/h_n)|\le q<1.       \tag{4.43q0}
\]
Indeed $X_n\Rightarrow Q_*$ and $h_n\to s_*$ give uniform convergence
on this compact interval, whereas
\[
 |\varphi_{Q_*}(2\pi t/s_*)|^2
 =1-2p_*q_*\{1-\cos(2\pi t)\}<1
 \quad(\delta\le t\le1-\delta).
\]
The same $q$ works for both kernels, so the normalized middle
contribution is $O(\sqrt n q^n)$, uniformly in the spatial phase.

For $I_1$, let $u_n$ be the unique critical spectral maximum from
Lemma~\ref{orig:lem:spectral} and put
$\tau_n=h_nu_n/(2\pi)\to1$.  The curvature estimate (4.43h1), the affine
change $u=2\pi t/h_n$, and $R_n(u_n)\le1$ give fixed $c_1>0$ such that
\[
 |\varphi_n(2\pi t/h_n)|^n
 \le e^{-c_1n(t-\tau_n)^2},\qquad t\in I_1.           \tag{4.43q1}
\]
Because $|K(1-y)|+|\overline K(1-y)|\le Cy$, either normalized endpoint
is at most
\[
 C\sqrt n\int_0^\delta y e^{-c_1n(y-e_n)^2}\,\dd y
 \le C\{n^{-1/2}+|e_n|\}=o(1),\qquad e_n=1-\tau_n.   \tag{4.43q}
\]
For the inequality use $y\le|y-e_n|+|e_n|$, extend the two nonnegative
Gaussian integrals to the real line, and evaluate them.  Notice that
$e_n$ may have either sign.  This argument uses no residual-power
conclusion and no estimate on the selected height other than
$|\varphi_n(u_n)|\le1$.

The three cells prove the two uniform formulas.  Their limiting plus
profile is maximized at zero and equals
$\phi_0(s_*/2+\gamma_*/6)/\rho_*=\CE$.  The reflected profile is also
maximized at zero because $s_*>\gamma_*$ and is smaller by
$\phi_0\gamma_*/(3\rho_*)$.
\end{proof}

\begin{lemma}[uniform marked lattice local limit]\label{orig:lem:markedllt-full}
For every row let $(J_{n,r},\eta_{n,r})$, $1\le r\le n$, be independent
copies of $(J_n,\eta_n)$, where $J_n$ is integer valued.  Put
\[
 m_n=\E J_n,\qquad Y_n=J_n-m_n,\qquad V_n=\E Y_n^2,
 \qquad K_n=\sum_{r=1}^nJ_{n,r},\qquad
 T_n=\sum_{r=1}^n\eta_{n,r}.
\]
Assume, with constants independent of the row, that
\[
 \E\eta_n=\E(Y_n\eta_n)=0,\qquad
 B_n:=\E\eta_n^2\le {C_B\over n},
 \qquad 0<v_-\le V_n\le v_+<\infty,                 \tag{4.48u0}
\]
that two adjacent values of $J_n$ each have probability at least $c_0>0$,
and that, for some $\bar d,C_0<\infty$ and some $d_n\in[0,\bar d]$,
\[
 \E\{|J_n|^r;|J_n|>H\}\le {C_0d_n\over H^{3-r}},
 \qquad r=0,1,2,\quad H\ge2.                         \tag{4.48u1}
\]
For fixed $T,M,X<\infty$, let
\[
 \mathscr G(T,M)=\left\{g\in C_c^2(\mathbb R):
 \operatorname {supp}g\subset[-T,T],\quad
 \max_{0\le r\le2}\|g^{(r)}\|_\infty\le M\right\}.
\]
Then, uniformly over all row laws satisfying the displayed bounds, all
$g_n\in\mathscr G(T,M)$, and all $k_n\in\mathbb Z$ for which
\[
 x_n={k_n-nm_n\over\sqrt n},\qquad |x_n|\le X,
\]
one has
\[
 \sqrt n\,\E\{g_n(T_n)1_{\{K_n=k_n\}}\}
 ={e^{-x_n^2/(2V_n)}\over\sqrt{2\pi V_n}}\,
      \E g_n(T_n)+o(1).                              \tag{4.48}
\]
In particular, the error in \emph{(4.48)} is one common error for the
whole $C^2$ class and for every bounded centering $x_n$; it does not depend
on a limiting law for $T_n$.
\end{lemma}

\begin{proof}
We give all estimates uniformly over the class in the statement.  First
the tail assumption controls the translation from $J_n$ to $Y_n$.  Indeed,
taking $H=2$ in (4.48u1) gives
\[
 \E|J_n|\le2+\E\{|J_n|;|J_n|>2\}\le2+C_0\bar d/4,
\]
so $|m_n|\le M_0$ for a fixed $M_0$.  Consequently there are fixed
$H_0,C_1$ such that, for $r=0,1,2$ and $H\ge H_0$,
\[
 \E\{|Y_n|^r;|Y_n|>H\}
 \le {C_1d_n\over H^{3-r}}.                          \tag{4.48u3}
\]
For if $H\ge2M_0$, the event $|Y_n|>H$ implies $|J_n|>H/2$, and
$|Y_n|^r\le C_r(|J_n|^r+M_0^r)$.  Applying (4.48u1) at $H/2$ proves
(4.48u3).  In particular, the family $Y_n^2$ is uniformly integrable:
\[
 \sup_n\E\{Y_n^2;|Y_n|>H\}\le {C_1\bar d\over H}
 \longrightarrow0.                                  \tag{4.48u4}
\]

Fix temporarily $U<\infty$ and, for $|u|\le U$, define
\[
 f_{n,u}=\E e^{\ii u\eta_n},\qquad
 c_{1,n}(u)=\E\{Y_ne^{\ii u\eta_n}\},\qquad
 c_{2,n}(u)=\E\{(Y_n^2-V_n)e^{\ii u\eta_n}\}.
\]
We first prove the two uniform decorrelation estimates
\[
 \sup_{|u|\le U}|c_{1,n}(u)|=o(n^{-1/2}),
 \qquad
 \sup_{|u|\le U}|c_{2,n}(u)|=o(1).                   \tag{4.48u5}
\]
Let $C_H=\{|Y_n|\le H\}$.  The two normal equations in (4.48u0) give the
exact identity
\begin{align*}
 c_{1,n}(u)={}&\E\{Y_n(e^{\ii u\eta_n}-1-\ii u\eta_n);C_H\}\\
 &+\E\{Y_n(e^{\ii u\eta_n}-1);C_H^c\}
       -\ii u\E\{Y_n\eta_n;C_H^c\}.
                                                               \tag{4.48u6}
\end{align*}
The inequalities $|e^{\ii y}-1-\ii y|\le y^2/2$ and
$|e^{\ii y}-1|\le2$, (4.48u3), and Cauchy--Schwarz bound the three terms by
\[
 C_UHB_n,\qquad {C_Ud_n\over H^2},\qquad
 C_U\left({B_nd_n\over H}\right)^{1/2}.
\]
Since
\[
 (B_nd_n/H)^{1/2}
 \le\tfrac12(HB_n+d_n/H^2),
\]
we have
\[
 \sup_{|u|\le U}|c_{1,n}(u)|
 \le C_U\{HB_n+d_n/H^2\}.                            \tag{4.48u7}
\]
If $nd_n\le H_0^3$, use $H=H_0$; then the right side is $O(n^{-1})$.
If $nd_n>H_0^3$, use $H=(nd_n)^{1/3}$; after multiplication by $\sqrt n$
both terms are bounded by
$C_Ud_n^{1/3}n^{-1/6}\le C_U\bar d^{1/3}n^{-1/6}$.  This proves the first
part of (4.48u5).

For the second part, use $\E(Y_n^2-V_n)=0$ to write
\[
 c_{2,n}(u)=\E\{(Y_n^2-V_n)(e^{\ii u\eta_n}-1)\}.
\]
On $C_H$, Cauchy--Schwarz and
$Y_n^4 1_{C_H}\le H^2Y_n^2$ give
\[
 \E\{|Y_n^2-V_n|\,|\eta_n|;C_H\}\le CH\sqrt{B_n}.
\]
On $C_H^c$, (4.48u3) and $V_n\le v_+$ give
\[
 \E\{|Y_n^2-V_n|;C_H^c\}\le C d_n/H.
\]
Thus
\[
 \sup_{|u|\le U}|c_{2,n}(u)|
 \le C_U\{H\sqrt{B_n}+d_n/H\}.                      \tag{4.48u8}
\]
Taking $H=\max(H_0,n^{1/8})$ proves the second part of (4.48u5).  Finally,
because $\E\eta_n=0$,
\[
 f_{n,u}-1=\E(e^{\ii u\eta_n}-1-\ii u\eta_n)=O_U(B_n)=O_U(n^{-1}),
                                                               \tag{4.48u9}
\]
so $f_{n,u}$ is bounded away from zero, uniformly on $|u|\le U$.

Put
\[
 \psi_n(t,u)=\E e^{\ii tY_n+\ii u\eta_n}.
\]
For $q=v/\sqrt n$, Taylor's formula in the first coordinate gives the
exact expansion
\[
 \psi_n(q,u)=f_{n,u}+\ii q c_{1,n}(u)
 -{q^2\over2}\{V_nf_{n,u}+c_{2,n}(u)\}+R_n(q,u),     \tag{4.48u10}
\]
where
\[
 |R_n(q,u)|\le Cq^2\E\{Y_n^2\min(|qY_n|,1)\}.        \tag{4.48u11}
\]
The last expectation tends to zero uniformly for bounded $v$: for fixed
$H$ its part on $|Y_n|\le H$ is at most
$|v|H\E Y_n^2/\sqrt n$, while its complement tends uniformly to zero by
(4.48u4).  Divide (4.48u10) by $f_{n,u}$ and use
(4.48u5), (4.48u9), and (4.48u11).  Uniformly on every compact
$|u|\le U$, $|v|\le V_0$,
\[
 n\log{\psi_n(v/\sqrt n,u)\over f_{n,u}}
 =-{V_nv^2\over2}+o(1).
\]
Since $|f_{n,u}|^n\le1$, exponentiation yields
\[
 \psi_n(v/\sqrt n,u)^n
 =f_{n,u}^n e^{-V_nv^2/2}+o(1)                       \tag{4.48u12}
\]
with one uniform error on each such compact set.

We next prove an integrable bound valid outside compact $v$-sets.  This is
the step for which convergence (4.48u12) alone is insufficient.  Let
$(Y_n',\eta_n')$ be an independent copy.  Then
\[
 |\psi_n(t,u)|^2
 =\E e^{\ii t(Y_n-Y_n')+\ii u(\eta_n-\eta_n')}.
\]
Taylor expansion in $t$ gives
\begin{align*}
 |\psi_n(t,u)|^2
 ={}&|f_{n,u}|^2
 +\ii t\E\{(Y_n-Y_n')e^{\ii u(\eta_n-\eta_n')}\}\\
 &-{t^2\over2}\E\{(Y_n-Y_n')^2
                e^{\ii u(\eta_n-\eta_n')}\}+\mathcal R_n(t,u),
                                                               \tag{4.48u13}
\end{align*}
where
\[
 |\mathcal R_n(t,u)|
 \le Ct^2\E\{(Y_n-Y_n')^2
                    \min(|t(Y_n-Y_n')|,1)\}.          \tag{4.48u14}
\]
Put $r_{1,n}(U)=\sup_{|u|\le U}|c_{1,n}(u)|$.  The coefficient of $t$ in
(4.48u13) has absolute value at most $2r_{1,n}(U)$.  The coefficient inside
the quadratic expectation is
\begin{align*}
 &2\Re\{(V_nf_{n,u}+c_{2,n}(u))\overline{f_{n,u}}\}
       -2|c_{1,n}(u)|^2
   =2V_n+o(1),                                        \tag{4.48u15}
\end{align*}
uniformly for $|u|\le U$.  The uniform integrability needed for the
difference is a direct consequence of (4.48u3)--(4.48u4): with
$D_n=Y_n-Y_n'$,
\[
 \E\{D_n^2;|D_n|>H\}
 \le4\E\{Y_n^2;|Y_n|>H/2\}
      +4V_n\Pp\{|Y_n|>H/2\}\longrightarrow0           \tag{4.48u15b}
\]
uniformly in the row.  Hence one can first choose a fixed
$\delta>0$ such that the right side of (4.48u14) is at most
$v_-t^2/8$ whenever $|t|\le\delta$, uniformly in the row, and then choose
$n_U$ so large that the quadratic expectation in (4.48u15) is at least
$3v_-/2$ for $n\ge n_U$ and $|u|\le U$.  Therefore (4.48u13) gives
\[
 |\psi_n(t,u)|^2
 \le1+2|t|r_{1,n}(U)-{3v_-\over4}t^2+{v_-\over8}t^2.
                                                               \tag{4.48u15a}
\]
Young's inequality, in the explicit form
\[
 2|t|r_{1,n}(U)\le {v_-\over8}t^2+{8\over v_-}r_{1,n}(U)^2,
\]
now gives
\[
 |\psi_n(t,u)|^2
 \le1-{v_-\over2}t^2+{8\over v_-}r_{1,n}(U)^2,
 \qquad |t|\le\delta,\quad |u|\le U,\quad n\ge n_U.  \tag{4.48u16}
\]
The right side of (4.48u16) is nonnegative because it bounds the
nonnegative left side.  Thus $\log y\le y-1$, applied after raising the
squared modulus to the power $n/2$, yields
\[
 |\psi_n(t,u)|^n
 \le\exp\left\{-{v_-\over4}nt^2+{4\over v_-}n
                         r_{1,n}(U)^2\right\}.         \tag{4.48u16a}
\]
By (4.48u5), $nr_{1,n}(U)^2=o(1)$.  For the indices $n<n_U$, the trivial
bound $|\psi_n(t,u)|\le1$ is dominated on $|t|\le\delta$ by the same
Gaussian after choosing, for example,
$C_U\ge\exp\{v_-n_U\delta^2/4\}$.  Thus, for every row,
\[
 |\psi_n(t,u)|^n\le C_Ue^{-v_-nt^2/4},
 \qquad |t|\le\delta,\quad |u|\le U.                 \tag{4.48u17}
\]

It remains to obtain a common gap away from the origin.  If $j_n,j_n+1$
have masses at least $c_0$, the independent-copy identity gives, for every
$t\in[-\pi,\pi]$,
\[
 1-|\E e^{\ii tY_n}|^2
 =\E\{1-\cos(t(Y_n-Y_n'))\}
 \ge2c_0^2(1-\cos t).                                \tag{4.48u18}
\]
Moreover
\[
 |\psi_n(t,u)-\E e^{\ii tY_n}|
 \le |u|\E|\eta_n|\le U\sqrt{B_n}.                  \tag{4.48u19}
\]
For the fixed $\delta$ in (4.48u17), (4.48u18)--(4.48u19) imply that
there are $q_U<1$ and $n_U$ such that
\[
 |\psi_n(t,u)|\le q_U,\qquad
 \delta\le|t|\le\pi,\quad |u|\le U,\quad n\ge n_U.   \tag{4.48u20}
\]
Thus (4.48u12), (4.48u17), and (4.48u20) give respectively compact
convergence and domination in the three ranges
\[
 |v|\le V_0,\qquad
 V_0<|v|\le\delta\sqrt n,\qquad
 \delta\sqrt n<|v|\le\pi\sqrt n.                   \tag{4.48u21}
\]

Fourier inversion on the unit lattice, followed by $v=t\sqrt n$, is the
exact identity
\[
 \sqrt n\,\E\{e^{\ii uT_n}1_{\{K_n=k_n\}}\}
 ={1\over2\pi}\int_{-\pi\sqrt n}^{\pi\sqrt n}
 e^{-\ii vx_n}\psi_n(v/\sqrt n,u)^n\,\dd v.          \tag{4.48u22}
\]
On the first range in (4.48u21), use (4.48u12).  On the second use the
integrable Gaussian in (4.48u17); on the third, the integral is at most
$2\pi\sqrt n q_U^n$.  Let $n\to\infty$, then $V_0\to\infty$.  The bounds
are independent of $|x_n|\le X$, and
\[
 {1\over2\pi}\int_{\mathbb R}e^{-\ii vx_n}e^{-V_nv^2/2}\,\dd v
 ={e^{-x_n^2/(2V_n)}\over\sqrt{2\pi V_n}}.
\]
We have therefore proved, uniformly for $|u|\le U$ and $|x_n|\le X$,
\[
 \sqrt n\,\E\{e^{\ii uT_n}1_{\{K_n=k_n\}}\}
 ={e^{-x_n^2/(2V_n)}\over\sqrt{2\pi V_n}}f_{n,u}^n+o(1).
                                                               \tag{4.48u23}
\]

We finally pass from exponential marks to the whole class
$\mathscr G(T,M)$.  Put
\[
 \widehat g(u)=\int_{\mathbb R}e^{-\ii ut}g(t)\,\dd t.
\]
Two integrations by parts give the uniform estimate
\[
 |\widehat g(u)|\le C_{T,M}(1+u^2)^{-1},
 \qquad g\in\mathscr G(T,M).                          \tag{4.48u24}
\]
The unmarked case $u=0$ of (4.48u17)--(4.48u20), inserted into ordinary
lattice inversion, also gives
\[
 \sup_{n,k}\sqrt n\,\Pp\{K_n=k\}\le C.              \tag{4.48u25}
\]
Since
\[
 |\E\{e^{\ii uT_n}1_{\{K_n=k_n\}}\}|\le\Pp\{K_n=k_n\},
 \qquad |f_{n,u}|^n\le1,
\]
the contribution of $|u|>U$ to the Fourier inversion of either side of
(4.48) is bounded by
\[
 C\int_{|u|>U}(1+u^2)^{-1}\,\dd u,                   \tag{4.48u26}
\]
uniformly in $n,k_n,g_n$.  On $|u|\le U$ use (4.48u23), integrate the
uniform error against (4.48u24), and note the exact identity
$f_{n,u}^n=\E e^{\ii uT_n}$.  First let $n\to\infty$ and then
$U\to\infty$ in (4.48u26).  This proves (4.48).
\end{proof}

\begin{lemma}[smooth cardinal gap]\label{orig:lem:cardgap}
Suppose that, for a triangular array,
\[
 X=a+hJ+\varepsilon,\qquad J\in\mathbb Z,\qquad
 \E\varepsilon=\E(J\varepsilon)=0,\qquad
 A:=\E\varepsilon^2=O(n^{-1}),\qquad d=d_n\to0,       \tag{4.44}
\]
where $h$ stays in a compact subinterval of $(0,\infty)$, two adjacent
values of $J$ have probabilities bounded below, and
\[
 \E\{|J|^k;|J|>H\}\le {Cd\over H^{3-k}}
 \quad(k=0,1,2;\ H\ge2).                               \tag{4.45}
\]
Put
\[
 K_n=\sum_{j=1}^nJ_j,\qquad
 T_n={1\over h}\sum_{j=1}^n\varepsilon_j,\qquad
 z_n=\E e^{2\pi\ii\varepsilon/h}.
\]
If $V_n=\operatorname {Var}J\to V\in(0,\infty)$,
$|k_n-n\E J|\le X_0\sqrt n$, and
$|z_n|^n\le1-\kappa$ for some fixed $\kappa>0$, then, uniformly for
$|\theta_n|\le1/2$,
\[
 \liminf_n\sqrt n\,
 \E e_B(K_n-k_n+T_n-\theta_n)>0,                       \tag{4.46}
\]
where $e_B$ is the excess of the one-sided Beurling majorant over the
indicator of the closed negative half-line.
\end{lemma}

\begin{proof}
Put
\[
 Y=J-\E J,\qquad \eta={\varepsilon\over h},\qquad
 B=\E\eta^2.
\]
The normal equations in (4.44) give
\[
 \E\eta=\E(Y\eta)=0,\qquad
 B={A\over h^2}=O(n^{-1}).
\]
After discarding finitely many rows, \(V/2\le V_n\le2V\), while
\(d_n\) is bounded.  Thus (4.45), the two adjacent masses, and the compact
range of \(h\) verify every hypothesis of
Lemma~\ref{orig:lem:markedllt-full}.  In particular, for fixed \(T,M<\infty\),
let
\[
 \mathscr G(T,M)=\left\{g\in C_c^2(\mathbb R):
 \operatorname {supp}g\subset[-T,T],\quad
 \max_{0\le j\le2}\|g^{(j)}\|_\infty\le M\right\}.     \tag{4.47b}
\]
For \(x_n=(k_n-n\E J)/\sqrt n\), \(|x_n|\le X_0\), equation (4.48)
therefore holds with one \(o(1)\), uniformly over
\(g_n\in\mathscr G(T,M)\), the phase, the bounded center, and all row
parameters occurring here.  No limiting distribution for \(T_n\) is
assumed.

The elementary cardinal subkernel is
\[
 e_B(x)\ge {\sin^2(\pi x)\over4\pi^2(1+x^2)}.         \tag{4.52}
\]
Here are the partial-fraction details.  Inserting the paired expansion of
$\pi\cot\pi x$ in the Beurling formula gives, for $x>0$,
\begin{align}
 e_B(x)&=\left({\sin\pi x\over\pi}\right)^2
 \left\{\sum_{j=0}^\infty{1\over(x+j)^2}-{1\over x}\right\},
                                                               \tag{4.52a}\\
 e_B(-x)&=\left({\sin\pi x\over\pi}\right)^2
 \left\{{1\over x}-\sum_{j=1}^\infty{1\over(x+j)^2}\right\}.
                                                               \tag{4.52b}
\end{align}
Let $q(t)=(x+t)^{-2}$.  Since $q$ is decreasing and convex, the
trapezoidal inequality and telescoping give
\begin{align*}
 \sum_{j=0}^\infty q(j)-\int_0^\infty q(t)\,\dd t
 &\ge {1\over2}\sum_{j=0}^\infty\{q(j)-q(j+1)\}
 ={1\over2x^2}.
\end{align*}
For the other sign, the tangent inequality at the right endpoint of each
unit interval gives
\begin{align*}
 \int_{j-1}^{j}q(t)\,\dd t-q(j)&\ge-{q'(j)\over2}
 ={1\over(x+j)^3},\\
 {1\over x}-\sum_{j=1}^\infty{1\over(x+j)^2}
 &\ge\sum_{j=1}^\infty{1\over(x+j)^3}
 \ge\int_1^\infty{\dd t\over(x+t)^3}
 ={1\over2(x+1)^2}.
\end{align*}
Now $1/(2x^2)\ge1/\{4(1+x^2)\}$ and
$1/\{2(x+1)^2\}\ge1/\{4(1+x^2)\}$.  Multiplication by
$(\sin\pi x/\pi)^2$ proves (4.52) for both signs.  At nonzero integers it
follows by continuity.  At $x=0$ the right side is zero and
$e_B(0)=M(0)-1=0$ by (P3), so no continuity of
the positive-side auxiliary expression at zero is being assumed.  Now
\[
 \E\sin^2\{\pi(T_n-\theta_n)\}
 ={1-\Re(e^{-2\pi\ii\theta_n}z_n^n)\over2}
 \ge {\kappa\over2}.                                  \tag{4.53}
\]
Moreover $\E T_n^2=nA/h^2\le C_T$ for a fixed $C_T$.  Choose \(R\) so
large that \(C_T/R^2\le\kappa/4\), and choose one fixed
\(\chi\in C_c^\infty(\mathbb R)\) with \(0\le\chi\le1\), equal to one on
\([-R,R]\), and supported in \([-R-1,R+1]\).  Put

\[
 g_n(t)=\chi(t)\sin^2\{\pi(t-\theta_n)\}.
\]
Equations (4.53) and Chebyshev give
\(\E g_n(T_n)\ge\kappa/4\), uniformly in the phase.  This family belongs
to one class \(\mathscr G(R+1,M)\).  On its support, (4.52) gives the
explicit phase-uniform bound
\[
 e_B(t-\theta_n)\ge c_Rg_n(t),\qquad
 c_R={1\over4\pi^2\{1+(R+3/2)^2\}}.                  \tag{4.54}
\]
Since $e_B\ge0$, retain only the event $K_n=k_n$ and apply the uniform
version of (4.48) to $g_n$.  If
\(x_n=(k_n-n\E J)/\sqrt n\), then
\[
 \begin{split}
 \sqrt n\,\E e_B(K_n-k_n+T_n-\theta_n)
 &\ge c_R\sqrt n\,
       \E\{g_n(T_n)1_{\{K_n=k_n\}}\}\\
 &=c_R{e^{-x_n^2/(2V_n)}\over\sqrt{2\pi V_n}}
       \E g_n(T_n)+o(1).
 \end{split}                                           \tag{4.54a}
\]
The Gaussian factor has a positive lower bound because
\(|x_n|\le X_0\) and \(V_n\to V\in(0,\infty)\); the expectation is at
least \(\kappa/4\).  This proves (4.46), uniformly in the phase, without
an infinite-cell interchange.
\end{proof}

\begin{corollary}[residual-power selection]\label{orig:cor:respower}
In addition to the hypotheses of Lemma~\ref{orig:lem:cardgap}, suppose
\((x_n)\) is an active near-extremal positive-branch sequence in the sense
fixed at the start of this subsection, and coordinate its thresholds by
\[
 x_n\sqrt n=na_n+h_n(k_n+\theta_n),\qquad |\theta_n|\le\tfrac12.
\]
Suppose also that $\sup_n\rho_n<\infty$ and that the adaptive majorant
 $\mathcal U^+_{X_n,n,h_n}$ satisfies
\[
 {\sqrt n\over\rho_n}\sup_x
       \{\mathcal U^+_{X_n,n,h_n}(x)-\Phi(x)\}
 \le\CE+o(1),                                         \tag{4.47}
\]
while the reflected branch has a fixed gap.  Then $|z_n|^n\to1$.
\end{corollary}

\begin{proof}
Write
\[
 S_n={\sqrt n\over\rho_n}\{F_{n,X_n}(x_n)-\Phi(x_n)\},\qquad
 M_n={\sqrt n\over\rho_n}
       \{\mathcal U^+_{X_n,n,h_n}(x_n)-\Phi(x_n)\}.
\]
The upper identity in (4.24w) and the exact threshold coordinate give,
with no inequality,
\[
 S_n=M_n-{\sqrt n\over\rho_n}
       \E e_B(K_n-k_n+T_n-\theta_n).                  \tag{4.55}
\]
By near extremality, \(S_n\to\CE\), whereas (4.47) gives
\(M_n\le\CE+o(1)\).  If \(|z_n|^n\not\to1\), then
\(0\le |z_n|^n\le1\) supplies a subsequence and a fixed \(\kappa>0\) on
which \(|z_n|^n\le1-\kappa\).  Lemma~\ref{orig:lem:cardgap} gives a fixed
\(c_\kappa>0\) such that
\[
 \liminf_n\sqrt n\,
 \E e_B(K_n-k_n+T_n-\theta_n)\ge c_\kappa.
\]
Choose \(R_\rho<\infty\) with \(\rho_n\le R_\rho\).  Then (4.55) would imply
\(\limsup S_n\le\CE-c_\kappa/R_\rho\), a contradiction.  Hence
\(|z_n|^n\to1\).  The reflected fixed gap is used only one step upstream,
to ensure that the near-extremal absolute score has the positive
orientation; no lower cardinal estimate is hidden in (4.55).
\end{proof}

\begin{lemma}[active-point exact-lattice envelope]\label{orig:lem:anchored}
Let \(V_n\) be centered variance-one laws supported on maximal lattices
\(a_n+b_n\mathbb Z\), where \(b_n\) stays in a compact subset of
\((0,\infty)\).  Put
\(\gamma_n=\E V_n^3\), \(\rho_n=\E|V_n|^3\), let $p_n$ be the matching
Bernoulli parameter, and set $d_n=\rho_n-\rho_2(\gamma_n)$.  Assume that
two adjacent sites have probabilities at least \(c_0>0\),
\[
 p_n\to p_*,\qquad d_n\to0,\qquad
 0<c\le\sqrt n\,d_n\le C.
\]
If \(x_n\to0\), then, uniformly in the lattice phase of \(x_n\),
\[
 {\sqrt n\over\rho_n}
 \{\mathcal U^+_{V_n,n,b_n}(x_n)-\Phi(x_n)\}
 \le {\phi_0\over\rho_n}\left({b_n\over2}+{\gamma_n\over6}\right)
 +\{\Lambda_{\rm abs,*}+o(1)\}d_n+O(n^{-1}).          \tag{4.58}
\]
Consequently
\[
 {\sqrt n\over\rho_n}
 \{\mathcal U^+_{V_n,n,b_n}(x_n)-\Phi(x_n)\}
 \le\CE-\{G_{\rm span,*}-\Lambda_{\rm abs,*}-o(1)\}d_n
       +O(n^{-1})<\CE.                                \tag{4.58a}
\]
The reflected operator has a fixed leading gap.
\end{lemma}

\begin{proof}
All constants in this proof are uniform after the row index is sufficiently
large.  We suppress that index and write
\(\gamma=\gamma_n\), \(\rho=\rho_n\), \(d=d_n\), and \(b=b_n\).
The assumptions put \(\gamma\) in a fixed compact neighbourhood
\(\Gamma\) of \(\gamma_*\), put \(b\) in a compact subset of
\((0,\infty)\), and give
\[
             d\asymp n^{-1/2},\qquad
             \rho=\rho_2(\gamma)+d\asymp1.            \tag{4.58b0}
\]

\emph{The actual span.}
Let \(w_\gamma\) be the nonnegative defect density from the proof of
Theorem~\ref{orig:thm:defect}.  If an atom \(v\) of \(V_n\) has mass at least
\(c_0\), then
\[
       0\le w_\gamma(v)\le c_0^{-1}\E w_\gamma(V_n)
                         =c_0^{-1}d.                 \tag{4.58b1}
\]
Uniformly for \(\gamma\in\Gamma\), the density \(w_\gamma\) is coercive
and its only zeros are
\(-a_\gamma\) and \(b_\gamma\), whose distance is
\(s_\gamma(\gamma)\).  Thus each of the two adjacent atoms in the
hypothesis has distance \(o(1)\) from this two-point zero set.  The
distance between the two atoms is the lattice span \(b_n\), which is
bounded away from zero; consequently the two atoms cannot converge to
the same zero.  Their order then forces the lower atom to converge to
\(-a_{\gamma_n}\) and the upper atom to \(b_{\gamma_n}\).  Hence
\[
                 b_n-s_\gamma(\gamma_n)\longrightarrow0,
        \qquad b_n\longrightarrow s_*.               \tag{4.58b}
\]

\emph{Exact folding.}
Choose the unique \(k\in\mathbb Z\) and
\(\theta\in[-1/2,1/2)\) such that
\[
 x={na_n+b_n(k+\theta)\over\sqrt n},
\]
put, for \(0\le q\le1/2\),
\[
 f(q)=\E e^{2\pi\ii qV_n/b_n},\qquad
 g(q)=e^{-2\pi\ii qx\sqrt n/b_n}f(q)^n.
\]
Since \(V_n\in a_n+b_n\mathbb Z\),
\(f(1-q)=e^{2\pi\ii a_n/b_n}\overline{f(q)}\).  Substitution of the
displayed coordinate for \(x\) therefore gives, with no approximation,
\[
              g(1-q)=e^{-2\pi\ii\theta}\overline{g(q)}.
\]
It follows that
\[
 \mathcal U^+_{V_n,n,b_n}(x)={1\over2}+2\Re\int_0^{1/2}
 \{K(q)g(q)+K(1-q)e^{-2\pi\ii\theta}\overline{g(q)}\}\,\dd q.
                                                               \tag{4.58c}
\]
In (4.58c) replace \(f(q)\) by
\(\chi_\gamma(2\pi q/b)\), and in the returned term replace
\(\overline{f(q)}\) by the corresponding conjugate.  Denote the resulting
analytic expression by \(\mathscr S_n^{\rm lat}(x,\theta)\).  This
definition uses the actual interval and the actual folding phase; it does
not assert that the matching two-point law has span \(b\).

\emph{A uniform second-order expansion of the surrogate.}
The support of \(Q_\gamma\) is uniformly bounded for
\(\gamma\in\Gamma\).  Hence there are \(\delta,C>0\) such that the
principal logarithm of \(\chi_\gamma(t)\) is defined for
\(|t|\le\delta\) and, uniformly in \(\gamma\in\Gamma\),
\[
 \log\chi_\gamma(t)=-{t^2\over2}-{\ii\gamma t^3\over6}
       +{\gamma^2-2\over24}t^4+r_\gamma(t),
       \qquad |r_\gamma(t)|\le C|t|^5.               \tag{4.58c0}
\]
Indeed, Taylor's formula for \(\chi_\gamma\), followed by Taylor's
formula for \(\log(1+z)\), is uniform on this compact family, and
\(\E Q_\gamma^4=\gamma^2+1\).  For \(t=w/\sqrt n\), put
\[
 A=-{\ii\gamma w^3\over6},\qquad
 B={\gamma^2-2\over24}w^4,\qquad
 E=n r_\gamma(w/\sqrt n).
\]
Then (4.58c0) says
\[
 n\log\chi_\gamma(w/\sqrt n)
 =-{w^2\over2}+{A\over\sqrt n}+{B\over n}+E,
 \qquad
 |A|\le Cw^3,\quad |B|\le Cw^4,\quad
 |E|\le {Cw^5\over n^{3/2}}.                        \tag{4.58c0a}
\]
For all sufficiently large \(n\), on \(0\le w\le n^{1/16}\), the sum
\(|A|/\sqrt n+|B|/n+|E|\) tends to zero uniformly.  Taylor's formula
for the exponential, with its integral remainder, therefore gives
\[
 \left|e^{A/\sqrt n+B/n+E}
 -1-{A\over\sqrt n}-{B+A^2/2\over n}\right|
 \le C\left\{{w^5+w^7+w^9\over n^{3/2}}
                  +{w^8\over n^2}\right\}.          \tag{4.58c0b}
\]
Here the four displayed monomials respectively dominate \(E\), the
product \(AB/n^{3/2}\), the cubic term \(A^3/n^{3/2}\), and
\(B^2/n^2\); every product containing \(E\), and every higher power, is
smaller on the stated window.  Since
\(B+A^2/2=R_{\rm even,\gamma}(w)\), (4.58c0a)--(4.58c0b) give
\[
 \chi_\gamma(w/\sqrt n)^n=e^{-w^2/2}
 \left[1-{\ii\gamma w^3\over6\sqrt n}
       +{R_{\rm even,\gamma}(w)\over n}
       +\mathcal R_{n,\gamma}(w)\right],             \tag{4.58c1}
\]
where
\[
 R_{\rm even,\gamma}(w)={\gamma^2-2\over24}w^4
                    -{\gamma^2\over72}w^6,           \tag{4.58c2}
\]
and, for example,
\[
 |\mathcal R_{n,\gamma}(w)|
       \le {C\over n^{3/2}}w^5(1+w)^{10}.             \tag{4.58c2a}
\]
Indeed (4.58c2a) follows term by term from (4.58c0b), using
\(n^{-2}w^8\le n^{-3/2}w^5(1+w)^{10}\).
The factor \(w^5\) in (4.58c2a) is recorded because the direct kernel is
singular at zero.  Notice also that the conjugate power has the same real
even term and the opposite cubic term.

We next justify extending the root-scale integrals to infinity.  By
symmetrisation there are \(c,\delta_1>0\), uniform in
\(\gamma\in\Gamma\), such that
\(|\chi_\gamma(u)|\le e^{-cu^2}\) for \(|u|\le\delta_1\).
Shrink \(\delta_1\), if necessary, so that
\(s_-\delta_1<\pi/4\), where
\(s_-:=\inf_{\gamma\in\Gamma}s_\gamma(\gamma)>0\).
By (4.58b), after removing finitely many rows,
\(3/4\le b/s_\gamma(\gamma)\le5/4\).  The explicit identity
\[
 |\chi_\gamma(u)|^2
 =1-2p(1-p)\{1-\cos(s_\gamma(\gamma)u)\}
\]
has \(p(1-p)\) bounded away from zero on \(\Gamma\).  On the rest of
the moving interval,
\[
 \delta_1\le u\le{\pi\over b}
 \quad\Longrightarrow\quad
 s_\gamma(\gamma)u\in[s_-\delta_1,4\pi/3].           \tag{4.58c2d}
\]
The compact interval on the right is disjoint from \(2\pi\mathbb Z\).
The explicit identity therefore gives one common modulus gap on
\(\delta_1\le u\le\pi/b\).  Consequently, under
\(u=w/\sqrt n\),
\[
 |\chi_\gamma(w/\sqrt n)|^n
       \le C e^{-cw^2}\quad(0\le w\le\delta_1\sqrt n),
\]
and it is at most \(C e^{-cn}\) beyond that interval.  In particular the
part \(w\ge n^{1/16}\), including the factor \(w^{-1}\) from the
singular kernel, is \(O(e^{-c n^{1/8}})\).  This proves the required tail
claim on the moving actual-span interval rather than on a fixed
surrogate interval.

Set \(q=bw/(2\pi\sqrt n)\).  Taylor expansion of the displayed Prawitz
kernel, with uniform remainders on \(w\le n^{1/16}\), gives
\[
 \begin{split}
 2K(q)\,\dd q={}&\left\{{\ii\over\pi w}
 +{b\over2\pi\sqrt n}-{b^2w\over4\pi^2n}
 -{\ii b^2w\over12\pi n}
 +O\!\left({w^2(1+w)\over n^{3/2}}\right)\right\}\dd w,\\
 2K(1-q)\,\dd q={}&\left\{{b^2w\over4\pi^2n}
 +O\!\left({w^2(1+w)\over n^{3/2}}\right)\right\}\dd w.
 \end{split}                                           \tag{4.58c2b}
\]
The following identities fix all signs in the calculation:
\[
 {1\over2}+{1\over\pi}\int_0^\infty
       {\sin(xw)\over w}e^{-w^2/2}\dd w=\Phi(x),
\]
\[
 {1\over\pi}\int_0^\infty e^{-w^2/2}\cos(xw)\dd w=\phi(x),
 \qquad
 {1\over\pi}\int_0^\infty w^2e^{-w^2/2}\cos(xw)\dd w
       =(1-x^2)\phi(x).                              \tag{4.58c2c}
\]
The first identity follows by differentiating in \(x\) and checking
\(x=0\); the other two are the Gaussian Fourier transform and its second
derivative.

Insert (4.58c1)--(4.58c2b), retain every product through raw order
\(n^{-1}\), and keep the spatial factors \(e^{-\ii xw}\) and
\(e^{\ii xw}\) until real parts are taken.  The complete direct ledger is
\[
\begin{array}{c|c|c}
\text{kernel factor}&\text{power factor}&\text{direct real contribution}
\\ \hline
\ii/(\pi w)&1&\Phi(x)-1/2\\
b/(2\pi\sqrt n)&1&b\phi(x)/(2\sqrt n)\\
\ii/(\pi w)&-\ii\gamma w^3/(6\sqrt n)&
       \gamma(1-x^2)\phi(x)/(6\sqrt n)\\ \hline
\ii/(\pi w)&R_{\rm even,\gamma}(w)/n&O(|x|/n)\\
b/(2\pi\sqrt n)&-\ii\gamma w^3/(6\sqrt n)&O(|x|/n)\\
-\ii b^2w/(12\pi n)&1&O(|x|/n)\\
-b^2w/(4\pi^2n)&1&-b^2/(4\pi^2n)+O(x^2/n).
\end{array}                                                     \tag{4.58c3}
\]
For each of the three \(O(|x|/n)\) entries the coefficient before
\(e^{-\ii xw}\) is purely imaginary and odd; its real part therefore
contains \(\sin(xw)\), and
\(|\sin(xw)|\le |x|w\) proves the stated bound by Gaussian integration.
For the last entry use
\(|1-\cos(xw)|\le x^2w^2/2\) and
\(\int_0^\infty we^{-w^2/2}\dd w=1\).

The returned kernel has neither a singular term nor a term of order
\(n^{-1/2}\).  Its only contribution of raw order \(n^{-1}\) is
\[
 {b^2\over4\pi^2n}\Re\left\{e^{-2\pi\ii\theta}
   \int_0^\infty we^{-w^2/2}e^{\ii xw}\,\dd w\right\}
 ={b^2\over4\pi^2n}\cos(2\pi\theta)+O(|x|/n),       \tag{4.58c4}
\]
uniformly for \(|\theta|\le1/2\).  Products not displayed in
(4.58c3)--(4.58c4) are integrable, by (4.58c2a), against the singular
kernel and are bounded by \(Cn^{-3/2}\) after integration.  The already
estimated tail contributes exponentially less.  We have therefore proved,
with a constant independent of \(|x|\le1\) and
\(|\theta|\le1/2\),
\[
 {\sqrt n\over\rho}\{\mathscr S_n^{\rm lat}(x,\theta)-\Phi(x)\}
 ={\phi(x)\over\rho}\left\{{b\over2}
             +{\gamma\over6}(1-x^2)\right\}
 +{b^2\over4\pi^2\rho\sqrt n}\{\cos(2\pi\theta)-1\}
 +O\left({|x|\over\sqrt n}+{1\over n}\right).       \tag{4.58d}
\]

\emph{Uniform comparison of the actual and surrogate powers.}
Let \(p=p(V_n)\), \(s=s_\gamma(\gamma)=s(p)\),
\(\varphi=\varphi_{V_n}\), \(\chi=\chi_\gamma\), and
\(\Delta(u)=\varphi(u)-\chi(u)\).  In the Bernoulli coordinate
\(Y=p+V_n/s\), cubic Hermite interpolation at \(0,1\), together with the
matching moments of \(Y\) and the Bernoulli law, gives the exact identity
\[
 \Delta(u)=e^{-\ii spu}d\int {E_{su}(y)\over W_p(y)}\,\dd\nu(y),
 \qquad \dd\nu(y)={W_p(y)\,\dd\Pp_Y(y)\over d}.       \tag{4.58e}
\]
Here \(\nu\) is a probability measure because
\(\E W_p(Y)=d\).

We record all uniformity needed below.  Every convex combination
\[
 \psi_a(u)=(1-a)\chi(u)+a\varphi(u),\qquad0\le a\le1,
\]
is the characteristic function of a mixture of centered variance-one
laws.  Their third absolute moments are uniformly bounded, since
\(p\to p_*\) and \(d\to0\).  Applying
\[
 1-|\E e^{\ii uZ}|^2
   =\E\{1-\cos(u(Z-Z'))\},\qquad
 1-\cos v\ge v^2/2-|v|^3/6,
\]
to these mixtures gives fixed \(c,\delta_0>0\) such that
\[
             |\psi_a(u)|\le e^{-cu^2}
       \quad(|u|\le\delta_0,\ 0\le a\le1).           \tag{4.58e1}
\]
Shrink \(\delta_0\), if necessary, so that
\(b\delta_0/(2\pi)<1/2\) for every \(b\) in the fixed span compactum.
Thus the exact divided power
\[
 \mathscr P_n(u):={\varphi(u)^n-\chi(u)^n
                    \over n\{\varphi(u)-\chi(u)\}}
 =\int_0^1\psi_a(u)^{\,n-1}\dd a                    \tag{4.58e2}
\]
(with its continuous value if the denominator vanishes) satisfies
\[
 |\mathscr P_n(w/\sqrt n)|\le Ce^{-cw^2}
       \quad(0\le w\le\delta_0\sqrt n).              \tag{4.58e3}
\]
Moreover, uniformly for \(w\) in each fixed compact interval and
\(a\in[0,1]\),
\[
 \psi_a(w/\sqrt n)=1-{w^2\over2n}+O(w^3n^{-3/2}),
\]
so (4.58e2) converges there uniformly to \(e^{-w^2/2}\).  No
linearisation in \(d\) has been made.

Put \(H=n^{1/8}\).  Equations (4.21)--(4.21a), already proved uniformly
for \(p\) in a compact subset of \((0,1)\), say
\[
 |E_z(y)-g(zy)|\le C|z|^4(1+|y|^3),\quad
 g(v)=e^{\ii v}-1-\ii v+v^2/2+\ii v^3/6,
\]
\[
 |E_z(y)|\le {|z|^4\over24}y^2(1-y)^2,\qquad
 {y^2(1-y)^2\over W_p(y)}\le C(1+|y|).
                                                               \tag{4.58e4}
\]
On \(|y|\le H\), the second line of (4.58e4) bounds the integral by
\(C|z|^4(H+1)\).  On \(|y|>H\), use the first line,
\(|g(v)|\le|v|^3/6\), and the explicit cubic tails of \(W_p\):
\[
 { |y|^3\over W_p(y)}\le
 \max\{C_+(p)^{-1},C_-(p)^{-1}\}\{1+O(H^{-1})\}.    \tag{4.58e5}
\]
Since \(\nu\) is a probability measure, these two estimates yield,
uniformly for \(0\le w\le\delta_0\sqrt n\),
\[
 {n^{3/2}|\Delta(w/\sqrt n)|\over d}
 \le {s^3\over6}\max\{C_+(p)^{-1},C_-(p)^{-1}\}
       \{1+O(H^{-1})\}w^3
       +C(H+1){w^4\over\sqrt n}.                    \tag{4.58f}
\]

We now insert (4.58f) into the folded operator, rather than merely invoke
a pointwise \(o(1)\).  On the direct root-scale cell, the singular part
of (4.58c2b), (4.58e2), and (4.58e3) give
\[
 \begin{split}
 {\sqrt n\over\rho}|\hbox{direct singular defect}|
 &\le {d\over\pi\rho}\int_0^{\delta_0\sqrt n}
 {|\mathscr P_n(w/\sqrt n)|\over w}
       {n^{3/2}|\Delta(w/\sqrt n)|\over d}\,\dd w\\
 &\le {d s^3\over6\pi\rho}
       \max\{C_+(p)^{-1},C_-(p)^{-1}\}
       \int_0^\infty w^2e^{-w^2/2}\dd w+o(d)\\
 &=\{\Lambda_{\rm abs}(p,d)+o(1)\}d
  =\{\Lambda_{\rm abs,*}+o(1)\}d.
 \end{split}                                           \tag{4.58g}
\]
 To justify the second line with its sharp Gaussian coefficient, first
 restrict to \(0\le w\le M\).  There (4.58f) and the uniform compact
 convergence following (4.58e3) give the displayed limsup with
 \(e^{-w^2/2}\).  On \(M\le w\le\delta_0\sqrt n\), (4.58e3) and
 (4.58f) give the row-independent majorant
 \(C(w^2+w^3)e^{-cw^2}\); let \(M\to\infty\).  The coefficient errors
 in (4.58f) are \(O(H^{-1})\) and \(O((H+1)/\sqrt n)\), both tending to
 zero.  The
identity
\(\pi^{-1}\int_0^\infty w^2e^{-w^2/2}\dd w=\phi_0\)
then gives exactly the definition (4.18a) of \(\Lambda_{\rm abs}\).

For completeness, all nonsingular pieces are smaller on the asserted
scale.  Since \((H+1)/\sqrt n\le1\) eventually, (4.58f) and
(4.58e3) imply
\[
 |\varphi(w/\sqrt n)^n-\chi(w/\sqrt n)^n|
 \le {Cd\over\sqrt n}(w^3+w^4)e^{-cw^2}.             \tag{4.58g1}
\]
The bounded part of the direct kernel is \(O(n^{-1/2})\dd w\), so its
normalized contribution is \(O(d/\sqrt n)\).  On the returned copy,
\(|2K(1-q)\dd q|\le Cw\dd w/n\), and hence
\[
 {\sqrt n\over\rho}|\hbox{returned defect}|
 \le {Cd\over n}\int_0^\infty
       w(w^3+w^4)e^{-cw^2}\dd w=O(d/n).              \tag{4.58g2}
\]

It remains to cover, uniformly, the complement of the root-scale cell.
Choose a fixed
\(0<u_0<\min\{\delta_0,\pi/(2s_*)\}\).  For the actual law,
the two adjacent atoms of mass at least \(c_0\) give
\[
 1-|\varphi(u)|^2
 =\E\{1-\cos(u(V_n-V_n'))\}
 \ge2c_0^2\{1-\cos(bu)\}.                            \tag{4.58g3}
\]
For \(u_0\le u\le\pi/b\), (4.58b) implies, for all sufficiently large
rows,
\[
             bu\in[s_*u_0/2,\pi],\qquad
             s(p)u\in[s_*u_0/2,3\pi/2].              \tag{4.58g4}
\]
Both intervals in (4.58g4) are compact and avoid \(2\pi\mathbb Z\).
Hence (4.58g3) gives a fixed positive lower bound for
\(1-|\varphi(u)|^2\).  The explicit identity
\[
 1-|\chi(u)|^2=2p(1-p)\{1-\cos(s(p)u)\},
\]
together with \(p\to p_*\in(0,1)\), gives another fixed positive lower
bound for \(1-|\chi(u)|^2\).  Thus both powers are at most \(r^n\), for
one row-independent \(r<1\), on the remainder of
\(0\le q\le1/2\).  The kernel is bounded there, so the normalized
contribution is \(O(\sqrt n r^n)=o(d)\) by (4.58b0).  Equations
(4.58g)--(4.58g4) therefore compare the actual folded operator with the
surrogate on the whole actual fundamental interval.

Now take the sequence \(x=x_n\to0\).  The cosine term in (4.58d) is
nonpositive.  For all sufficiently large rows the positive leading
profile is at most its value at zero.  Furthermore
\(|x_n|/\sqrt n=o(d)\), and the nonsingular comparison error
\(d/\sqrt n\) is \(O(n^{-1})\), by (4.58b0).  Combining
(4.58d)--(4.58g4) gives
\[
 {\sqrt n\over\rho_n}
 \{\mathcal U^+_{V_n,n,b_n}(x_n)-\Phi(x_n)\}
 \le {\phi_0\over\rho_n}\left({b_n\over2}
                       +{\gamma_n\over6}\right)
       +\{\Lambda_{\rm abs,*}+o(1)\}d_n+O(n^{-1}),
\]
uniformly in \(\theta\).

\emph{The signed denominator loss, with the current \(p\) frozen.}
If \(V'\) is an independent copy of \(V\), then
\(V-V'\in b\mathbb Z\), and hence
\[
 2b=b\E(V-V')^2\le\E|V-V'|^3
       \le2\{\rho+\E|V|\}.
\]
The last inequality follows by expanding
\(|x-y|^3\le(x-y)^2(|x|+|y|)\) and using centering.  Proposition
\ref{orig:prop:m1} now gives
\[
                       b\le s(p)+K_1(\gamma(p))d.     \tag{4.58h}
\]
Define, at the current value of \(p\),
\[
 \begin{gathered}
 N_0(p)=\phi_0\{s(p)/2+\gamma(p)/6\},\qquad
 N_1(p)=\phi_0K_1(\gamma(p))/2,\\
 D_0(p)=\rho_2(p),\qquad A(p)=N_0(p)/D_0(p).
 \end{gathered}
\]
Then (4.58h) and the exact algebraic identity
\[
 {N_0(p)+N_1(p)d\over D_0(p)+d}
 =A(p)-{A(p)-N_1(p)\over D_0(p)+d}\,d               \tag{4.58i}
\]
show, without any joint expansion in \((p,d)\), that
\[
 {\phi_0\over\rho}\left({b\over2}+{\gamma\over6}\right)
 \le A(p)-\{G_{\rm span,*}+o(1)\}d
 \le \CE-\{G_{\rm span,*}+o(1)\}d.                 \tag{4.58j}
\]
Indeed \(A(p)\le\CE\) by Lemma~\ref{orig:lem:tangentexp}, while continuity
and \(p\to p_*\), \(d\to0\), give
\[
 {A(p)-N_1(p)\over D_0(p)+d}
 \longrightarrow {\CE-\phi_0K_1(\gamma_*)/2\over\rho_*}
 =G_{\rm span,*}.
\]
Thus (4.58i) freezes \(p\) and expands only in \(d\); in particular no
unproved relation such as \(|p-p_*|=o(d)\) is used.  Combining (4.58)
and (4.58j), and then using (4.24), gives
\[
 {\sqrt n\over\rho_n}
 \{\mathcal U^+_{V_n,n,b_n}(x_n)-\Phi(x_n)\}
 \le\CE-\{G_{\rm span,*}-\Lambda_{\rm abs,*}-o(1)\}d_n
       +O(n^{-1})<\CE,
\]
because \(n^{-1}=o(d_n)\).

Finally, the reflected assertion is not merely a change of notation.
Directly from (4.24u), conjugation gives
\[
 \mathcal U^+_{-V,n,b}(-x)=1-\mathcal U^-_{V,n,b}(x),
 \qquad
 \Phi(x)-\mathcal U^-_{V,n,b}(x)
 =\mathcal U^+_{-V,n,b}(-x)-\Phi(-x).                \tag{4.58k}
\]
The reflected law has the same \(b,\rho,d\), skewness \(-\gamma\), and
the same adjacent-atom and comparison bounds.  Repeating the already
proved surrogate calculation on the reflected compact parameter set (the
defect term is \(O(d)=o(1)\)) yields, uniformly in phase,
\[
 {\sqrt n\over\rho_n}
 \{\Phi(x_n)-\mathcal U^-_{V_n,n,b_n}(x_n)\}
 ={\phi_0\over\rho_*}\left({s_*\over2}-{\gamma_*\over6}\right)
       +o(1)
 =\CE-{\phi_0\gamma_*\over3\rho_*}+o(1).             \tag{4.58l}
\]
Since \(\gamma_*>0\), this is bounded above by \(\CE\) minus a fixed
positive constant for all sufficiently large \(n\), which proves the
claimed reflected gap.
\end{proof}

\begin{lemma}[near-lattice transfer]\label{orig:lem:latticetransfer}
Let $X_n$ be centered, variance-one laws, let $p_n$ be their matching
Bernoulli parameters, and put
$d_n=\E|X_n|^3-\rho_2(\E X_n^3)$.  Assume $p_n\to p_*$, $d_n\to0$, and
\[
                 0<\liminf\sqrt n\,d_n
                 \le\limsup\sqrt n\,d_n<\infty.       \tag{4.56}
\]
Let \((x_n)\) be an active near-extremal positive-branch sequence.  If
$\liminf|\varphi_{X_n}(u_n)|^n>0$ at the selected critical resonance,
then there is a centered,
variance-one maximal-span exact lattice projection $L_n^\circ$, physical
spans $h_n$ before and $b_n$ after standardization, and points $y_n\to0$
such that, for both one-sided Prawitz
operators,
\[
 \sqrt n\,|\mathcal U^\pm_{X_n,n,h_n}(x_n)
       -\mathcal U^\pm_{L_n^\circ,n,b_n}(y_n)|=o(d_n),
 \qquad \rho(L_n^\circ)=\rho(X_n)+o(d_n),
 \qquad d(L_n^\circ)=d_n+o(d_n).                       \tag{4.57}
\]
Consequently the positive active score satisfies
\[
 {\sqrt n\over\rho(X_n)}
 \{\mathcal U^+_{X_n,n,h_n}(x_n)-\Phi(x_n)\}
 \le\CE-\{G_{\rm span,*}-\Lambda_{\rm abs,*}-o(1)\}d_n
               +O(n^{-1}).                            \tag{4.57a}
\]
\end{lemma}

\begin{proof}
We suppress row subscripts until the last step.  Let $u_0$ be the unique
critical spectral maximum from Lemma~\ref{orig:lem:spectral}, write
\[
 \varphi_X(u_0)=r e^{\ii\alpha},\qquad h_0={2\pi\over u_0},
\]
and choose $a_0=\alpha/u_0$ modulo $h_0$.  Let $J$ be the nearest integer
to $(X-a_0)/h_0$ and put
\[
 X=a_0+h_0J+\varepsilon_0,\qquad |\varepsilon_0|\le h_0/2.
\]
This phase convention gives
\[
 \E e^{\ii u_0\varepsilon_0}=r.
\]
Since $1-\cos(2\pi v)\ge8v^2$ for $|v|\le1/2$, the assumed lower bound
for $r^n$ implies
\[
 A_0:=\E\varepsilon_0^2
 \le {h_0^2\over8}(1-r)=O(n^{-1}).                    \tag{4.59}
\]
We now derive, rather than assume, the adjacency and tail assertions used
below.  Uniform convergence of $R_n'$ in Lemma~\ref{orig:lem:spectral} puts its
unique critical point at
$u_0-2\pi/s_\gamma(\gamma)=o(1)$; hence
\[
                         h_0={2\pi\over u_0}\longrightarrow s_* .   \tag{4.59a0}
\]
Let $-a_\gamma,b_\gamma$ be the two zeros of $w_\gamma$, so that
$a_\gamma+b_\gamma=s_\gamma(\gamma)$, and set
\[
 \delta_n=(d+A_0)^{1/4}.
\]
Corollary~\ref{orig:cor:transport} and Chebyshev's inequality give
\[
 \Pp\{\operatorname {dist}(X,\{-a_\gamma,b_\gamma\})>\delta_n\}=o(1),
 \qquad \Pp\{|\varepsilon_0|>\delta_n\}=o(1).          \tag{4.59a1}
\]
If $Z$ is a nearest one of these two roots and
$\pi_n=\Pp\{Z=b_\gamma\}$, then
$\E|X-Z|=O(\sqrt d)$ and centering gives
\[
 0=\E X=s_\gamma(\gamma)\pi_n-a_\gamma+O(\sqrt d),
 \qquad \pi_n={a_\gamma\over s_\gamma(\gamma)}+O(\sqrt d). \tag{4.59a2}
\]
Thus both root clusters have masses bounded below.

Intersect each cluster with the two good events in (4.59a1).  On either
intersection every lattice site $a_0+h_0J=X-\varepsilon_0$ lies within
$2\delta_n$ of the corresponding root.  Since $h_0$ is bounded below,
two different labels cannot occur in one such cluster once
$4\delta_n<h_0$.  Call the resulting labels $j_-$ and $j_+$.  They have
uniformly positive masses, and
\[
 h_0(j_+-j_-)=a_\gamma+b_\gamma+O(\delta_n)
              =s_\gamma(\gamma)+o(1).                 \tag{4.59a3}
\]
The integer $j_+-j_-$ therefore tends to one and equals one for all
sufficiently large rows.  Translate $J$ so that these labels are $0,1$;
this only changes $a_0$ by an integral multiple of $h_0$ and preserves the
phase convention.

Finally $h_0$ is in a fixed compact span range, $|\varepsilon_0|\le h_0/2$,
and the labels $0,1$ are the only cells containing the two defect roots.
The explicit cubic tails of $w_\gamma$ consequently give, uniformly in
the row,
\[
 |J|^3 1_{\{|J|>H\}}\le Cw_\gamma(X)1_{\{|J|>H\}}
 \qquad(H\ge2).
\]
Since $\E w_\gamma(X)=d$, multiplication by
$|J|^{k-3}\le H^{k-3}$ proves, simultaneously for every $H\ge2$,
\[
 \E\{|J|^k;|J|>H\}\le {Cd\over H^{3-k}},
 \qquad k=0,1,2.                                      \tag{4.59a}
\]

We give the spectral-to-least-squares bridge explicitly.  Since the phase
choice gives \(\E e^{\ii u_0\varepsilon_0}=r\in\mathbb R\), Taylor's
theorem gives
\[
 |\E\varepsilon_0|\le CA_0.                           \tag{4.59b}
\]
Stationarity of $|\varphi_X|^2$ at $u_0$ gives
$\Im\E\{Xe^{\ii u_0\varepsilon_0}\}=0$.  We record the quantitative
argument rather than hide it in this stationarity statement.  Since
$u_0$ is bounded above and below,
\[
 u_0\E(X\varepsilon_0)
 =\E X\{u_0\varepsilon_0-\sin(u_0\varepsilon_0)\}.    \tag{4.59b1}
\]
On $C_H=\{|J|\le H\}$, one has $|X|\le C(1+H)$ and
$|y-\sin y|\le Cy^2$, whence the right side of (4.59b1), restricted to
$C_H$, is at most $CHA_0$.  On $C_H^c$, use
$|y-\sin y|\le C|y|$, Cauchy--Schwarz, and (4.59a):
\[
\begin{split}
 \E\{|X\varepsilon_0|;C_H^c\}
 &\le C\{\E[(1+J^2);C_H^c]\}^{1/2}A_0^{1/2}\\
 &\le C(A_0d/H)^{1/2}
 \le C\{HA_0+d/H^2\}.
\end{split}                                           \tag{4.59b2}
\]
All constants here are uniform for $H\ge2$ and for every row under
consideration.  Equations (4.59b1)--(4.59b2) prove
\[
 |\E(X\varepsilon_0)|
 \le C\{HA_0+d/H^2+A_0\}.                            \tag{4.59c}
\]
If $A_0=0$ the bridge is immediate.  Otherwise take $H=A_0^{-1/3}$,
replaced by $2$ for the finitely many rows for which it is smaller.
Then $HA_0=A_0^{2/3}$ and $d/H^2=dA_0^{2/3}$.
Because $d$ is bounded, (4.59b)--(4.59c) and
$X=a_0+h_0J+\varepsilon_0$ imply
\[
 |\E(J\varepsilon_0)|+|\E\varepsilon_0|
       =O(A_0^{2/3}).                                  \tag{4.59d}
\]
The two core labels give $\operatorname {Var}J\ge c>0$.  Hence the affine
least-squares projection
\[
 L=a+hJ,\qquad \varepsilon=X-L
\]
satisfies
\[
 \E\varepsilon=\E(J\varepsilon)=0,\qquad
 A:=\E\varepsilon^2\le A_0,\qquad
 |a-a_0|+|h-h_0|=O(A_0^{2/3}).                        \tag{4.60a0}
\]
Put $u_L=2\pi/h$ and $z=\E e^{\ii u_L\varepsilon}$.  Exact lattice
phase cancellation gives $|z|=|\varphi_X(u_L)|$.  On the selected window
$|\varphi_X|$ has a uniformly bounded second derivative, while its first
derivative vanishes at $u_0$.  Thus
\[
 n\big||z|-r\big|
 \le Cn|u_L-u_0|^2
 =O(nA_0^{4/3})=o(1).                                  \tag{4.60}
\]
This is the required bridge: the selected spectral power and the
least-squares residual power have the same subsequential limit.

Since $u_0\to2\pi/s_*$ and (4.60a0) gives $h-h_0=o(1)$, one has
$h\to s_*$.  Lemma~\ref{orig:lem:adaptiveupstream} now gives exactly the
upper-majorant hypothesis (4.47) in Corollary~\ref{orig:cor:respower}, without
using residual power or $A=o(n^{-1})$; it also gives the reflected fixed
gap.  The middle has the adjacent-mass Cram\'er gap.  Moreover
$A=O(n^{-1})$, so
$\operatorname {Var}J=(1-A)/h^2\to s_*^{-2}\in(0,\infty)$, while outer
localization gives $x_n=o(1)$.  Define the integer $k_n$ and
$\theta_n\in[-1/2,1/2)$ by the exact threshold identity
\[
 x_n\sqrt n=na+h(k_n+\theta_n).                        \tag{4.60a00}
\]
The first normal equation gives $a=-h\E J$, and hence
\[
 {k_n-n\E J\over\sqrt n}
 ={x_n\over h}-{\theta_n\over\sqrt n}\longrightarrow0. \tag{4.60a01}
\]
Moreover, with $K_n=\sum J_i$ and
$T_n=h^{-1}\sum\varepsilon_i$, the closed-half-line convention is exactly
\[
 \sum_{i=1}^nX_i\le x_n\sqrt n
 \quad\Longleftrightarrow\quad
 K_n-k_n+T_n-\theta_n\le0.                             \tag{4.60a02}
\]
Thus the cardinal-gap hypotheses are now all explicit: the normal
equations are (4.60a0), $A=O(n^{-1})$, the variance has the displayed
positive limit, (4.59a) supplies adjacent masses and uniform tails,
(4.60a01) supplies the centering, (4.60a02) fixes the sign and index, and
\(\rho_n=\rho_2(p_n)+d_n=O(1)\) follows from \(p_n\to p_*\) and
\(d_n\to0\).
Consequently every hypothesis of Corollary~\ref{orig:cor:respower} is
satisfied.  Since
the actual score tends to $\CE$, that corollary yields
\[
                              |z|^n\longrightarrow1.  \tag{4.60a1}
\]
Only \(A=O(n^{-1})\) was used to obtain (4.60a1).  The stronger residual
estimate \(A=o(n^{-1})\) is deduced below from (4.60a1), so the argument is
not circular.
Choose $\beta\in[-1/2,1/2)$ so that
$e^{-2\pi\ii\beta}z=|z|$, let $N$ be a nearest integer to
$\varepsilon/h-\beta$, and write exactly
\[
 X=(a+h\beta)+h(J+N)+\widetilde\varepsilon,\qquad
 \widetilde\varepsilon=\varepsilon-h(\beta+N).       \tag{4.60a2}
\]
The chord-to-coset inequality and (4.60a1) give
\[
 \E\widetilde\varepsilon^2
 \le {h^2\over8}\{1-|z|\}=o(n^{-1}).                 \tag{4.60a}
\]
Here there is no hidden compactness step.  By the definition of $N$,
$v:=\varepsilon/h-\beta-N$ lies in $[-1/2,1/2]$, and therefore
\[
\begin{split}
 1-|z|
 &=\E\{1-\cos(2\pi(\varepsilon/h-\beta))\}\\
 &=\E\{1-\cos(2\pi v)\}
 \ge8\E v^2={8\over h^2}\E\widetilde\varepsilon^2.
\end{split}                                           \tag{4.60a3}
\]
Moreover $|z|^n\to1$ implies $n(1-|z|)\to0$: otherwise a subsequence
with $n(1-|z|)\ge c>0$ would satisfy
$|z|^n\le e^{-n(1-|z|)}\le e^{-c}<1$.  Thus (4.60a) follows from
(4.60a3), with the explicit constant $1/8$.
Take the affine least-squares projection once more onto $J+N$ and rename
it $L=a+hJ$, with residual $\varepsilon$.  Least squares cannot increase
the last squared distance, so
\[
 \E\varepsilon=\E(J\varepsilon)=\E(L\varepsilon)=0,
 \qquad A=\E\varepsilon^2=o(n^{-1}).                  \tag{4.61}
\]

We verify here that the lattice hypotheses survive both the reindexing and
this second projection.  Before the renaming, let $j_0,j_0+1$ be the two
core labels constructed immediately before (4.59a), put $J'=J+N$, and let
$\varepsilon^{(1)}$ denote
the first least-squares residual in (4.60a0).  Since
$\E\varepsilon^{(1)}=0$, $u_L=2\pi/h$ stays bounded, and
$A_1:=\E(\varepsilon^{(1)})^2=O(n^{-1})$,
\[
 1-\Re z\le {u_L^2A_1\over2},\qquad
 |\Im z|=|\E(\sin(u_L\varepsilon^{(1)})
                    -u_L\varepsilon^{(1)})|
       \le C u_L^2A_1.                                \tag{4.61a}
\]
Thus $\arg z=O(A_1)$ and the representative chosen in (4.60a2) satisfies
$\beta=O(A_1)$.  If $N\ne0$, the nearest-integer definition implies
$|\varepsilon^{(1)}/h-\beta|\ge1/2$.  Since $h$ is bounded away from
zero, for all sufficiently large $n$ this entails
$|\varepsilon^{(1)}|\ge c>0$; hence
\[
 \Pp\{N\ne0\}\le C\E(\varepsilon^{(1)})^2=O(n^{-1}). \tag{4.61b}
\]
If the two old core masses are at least $c_0$, then
\[
 \Pp\{J'=j_0+r\}\ge
 \Pp\{J=j_0+r,N=0\}\ge c_0-\Pp\{N\ne0\}
 \ge {c_0\over2},\qquad r=0,1.                        \tag{4.61c}
\]
In particular $\operatorname {Var}J'$ is bounded below.  Write the second projection,
before renaming, as $L=a'+h'J'$.  From the preliminary representation
(4.60a2) and the normal equations,
\[
 |h'-h|\le {\|\widetilde\varepsilon\|_2\over
                    \sqrt{\operatorname {Var}J'}}=o(n^{-1/2}).       \tag{4.61d}
\]
Consequently $h'>0$ and remains in the same compact span range.  Projection
changes neither $J'$ nor its probabilities, so (4.61c) still holds for the
renamed lattice law.  Its support is contained in $a'+h'\mathbb Z$ and
contains both $a'+h'j_0$ and $a'+h'(j_0+1)$.  Any lattice containing the
whole support must therefore have span at most $h'$, whereas $h'$ itself is
a span.  Hence $h'$ is the maximal span.  Centering and standardizing later
in (4.68a) merely replace it by $h'/\sigma$ and preserve both adjacent
masses and maximality.

We next prove the full operator transfer.  Put
$\alpha_0=(\E\widetilde\varepsilon^2)^{1/2}=o(n^{-1/2})$ for the
preliminary reindexed residual in (4.60a2).  The normal equations, the two
core labels, and (4.59a) give the pointwise least-squares geometry as
follows.  Before the second projection write
$X=c+\bar hJ+\widetilde\varepsilon$, choosing $c$ modulo $\bar h$ so
$|c|\le C$.  The span $\bar h$ is bounded above and below and
$|\widetilde\varepsilon|\le\bar h/2$.  Centering gives $|\E J|\le C$,
while the adjacent core labels give $\operatorname {Var}J\ge c_0$.  The least-squares
normal equations therefore imply
\[
 |h-\bar h|={|\operatorname {Cov}(J,\widetilde\varepsilon)|
              \over\operatorname {Var}J}
 \le C\alpha_0,\qquad
 |a-c|\le|\E\widetilde\varepsilon|
       +|h-\bar h|\,|\E J|\le C\alpha_0.
\]
Moreover $|J|\le C(1+|X|)$.  Hence
\[
 |X-L|\le C\{1+\alpha_0(1+|X|)\}.                    \tag{4.62}
\]
For the exact defect density, uniformly on the critical compactum,
\[
 w_\gamma(x)\ge c|x|^3\quad(|x|\ge R_0),\qquad
 \E w_\gamma(X)=d.                                    \tag{4.63}
\]
Consequently, for $R\ge R_0$,
$\E\{|X|^k;|X|>R\}\le Cd/R^{3-k}$, $k=0,1,2$.  From
\[
 \big||x|^3-|l|^3\big|+|x^3-l^3|
 \le C(x^2+l^2)|x-l|
\]
and (4.62)--(4.63), we now record the split which also fixes the
dependence on the preliminary residual $\alpha _0$.  Write
$C_R=\{|X|\le R\}$.  On $C_R$, $L=X-(X-L)$ and (4.62) give
\[
 |L|\le C(1+R),\qquad |X-L|\le C(1+\alpha _0R),
\]
and hence, using $\E|X-L|\le\sqrt A$,
\[
 \E\{(X^2+L^2)|X-L|;C_R\}
 \le C\{R^2\sqrt A+(1+\alpha _0R)A\}.                \tag{4.63a}
\]
Indeed, with $e=X-L$, the elementary bound
$X^2+L^2\le C(X^2+e^2)$ makes the two terms on the left at most
$CR^2\E|e|$ and $\E|e|^3$.  The first is at most $CR^2\sqrt A$, while
(4.62) gives
$\E\{|e|^3;C_R\}\le\sup_{C_R}|e|\,\E e^2
\le C(1+\alpha _0R)A$.  Thus (4.63a) uses no fourth moment.
On $C_R^c$, after increasing $R_0$ and taking $n$ large, (4.62) also gives
$|L|\le C(1+|X|)$ and
\[
 (X^2+L^2)|X-L|
 \le C\{1+X^2+\alpha _0(1+|X|^3)\}.
\]
The three tail estimates following (4.63), together with
$\E\{|X|^3;|X|>R\}\le Cd$, therefore yield
\[
 \E\{(X^2+L^2)|X-L|;C_R^c\}
      \le C\{d/R+\alpha _0d\}.                       \tag{4.63b}
\]
More explicitly, the part without $\alpha _0$ is controlled by
$\Pp(C_R^c)+\E\{X^2;C_R^c\}\le C(d/R^3+d/R)$, and the
$\alpha _0$ part by
$\alpha _0\E\{1+|X|^3;C_R^c\}\le C\alpha _0d$.
All constants in (4.62)--(4.63b) are independent of $R\ge R_0$ and of
the row once $p$ lies in the fixed critical compactum.
Combining (4.63a)--(4.63b) proves
\[
 |\E X^3-\E L^3|+|\E|X|^3-\E|L|^3|
 \le C\{R^2\sqrt A+(1+\alpha_0R)A+d/R+\alpha_0d\}. \tag{4.64}
\]
The relation needed here is
\[
 \sqrt A=o(n^{-1/2}),\qquad n^{-1/2}=O(d),\qquad
 \sqrt A=o(d);                                         \tag{4.64a}
\]
the previously suggested relation $n^{-1/2}=o(d)$ would be false on the
scale (4.56).  We choose all later free cutoffs simultaneously.  With a vanishing
denominator interpreted as $+\infty$, put
\[
 H_n^{\rm cut}=\min\left\{
 \left({d\over\sqrt A}\right)^{1/4},
 \left({d\over A}\right)^{1/2},
 (nd)^{1/2},\alpha_0^{-1/2}\right\},
 \qquad\text{and below we take }R=H=H_n^{\rm cut}.     \tag{4.64b}
\]
Every entry in the minimum tends to infinity.  Moreover
\[
 {(H_n^{\rm cut})^2\sqrt A\over d}\le
       \sqrt{{\sqrt A\over d}}\to0,
 \quad {H_n^{\rm cut}A\over d}\le\sqrt{{A\over d}}\to0,
 \quad {H_n^{\rm cut}\over nd}\le{1\over\sqrt{nd}}\to0,
 \quad \alpha_0H_n^{\rm cut}\le\sqrt{\alpha_0}\to0.    \tag{4.64c}
\]
It follows at once that
\[
 R^2\sqrt A=o(d),\quad (1+\alpha_0R)A=o(d),\quad
 d/R=o(d),\quad HA=o(d),\quad H=o(nd),                 \tag{4.64d}
\]
as well as $d/H^2=o(d)$ and $\alpha_0d/H=o(d)$.  Thus the
right side of (4.64) is $o(d)$ and the same single cutoff will be used in
(4.67b)--(4.67d).
For later reference, the normalized assertions just used are
\[
 {R^2\sqrt A\over d}\to0,\quad {A\over d}\to0,\quad
 \alpha_0R\to0,\quad {1\over R}\to0,\quad
 {HA\over d}\to0,\quad {H\over nd}\to0.              \tag{4.64e}
\]
They follow respectively from (4.64c), (4.64a), (4.64c), $R\to\infty$,
(4.64c), and (4.64c).  This list is the complete scale input used below;
in particular no step assumes the false relation $n^{-1/2}=o(d)$.

Here is the zero-cell comparison, including its full frequency range.
Set
\[
 G(y)=e^{\ii y}-1-\ii y+{y^2\over2}+{\ii y^3\over6},
 \qquad |G'(y)|\le C\min\{|y|^3,1+y^2\}.             \tag{4.65}
\]
Put $\kappa=2\pi/h$; it ranges in the compact set of scaling factors
generated by the admissible spans.  Put $u=\kappa v/\sqrt n$, and initially
restrict
\[
                         0\le v\le\delta\sqrt n.      \tag{4.65a}
\]
For $|X|\ge R_0$, write $q_X=|X-L|/|X|$ and
$s_X=|X|/\sqrt n$.  Increase $R_0$ once so that
$|L|\le C|X|$ on this set.  The mean-value theorem and (4.65) then give
\[
 |G(uX)-G(uL)|
 \le Cq_X\min\{|u|^4|X|^4, |u||X|+|u|^3|X|^3\}.
\]
Dividing by the cubic lower bound in (4.63), and using
$u=\kappa v/\sqrt n$, proves the first inequality in
\[
 {n^{3/2}|G(uX)-G(uL)|\over w_\gamma(X)}
 \le Cq_X\min\{v^4s_X,v/s_X^2+v^3\}
 \le C(\alpha_0+n^{-1/2})(v^3+v^4).                  \tag{4.65b}
\]
Indeed (4.62) implies
$q_X\le C(\alpha _0+|X|^{-1})$.  If $s_X\le1$, use the first member
of the minimum and
$q_Xs_X\le C(\alpha_0+n^{-1/2})$.  If $s_X\ge1$, then
$q_X\le C(\alpha_0+n^{-1/2})$, while
\[
 \min\{v^4s_X,v/s_X^2+v^3\}
 \le \min\{v^4s_X,v/s_X^2\}+v^3\le2v^3;
\]
the last inequality follows by separating $vs_X\le1$ and
$vs_X\ge1$.  This proves the second inequality in (4.65b).  On
$|X|<R_0$, no division by
$w_\gamma$ is made (so both zeros of the defect are included), and
\[
 \E\{|G(uX)-G(uL)|;|X|<R_0\}
 \le C{v^4\over n^2}\sqrt A.                         \tag{4.65c}
\]
After integrating (4.65b) against $w_\gamma(X)\,\dd\Pp/d$ and adding
(4.65c),
\[
 {n^{3/2}\over d}\E|G(uX)-G(uL)|
 \le C\vartheta_n(v^3+v^4),
 \quad
 \vartheta_n=\alpha_0+n^{-1/2}+{\sqrt A\over\sqrt n\,d}=o(1). \tag{4.66}
\]
This single estimate covers compact, root-scale, mesoscopic, and
super-root atoms.

Let $R(u)=\varphi_X(u)-\varphi_L(u)$ and
$B_3=\E X^3-\E L^3=o(d)$.  Since both variables are centered and
$\E X^2-\E L^2=A$, the identity defining $G$ gives
\[
 R(u)=-{Au^2\over2}-{\ii B_3u^3\over6}
       +\E\{G(uX)-G(uL)\}.                            \tag{4.66a}
\]
Thus, with
$\eta_n=|B_3|/d+\vartheta_n\to0$ and
$P(v)=v^3+v^4$,
\[
 |R(u)|\le C\left\{{Av^2\over n}
              +{d\eta_nP(v)\over n^{3/2}}\right\}.  \tag{4.66b}
\]
The variances of $X,L$ are bounded below and their third absolute moments
are bounded above.  Applying
$1-\cos y\ge y^2/2-|y|^3/6$ to an independent difference gives fixed
$c,\delta>0$ such that
\[
 \max\{|\varphi_X(\kappa t)|,|\varphi_L(\kappa t)|\}^{n-j}
 \le Ce^{-cnt^2},\qquad 0\le t\le\delta,\quad j=0,1,2. \tag{4.66c}
\]
For complex $a,b$,
\[
 |a^n-b^n-nb^{n-1}(a-b)|
 \le {n(n-1)\over2}|a-b|^2\max(|a|,|b|)^{n-2}.       \tag{4.66d}
\]
Equations (4.66b)--(4.66d), only on the range (4.65a), imply
\[
 |\varphi_X(u)^n-\varphi_L(u)^n|
 \le Ce^{-cv^2}\{n|R(u)|+n^2|R(u)|^2\}.              \tag{4.66e}
\]
Finally
\[
 {2\over\sqrt n}|K(v/\sqrt n)|\le {C_\delta\over v},
 \qquad 0<v\le\delta\sqrt n.                         \tag{4.66f}
\]
Write $I_0=[0,\delta]$ and $I_1=[1-\delta,1]$ for the zero and return
cells.  After $t=v/\sqrt n$, the zero-cell integral is bounded by
\begin{align}
 &{\sqrt n\over d}\max_{\sigma\in\{+,-\}}\sup_x
 |\mathcal J^\sigma_{X,n,h;I_0}(x)
       -\mathcal J^\sigma_{L,n,h;I_0}(x)|                         \notag\\
 &\quad\le C\int_0^{\delta\sqrt n}{e^{-cv^2}\over v}
 \Big\{{\sqrt nA\over d}v^2+\eta_nP(v)
       +{\sqrt nA^2\over d}v^4                         \notag\\
 &\hspace{42mm}+A\eta_nv^2P(v)+{d\eta_n^2\over\sqrt n}P(v)^2
 \Big\}\,\dd v=o(1).                                 \tag{4.67}
\end{align}
Every coefficient in the last integral tends to zero:
\[
 {\sqrt nA\over d}=o(1),\qquad \eta_n=o(1),\qquad
 {\sqrt nA^2\over d}=A\,{\sqrt nA\over d}=o(1),
 \quad A\eta_n=o(1),\quad {d\eta_n^2\over\sqrt n}=o(1). \tag{4.67a0}
\]
Here the first relation is exactly $A=o(n^{-1})$ together with
$d\asymp n^{-1/2}$; all Gaussian polynomial integrals are finite.
Thus (4.67) says precisely that the raw zero-cell difference is
$o(d/\sqrt n)$, uniformly in the cdf coordinate and for both signs.
Only the nonnegative majorant in (4.67) is extended to $[0,\infty)$.
The estimate (4.66c) is not asserted beyond the fixed zero cell.  On the
fixed physical middle interval, the adjacent masses give
$|\varphi_L(u)|\le1-c$, while
\[
 |\varphi_X(u)-\varphi_L(u)|\le |u|\E|X-L|\le C\sqrt A=o(1).
\]
Thus both powers are exponentially small there.  The remaining endpoint
is handled below.

For that endpoint put $t_0=2\pi/h$,
\[
 z=\E e^{\ii t_0\varepsilon}=r_0\omega,\qquad
 \lambda=e^{\ii t_0a},\qquad
 M_1=\E\{Xe^{\ii t_0\varepsilon}\}.
\]
Since $e^{\ii t_0L}=\lambda$, the marked first moment below satisfies
\[
 m=\E\{Xe^{\ii t_0X}\}=\lambda M_1,\qquad |m|=|M_1|.   \tag{4.67a00}
\]
Also put
\[
 q_2=\E\{(X^2-1)e^{\ii t_0X}\},\quad
 q_3=\E\{(X^3-\gamma)e^{\ii t_0X}\}.                \tag{4.67a}
\]
Since $\E\varepsilon=0$,
\[
 1-r_0\le1-\Re z\le {t_0^2A\over2},\qquad
 |\arg z|\le C|\Im z|
 \le C\E|\sin(t_0\varepsilon)-t_0\varepsilon|\le CA.
\]
Consequently
\[
 {1-r_0^n+|\omega^n-1|\over\sqrt n}
 \le C\sqrt n A=o(d).                                \tag{4.67a1}
\]
The common unit phase relates $m$ and $M_1$.  Since
$\E\varepsilon=\E(L\varepsilon)=0$, one has
$\E(X\varepsilon)=A$.  Because also $\E X=0$, the exact centered Taylor
identity is
\[
 M_1=\ii t_0A+
      \E\{X(e^{\ii t_0\varepsilon}-1-\ii t_0\varepsilon)\}.       \tag{4.67a2}
\]
On $|X|\le H$, the remainder bound $|e^{\ii y}-1-\ii y|\le Cy^2$
makes the last expectation at most $CHA$.  On $|X|>H$, use instead
$|e^{\ii y}-1-\ii y|\le C|y|$, followed by (4.62), to get
\[
 \E\{|X||\varepsilon|;|X|>H\}
 \le C\{d/H^2+\alpha _0d/H\}.
\]
Thus, for every $H\ge\max\{2,R_0\}$,
\[
 |M_1|\le C\{A+HA+d/H^2+\alpha_0d/H\}.                \tag{4.67b0}
\]
Taking the single common choice $H=H_n^{\rm cut}$ from
(4.64b)--(4.64d) gives
\[
 |M_1|=o(d),\qquad \sqrt n|M_1|^2=o(d),               \tag{4.67b}
\]
the second conclusion using \(\sqrt n\,d=O(1)\).
For completeness we prove the two marked estimates used next without any
moment beyond the third.  Put
\[
 P_2(x)=x^2-\gamma x-1,\qquad R_3(x)=(x+\gamma)P_2(x).
\]
Uniformly for $\gamma$ in the critical compactum, the factorizations
\eqref{orig:eq:wplus}--\eqref{orig:eq:wminus} imply, for $H\ge2$,
\begin{equation}\tag{4.67b1}
\begin{aligned}
 P_2(x)^2 1_{\{|x|\le H\}}&\le CHw_\gamma(x),\\
 |P_2(x)|1_{\{|x|>H\}}&\le {C\over H}w_\gamma(x),\\
 |R_3(x)|&\le C\{w_\gamma(x)^{1/2}+w_\gamma(x)\}.
\end{aligned}
\end{equation}
Indeed the quotients extend through the two zeros of $w_\gamma$, and
their limits at infinity give the displayed powers of $H$.  Since
$z=\E e^{\ii t_0\varepsilon}=r_0\omega$, $|\omega|=1$,
\[
 \E|e^{\ii t_0\varepsilon}-\omega|^2=2(1-r_0).        \tag{4.67b2}
\]
Moreover $\E P_2(X)=0$, while $e^{\ii t_0L}$ is a constant unit phase.
Subtracting $\omega$ inside the expectation, applying Cauchy--Schwarz on
$|X|\le H$, and using the second bound in (4.67b1) on its complement
gives
\[
 |q_2-\gamma m|
 \le C\{\sqrt{Hd(1-r_0)}+d/H\}.                       \tag{4.67c}
\]
Finally
\[
 x^3-\gamma=(\gamma^2+1)x+R_3(x),\qquad \E R_3(X)=0.
\]
One more subtraction of $\omega$, followed by (4.67b1)--(4.67b2), yields
\[
 |q_3-(\gamma^2+1)m|
 \le C\{\sqrt{d(1-r_0)}+d\}.
\]
Since $|m|=|M_1|=o(d)$ by (4.67b), this proves
\[
                     |q_3|\le C\{\sqrt{1-r_0}+d\}.     \tag{4.67c1}
\]
Use the common $H=H_n^{\rm cut}$ from (4.64b)--(4.64d).  Then
(4.67a1)--(4.67c) give
\[
 q_2=o(d),\qquad q_3/\sqrt n=o(d),\qquad
 {1-r_0^n\over\sqrt n}=o(d).                         \tag{4.67d}
\]
Indeed the three normalized estimates, including the terms suppressed in
the preceding sentence, are
\[
\begin{split}
 {|q_2|\over d}
 &\le C\left\{\sqrt{HA/d}+{1\over H}+{|M_1|\over d}\right\}=o(1),\\
 {|q_3|\over\sqrt n\,d}
 &\le C\left\{{\sqrt A\over\sqrt n\,d}+{1\over\sqrt n}
               +{|M_1|\over\sqrt n\,d}\right\}=o(1),\\
 {1-r_0^n+|\omega^n-1|\over\sqrt n\,d}
 &\le C{\sqrt nA\over d}=o(1).
\end{split}                                           \tag{4.67d1}
\]
The first line uses $1-r_0\le CA$ and (4.64e); in the second,
$\sqrt n\,d$ is bounded below and $\sqrt A=o(n^{-1/2})$; the last is
(4.67a1).  This verifies all three $o(d)$ statements on the full scale
$d\asymp n^{-1/2}$.
The same tail comparison makes $|X|^3$ uniformly integrable as $d\to0$.
Taylor's Peano remainder is therefore uniform along the triangular array,
and the marked factorization at $t_0+s$ is
\[
 \varphi_X(t_0+s)-\varphi_X(t_0)\varphi_X(s)
 =\ii s m-{s^2\over2}q_2-{\ii s^3\over6}q_3+o(|s|^3). \tag{4.67e}
\]
Recall that
\(e^{\ii t_0L}=\lambda\), \(\varphi_X(t_0)=\lambda z\), and
\(m=\lambda M_1\).  If
\[
 D_n(s)=\varphi_X(t_0+s)-\varphi_X(t_0)\varphi_X(s),
\]
uniform third-moment Peano expansion gives
\[
 D_n(s)=\ii sm-{s^2q_2\over2}-{\ii s^3q_3\over6}
                     +s^3\epsilon_n(s),\qquad
 \lim_{\delta\downarrow0}\sup_n\sup_{|s|\le\delta}|\epsilon_n(s)|=0. \tag{4.67e1}
\]
For \(\mathscr F_n(s)=\varphi_X(t_0+s)\),
\[
 \mathscr F_n(0)=\lambda z,\qquad \mathscr F_n'(0)=\ii m,\qquad
 \mathscr F_n''(0)=-\mathscr F_n(0)-q_2.
\]
Consequently
\[
 (|\mathscr F_n|^2)''(0)
 =-2r_0^2-2\Re\{q_2\overline{\lambda z}\}+2|m|^2=-2+o(1).          \tag{4.67e2}
\]
Uniform integrability of \(X_n^2\) makes these second derivatives
equicontinuous.  Thus \((|\mathscr F_n|^2)''\le-c\) on a fixed endpoint
cell, and
Taylor's theorem plus completion of the square gives
\[
 |\mathscr F_n(s)|^n\le C\exp\{-cns^2+Cn|m|^2\}
       \le C'e^{-c'ns^2}.                                  \tag{4.67e3}
\]
We now verify all bases which occur in the divided powers.  Write
\[
 \chi(s)=\varphi_X(s),\quad \psi(s)=\varphi_L(s),\quad
 F(s)=\varphi_X(t_0+s),\quad B(s)=\lambda z\chi(s),\quad
 C(s)=\lambda\psi(s),
\]
and put $D=F-B$, $R=\chi-\psi$.  For
\(\alpha,\beta,\kappa\ge0\), \(\alpha+\beta+\kappa=1\), set
\(G_{\alpha\beta\kappa}=\alpha F+\beta B+\kappa C\), and abbreviate this
function by \(G\) in the next two displays.  At the endpoint,
\[
 \begin{split}
 G(0)&=\lambda\{(\alpha+\beta)z+\kappa\},\qquad
 G'(0)=\ii\alpha m,\\
 G''(0)&=-\lambda\{(\alpha+\beta)z+\kappa(1-A)\}-\alpha q_2.
 \end{split}                                                    \tag{4.67e4}
\]
It follows, uniformly on this compact simplex, that
\[
 (|G|^2)'(0)=O(|m|),\qquad (|G|^2)''(0)=-2+o(1).          \tag{4.67e5}
\]
Uniform integrability of $X^2$ and $L^2$ makes these second derivatives
equicontinuous.  Taylor's theorem and completion of the square therefore
give, for $j=0,1,2$,
\[
 |G_{\alpha\beta\kappa}(s)|^{\,n-j}
 \le C\exp\{-cns^2+Cn|m|^2\}\le C'e^{-c'ns^2}.          \tag{4.67e6}
\]
The last inequality uses
\(n|m|^2=o(1)\), which follows from (4.67b) and
\(\sqrt n\,d=O(1)\).  The zero-cell divided power uses
 \(Z_\tau=(1-\tau)\psi+\tau\chi\); here
\[
 Z_\tau(0)=1,\quad Z_\tau'(0)=0,\quad
 Z_\tau''(0)=-\{1-(1-\tau)A\},
\]
so the same Gaussian bound holds uniformly in \(0\le\tau\le1\).

For completeness, these local bounds have a common continuation over the
whole endpoint cell.  If two adjacent lattice sites have probabilities at
least $c_0$, then
\[
 1-|\psi(s)|^2\ge2c_0^2\{1-\cos(hs)\}.                 \tag{4.67e7}
\]
Away from $s=0$ this is a fixed gap.  Exact lattice periodicity and the
residual bound give
\[
 |\chi(s)-\psi(s)|\le |s|\sqrt A,\qquad
 |F(s)-C(s)|\le |t_0+s|\sqrt A,\qquad |z-1|=O(A).
\]
The triangle inequality therefore transfers the gap to all three vertices
\(F,B,C\), and hence to their convex hull.  This proves the claimed
power domination without interpolating unproved endpoint estimates.

The algebraic ledger is also exact.  With
\(E_\tau=(1-\tau)B+\tau F\),
\begin{align}
 F^n-C^n={}&\lambda^nz^n nR\int_0^1Z_\tau^{\,n-1}\,\dd\tau
 +\lambda^n(z^n-1)\psi^n+nB^{\,n-1}D                         \notag\\
 &+n(n-1)D^2\int_0^1(1-\tau)E_\tau^{\,n-2}\,\dd\tau . \tag{4.67e8}
\end{align}
For the first term in the first row, put
$s=-t_0v/\sqrt n$, so that
$t=1-v/\sqrt n$, $\dd t=-\dd v/\sqrt n$, and
$|K(1-v/\sqrt n)|\le Cv/\sqrt n$.  The zero-cell divided-power bound for
$Z_\tau$ following (4.67e6) gives
\[
 \left|\int_0^1Z_\tau(s)^{n-1}\,\dd\tau\right|
      \le Ce^{-cv^2}.
\]
Consequently, since the phases and $z^n$ have modulus at most one,
(4.66b) gives
\begin{align}
 &\sqrt n\left|2\int_{1-\delta}^1K(t)\,
   \lambda^nz^n nR(s)
       \int_0^1Z_\tau(s)^{n-1}\,\dd\tau\,\dd t\right|\notag\\
 &\quad\le C\sqrt n\int_0^{\delta\sqrt n}
      v e^{-cv^2}|R(-t_0v/\sqrt n)|\,\dd v\notag\\
 &\quad\le C\left\{{A\over\sqrt n}+{d\eta_n\over n}\right\}.
                                                                   \tag{4.67e8a}
\end{align}
Thus the first term has the explicit normalized-score bound
\[
 C\left\{{A\over\sqrt n}+{d\eta_n\over n}\right\}=o(d).             \tag{4.67e9}
\]
The second term in that row costs
\(C\{1-r_0^n+|\omega^n-1|\}/\sqrt n\).  Insert (4.67e1) in the
linear term and use (4.67e6) in the quadratic term.  With
\(s=-t_0v/\sqrt n\) and \(|K(1-v/\sqrt n)|\le Cv/\sqrt n\), their costs
are respectively
\[
 O\left(|m|+{|q_2|\over\sqrt n}+{|q_3|\over n}\right)+o(d),
 \qquad O(\sqrt n|m|^2)+o(d).
\]
Here is the full content of the second $o(d)$.  Squaring the four terms in
(4.67e1), absorbing cross products by $2ab\le a^2+b^2$, and using the
returned-kernel factor \(v/n\) gives the normalized bound
\[
 C\left\{\sqrt n|m|^2+{|q_2|^2\over\sqrt n}
       +{|q_3|^2\over n^{3/2}}\right\}+o(n^{-1}).     \tag{4.67e10}
\]
For the Peano terms choose
\(B_n=C_0\sqrt{\log(n+2)}\), with $C_0$ large enough that
$n^3e^{-cB_n^2}\to0$.  Then $B_n/\sqrt n\to0$, so the uniform modulus in
(4.67e1) tends to zero on $v\le B_n$; on $v>B_n$, (4.67e6) and the
choice of $C_0$ dominate every polynomial prefactor.  This proves the
$o(n^{-1})$ in (4.67e10), rather than merely asserting a diagonal
cutoff.  Moreover
\[
\begin{gathered}
 |m|=o(d),\quad |q_2|/\sqrt n=o(d),\quad |q_3|/n=o(d),\\
 \sqrt n|m|^2=o(d),\quad |q_2|^2/\sqrt n=o(d),\quad
 |q_3|^2/n^{3/2}=o(d),\quad n^{-1}=o(d),              \tag{4.67e11}
\end{gathered}
\]
by (4.56), (4.67b), and (4.67d); boundedness of the third absolute
moments suffices for the penultimate relation.  The fixed physical
middle is \(O(\operatorname {poly}(n)e^{-cn})=o(d)\) by (4.56).
Notice in particular that the cubic marked term costs $|q_3|/n$ after
the outer $\sqrt n$ normalization: $s^3$ supplies $n^{-3/2}$, the
linear divided power supplies $n$, and the returned kernel together with
$\dd t$ supplies $n^{-1}$.  The stronger estimate
$q_3/\sqrt n=o(d)$ in (4.67d) is therefore more than is needed.
Consequently,
uniformly in the cdf coordinate and in $\sigma\in\{+,-\}$,
\begin{align}
 &\sqrt n\max_{\sigma\in\{+,-\}}\sup_x
 |\mathcal J^\sigma_{X,n,h;I_1}(x)
       -\mathcal J^\sigma_{L,n,h;I_1}(x)|                         \notag\\
 &\quad\le C\Big\{{1-r_0^n+|\omega^n-1|\over\sqrt n}+|M_1|
       +{|q_2|\over\sqrt n}                            \notag\\
 &\hspace{42mm}+\sqrt n|M_1|^2+{|q_3|\over n}\Big\}
 =o(d).                                                \tag{4.67f}
\end{align}
Thus (4.67), (4.67e4)--(4.67e8), (4.67f), and the adjacent-mass middle
Cram\'er gap transfer both one-sided operators with $o(d)$ cost.  The
defining outer minus in the reflected contribution (4.24v) is already
included and is immaterial in these absolute-value estimates.

It remains to center and standardize the lattice comparison.  Least
squares gives
\[
 \E L=0,\qquad\sigma_L^2:=\E L^2=1-A,\qquad
 L^\circ={L\over\sigma_L},\qquad b={h\over\sigma_L}.   \tag{4.68a}
\]
The moment estimate (4.64), followed by
\(\sigma_L^{-3}=1+O(A)\), yields
\[
 \gamma(L^\circ)=\gamma(X)+o(d),\qquad
 \rho(L^\circ)=\rho(X)+o(d),\qquad
 d(L^\circ)=d(X)+o(d).                                \tag{4.68}
\]
For \(\Phi_{\sigma_L}(x)=\Phi(x/\sigma_L)\),
\[
 \sup_x|\Phi_{\sigma_L}(x)-\Phi(x)|\le CA,\qquad
 \sqrt nA=o(d).                                       \tag{4.68b}
\]
Exact scaling gives
\[
 \mathcal U^\pm_{L,n,h}(x)=
 \mathcal U^\pm_{L^\circ,n,b}(x/\sigma_L).            \tag{4.68c}
\]
At the active point put \(y_n=x_n/\sigma_L\to0\), with the row index on
\(\sigma_L\) suppressed.  Equations
(4.67), (4.67f), and (4.68b)--(4.68c) prove (4.57), including the true
third-absolute-moment denominators.  Before invoking the exact-lattice
lemma, its hypotheses are as follows: (4.61) gives centering and variance
(1-A); (4.61c) preserves two adjacent masses; (4.61d) keeps the span in
a fixed compact subset of $(0,\infty)$; the support argument following
(4.61d) proves maximality; and standardization preserves both adjacency
and maximality.  Finally (4.64) and (4.68) give
\[
 p(L^\circ)\to p_*,\quad d(L^\circ)=d+o(d),\quad
 \rho(L^\circ)=\rho(X)+o(d),\quad
 0<c'\le\sqrt n\,d(L^\circ)\le C',\quad y_n\to0.        \tag{4.68d}
\]
Thus every hypothesis is checked before Lemma~\ref{orig:lem:anchored} is applied
directly at \(y_n\); that lemma gives
\[
 {\sqrt n\over\rho(X_n)}
 \{\mathcal U^+_{X_n,n,h_n}(x_n)-\Phi(x_n)\}
 \le\CE-\{G_{\rm span,*}-\Lambda_{\rm abs,*}-o(1)\}d_n
       +O(n^{-1})+o(d_n).
\]
Since \(n^{-1}=o(d_n)\) in (4.56), this proves (4.57a).
The dependency order is therefore
\[
\begin{aligned}
 \text{spectral selection}&\ \longrightarrow\
 \text{adaptive majorant}\ \longrightarrow\
 \text{cardinal gap}\\
 &\ \longrightarrow\ A=o(n^{-1})\ \longrightarrow\
 \text{operator transfer}\ \longrightarrow\
 \text{exact-lattice envelope}.
\end{aligned}
\]
In particular, neither the adaptive majorant nor the residual-power
selection uses the near-lattice transfer that it is later used to prove.
\end{proof}

\begin{lemma}[moving-bandwidth low-frequency expansion]
\label{orig:lem:movinglow}
Let $n_j\to\infty$, and for each $j$ let $X_j$ be centered and
variance one.  Denote its third moment and third absolute moment by
$\gamma_j$ and $\rho_j$, let $p_j$ be the matching Bernoulli parameter,
and put
\[
 d_j=\rho_j-\rho_2(p_j),\qquad p_j\to p_*,\qquad d_j\to0,
\]
and
\[
 \lambda_j=\rho_j+\E|X_j|.
\]
For the pointwise Prawitz operator at physical bandwidth $\lambda_j$ one
has, uniformly for $x\in\mathbb R$,
\begin{equation}\tag{4.69}
\begin{split}
 {\sqrt {n_j}\over\rho_j}
   \{\mathcal U^+_{X_j,n_j,\lambda_j}(x)-\Phi(x)\}
 \le{}&{\phi(x)\over\rho_j}
       \left\{{\lambda_j\over2}+{\gamma_j\over6}(1-x^2)\right\}\\
 &+\{\Lambda_{\rm abs}(p_j,d_j)+o(1)\}d_j
       +O(n_j^{-1/2}).
\end{split}
\end{equation}
The $o(1)$ and $O(n_j^{-1/2})$ are uniform in $x$.  For the conjugate
kernel the sign of the skewness term is reversed.  The assertion is
pointwise and uniform in the cdf coordinate; it is not notation for a
supremum over that coordinate.  More precisely, along every array in the
hypotheses there are numbers $\varepsilon_j\downarrow0$ and a constant
$C<\infty$, independent of $x$, such that the two remainders on the
right of \textup{(4.69)} are bounded by
$\varepsilon_jd_j$ and $Cn_j^{-1/2}$, respectively.
\end{lemma}

\begin{proof}
The rows with $d=0$ are separate: Theorem~\ref{orig:thm:defect} gives
$X\stackrel d=Q_p$, so the defect contribution (4.69i) is zero.  For the defect
calculation below pass to rows with $d>0$.  The matching-law expansion and
the three frequency-region estimates below apply unchanged to the rows
$d=0$, where $\lambda=s$.  Suppress the row index; in either case
Proposition~\ref{orig:prop:m1} gives
\[
 0\le\lambda-s\le K_1(\gamma)d,                       \tag{4.69a}
\]
so $\lambda$ stays in a fixed compact subinterval of $(0,\infty)$.
Choose a fixed zero cell on which the physical argument is within the
range of (4.7).  The substitution
\[
 t={\lambda v\over2\pi\sqrt n}                        \tag{4.69b}
\]
has two exact consequences which are important here: the characteristic
function is evaluated at $v/\sqrt n$, and
\[
 2K(t)\,\dd t={\ii\over\pi v}\,\dd v
       +{\lambda\over2\pi\sqrt n}\,\dd v
       +O(v/n)\,\dd v.                                \tag{4.69c}
\]
Thus neither the singular coefficient nor the physical Fourier scale
changes when the bandwidth moves.

We first expand the matching two-point transform.  Uniformly on
$0\le v\le n^{1/21}$,
\[
 \chi_p(v/\sqrt n)^n=e^{-v^2/2}
 \left(1-{\ii\gamma v^3\over6\sqrt n}\right)
 +O\left({(v^4+v^6)e^{-cv^2}\over n}\right),          \tag{4.69d}
\]
and both sides have a common Gaussian bound up to the end of the fixed
zero cell.  Insert (4.69c)--(4.69d), retain
$e^{-\ii vx}$ exactly, and use normal inversion.  The singular skewness
term and the bounded kernel term give respectively
\[
 {\gamma\over6\rho}(1-x^2)\phi(x),\qquad
 {\lambda\over2\rho}\phi(x),                         \tag{4.69e}
\]
with an $O(n^{-1/2})$ normalized remainder.  Its dominator is a fixed
Gaussian polynomial and is independent of $x$.

We next calculate the defect term at this bandwidth, rather than invoking
Lemma~\ref{orig:lem:abslow} by analogy.  In Bernoulli coordinates,
\[
 \varphi_X(v/\sqrt n)-\chi_p(v/\sqrt n)
 =e^{-\ii spv/\sqrt n}d\int
       {E_{sv/\sqrt n}(y)\over W_p(y)}\dd\nu(y),
 \qquad d\,\dd\nu=W_p\,\dd\Pp_Y.                     \tag{4.69f}
\]
With $H=n^{1/4}$, Hermite--Genocchi gives on $|y|\le H$
\[
 { |E_{sv/\sqrt n}(y)|\over W_p(y)}
 \le {C(1+|y|)v^4\over n^2},                         \tag{4.69g}
\]
whose normalized singular-kernel contribution is $O(dH/\sqrt n)=o(d)$.
On $|y|>H$, subtract
$g(svy/\sqrt n)$ and use (4.21); that correction is $O(d/\sqrt n)$.
The remaining term satisfies
\[
 {n^{3/2}\over d}
 \left|\E\{g(svY/\sqrt n);|Y|>H\}\right|
 \le {s^3v^3\over6}
 \max\{C_+(p)^{-1},C_-(p)^{-1}\}
        \{1+O(H^{-1})\}.                               \tag{4.69h}
\]
For reference, the exact divided power used at this point is
\[
 \mathscr P_n(v):={\varphi_X(v/\sqrt n)^n-\chi_p(v/\sqrt n)^n
       \over n\{\varphi_X(v/\sqrt n)-\chi_p(v/\sqrt n)\}}
 =\int_0^1\{(1-a)\chi_p(v/\sqrt n)
                 +a\varphi_X(v/\sqrt n)\}^{n-1}\dd a,              \tag{4.69h1}
\]
with the continuous value when the denominator vanishes.  Every transform
inside the braces is that of a centered variance-one mixture whose third
absolute moment is bounded uniformly along the array.  Hence (4.7), first
on $|v|\le\delta\sqrt n$ and then the fixed Cram\'er gap on the rest of
the zero cell, gives the required common power bound.  Explicitly, on
every fixed annulus \(\delta\le |u|\le u_0\) contained in that cell,
\[
 \sup_{p\ {\rm near}\ p_*}|\chi_p(u)|\le q<1,\qquad
 |\varphi_X(u)-\chi_p(u)|\le C d,
\]
so, for all sufficiently large rows, both endpoints and every convex
combination in (4.69h1) have modulus at most
\((1+q)/2<1\).  Consequently
\[
 |\mathscr P_n(v)|\le Ce^{-cv^2},\qquad
 \mathscr P_n(v)\longrightarrow e^{-v^2/2}             \tag{4.69h2}
\]
locally uniformly in $v$.  The same statements hold after multiplication
by any spatial phase, since that phase has modulus one.  In particular,
dominated convergence applied to the nonnegative majorant
$C(v+v^2+v^3+v^4)e^{-cv^2}$ supplies a number
$\varepsilon_n\downarrow0$, independent of $x$, which bounds by
$\varepsilon_nd$ the sum of: the compact part in (4.69g), the
interpolant correction in (4.21), the $O(H^{-1})$ term in (4.69h), and
the replacement of $\mathscr P_n$ by its Gaussian limit.  Equations
(4.69c), (4.69f)--(4.69h2), and
$\pi^{-1}\int_0^\infty v^2e^{-v^2/2}\dd v=\phi_0$ therefore bound the
singular defect by
\[
 \{\Lambda_{\rm abs}(p,d)+o(1)\}d.                   \tag{4.69i}
\]
The bounded-kernel part of the defect is $O(d/\sqrt n)$ in normalized
score.  Its change caused by $\lambda-s=O(d)$ is no larger; the matching
smoothing change itself is not discarded, but is retained exactly as
$\lambda/2$ in (4.69e).  Every estimate above is unchanged after taking a
supremum over $x$, because the only omitted spatial factor has modulus one.

It remains to justify that no other frequency cell contributes at the
displayed scale.  Choose $0<\delta_1<1/4$ so small that (4.7) applies
whenever $0\le t\le2\delta_1$, uniformly for all large rows.  Near zero
the resulting Gaussian envelope comes from (4.7), after
$u=2\pi t/\lambda$; it does not come from the quartic lower bound below.
For $\delta_1\le t\le1-\delta_1$, (4.6) and compactness of the bandwidth
range give constants $q<1$ and $C<\infty$ such that the whole integral is
at most $Cq^n$ in raw units.  Finally, at
$t=1-h$, $0\le h\le\delta_1$, the same inequality gives
$|\varphi_X(2\pi t/\lambda)|^n\le e^{-cnh^2}$, while
$|K(1-h)|\le Ch$.  Hence this endpoint has raw size $O(n^{-1})$ and
normalized size $O(n^{-1/2})$.  These bounds also ignore only a unit
spatial phase.  Combining the three frequency regions proves (4.69).
Conjugation reverses the singular skewness sign and leaves all absolute
defect estimates unchanged.
\end{proof}

\begin{lemma}[moment-band signed gap]\label{orig:lem:momentbandgap}
Let $X=X_n$ be a triangular array of centered variance-one laws, let
$p=p_n$ be their matching Bernoulli parameters, and put
$d=d_n=\E|X_n|^3-\rho_2(p_n)$.  Suppose
$n\to\infty$, $p\to p_*$ and $d\to0$, and let
$\rho=\rho_2(p)+d$ and $\lambda=\rho+\E|X|$.  Uniformly at every
cdf coordinate $x\in\mathbb R$, the pointwise positive one-sided operator
at physical bandwidth $\lambda$ satisfies
\[
 {\sqrt n\over\rho}\{\mathcal U^+_{X,n,\lambda}(x)-\Phi(x)\}
 \le \CE-\{G_{\rm span,*}-\Lambda_{\rm abs,*}-o(1)\}d
             +O(n^{-1/2}).                              \tag{4.70a}
\]
For the conjugate operator the same estimate holds with a strictly smaller
leading constant: there is a fixed $c_->0$ such that its right side may be
replaced by $\CE-c_-+o(1)$.
\end{lemma}

\begin{proof}
The three separate envelope sources are, explicitly,
\[
 \begin{array}{c|c}
 \text{frequency region}&\text{owner of the bound}\\ \hline
 t\simeq0&(4.7),\ \text{Gaussian in }\sqrt n\,t,\\
 \delta\le t\le1-\delta&(4.6),\ \text{a fixed Cram\'er gap},\\
 t=1-h&(4.6),\ \text{Gaussian in }\sqrt n\,h.
 \end{array}                                             \tag{4.70b}
\]
In particular, the quartic behavior of the right side of (4.6) at zero
is not used as a Gaussian estimate.  Lemma~\ref{orig:lem:movinglow}, whose
proof carries out the moving-bandwidth calculation, gives pointwise
\[
 {\sqrt n\over\rho}\{\mathcal U^+_{X,n,\lambda}(x)-\Phi(x)\}
 \le {\phi_0\over\rho}\left({\lambda\over2}+{\gamma\over6}\right)
       +\{\Lambda_{\rm abs,*}+o(1)\}d+O(n^{-1/2}).     \tag{4.70d}
\]
Here $\lambda>|\gamma|$ and the positive profile in (4.69) is maximized at
$x=0$, which justifies the replacement of $\phi(x)$ by $\phi_0$ in
(4.70d).
Proposition~\ref{orig:prop:m1} gives $\lambda\le s+K_1(\gamma)d$.  Thus the
actual profile in (4.70d) is first bounded above by the numerator displayed
below.  Expanding that upper envelope with the true denominator, before
replacing any parameter by its limit, gives
\[
 {\phi_0\{s/2+\gamma/6+(K_1(\gamma)/2)d\}\over\rho_2(p)+d}
 =A(p)-\{G_{\rm span,*}+o(1)\}d\le
 \CE-\{G_{\rm span,*}+o(1)\}d.
\]
The equality is simply the first-order identity
\[
 {N_0(p)+N_1(p)d\over D_0(p)+d}
 ={N_0(p)\over D_0(p)}
 -{N_0(p)/D_0(p)-N_1(p)\over D_0(p)}d+O(d^2),           \tag{4.70e}
\]
used with
\[
 N_0(p)=\phi_0(s/2+\gamma/6),\quad
 N_1(p)=\phi_0K_1(\gamma)/2,\quad D_0(p)=\rho_2(p).
\]
At $p=p_*$, $N_0/D_0=\CE$ and the coefficient of $d$ in
(4.70e) is exactly $G_{\rm span,*}$; continuity makes the discrepancy
$o(1)$ along the array.  Thus no denominator has been replaced by its
limit before the signed loss is extracted.
The elementary bounds (4.22)--(4.23) give
$G_{\rm span,*}-\Lambda_{\rm abs,*}>19/200>0$, proving (4.70a).
Reflection reverses the
skewness term at the equality parameters and so leaves a strict leading
gap.
\end{proof}

\subsection*{Uniform signed local stability and completion}

\begin{proposition}[uniform signed local stability]\label{orig:prop:qlocal}
Let $n_j\to\infty$.  Suppose that laws $X_j$ satisfy
\[
 p_j\to p_*,\qquad d_j\to0,                            \tag{LS0}
\]
and that $\mathcal C_{n_j}(X_j)>\CE$ for every $j$.  Choose a one-sided
cdf point whose normalized score is still larger than $\CE$.  After
reflection and passage to a subsequence it lies in the positive central
branch.  Choose its unique half-open Bernoulli-cell representation
\[
 x_j={s(p_j)(\eta_j+\theta_j)\over\sqrt {n_j}},\qquad
 n_jp_j+\eta_j\in\mathbb Z,\qquad
 -\tfrac12\le\theta_j<\tfrac12.                         \tag{LS1}
\]
It is impossible that the corresponding one-sided scores satisfy
\[
 \mathcal C_{n_j}^+(X_j;\eta_j,\theta_j)>\CE            \tag{LS}
\]
for every $j$.
\end{proposition}

\begin{proof}
If $d_j=0$ on an infinite subsequence, Theorem~\ref{orig:thm:defect} makes
$X_j$ exactly its matching two-point law there, and
$C_Q^*(n_j)=\CE-\kappa_\infty/n_j+o(n_j^{-1})<\CE$, a contradiction.
After discarding finitely many rows we may therefore take $d_j>0$; every
later normalization by $d_j$ is confined to these rows.
Fix \(c_{\rm gap}>0\) smaller than both positive margins in (4.24).
Put $\tau_j=\sqrt{n_j}d_j$.  The compactification $[0,+\infty]$ is
sequentially compact, so every subsequence has a further subsequence on
which either $\tau_j\to0$, $\tau_j\to+\infty$, or
$0<c\le\tau_j\le C<\infty$.  It is therefore enough to contradict (LS)
on each of these three alternatives.  In the first alternative the same
argument, applied to the positive sequence $n_jd_j$, permits exactly the
three further cases used below: convergence to zero, confinement to a
compact subinterval of $(0,\infty)$, or convergence to infinity.

Suppose first that $\tau_j\to\infty$.  Lemma~\ref{orig:lem:momentbandgap} gives
\[
 \mathcal C_{n_j}^+(X_j;\eta_j,\theta_j)
 \le\CE-\{G_{\rm span,*}-\Lambda_{\rm abs,*}-o(1)\}d_j
             +O(n_j^{-1/2}).                           \tag{4.71}
\]
The coefficient is positive by (4.22)--(4.23), as recorded in the proof of
(4.70a), and $n^{-1/2}=o(d_j)$, so (4.71)
contradicts (LS).

Next suppose $\tau_j\to0$.  Pass again to a subsequence according to
$n_jd_j$.  Outer stability gives $\xi_j=\eta_j/\sqrt {n_j}\to0$.
Lemma~\ref{orig:lem:coarsephase}, proved without (4.31)--(4.33), therefore gives
\[
                         \eta_j^2=O(n_jd_j+1).          \tag{4.72}
\]
This is the non-circular owner of the tangent window.

If $n_jd_j\to0$, (4.72) makes $\eta_j$ bounded.  Use
Lemma~\ref{orig:lem:pertendpoint} with any fixed $D$.  The value/radial mode is
nonpositive, the shape and low-frequency costs are
$O(d_j)=o(n_j^{-1})$, and (4.43a) is $o(n_j^{-1})$.  Therefore
\[
 \mathcal C_{n_j}^+(X_j;\eta_j,\theta_j)
 \le C_Q^*(n_j)+o(n_j^{-1})
 =\CE-{\kappa_\infty\over n_j}+o(n_j^{-1})<\CE.        \tag{4.73}
\]

If $n_jd_j$ stays bounded away from both zero and infinity, choose a fixed
$D$ containing the subsequence.  The complete node and phase have already
been joined in the nonpositive square (4.43e), and (4.43f) gives
\[
 \mathcal C_{n_j}^+(X_j;\eta_j,\theta_j)
 \le C_Q^*(n_j)-\{c_{\rm gap}+o(1)\}d_j+o(n_j^{-1})<\CE \tag{4.74}
\]
because its linear coefficient is positive and $d_j\asymp n_j^{-1}$.

It remains within $\tau_j\to0$ that $n_jd_j\to\infty$.  The window
(4.72) is already available, so the paired growing-window conclusion
(4.43g) applies.  For reference, its normalized scale ledger is
\[
 {\eta^4\over n^{3/2}}=O(\sqrt n\,d^2)=o(d),\qquad
 {|\eta|^5\over n^2}=O(\sqrt n\,d^{5/2})=o(d),
\]
\[
 {1+\eta^2\over n^{3/2}}=o(d),\qquad
 {1+|\eta|\over n}=o(d),\qquad {1\over n}=o(d).      \tag{4.77}
\]
Indeed division of the first two bounds by $d$ leaves respectively
$\sqrt n\,d$ and $\sqrt n\,d^{3/2}$, both tending to zero.  Equation
(4.43g), with its coefficient converging to the positive first margin in
(4.24), gives
\[
 \max\{\mathcal C_{n_j}^+(X_j;\eta_j,\theta_j),
          \mathcal C_{n_j}^-(X_j;\eta_j,\theta_j)\}
 \le C_Q^*(n_j)-\{c_{\rm gap}+o(1)\}d_j+o(d_j)<\CE.   \tag{4.78}
\]
This contradicts (LS) throughout the scale $\tau_j\to0$.

Finally suppose that $\tau_j$ stays in a compact subinterval of
$(0,\infty)$.  Strict resonance curvature selects the critical maximum
$u_j$ used in Lemma~\ref{orig:lem:latticetransfer}.  The selected powers belong
to $[0,1]$; pass once more so that they converge.  If their limit is zero,
Lemma~\ref{orig:lem:spectral}, specifically the strict bound (4.43h),
contradicts (LS).  Pass therefore to a subsequence on which the power is
bounded below.  Proposition~\ref{orig:prop:qouter} implies that the scores tend
to $\CE$ and its active points tend to zero.  Lemma~\ref{orig:lem:latticetransfer}
now performs the preliminary regression, invokes the upstream adaptive
Lemma~\ref{orig:lem:adaptiveupstream} and the cardinal gap, and only then
reindexes to the $o(n_j^{-1})$ lattice tube.  Its active-point transfer and
Lemma~\ref{orig:lem:anchored} give
\[
 \mathcal C_{n_j}^+(X_j;\eta_j,\theta_j)
 \le\CE-\{G_{\rm span,*}-\Lambda_{\rm abs,*}-o(1)\}d_j
       +O(n_j^{-1})+o(d_j).                            \tag{4.79}
\]
Here $n_j^{-1}=o(d_j)$ by (4.56), and the absolute span margin in (4.24) is
positive.  Thus (4.79) is strictly below $\CE$, the final contradiction.
Since every putative counterexample sequence has a further subsequence in
one of the alternatives just excluded, no such sequence exists.  This
proves the proposition.  Reflection supplies the identical conclusion when
the other one-sided error is competitive.
\end{proof}

\begin{proof}[Proof of Theorem~\ref{orig:thm:eventual}]
If the assertion were false, for every $j$ one could choose an integer
$n_j\ge j$ and a centered standardized law $X_j$ such that
\[
                         \mathcal C_{n_j}(X_j)>\CE.
\]
Proposition~\ref{orig:prop:qouter}, after reflection and passage to a
subsequence, supplies (LS0) and the positive central branch.
Proposition~\ref{orig:prop:qlocal} then gives a contradiction.  This proves
exactly
\[
 \boxed{\qquad \exists N_0\ \forall n\ge N_0:\ C_n\le\CE.\qquad}
\]
No numerical value of $N_0$ is asserted or needed.
\end{proof}


\section{Eventual two-point optimality and quantitative stability}\label{stab:section}
\subsection{The quantitative theorem}

For a law $P$ of mean zero, variance one and finite third absolute
moment, put
\[
 \beta=\E|X|^3,\quad \gamma=\E X^3,\quad
 f(\gamma)=\frac{\gamma^2+2}{\sqrt{\gamma^2+4}},\quad
 d(P)=\beta-f(\gamma),\quad q(P)=\frac{d(P)}{\beta}.
\]
The moment-defect theorem of Part~\ref{part:structural} gives $d\ge0$, with equality
exactly for the standardized two-point law $Q_\gamma$ having third
moment $\gamma$. In particular $0\le q<1$. Let $\Phi$ and $\phi$
be the standard normal distribution function and density, and let
$F_{n,P}(t)=P^{*n}(({-\infty},t])$ be the distribution function
of the unnormalized sum of $n$ independent variables with law $P$.
Write
\[
 R_n(P)=\frac{\sqrt n}{\beta(P)}
    \sup_t|F_{n,P}(t)-\Phi(t/\sqrt n)|,\qquad
 \CE=\frac{3+\sqrt{10}}{6\sqrt{2\pi}}.
\]
\[
 C_n=\sup_P R_n(P),\quad
 B_n=\sup_{0<p<1}R_n(\operatorname{Bernoulli}(p)).
\]
The supremum defining $C_n$ runs over standardized laws with finite
third absolute moment. All moment parameters and $R_n$ for an
unstandardized nondegenerate law refer to its standardization.

\begin{theorem}\label{stab:thm:penalty}
There are constants $\theta>0$ and $N<\infty$ such that
\begin{equation}\label{stab:eq:penalty-bound}
 R_n(P)+\theta\frac{d(P)}{\beta(P)}\le B_n
                 \qquad(n\ge N)
\end{equation}
for every standardized law $P$ with finite third absolute moment.
Consequently, for some $a>0$ and all sufficiently large $n$,
\begin{equation}\label{stab:eq:linear-bound}
 C_n-R_n(P)\ge a\min\{d(P),1\},\qquad C_n=B_n<\CE.
\end{equation}
One may take $\theta=1/10000$ and $a=1/18500$.
Every maximizing law is two-point. Moreover,
\[
 C_n=\CE-\frac{\kappa_\infty}{n}+o(n^{-1}),\qquad
 \kappa_\infty=\frac{867+287\sqrt{10}}{25920\sqrt{2\pi}}.
\]
\end{theorem}

The inputs from Part~\ref{part:structural} are the exact moment defect, the optimized
two-point expansion $B_n=\CE-\kappa_\infty/n+o(n^{-1})$, and the
uniform outer bound (4.10i)--(4.14). From He--Cheng \cite{HC} we use
the sequential smoothing Lemma 2.2 and its Corollary 2.4, the local
comparison Lemmas 3.3 and 3.5, and Gaussian cancellation Lemma 4.4.
The contact argument is derived below for the penalized functional.
Neither earlier eventual exactness theorem is assumed.

\subsection{The outer estimate with a penalty}

Let $D=37/20$ and $\phi_0=\phi(0)$. The structural estimate used in
Part~\ref{part:structural} gives, whenever $\beta\ge D$,
\begin{equation}\label{stab:eq:cap}
 R_n(P)\le h_*:=\frac{1661}{5000}
             \left(1+\frac{429}{1850}\right)<\CE.
\end{equation}
For $\beta\le D$, reflect the law to have $\gamma\ge0$ and put
\[
 r=\frac{\gamma}{\sqrt{4+\gamma^2}},\quad u=\sqrt{1-r^2},\quad
 \rho_2(r)=\frac{1+r^2}{u}.
\]
\[
 H(r)=\frac{3+r}{3u},\qquad
 L(r)=\frac1{(1-r)(2+r)}+\frac1{3(1-r)^2(2+r)}.
\]
The compact parameter region is
\[
 \mathcal K=\{(r,d):0\le r\le2/3,\quad
                    0\le d\le D-\rho_2(r)\}.
\]
For any sequence with $n_j\to\infty$, $\beta_j\le D$ and
$(r_j,d_j)\to(r,d)$, the uniform Fourier estimate of Part~\ref{part:structural} is
\begin{equation}\label{stab:eq:outer}
 \limsup_j R_{n_j}(P_j)
       \le\frac{\phi_0 H(r)+\phi_0 L(r)d}{\rho_2(r)+d}.
\end{equation}
Its proof requires only the moment cap, not $R_n(P)>\CE$.
The estimates of the nonzero frequency cells and of the divided-power
remainder are uniform on this cap. Thus no uniform integrability
of third moments is being assumed in \eqref{stab:eq:outer}.

At $d=0$ the right side has unique maximum $\CE$ at
$r_E=\sqrt{10}-3$. At $d=D-\rho_2(r)$ the scalar estimate (4.14)
of Part~\ref{part:structural} bounds it strictly by $41\phi_0/40$. Define
\[
                   \mu=\CE-41\phi_0/40>0.
\]
For $0<\theta<\mu$, the penalized scalar expression
\begin{equation}\label{stab:eq:penalized-outer}
 \frac{\phi_0 H(r)+(\phi_0 L(r)+\theta)d}{\rho_2(r)+d}
\end{equation}
is still at most $\CE$, with equality only at $(r_E,0)$.
Indeed it is fractional linear in $d$; its lower endpoint is unchanged,
and its upper endpoint is at most $41\phi_0/40+\theta<\CE$.

Set
\[
 J_n(P)=R_n(P)+\theta q(P),\qquad A_n=\sup_P J_n(P).
\]
Choose also $\theta<\CE-h_*$. Since $J_n(P)\le h_*+\theta<\CE$
when $\beta\ge D$, any asymptotically maximizing sequence for
$A_n$ has $\beta<D$: here $A_n\ge B_n\to\CE$.
Applying \eqref{stab:eq:outer} and \eqref{stab:eq:penalized-outer} along
convergent moment parameters proves
\begin{equation}\label{stab:eq:A-limit}
 A_n\longrightarrow\CE.
\end{equation}
The same argument shows that if $n_j\to\infty$ and
$A_{n_j}-J_{n_j}(P_j)\to0$, then
\begin{equation}\label{stab:eq:penalty-outer-stability}
 d(P_j)\to0,\quad |\gamma(P_j)|\to\gamma_E,\quad
 \beta(P_j)\to\beta_E,\quad R_{n_j}(P_j)\to\CE.
\end{equation}
Here $\gamma_E,\beta_E$ are the moments of the positive-skew
Esseen law $\PE$.

\subsection{Penalized extremizers}

\begin{lemma}[Attainment]\label{stab:lem:attainment}
For all sufficiently large $n$, the value $A_n$ is attained by a law
$P$ and a positive closed-endpoint discrepancy at some finite $t$.
\end{lemma}
\begin{proof}
The penalty $q$ is invariant under reflection, so the supremum also
equals the supremum of the penalized positive signed discrepancy.
Fix a large $n$ with $A_n>h_*+\theta$ and $A_n>2\theta$.
Choose signed maximizing pairs $(P_k,t_k)$, with penalized signed
values $j_k\to A_n$. After discarding finitely many terms,
$\beta_k<D$ by \eqref{stab:eq:cap}. Their unpenalized signed values
are at least $j_k-\theta$ and hence stay bounded away from zero.
The same normal-tail and Chebyshev argument
as in an ordinary extremal problem bounds $t_k$ for this fixed $n$.

Pass to a subsequence with $P_k\Rightarrow P$, $t_k\to t$,
$\beta_k\to b$ and $\gamma_k\to g$. Uniform integrability of
squares preserves mean zero and variance one; let
$\beta=\beta(P)\le b$ and $\gamma=\gamma(P)$.
Lower semicontinuity of $|x|^3+x^3$ and $|x|^3-x^3$ gives
\begin{equation}\label{stab:eq:escape}
                         |g-\gamma|\le b-\beta.
\end{equation}
The derivative
\[
 f'(v)=\frac{v(v^2+6)}{(v^2+4)^{3/2}},\qquad
 1-f'(v)^2=\frac{4(3v^2+16)}{(v^2+4)^3}>0
\]
shows that $f$ is $1$-Lipschitz. Portmanteau at the moving closed
threshold, followed by continuity of $\Phi$, gives
\[
 \sqrt n\{F_{n,P}(t)-\Phi(t/\sqrt n)\}
                  \ge(A_n-\theta)b+\theta f(g).
\]
The resulting signed penalized value at $P,t$ is at least
\[
 \begin{split}
 \theta+\frac{(A_n-\theta)b+\theta f(g)-\theta f(\gamma)}\beta
 &\ge A_n+(A_n-2\theta)\frac{b-\beta}\beta\ge A_n.
 \end{split}
\]
Maximality forces $b=\beta$ and equality of the value with $A_n$.
The limiting signed discrepancy is positive, since
$A_n-\theta q(P)>0$. This proves the claim.
\end{proof}

At such a pair, write $S=A_n$, $R=R_n(P)=S-\theta q(P)$,
$z=t/\sqrt n$, $M=\E X|X|$ and $G=F_{n-1,P}$. Under raw
contamination by a point mass at $y$, followed by standardization,
the moment derivatives are
\[
 U_\beta(y)=|y|^3-\beta-3My-\tfrac32\beta(y^2-1),\qquad
 U_\gamma(y)=y^3-\gamma-3y-\tfrac32\gamma(y^2-1).
\]
Differentiating the penalized signed discrepancy at the same raw
threshold gives
\begin{equation}\label{stab:eq:penalized-contact}
 \begin{split}
 \beta I(y)={}&n^{3/2}\{G(t-y)-F_{n,P}(t)\}
       +ny\phi(z)+\tfrac{\sqrt n}{2}z\phi(z)(y^2-1)\\
       &-(S-\theta)U_\beta(y)-\theta f'(\gamma)U_\gamma(y).
 \end{split}
\end{equation}
For clarity, the penalty derivative is
$\beta\theta q'=\theta\{f(\gamma)U_\beta/\beta-f'(\gamma)U_\gamma\}$;
combining it with $-R U_\beta$ yields \eqref{stab:eq:penalized-contact}.
Maximality gives $I(y)\le0$ for every $y$. The integral of its right
side under $P$ is zero, since $\E U_\beta=\E U_\gamma=0$ and the
convolution term also integrates to zero. Upper semicontinuity then
gives
\begin{equation}\label{stab:eq:whole-contact}
                            I(y)=0\quad(y\in\supp P).
\end{equation}
No point mass at an individual support point is required.

\subsection{Selection and support control}

Suppose $A_n>B_n$ for arbitrarily large $n$. Apply harmonic
selection to $a_n=A_n-B_n\ge0$, which tends to zero. It gives
$n_j\uparrow\infty$ with $a_{n_j}>0$ and
$n_j(a_{n_j-1}-a_{n_j})_+\to0$. This elementary selection follows
by summing $a_{n-1}-a_n\ge\epsilon/n$ backwards along a hypothetical
arbitrarily late run of positive terms.
Writing $B_n=\CE-\kappa_\infty/n+\epsilon_n/n$, $\epsilon_n\to0$,
we have the exact identity
\[
 n(B_{n-1}-B_n)=-\frac{\kappa_\infty}{n-1}
                +\frac n{n-1}\epsilon_{n-1}-\epsilon_n\to0.
\]
Hence
\begin{equation}\label{stab:eq:selected}
 A_{n_j}>B_{n_j},\qquad
                  n_j(A_{n_j-1}-A_{n_j})_+\longrightarrow0.
\end{equation}
Choose positive penalized maximizing pairs by
Lemma~\ref{stab:lem:attainment}.

Suppress the index $j$. Put $\alpha_n=\sqrt{n/(n-1)}$ and let
$H_n(z,y)$ denote the Gaussian cancellation expression
\[
 n^{3/2}\left\{\Phi\left(\frac{z-y/\sqrt n}{\sqrt{1-1/n}}\right)
                    -\Phi(z)\right\}
      +ny\phi(z)+\tfrac{\sqrt n}{2}z\phi(z)(y^2-1).
\]
The definition of $A_{n-1}$, for the same law $P$, gives
\[
 G(t-y)-\Phi\left(\frac{z-y/\sqrt n}{\sqrt{1-1/n}}\right)
       \le\frac{\beta(A_{n-1}-\theta q(P))}{\sqrt{n-1}}.
\]
Together with the exact active value at $n$, this implies
\begin{equation}\label{stab:eq:penalized-majorant}
 \begin{split}
 \beta I(y)\le{}&H_n(z,y)+\beta n(A_{n-1}-S)\\
   &+\beta n(\alpha_n-1)(A_{n-1}-\theta q(P))\\
   &-(S-\theta)U_\beta(y)-\theta f'(\gamma)U_\gamma(y).
 \end{split}
\end{equation}
All moments in this expression are bounded on the selected sequence.
The recursive positive part tends to zero and $n(\alpha_n-1)$
is bounded. The Gaussian bound from \cite[Lemma 4.4]{HC} gives
$|H_n(z,y)|\le\phi_0|y|^3/6+C(1+|y|)$.
Choose $\theta<\CE/8$. Since $S\to\CE$ and $|f'|<1$, the cubic
coefficient in the upper bound is at most
$-S+2\theta+\phi_0/6$, bounded away from zero on the negative side.
Thus \eqref{stab:eq:whole-contact} forces one common compact support.

By \eqref{stab:eq:penalty-outer-stability}, $d_j\to0$ and $R_j\to\CE$.
Extract a $W_3$-convergent subsequence using the common compact support.
The moment-defect equality identifies its limit as the Esseen law
or its reflection. The positive envelope in \cite[Corollary 2.4]{HC}
excludes the reflected law and forces $z_j\to0$: the reflected
envelope has strictly smaller supremum, while the positive-skew
envelope has unique maximum at zero. Thus
\[
                       P_j\to\PE\text{ in }W_3,\qquad z_j\to0.
\]

Write $p_E=(4-\sqrt{10})/2$, $q_E=1-p_E$,
$\sigma_E=\sqrt{p_Eq_E}$, $a_E=p_E/\sigma_E$,
$b_E=q_E/\sigma_E$, $h_E=a_E+b_E$, and $M_E=q_E-p_E$.
For any convergent support sequence $y_j\to y$, take upper limits
in \eqref{stab:eq:penalized-majorant}, using contact. The recursive
positive part vanishes, $q_j\to0$, and $n(\alpha_n-1)\to1/2$.
We obtain
\begin{equation}\label{stab:eq:polynomial}
                           0\le V(y)+\theta w_E(y),
\end{equation}
where
\[
 V(y)=\frac{\phi_0}{6}(y^3-3y)-\CE|y|^3
                       +\frac32\CE\beta_Ey^2+3\CE M_Ey
\]
and $w_E$ is the exact defect polynomial at $\gamma_E$.
Indeed the additional term is
$U_{\beta,E}(y)-f'(\gamma_E)U_{\gamma,E}(y)=w_E(y)$;
this follows either by the defect factorization or by direct expansion.

This perturbation preserves a sufficient support interval for small
fixed $\theta$. With $c_*=3+\sqrt{10}$, the identities
$c_*M_E=1$ and $c_*\beta_E=\sqrt{10}/\sigma_E$ give
\[
 V(y)=-\frac{\phi_0}{6}y^2
       \left\{(c_*-\sgn y)|y|-\tfrac32c_*\beta_E\right\}.
\]
Its nonzero roots are $-\sqrt{10}a_E/2$ and $\sqrt{10}b_E/2$.
Under $v=p_E+\sigma_Ey$, they lie strictly inside $[-1/2,3/2]$.
Therefore $V<0$ on the closed set
\[
                    \mathcal O=\{y:p_E+\sigma_Ey\notin(-1/2,3/2)\}.
\]
Also $w_E>0$ on this set. The ratio $-V/w_E$ is continuous there
and has positive limits at both infinities:
\[
 \frac{\CE-\phi_0/6}{1-f'(\gamma_E)},\qquad
 \frac{\CE+\phi_0/6}{1+f'(\gamma_E)}.
\]
Thus
\begin{equation}\label{stab:eq:theta-support}
                  \theta_*:=\inf_{y\in\mathcal O}\frac{-V(y)}{w_E(y)}>0.
\end{equation}
Choosing $\theta<\theta_*$ makes \eqref{stab:eq:polynomial} impossible
on $\mathcal O$.

It remains to separate distinct support limits. The bounded polynomial
terms in \eqref{stab:eq:penalized-contact} give, uniformly on the support,
\[
 G_j(t_j-v)=F_{n_j,P_j}(t_j)
                  -\phi(z_j)v/\sqrt{n_j}+O(n_j^{-1}).
\]
Because the unpenalized active values $R_j$ tend to $\CE$, each
such residual threshold has the first-order discrepancy
$\CE\beta_E$ at sample size $m_j=n_j-1$. Apply \cite[Lemma 2.2]{HC}
to that convolution, with uniform jitter of width $h_E$.
If $J_j$ is its jittered CDF and $u_j=t_j-v_j$, this gives
\[
 \sqrt{m_j}\{J_j(u_j+h_E/2)-G_j(u_j)\}\to0,
\quad
 J_j(u_j+h_E/2)-G_j(u_j)
   =\frac1{h_E}\int_0^{h_E}[G_j(u_j+s)-G_j(u_j)]\,ds.
\]
Consequently, for $\ell_j\to\ell\in(0,h_E)$, monotonicity gives
$\sqrt{m_j}[G_j(u_j+\ell_j)-G_j(u_j)]\to0$.
Two support limits $x<y$ with $y-x<h_E$ would contradict this:
subtracting their contact expansions makes the same increment tend
to $\phi_0(y-x)>0$. Distinct support limits are therefore separated
by at least $h_E$.

In the coordinate $v=p_E+\sigma_E y$, all support limits lie in
$(-1/2,3/2)$ and are separated by at least one. Weak convergence
supplies the limits zero and one. No other support limit is possible.
The common compact support now gives
\begin{equation}\label{stab:eq:clusters}
           \sup_{x\in\supp P_j}\dist(x,\{-a_E,b_E\})\longrightarrow0.
\end{equation}

\subsection{Absorption of the penalty}

We record the local estimate that connects the moment defect to
the conditional noise. Suppose $Y=B_p+U$, $\E(U\mid B_p)=0$,
$|U|\le\varepsilon$, with $p$ in a fixed sufficiently small compact
interval about $p_E$. Put $s=\E U^2$, $v=p(1-p)$ and
$\tau=p^2+(1-p)^2$. For sufficiently small fixed $\varepsilon$,
the signs of $Y-p$ on the two clusters are fixed, so
\[
 \E Y=p,\quad \operatorname{Var}Y=v+s,\quad
 |\E|Y-p|^3-v\tau|\le C s,\quad
 |\E(Y-p)^3-v(1-2p)|\le C s.
\]
The two last bounds follow by cubic expansion, conditional centering,
and $\E|U|^3\le\varepsilon s$. Since $v$ is bounded away from zero,
standardization and $|f'|<1$ give a constant $K<\infty$ with
\begin{equation}\label{stab:eq:local-defect}
                  0\le d(Y)\le K s,\qquad q(Y)\le K s.
\end{equation}
The second inequality uses the standardized bound $\beta(Y)\ge1$.

From \eqref{stab:eq:clusters}, recenter and rescale the conditional
means of the two clusters. This writes an affine image of $P_j$ as
$Y_j=B_{p_j}+U_j$, where $p_j\to p_E$,
$\E(U_j\mid B_{p_j})=0$ and $\|U_j\|_\infty\to0$.
Explicitly, if $Z_j=p_E+\sigma_E X_j$ and the two cluster radii are
$\delta_j\to0$, their conditional-mean difference lies in
$[1-2\delta_j,1+2\delta_j]$ and the resulting noise is bounded by
$2\delta_j/(1-2\delta_j)$.
Affine invariance and \eqref{stab:eq:penalty-outer-stability} give
$R_{n_j}(Y_j)\to\CE$. Lemma 3.5 of \cite{HC} therefore gives
$\lambda_j=n_j\E U_j^2\to0$.

The proof of \cite[Lemma 3.3, equation (38)]{HC} supplies fixed
$k>0$, $\lambda_0>0$, $\varepsilon_0>0$ and a sample-size threshold
such that
\begin{equation}\label{stab:eq:local-loss}
                 R_n(B_p+U)\le R_n(B_p)-k n\E U^2
\end{equation}
whenever its local hypotheses hold and $n\E U^2\le\lambda_0$.
Indeed its loss is
$\min\{c_0\lambda/\sqrt{C_0\lambda+\varepsilon_0^2},\eta_0\}$,
which is at least $k\lambda$ on $[0,\lambda_0]$ for
$k=\min\{c_0/\sqrt{C_0\lambda_0+\varepsilon_0^2},\eta_0/\lambda_0\}>0$.
Shrink $\varepsilon_0$ and the parameter interval as necessary for
\eqref{stab:eq:local-defect}. All these hypotheses hold eventually along
the selected sequence. Combining the two estimates gives
\[
 A_{n_j}=J_{n_j}(Y_j)
       \le R_{n_j}(B_{p_j})-(k n_j-\theta K)\E U_j^2
       \le B_{n_j}
\]
for all sufficiently large $j$, contradicting \eqref{stab:eq:selected}.

Choose once and for all
\[
 0<\theta<\min\{\mu,\CE-h_*,\CE/8,\theta_*\}.
\]
Every restriction used above is then satisfied. We have proved
$A_n=B_n$ for all sufficiently large $n$, which is
\eqref{stab:eq:penalty-bound}. In particular $C_n=B_n$, and $d=0$
for any maximizing law. A binary maximizer exists: its score is
continuous in $p$, tends to zero at $p=0,1$, and has a positive
interior maximum. Schulz's strict inequality \cite{Schulz2016} gives
$B_n<\CE$ for each fixed $n$. The two-point expansion proves the
last assertion of Theorem~\ref{stab:thm:penalty}.

Finally increase $N$ so that $B_n\ge(\CE+h_*)/2$ for $n\ge N$.
When $\beta\le D$, \eqref{stab:eq:penalty-bound} gives a loss at least
$\theta d/D$. When $\beta>D$, \eqref{stab:eq:cap} gives a loss at least
$(\CE-h_*)/2$. Thus \eqref{stab:eq:linear-bound} holds with
\[
                 a=\min\{\theta/D,(\CE-h_*)/2\}>0.
\]
This completes the proof.

\subsection{An explicit choice of constants}

We verify the constants stated in Theorem~\ref{stab:thm:penalty}. In this
section write $s=\sqrt{10}$. The elementary bounds
$79/25<s<16/5$ and $\pi<22/7$ give
\[
 \phi_0>997/2500>1/3,\qquad
 \CE>\frac{76769}{187500}>\frac25.
\]
Indeed $44\cdot997^2<7\cdot2500^2$ and $(3+s)/6>77/75$.
Since $h_*=3785419/9250000$, subtraction yields
\[
 \CE-h_*>\frac{1111}{5550000}>\frac1{5000},\qquad
 \mu=\phi_0\frac{20s-63}{120}>\frac1{1800},\qquad
 \CE/8>\frac1{20}.
\]

It remains to bound \eqref{stab:eq:theta-support}. Put $p=p_E$, $q=q_E$
and $\sigma=\sigma_E$. The moduli of the nonzero roots of $V$ are
$\ell_-=s a_E/2$ and $\ell_+=s b_E/2$. On the negative component
of $\mathcal O$, $|y|\ge(p+1/2)/\sigma>1$ and
\[
 \frac{\ell_-}{|y|}\le\frac s2\frac p{p+1/2}
                  <\frac45<\frac78.
\]
On the positive component, $y\ge(q+1/2)/\sigma>1$ and, since $q<3/5$,
\[
 \frac{\ell_+}{y}\le\frac s2\frac q{q+1/2}
                  <\frac{48}{55}<\frac78.
\]
The factorization of $V$, together with $c_*-\sgn y>5$, now gives
\[
                         -V(y)>\frac5{144}|y|^3\quad(y\in\mathcal O).
\]
For $S_E=\sqrt{\gamma_E^2+4}\ge2$, the exact defect polynomial is
\[
 w_E(y)=|y|^3-f'(\gamma_E)y^3-\frac{12y^2}{S_E^3}
                         +\frac{6\gamma_Ey+4}{S_E^3}.
\]
Its two zeros are outside $\mathcal O$. Using $|f'|<1$,
$|\gamma_E|\le S_E$ and $|y|>1$ gives there
\[
 0<w_E(y)\le2|y|^3+\tfrac32|y|+\tfrac12\le4|y|^3.
\]
Consequently $\theta_*\ge5/576$. All four restrictions on $\theta$
hold for $\theta=1/10000$, and
\[
 \min\{\theta/D,(\CE-h_*)/2\}=\theta/D=1/18500.
\]
This makes the coefficients explicit, but not the sample-size threshold.

\subsection{Cubic transport stability}

\begin{lemma}\label{stab:lem:transport}
If $\beta(P)\le37/20$, then
\begin{equation}\label{stab:eq:W3}
                         W_3(P,Q_\gamma)^3\le1400\,d(P).
\end{equation}
\end{lemma}
\begin{proof}
Let $S=\sqrt{\gamma^2+4}$, $a=(S-\gamma)/2$,
$b=(S+\gamma)/2$, $m=\min(a,b)$, $M=\max(a,b)$, and
$A=S^3/(4m^2)$. Round $X$ by sign to $Y\in\{-a,b\}$.
The defect factorization in Part~\ref{part:structural} is
\[
 w_\gamma(x)=\frac{2a^2}{S^3}(x-b)^2\{2+(a+3b)x\}\quad(x\ge0),
\]
with $a,b$ interchanged on the negative half-line; $\E w_\gamma(X)=d$.
Since $ab=1$, $M\ge1$ and $a+3b\ge2\sqrt3$,
\[
 \frac{x+b}{2+(a+3b)x}\le M/2\qquad(x\ge0).
\]
The same bound holds on the other half-line. Consequently
\[
                            \E|X-Y|^3\le A M d.
\]
To correct the mass of $Y$, use the cubic $H(x)=3t^2-2t^3$,
$t=(x+a)/S$. Moment matching gives $\E H(X)=a/S$.
Direct factorization yields
\[
 |1_{\{x\ge0\}}-H(x)|
 \le\frac{3S}{4m^2}w_\gamma(x).
\]
For example, on $x\ge0$ the ratio is
$(2x+3a+b)/(2a^2\{(a+3b)x+2\})$, bounded by $3S/(4m^2)$;
the negative half-line is symmetric.
Thus the mass correction needed to turn $Y$ into $Q_\gamma$ is
at most $3S d/(4m^2)$. Move exactly that mass between the two nodes
and call the corrected variable $Y_0$. Gluing the coupling gives
\[
 \E|Y-Y_0|^3\le3S A d,\qquad
 \E|X-Y_0|^3\le4A(M+3S)d.
\]
Here $|\gamma|\le\beta\le37/20$ implies
$S<11/4$, $M<5/2$, $m>2/5$, and
$A<33275/1024$. Therefore
\[
             4A(M+3S)<\frac{1430825}{1024}<1400.
\]
Taking the infimum over couplings proves \eqref{stab:eq:W3}.
\end{proof}

\begin{corollary}\label{stab:cor:stability}
There is $N<\infty$ such that, for every $n\ge N$ and every standardized
law $P$ with finite third absolute moment, if
$\delta=C_n-R_n(P)<1/10000$, then
\[
 \beta(P)<37/20,\qquad
 d(P)\le18500\delta,\qquad
 W_3(P,Q_{\gamma(P)})^3\le25900000\delta.
\]
In particular $W_3(P,Q_{\gamma(P)})\le300\delta^{1/3}$.
No bounded-support or fourth-moment hypothesis is imposed.
\end{corollary}
\begin{proof}
Increase $N$ so that
$B_n\ge(\CE+h_*)/2$. Equation \eqref{stab:eq:cap} excludes
$\beta\ge37/20$ when $\delta<1/10000$, since
$(\CE-h_*)/2>1/10000$.
Then \eqref{stab:eq:penalty-bound} gives
$d\le18500\delta$; apply Lemma~\ref{stab:lem:transport} and
$1400\cdot18500=25900000<300^3$.
\end{proof}

The comparison law in this statement is the matching two-point law,
not necessarily the optimizing Bernoulli parameter at that $n$.
The conclusion controls moment-weighted transport, not the Hausdorff
distance of supports. An arbitrarily small remote atom prevents such
a support-distance statement. Only the threshold $N$ remains qualitative;
the inequality for every sample size remains open.

\subsection{Sharpness of the stability exponents}

Let $\mathcal Q_2$ be the class of probability laws supported on at most
two points. The exponent $1/3$ cannot be increased, even if the comparison
law is freely chosen in $\mathcal Q_2$.

\begin{theorem}\label{stab:thm:sharpness}
For each fixed sufficiently large $n$, there are standardized three-point
laws $P_\varepsilon$, with uniformly bounded supports, such that as
$\varepsilon\downarrow0$,
\[
 \begin{gathered}
 d(P_\varepsilon)\asymp\varepsilon,\qquad
 C_n-R_n(P_\varepsilon)\asymp\varepsilon,\\
 \inf_{Q'\in\mathcal Q_2}W_3(P_\varepsilon,Q')\asymp\varepsilon^{1/3}.
 \end{gathered}
\]
The comparison constants may depend on $n$. Thus neither the linear
defect loss nor the cubic transport exponent admits a stronger power
uniformly near the maximizing class, even at one such fixed index.
\end{theorem}
\begin{proof}
Fix a standardized binary maximizer $Q$ at $n$ and a point $y$ outside
its two-point support. Let $P_\varepsilon$ be the standardization of
\[
                 Z_\varepsilon\sim(1-\varepsilon)Q+\varepsilon\delta_y.
\]
After reflecting $Q$ and $y$ together if necessary, $Q$ has an active positive closed
threshold $t_0$ for its unnormalized sum. A negative active left limit
becomes precisely such a threshold under reflection. At the same raw
threshold $t_0$, the distribution functions of the $n$-fold sums of
$Z_\varepsilon$ and $Q$ differ by at most $n\varepsilon$.
Writing $m_\varepsilon=\varepsilon y$ and
$\sigma_\varepsilon^2=1+\varepsilon(y^2-1)-\varepsilon^2y^2$, the
corresponding standardized threshold is
$(t_0-nm_\varepsilon)/\sigma_\varepsilon$.
It changes the normal term by $O_n(\varepsilon)$, and
$\beta(P_\varepsilon)=\beta(Q)+O(\varepsilon)$. Evaluating the signed
score there gives
\[
                 R_n(P_\varepsilon)\ge C_n-O_n(\varepsilon).
\]
The moment derivatives used in \eqref{stab:eq:penalized-contact}, at a
two-point equality law, give
\[
 d(P_\varepsilon)=\varepsilon w_{\gamma(Q)}(y)+O(\varepsilon^2),
 \qquad w_{\gamma(Q)}(y)>0.
\]
Equation \eqref{stab:eq:penalty-bound} supplies the matching positive lower
bound for $C_n-R_n(P_\varepsilon)$.

The three support points of $P_\varepsilon$ have a common positive
separation, and each has mass at least $\varepsilon$ for small
$\varepsilon$. Choose disjoint balls of a fixed radius $r>0$ around
these points. Any law in $\mathcal Q_2$ omits at least one of the balls;
every coupling therefore has cubic cost at least $r^3\varepsilon$.
Conversely, the cap estimate gives $\beta(Q)<37/20$ for large $n$,
and hence $\beta(P_\varepsilon)<37/20$ for small $\varepsilon$.
Lemma~\ref{stab:lem:transport}, with the matching law
$Q_{\gamma(P_\varepsilon)}$, gives a cubic cost $O(\varepsilon)$.
This proves all three comparisons.
\end{proof}

\section{The oscillatory second-order correction}\label{bin:section}
\subsection{Set-up and statement}

Let $0<p\le 1/2$, $q=1-p$, and let $Q_p$ be the standardized two-point law
which takes the values $\sqrt{q/p}$ and $-\sqrt{p/q}$ with probabilities
$p$ and $q$.  If $K\sim\operatorname{Bin}(n,p)$, then
\[
 \frac{Q_{p,1}+\cdots+Q_{p,n}}{\sqrt n}
 =\frac{K-np}{\sqrt{npq}},
 \qquad
 \rho(p)=\frac{p^2+q^2}{\sqrt{pq}}.
\]
For an integer $k$, put
\[
 \eta=k-np,
 \qquad
 B_{n,p}(k)=\Pp\{K\le k\},
\]
and define the normalized positive right-edge score
\begin{equation}\label{bin:eq:score}
 S_n(p,\eta)
 :=\frac{\sqrt n}{\rho(p)}
 \left\{B_{n,p}(k)-\Phi\left(\frac{\eta}{\sqrt{npq}}\right)\right\}.
\end{equation}
The notation presupposes that $k=np+\eta$ is an integer.

Set
\[
 \phiZ=(2\pi)^{-1/2},\qquad D=p^2+q^2,
\]
\begin{equation}\label{bin:eq:G}
 G(p)=\phiZ\frac{2-p}{3D},
\end{equation}
and
\begin{equation}\label{bin:eq:F}
 F(p,\eta)=-\frac{\phiZ}{540pqD}
 \left\{
 270q\eta^2+45(p^2-7p+5)\eta+p^3+21p^2-69p+46
 \right\}.
\end{equation}
The first-order constants from Section~\ref{orig:sec:binary} are
\begin{equation}\label{bin:eq:starfirst}
 p_* = \frac{4-\sqrt{10}}2,
 \qquad
 \eta_*=-\frac{20+\sqrt{10}}{72},
 \qquad
 \CE=G(p_*)=\frac{\sqrt{10}+3}{6\sqrt{2\pi}},
\end{equation}
\begin{equation}\label{bin:eq:kappa}
 \kappa_\infty=
 \frac{867+287\sqrt{10}}{25920\sqrt{2\pi}},
 \qquad
 F(p_*,\eta_*)=-\kappa_\infty.
\end{equation}

The new constants are
\begin{align}
 H_*&=-\phiZ\frac{653866157+211759711\sqrt{10}}{25395793920},
 \label{bin:eq:Hstar}\\
 F_{p,*}&:=\partial_pF(p_*,\eta_*)
 =-\phiZ\frac{1483+463\sqrt{10}}{19440},
 \label{bin:eq:Fpstar}\\
 G_*''&:=G''(p_*)
 =-\phiZ\frac{60+19\sqrt{10}}{15},
 \label{bin:eq:Gppstar}\\
 \delta_0&:=-\frac{F_{p,*}}{G_*''}
 =\frac{1010-397\sqrt{10}}{12960},
 \label{bin:eq:delta0}\\
 A_2&:=-\frac12G_*''
 =\phiZ\frac{60+19\sqrt{10}}{30},
 \label{bin:eq:A2}\\
 \kappa_2&:=-H_*+\frac{F_{p,*}^2}{2G_*''}
 =\phiZ\frac{3183580225+1028276063\sqrt{10}}
 {126978969600}.
 \label{bin:eq:kappa2}
\end{align}
Numerically,
\begin{align*}
 \delta_0&=-0.0189370548678122\ldots,\\
 \kappa_2&=0.0202183334057192\ldots,\\
 A_2&=1.59687652610925\ldots,\\
 G_*''&=-3.19375305221850\ldots.
\end{align*}

\begin{theorem}[second-order lattice oscillation]\label{bin:thm:main}
Let $C_Q^*(n)$ be the optimal fixed-$n$ Berry--Esseen constant restricted
to standardized two-point laws, with the normalization of Section~\ref{orig:sec:binary}.
Define
\begin{equation}\label{bin:eq:omega}
 \omega_n:=\dist\bigl(np_*+\eta_*+\delta_0,\mathbb Z\bigr)
 \in[0,1/2].
\end{equation}
Then
\begin{equation}\label{bin:eq:main}
 \boxed{
 C_Q^*(n)=\CE-\frac{\kappa_\infty}{n}
 -\frac1{n^2}
 \left\{\kappa_2+A_2\omega_n^2\right\}
 +o(n^{-2}).
 }
\end{equation}
The $o(n^{-2})$ term is uniform over the compact local parameter family
associated with the finitely many central-edge types retained by Schulz's
localization theorem.
\end{theorem}

Because $A_2>0$ and $\kappa_2>0$, the second correction is always strictly negative.  It
varies between approximately $-0.0202/n^2$ and $-0.4194/n^2$.

\subsection{The fifth-order central saddle}

The main calculation is a one-order extension of the incomplete-beta saddle
in the second note.

\begin{lemma}[uniform fifth-order saddle]\label{bin:lem:fifth}
Fix a compact interval $U\Subset(0,1)$ and $M<\infty$.  Uniformly for
$p\in U$, $|\eta|\le M$, and integers $k=np+\eta$,
\begin{equation}\label{bin:eq:fifthexp}
 S_n(p,\eta)
 =G(p)+\frac{F(p,\eta)}n+\frac{H(p,\eta)}{n^2}
 +o_{U,M}(n^{-2}),
\end{equation}
where
\begin{equation}\label{bin:eq:H}
 H(p,\eta)=\frac{\phiZ}{30240p^2q^2D}\,P(p,\eta)
\end{equation}
and
\begin{align}
 P(p,\eta)={}&1260(4-5p)\eta^4
 +420(25p^2-55p+28)\eta^3 \notag\\
 &+1260(-p^3+14p^2-21p+9)\eta^2 \notag\\
 &+21(5p^4-70p^3+355p^2-470p+189)\eta \notag\\
 &+5(25p^5-52p^4-17p^3+88p^2-115p+46).
 \label{bin:eq:P}
\end{align}
One may replace the little-oh in \eqref{bin:eq:fifthexp} by
$O_{U,M}(n^{-181/80})$.
\end{lemma}

\begin{proof}
Write $\eps=n^{-1/2}$ and $r=\sqrt{pq}$.  The incomplete-beta identity is
\[
 B_{n,p}(k)
 =\frac{\int_0^q t^{n-k-1}(1-t)^k\dd t}
 {\int_0^1 t^{n-k-1}(1-t)^k\dd t}.
\]
With $t=q+\eps r y$, $A=p/r$ and $B=q/r$, the common multiplicative
constant cancels and the logarithm of the remaining density is
\begin{equation}\label{bin:eq:L}
 L(y)=(nq-\eta-1)\log(1+\eps Ay)
 +(np+\eta)\log(1-\eps By).
\end{equation}
Expanding the logarithms gives
\begin{equation}\label{bin:eq:ell}
 L(y)=-\frac{y^2}{2}+\sum_{j\ge1}\eps^j\ell_j(y),
 \qquad
 \ell_j(y)=c_jy^{j+2}+d_jy^j,
\end{equation}
where
\[
 c_j=\frac{q(-1)^{j+3}A^{j+2}-pB^{j+2}}{j+2},
 \qquad
 d_j=\frac{-(\eta+1)(-1)^{j+1}A^j-\eta B^j}{j}.
\]
Define $E_0=1$ and
\begin{equation}\label{bin:eq:Erec}
 E_m=\frac1m\sum_{j=1}^m j\ell_jE_{m-j}.
\end{equation}
Then
\[
 \exp\left(\sum_{j\ge1}\eps^j\ell_j\right)
 \sim\sum_{m\ge0}\eps^mE_m.
\]
Since $\ell_j(-y)=(-1)^j\ell_j(y)$, recursion
\eqref{bin:eq:Erec} gives $E_m(-y)=(-1)^mE_m(y)$.

For a uniform remainder, restrict first to
$|y|\le R_\eps:=\eps^{-1/40}$.  On $U$ and for bounded $\eta$, Taylor's
formula for the two logarithms gives the uniform remainder
\[
 O_{U,M}(\eps^6R_\eps^8)
\]
after $\ell_5$.  The part omitted after total $\eps$-degree five in the
exponential is
\[
 O_{U,M}(\eps^6R_\eps^{18}).
\]
After integration over the central interval this is, conservatively,
\[
 O_{U,M}(\eps^6R_\eps^{19})
 =O_{U,M}(\eps^{221/40})=o(\eps^5).
\]

Here is a uniform tail justification.  Put
\[
 \alpha=\min\{\inf U,1-\sup U\}>0,
 \qquad t=q+\eps r y.
\]
On the beta integration interval $0<t<1$, exact differentiation of
\eqref{bin:eq:L} gives
\begin{equation}\label{bin:eq:Lpp}
 L''(y)=-pq\left\{
 \frac{q-\eps^2(\eta+1)}{t^2}
 +\frac{p+\eps^2\eta}{(1-t)^2}\right\}.
\end{equation}
For all sufficiently small $\eps$, the two coefficients in braces are at
least $\alpha/2$.  Since
$t^{-2}+(1-t)^{-2}\ge8$ and $pq\ge\alpha^2$, it follows that
\[
 L''(y)\le-4\alpha^3=:-c_U
\]
throughout the whole beta interval.  Moreover, cancellation of the
$\eps^{-1}$ terms gives
\[
 L'(0)=-\eps\{(\eta+1)A+\eta B\}=O_{U,M}(\eps).
\]
If $y_0$ is the unique maximizer of $L$, then
$|y_0|=O_{U,M}(\eps)$ and
\[
 L(y)\le O_{U,M}(\eps^2)-\frac{c_U}{2}(y-y_0)^2.
\]
Consequently the beta tails outside $[-R_\eps,R_\eps]$ are
$O_{U,M}(\exp(-c_UR_\eps^2/8))$, hence smaller than every power of
$\eps$.  The polynomials $E_0,\ldots,E_5$ have uniformly bounded
coefficients and degree at most $15$, so their Gaussian tails satisfy the
same conclusion.  Thus the exact denominator and numerator have the
expansions
\begin{align}
 \mathcal D
 &=D_0+\eps^2D_2+\eps^4D_4+o_{U,M}(\eps^5),
 \label{bin:eq:denexp}\\
 \mathcal N
 &=\frac12(D_0+\eps^2D_2+\eps^4D_4)
 +\eps A_1+\eps^3A_3+\eps^5A_5
 +o_{U,M}(\eps^5),
 \label{bin:eq:numexp}
\end{align}
where
\[
 D_m=\int_{\mathbb R}e^{-y^2/2}E_m(y)\dd y,
 \qquad
 A_m=\int_{-\infty}^0e^{-y^2/2}E_m(y)\dd y,
 \qquad D_0=\sqrt{2\pi}.
\]
Put $\delta_j=D_j/D_0$.  Division of \eqref{bin:eq:numexp} by
\eqref{bin:eq:denexp} yields
\begin{align}
 B_{n,p}(k)
 ={}&\frac12+\phiZ\eps A_1
 +\phiZ\eps^3(A_3-A_1\delta_2) \notag\\
 &+\phiZ\eps^5
 \{A_5-A_3\delta_2+A_1(\delta_2^2-\delta_4)\}
 +o_{U,M}(\eps^5).
 \label{bin:eq:quotient}
\end{align}
Here $A_1,A_3,A_5$ denote the raw half-Gaussian integrals above; the
factor $D_0^{-1}=\phiZ$ is displayed explicitly.

The normal cdf contributes
\begin{equation}\label{bin:eq:normalexp}
 \Phi(\eta\eps/r)
 =\frac12+\phiZ\frac{\eta\eps}{r}
 -\phiZ\frac{\eta^3\eps^3}{6r^3}
 +\phiZ\frac{\eta^5\eps^5}{40r^5}
 +O_{U,M}(\eps^7).
\end{equation}
The coefficients of $\eps$ and $\eps^3$ reproduce \eqref{bin:eq:G} and
\eqref{bin:eq:F}.  For the next coefficient, set
$\widetilde\ell_j=r^j\ell_j$ and
$\widetilde E_m=r^mE_m$.  If
\[
 \widetilde A_m=\int_{-\infty}^0e^{-y^2/2}\widetilde E_m(y)\dd y,
 \qquad
 \widetilde\delta_m=\frac1{D_0}
 \int_{\mathbb R}e^{-y^2/2}\widetilde E_m(y)\dd y,
\]
then exact Gaussian integration gives the coefficient ledger in
Appendix~\ref{bin:app:ledger}.  In particular,
\begin{equation}\label{bin:eq:Hledger}
 \frac{H(p,\eta)}{\phiZ}
 =\frac{
 \widetilde A_5-\widetilde A_3\widetilde\delta_2
 +\widetilde A_1(\widetilde\delta_2^2-\widetilde\delta_4)
 -\eta^5/40
 }{Dr^4}.
\end{equation}
Substitution simplifies exactly to \eqref{bin:eq:H}--\eqref{bin:eq:P}.  Notice
that the apparent fifth power of $\eta$ cancels identically.

The raw quotient error is $O(\eps^{221/40})$.  Dividing by
$\rho(p)\eps$, with $\rho$ uniformly bounded away from zero on $U$, gives
$O(\eps^{181/40})=O(n^{-181/80})$.  This proves the lemma.
\end{proof}

\subsection{Optimization with the integer node retained}

For later use define
\[
 \eta_0(p)=-\frac{p^2-7p+5}{12(1-p)}.
\]
Completing the square in \eqref{bin:eq:F} gives the exact identity
\begin{equation}\label{bin:eq:Fquad}
 F(p,\eta)=F(p,\eta_0(p))
 -\frac{\phiZ}{2p^{3/2}q^{1/2}\rho(p)}
 \{\eta-\eta_0(p)\}^2.
\end{equation}
In particular, $\eta_0(p_*)=\eta_*$.

\begin{proof}[Proof of Theorem~\ref{bin:thm:main}]
The binomial localization theorem used in the second note reduces the
global two-point supremum, up to alternatives separated by a fixed gap
from $\CE$, to the four central jump edges.  Reflection and the strict
gaps in their leading coefficients leave only positive right edges with
$p$ in a fixed neighborhood of $p_*$.  Their lattice coordinate satisfies
$|\eta|\le1$.  Lemma~\ref{bin:lem:fifth} is therefore uniform over every
candidate that can be asymptotically optimal.

For the lower bound, let $k_n$ be an integer nearest to
\[
 np_*+\eta_*+\delta_0
\]
and put
\[
 p_n=\frac{k_n-\eta_*}{n}.
\]
Then $k_n=np_n+\eta_*$ and
\[
 u_n:=n(p_n-p_*)=k_n-(np_*+\eta_*).
\]
Taylor expansion in Lemma~\ref{bin:lem:fifth} gives
\begin{align}
 S_n(p_n,\eta_*)
 ={}&\CE-\frac{\kappa_\infty}{n}
 +\frac1{n^2}
 \left\{H_*+F_{p,*}u_n+\frac12G_*''u_n^2\right\}
 +o(n^{-2}).
 \label{bin:eq:lowerquad}
\end{align}
By \eqref{bin:eq:delta0}--\eqref{bin:eq:kappa2},
\begin{equation}\label{bin:eq:completesquare}
 H_*+F_{p,*}u+\frac12G_*''u^2
 =-\kappa_2-A_2(u-\delta_0)^2.
\end{equation}
The choice of $k_n$ makes $|u_n-\delta_0|=\omega_n$, proving the lower
bound in \eqref{bin:eq:main}.

For the upper bound, choose a central right-edge candidate within
$o(n^{-2})$ of $C_Q^*(n)$ and denote its parameters by $p_n,\eta_n,k_n$.
The preceding lower bound and the strict quadratic maximum of $G$ first
give
\[
 p_n-p_*=O(n^{-1/2}).
\]
Using the $F/n$ term and the exact concavity in \eqref{bin:eq:Fquad} improves
this to
\begin{equation}\label{bin:eq:localrates}
 p_n-p_*=O(n^{-1}),
 \qquad
 \eta_n-\eta_*=O(n^{-1/2}).
\end{equation}
Indeed, comparison with the lower bound gives, for fixed positive
$c,b$ and some $C$,
\[
 c(p_n-p_*)^2
 \le \frac{C|p_n-p_*|}{n}+\frac{C}{n^2},
\]
and then
\[
 \frac{b}{n}\{\eta_n-\eta_0(p_n)\}^2\le\frac{C}{n^2}.
\]

Put
\[
 u_n=n(p_n-p_*),
 \qquad
 d_n=k_n-(np_*+\eta_*).
\]
By \eqref{bin:eq:localrates}, $u_n$ and $d_n$ are bounded and
$d_n-u_n=\eta_n-\eta_*=o(1)$.  Applying Lemma~\ref{bin:lem:fifth}, using
\eqref{bin:eq:Fquad} to discard its nonpositive square, and expanding at
$p_*$ gives
\begin{align*}
 C_Q^*(n)
 &\le \CE-\frac{\kappa_\infty}{n}
 +\frac1{n^2}
 \left\{H_*+F_{p,*}u_n+\frac12G_*''u_n^2\right\}
 +o(n^{-2})\\
 &=\CE-\frac{\kappa_\infty}{n}
 -\frac1{n^2}\left\{\kappa_2+A_2(d_n-\delta_0)^2\right\}
 +o(n^{-2}).
\end{align*}
Since $d_n\in\mathbb Z-(np_*+\eta_*)$ and $G_*''<0$,
\[
 -\kappa_2-A_2(d_n-\delta_0)^2
 \le -\kappa_2-A_2\omega_n^2.
\]
This proves the matching upper bound.
\end{proof}


\subsection{The full optimal constants}
\begin{corollary}\label{main:full-second-order}
With the constants in \eqref{bin:eq:kappa},
\eqref{bin:eq:delta0}, \eqref{bin:eq:A2} and \eqref{bin:eq:kappa2},
put
\[
 \omega_n=\dist(np_*+\eta_*+\delta_0,\Z).
\]
Then
\[
 C_n=C_{\mathrm E}-\frac{\kappa_\infty}{n}
 -\frac{\kappa_2+A_2\omega_n^2}{n^2}+o(n^{-2}).
\]
The complete set of subsequential limits of
\(n^2(C_n-C_{\mathrm E}+\kappa_\infty/n)\) is
\[
 \left[-\kappa_2-\frac{A_2}{4},-\kappa_2\right].
\]
In particular this sequence has no limit, and eventually
\[
 C_n\le C_{\mathrm E}-\frac{\kappa_\infty}{n}
                       -\frac{\kappa_2}{2n^2}.
\]
\end{corollary}
\begin{proof}
Theorem~\ref{stab:thm:penalty} gives \(C_n=C_Q^*(n)\) for all
sufficiently large \(n\). Substitute this equality into
Theorem~\ref{bin:thm:main}. Irrationality of \(p_*\) makes
\(np_*+\eta_*+\delta_0\) dense modulo one; consequently the limit
set of \(\omega_n^2\) is \([0,1/4]\). Finally
\(A_2\omega_n^2\ge0\) and \(\kappa_2>0\) absorb the \(o(n^{-2})\)
remainder into \(\kappa_2/(2n^2)\).
\end{proof}

\section{Proper convolution roots at every index}\label{conv:section}
For a probability law $P$ with mean zero, variance one and finite third
absolute moment, set
\[
 \beta(P)=\int |x|^3\,P(dx),\qquad
 R_n(P)=\frac{\sqrt n}{\beta(P)}
 \sup_x\left|P^{*n}(({-\infty},\sqrt n x])-\Phi(x)\right|.
\]
Write $T_k(Q)$ for the law of $k^{-1/2}(X_1+\cdots+X_k)$, where the
$X_i$ are independent with standardized common law $Q$. Recall
\[
 \CE=\frac{3+\sqrt{10}}{6\sqrt{2\pi}}.
\]

\begin{theorem}\label{conv:thm:root}
Suppose $P=T_k(Q)$ for an integer $k\ge2$. Then, for every $n\ge1$,
\begin{equation}\label{conv:eq:main}
 R_n(P)\le
 0.3031\frac{\beta(Q)+0.646}
 {\beta(Q)+\dfrac{3(k-1)}{2k\beta(Q)^3}}
 <\frac{39}{100}<\CE.
\end{equation}
The conclusion applies whenever $P$ is a standardized nonzero affine
image of a sum of $k$ iid nondegenerate random variables. Finiteness of
the third absolute moment of $P$ suffices; no moment assumption on the
root is needed separately.
\end{theorem}

The only Berry--Esseen estimate used below is Shevtsova's structural
bound \cite{Shevtsova2013}, in the normalization recorded in
\cite[Eq.~(6)]{HC}:
\begin{equation}\label{conv:eq:structural}
 R_j(Q)\le0.3031\left(1+\frac{0.646}{\beta(Q)}\right),
 \qquad j\ge1.
\end{equation}
All terminating decimals denote exact rational constants.

\begin{lemma}\label{conv:lem:surplus}
Let $A,B$ be independent centered real random variables with finite
third absolute moments. Put $u=\E|A|$ and $v=\E|B|$. If $u+v>0$, then
\begin{equation}\label{conv:eq:surplus}
 \E|A+B|^3-\E|A|^3-\E|B|^3\ge\frac{3u^2v^2}{u+v}.
\end{equation}
\end{lemma}
\begin{proof}
Define
\[
 g(x,y)=|x+y|^3-|x|^3-|y|^3-3x|x|y-3y|y|x.
\]
This is zero for $xy\ge0$. For opposite signs, put
$a=\max\{|x|,|y|\}$ and $b=\min\{|x|,|y|\}$. Then
\[
 g(x,y)=2b^2(3a-b)\ge\frac{6a^2b^2}{a+b},
\]
since the difference multiplied by $a+b$ is $2b^3(2a-b)\ge0$.
The mixed expectations in $g(A,B)$ vanish. With $D=\{AB<0\}$,
centering and independence give
\[
 \E[|AB|1_D]=\frac{uv}{2},\qquad
 \E[(|A|+|B|)1_D]
 =\frac u2\Prob(B\ne0)+\frac v2\Prob(A\ne0)
 \le\frac{u+v}{2}.
\]
For $uv>0$, Cauchy--Schwarz yields
\[
 \E g(A,B)\ge
 6\frac{\bigl(\E[|AB|1_D]\bigr)^2}
 {\E[(|A|+|B|)1_D]}
 \ge\frac{3u^2v^2}{u+v}.
\]
The case $uv=0$ is immediate.
\end{proof}

\begin{lemma}\label{conv:lem:iteration}
For iid standardized $X_i$ with $\beta=\E|X_1|^3<\infty$, put
$M_k=\E|X_1+\cdots+X_k|^3$. Then
\begin{equation}\label{conv:eq:iteration}
 M_k\ge k\beta+\frac{3(k-1)}{2\beta^3},\qquad k\ge1.
\end{equation}
\end{lemma}
\begin{proof}
Let $m=\E|X_1|$. Cauchy--Schwarz gives $m\beta\ge1$.
Conditional Jensen gives $\E|X_1+\cdots+X_j|\ge m$.
For $u\ge m$,
\[
 \frac{3u^2m^2}{u+m}\ge\frac32m^3.
\]
Apply Lemma~\ref{conv:lem:surplus} to the partial sum and the next independent
summand. Every increment $M_{j+1}-M_j-\beta$ is at least
$3m^3/2\ge3/(2\beta^3)$. Summation proves the claim.
\end{proof}

\begin{proof}[Proof of Theorem~\ref{conv:thm:root}]
First consider a possibly unnormalized iid root $Y_1,\ldots,Y_k$ whose
sum $S$ has a finite third absolute moment. Put
$T=Y_2+\cdots+Y_k$ and choose $L<\infty$ with
$r=\Prob(|T|\le L)>0$. Independence and Tonelli give
\[
 r\E|Y_1|^3
 =\E[|Y_1|^3 1_{\{|T|\le L\}}]
 \le4\E|S|^3+4L^3r<\infty.
\]
Thus the root has finite mean and variance. Its variance is positive,
and centering, rescaling and, if needed, reflection give $P=T_k(Q)$.

Write $\beta=\beta(Q)$. Grouping $nk$ iid root summands into $n$ blocks
gives the exact identities
\begin{equation}\label{conv:eq:transfer}
 \beta(P)=\frac{M_k}{k^{3/2}},\qquad
 R_n(P)=\frac{k\beta}{M_k}R_{nk}(Q).
\end{equation}
Equations~\eqref{conv:eq:structural} and \eqref{conv:eq:iteration} give the first
inequality in \eqref{conv:eq:main}. Since $3(k-1)/(2k)\ge3/4$, it remains to
bound
\[
 0.3031\frac{\beta+0.646}{\beta+3/(4\beta^3)}.
\]
This is less than $0.39$ for every $\beta>0$: after multiplying by the
denominator, the assertion is
\[
 0.0869\beta+0.2925\beta^{-3}>0.1958026.
\]
Indeed,
\[
 0.0869\beta+0.2925\beta^{-3}
 >\frac{3\beta}{35}+\frac{2}{7\beta^3}
 \ge4\left(\frac{2}{7\cdot35^3}\right)^{1/4}
 >\frac15>0.1958026.
\]
The penultimate comparison is equivalent to $320000>300125$.
Finally $\sqrt{10}>79/25$ and $\pi<22/7$ give
$\CE>76769/187500>2/5>39/100$.
\end{proof}

\begin{corollary}\label{conv:cor:infdiv}
For $P=T_k(Q)$ and every $n\ge1$,
\begin{equation}\label{conv:eq:order}
 R_n(P)\le0.3031+
 \frac{0.1958026}{\sqrt{k}\,\beta(P)}.
\end{equation}
If $P$ is infinitely divisible, then $R_n(P)\le0.3031$ for every $n\ge1$.
\end{corollary}
\begin{proof}
The weaker inequality $M_k\ge k\beta(Q)$ implies
$\beta(Q)\le\sqrt{k}\,\beta(P)$. Use \eqref{conv:eq:structural} in
\eqref{conv:eq:transfer} to obtain \eqref{conv:eq:order}. Infinite divisibility
provides a root for every $k$. The preceding moment argument applies to
each root; $P$ and $n$ remain fixed. Let $k\to\infty$.
\end{proof}

\clearpage
\part{Effective bounds and finite certificates}\label{part:effective}
\section{The refinement to thirty-one}\label{cut31:section}
The original effective argument below gives a self-contained threshold
of 49. The refinement uses the same contact, support and Fourier
principles, but retains the dependence between inward displacement,
conditional variance and threshold phase in the count comparison.
Its finite certificates and the full effective continuation are supplied
in \cite{EffectiveSupplement}. The conclusions concern unrestricted
laws, not optimizers in a moment-restricted class.

\begin{theorem}\label{cut31:global:main}
For every integer \(n\ge31\),
\[
 C_n<\CE.
\]
This is a computer-assisted statement, using the stated analytic
inputs and the finite-certificate implications below.
\end{theorem}

\paragraph{Proof inputs.}
The continuation uses unrestricted attainment and reflection,
same-law contact and support bounds, the charged Fourier enclosure,
nonnegativity of the gap defect, and the conditional count inequalities.
Their analytic foundation is given below and in
\cite{EffectiveSupplement}. The finite inputs are the positive-budget
certificates at 27--30, the complete sharp covers at 31--34, and the
all-index continuation from 35. The positive-budget enclosure is used
in its charged form, including at index 27. The theorem incorporates
these previously established implications at their stated analytic
and finite-arithmetic scope; the supplement records the complete premises.

\subsection{A positive budget before the sharp budget}
Put \(a_n=C_n-\CE\). The finite continuation first proves
\begin{equation}\label{cut31:positive-ladder}
 a_{27}<\frac2{45},\qquad
 a_{28}<\frac9{280},\qquad
 a_{29}<\frac3{145},\qquad
 a_{30}<\frac1{100}.
\end{equation}
These are positive ceilings, not sharp estimates at indices
27--30. Their role is to improve the available contact allowance
before attempting the zero budget at 31.
For a contrary attained maximum at index \(n\), with
\(na_n\ge K\) and predecessor bound \(a_{n-1}<A\),
\[
 n(C_{n-1}-C_n)_+\le\max\{0,nA-K\}.
\]
The precise allowances in the finite certificates are
\[
\begin{array}{c|ccccc}
 n&27&28&29&30&31\\\hline
 K&6/5&9/10&3/5&3/10&0\\
 q&813/2000&31/90&93/280&93/290&31/100.
\end{array}
\]
The first row uses \(C_{26}\le469/1000\) and \(\CE>819/2000\).
Every subsequent allowance follows from the preceding strict
ceiling. Thus the argument is a finite descent in the permissible
excess; it assumes no monotonicity of \(C_n\).

The charge used in a positive-budget contact estimate is conservative:
if \(n(C-\CE)\ge K\), then
\[
 n(C-\CE)\ge\max\{K,n(C-40974/100000)\}.
\]
Splitting at \(C=40974/100000+K/n\) makes this charge affine
on each side. This permits the same corner-enclosure argument
as at the zero budget, without discarding the score dependence.

\subsection{One contrary law and a closed moment cover}
Suppose \(C_{31}\ge\CE\). Lemma~\ref{cut49:global:attainment}
supplies one unrestricted global maximizing law and a positive
closed threshold; reflection is performed before acquiring any
resources. In particular \(\rho\ge1\). The preceding ceiling
at 30 gives contact allowance \(31/100\).
Lemma~\ref{cut49:global:smoothing} also gives a uniform score cap.
Indeed \(\CE<37/90\), \(\sqrt{31/30}<61/60\) and
\(\sqrt{30}>5\) imply
\[
 C_{31}<\frac{61}{60}
 \left(\frac{37}{90}+\frac1{100}+\frac1{450}+\frac2{3750}\right)
 =\frac{193919}{450000}<\frac{432313}{10^6}.
\]
No separate barrier theorem is required for this cap.

The sharp 31 certificate excludes this same law on each of the
eight closed intervals
\begin{equation}\label{cut31:moment-cover}
\begin{gathered}
 [1,11/10],\ [11/10,6/5],\ [6/5,13/10],\ [13/10,7/5],\\
 [7/5,3/2],\ [3/2,33/20],\ [33/20,7/4],\ [7/4,7361/4000].
\end{gathered}
\end{equation}
For every interval, all 100 closed probability panels are either
necessarily impossible or covered by an acquired resource and
count comparison. The support, moments, endpoint bounds and noise
budget in that comparison all belong to the fixed contrary law.
Shared endpoints remain included; there is no optimizer chosen
separately for each interval.

For the remaining moment range, Shevtsova's bound
\eqref{cut49:global:finitecaps} gives the decreasing envelope
\[
 T(\rho)=\min\{.469,\ .3322(1+.429/\rho),\ .3031(1+.646/\rho)\}.
\]
At \(\rho_*=7361/4000\),
\[
 T(\rho_*)=\frac{819}{2000}<\CE.
\]
Hence the closed cover and this tail exclude every possible
\(\rho\ge1\). This proves \(C_{31}<\CE\).

\subsection{The count inequality preserves one phase}
The decisive refinement of the finite comparison can be seen
before any arithmetic. Write the localized law as
\(X=\ell(B+U-p)\), with \(B\sim\operatorname{Bernoulli}(p)\) and
\(\E(U\mid B)=0\). Let \(W_j\) be the conditional noise sum
at count \(j\), \(\pi_j=\binom njp^j(1-p)^{n-j}\), and
\(F_k=\sum_{j\le k}\pi_j\). At phase \(u\), conditioning gives
\begin{align}\label{cut31:one-phase}
 &\Pp\{\textstyle\sum_i(B_i+U_i)\le k+u\}-\Phi(Z(u))\notag\\
 &\quad\le F_k-\Phi(Z(u))-\pi_k\Pp\{W_k>u\}
       +\sum_{j>k}\pi_j\Pp\{W_j\le k+u-j\},
\end{align}
where \(Z(u)=(k+u-np)/\sqrt{np(1-p)+\lambda}\) and
\(\lambda=n\E U^2\).
The Gaussian term, the outgoing own-count mass and the incoming
count tails must be evaluated at the same phase. Maximizing them
independently loses precisely the compensation needed near the
sharp level. Polynomial minorants for the outgoing mass and
Cantelli or exponential bounds for the incoming tails reduce
\eqref{cut31:one-phase} to an explicit polynomial on a closed
phase interval. The resource inequalities constrain its
coefficients before its maximum is enclosed.
The strict outgoing event is the complement of the closed
own-count event, including atoms. At \(\lambda=0\), the law is
Bernoulli and Schulz's strict theorem applies directly.

Here is the common-phase estimate in a form independent of the
chosen minorants. Fix a closed interval \([l,h]\), and nonnegative
smooth incoming envelopes \(B_j(u)\ge\Pp\{W_j\le u-(j-k)\}\).
Suppose, on the whole parameter cell,
\[
 0\le\underline\pi_k\le\pi_k,\quad \pi_j\le m_j,\quad
 B_j'\le L_j,\quad (\Phi\circ Z)'\ge g,
\]
and set
\[
 B_0\ge F_k-\Phi(Z(l)),\qquad
 I_0\ge\sum_{j>k}m_jB_j(l),\qquad L=\sum_{j>k}m_jL_j.
\]
For each shift \(a\) in a chosen nonempty finite subset of \((0,\infty)\),
let a complete finite list of
quartic cases satisfy, for every actual law and phase,
\[
 \Pp\{W_k>u\}\ge Q_{a,c}(u)/(M_a a^4),\qquad M_a>0,
\]
for at least one applicable case \(c\).
The shifted-quartic inequality of
Section~\ref{cut49:shortloc:section} supplies elementary minorants;
the refined cases also retain the acquired moment budgets.

\begin{lemma}[Common phase]\label{cut31:common-phase}
The signed error in \eqref{cut31:one-phase} is at most
\[
 \min_a\max_c\sup_{u\in[l,h]}
 \left\{B_0+I_0+(L-g)(u-l)
       -\frac{\underline\pi_k Q_{a,c}(u)}{M_a a^4}\right\}.
\]
\end{lemma}
\begin{proof}
Integrate the incoming and Gaussian derivative bounds from the
same \(l\) to \(u\). For the applicable quartic case,
\[
 \pi_k\Pp\{W_k>u\}
 \ge\pi_k\max\{0,Q_{a,c}(u)/(M_a a^4)\}
 \ge\underline\pi_k Q_{a,c}(u)/(M_a a^4).
\]
This is valid even when the polynomial is negative. The case may
depend on the law and phase, so take the maximum over cases and
the supremum over the whole interval. Each shift gives a valid
bound; take their minimum. Only smooth envelopes were
differentiated, never the actual distribution function.
\end{proof}

The remaining supremum is polynomial. For
\(P(u)=\sum_{r=0}^4q_ru^r\), put \(w=h-l\) and
\[
 c_s=w^s\sum_{r=s}^4\binom rsq_rl^{r-s},\qquad
 b_j=\sum_{s=0}^j\frac{\binom js}{\binom4s}c_s.
\]
Then
\[
 P(l+wt)=\sum_{j=0}^4 b_j\binom4j t^j(1-t)^{4-j},\qquad
 \sup_{[l,h]}P\le\max_{0\le j\le4}b_j.
\]
The basis is nonnegative and sums to one. Rational upper
enclosures of its coefficients therefore prove the inequality
on a whole closed interval; a subdivision retains every shared
endpoint. This is the finite algebraic step, not a numerical
search for a maximizing phase.

\subsection{Actual-index bridges and the infinite range}
The continuation supplies three separate certificate implications
\begin{equation}\label{cut31:bridges}
 C_{m-1}<\CE\ \Longrightarrow\ C_m<\CE,
 \qquad m=32,33,34.
\end{equation}
They have the same closed moment cover \eqref{cut31:moment-cover}.
In the resource transfer, the physical defect and noise prices
scale by \(31/m\); the binomial probabilities, threshold
normalization and count covers use the actual index \(m\).
Thus \eqref{cut31:bridges} consists of three actual-index proofs,
not relabeled 31 certificates. Discharge their predecessor
assumptions in the order \(31,32,33,34\).

The effective continuation in \cite{EffectiveSupplement} also
proves \(C_n<\CE\) for every \(n\ge35\). Its finite descent
uses excess budgets \(.54,.36,.18\) at indices 32,33,34,
then sharp count comparisons at 35--39; the remaining indices
use the retained all-index theorem from 40. The finite 36--39
range is handled by a least-failure argument. The same-law resource and
common-phase inequality above are the essential new estimates.
Joining this all-index tail with the four sharp finite indices
proves Theorem~\ref{cut31:global:main}. Neither this joining step
nor the finite descent asserts monotonicity or eventual
two-point optimality.

\clearpage
\section{The result and the organizing ideas}\label{cut49:intro:chapter}
Let $\mathcal P$ be the probability laws on $\mathbb R$ with mean zero,
variance one and finite third absolute moment $\rho(P)=\mathbb E_P|X|^3$.
For iid variables with law $P$, set $S_n=X_1+\cdots+X_n$ and
\[
 C_n(P)=\frac{\sqrt n}{\rho(P)}\sup_{H\in\mathbb R}
 \left|\Pr_P(S_n\le H)-\Phi(H/\sqrt n)\right|,
 \qquad C_n=\sup_{P\in\mathcal P}C_n(P).
\]
Here $H$ is a raw sum threshold, $\Phi$ is the standard normal CDF,
and $\phi=\Phi'$. All terminating decimals in the proof are exact rationals.

\begin{theorem}\label{cut49:global:main}
For every integer $n\ge49$,
\[
 C_n<C_E,\qquad C_E=\frac{3+\sqrt{10}}{6\sqrt{2\pi}}.
\]
\end{theorem}

The argument is easiest to read as three ideas.
\begin{enumerate}[leftmargin=*]
\item \emph{A first failure is easier than an arbitrary large index.}
Once the case $n=49$ is proved, the first later failure has a good
predecessor. Smoothing that predecessor automatically bounds its excess
by $.29/n$. This supplies the hypotheses of a local exclusion, so the
same argument proves the entire infinite tail.
\item \emph{Spend one budget only once.}
At an extremizer, excess score, threshold displacement, nonbinary
defect and probability leaking into a support gap are nonnegative
terms of one exact identity. Moment, support and concentration estimates
refine that common budget. They do not choose a different law for each
term. The many numerical passes are instances of one enclosure theorem.
\item \emph{Compare each noisy law with its own Bernoulli law.}
Writing $Y=B+U$ with $\mathbb E(U\mid B)=0$ gives a count and a centered
error. One exact score identity and two tail estimates cover every
noise and phase regime. Positive noise lowers the signed score;
zero noise is already Bernoulli. Schulz's theorem treats every
$0<p<1$, so no finite binomial subproblem is needed.
\end{enumerate}

Section~\ref{cut49:global:route} proves the global implication from explicit
finite and local propositions. The remaining sections prove those
propositions and specify their exact arithmetic. This order separates
the short mathematical mechanism from the estimates needed to certify
its numerical hypotheses. The proof does not invoke any earlier
project cutoff, an eventual-tail theorem, or monotonicity of $C_n$.

\subsection{External inputs}
Shevtsova's finite iid estimates~\cite{Shevtsova2013} give, in this normalization,
\begin{equation}\label{cut49:global:finitecaps}
 C_n(P)\le\min\{.469,\ .3322(1+.429/\rho),\ .3031(1+.646/\rho)\}.
\end{equation}
The elementary numerical enclosures used below are
\[
 .3988<\phi(0)<.4,\qquad .4095<C_E<.40974,
 \qquad 3.162277660<\sqrt{10}<3.162277661.
\]
Their rational checks are part of the arithmetic kernel. In particular,
a law attaining $C_n\ge C_E$ has $1\le\rho<1.85$ before any
higher-moment or bounded-support assumption is made.

Schulz's Theorem~1~\cite{Schulz2016}, for every $n\ge1$ and $0<p<1$,
states
\begin{equation}\label{cut49:global:schulz}
 \sup_{x\in\mathbb R}\left|\Pr(K\le x)
 -\Phi\left(\frac{x-np}{\sqrt{np(1-p)}}\right)\right|
 <\frac{C_E\{p^2+(1-p)^2\}}{\sqrt{np(1-p)}},
 \qquad K\sim\operatorname{Bin}(n,p).
\end{equation}
Thus every standardized Bernoulli law has score strictly below $C_E$.
The restriction $p\in[1/3,2/3]$ belongs to a sharper envelope in that
work; it is not a restriction of \eqref{cut49:global:schulz}.

From He--Cheng~\cite{HC} we use the finite Gaussian bound
(Appendix~A.1, (70)) in the effective large-index argument.
The fixed-index attainment and contact arguments are proved next;
the contact identity agrees with their Lemma~4.3.
For the unrestricted fourth index we prove a separate Gaussian estimate
in Section~\ref{n4all:section}. All remaining analytic reductions
in this manuscript are proved here.

\section{Fixed-index extremizers and the contact interface}
\subsection{Attainment, including the critical value}
\begin{lemma}\label{cut49:global:attainment}\label{cut49:found:orientation}
For every fixed $n\ge1$, $0<C_n\le\sqrt n$, and a law $P\in\cP$ and
threshold $H$ attain $C_n$ as a positive closed discrepancy.
\end{lemma}
\begin{proof}
A negative closed discrepancy reflects into a positive strict discrepancy.
Approach that threshold from below; continuity of $\Phi$ gives the same
supremum using positive closed discrepancies. Since $\rho\ge1$, $C_n\le\sqrt n$; the symmetric binary law at
$H=n$ gives positivity. Choose positive scores $r_j\to C_n$ with
$r_j\ge c>0$. Then $\rho_j\le\sqrt n/c=:M$. Tightness gives a subsequence
$P_j\Rightarrow P$, and we may also assume $\rho_j\to b$. The estimates
\[
 \int_{|x|>A}x^2\,dP_j\le M/A,\qquad
 \int_{|x|>A}|x|\,dP_j\le M/A^2
\]
preserve the first two moments under weak convergence. Lower
semicontinuity gives $\rho(P)\le b$, so $P\in\cP$.
Chebyshev's inequality on a negative threshold, and the Gaussian upper
tail on a positive threshold, give $|H_j|\le n^{3/4}/\sqrt c$.
Pass to $H_j\to H$. Since $n$ is fixed, weak convergence of products
and the continuous sum map give $P_j^{*n}\Rightarrow P^{*n}$.
For every $\varepsilon>0$, eventually $H_j\le H+\varepsilon$;
closed-set Portmanteau followed by $\varepsilon\downarrow0$ gives
\[
 F_{n,P}(H)-\Phi(H/\sqrt n)\ge C_nb/\sqrt n.
\]
Consequently $C_n\ge C_nb/\rho(P)\ge C_n$. Equality proves attainment
and $\rho(P)=b$. Thus third-moment convergence is a conclusion of
extremality, not an unjustified uniform-integrability assumption.
This proof includes $C_n=\CE$; the strict-violation version is
Lemma 4.1 of \cite{HC}.
\end{proof}

\subsection{Contact at every support point}
For $n\ge2$, fix an attaining pair $(P,H)$, and write
\[
 C=C_n,\quad F=F_{n,P}(H),\quad z=H/\sqrt n,\quad
 \rho=\E|X|^3,\quad\gamma=\E X^3,\quad\beta=\E[X|X|].
\]
Here $\beta$ is a signed second moment; it is not the third absolute
moment denoted by $\beta$ in \cite{HC}.

\begin{lemma}[Contact]\label{cut49:global:contact}\label{cut49:found:majorant}
Define
\begin{align}
 \rho I(y)={}&n^{3/2}\{F_{n-1,P}(H-y)-F\}+ny\phi(z)
 +\tfrac12H\phi(z)(y^2-1)\notag\\
 &+\tfrac32C\rho(y^2-1)-C\{|y|^3-\rho-3y\beta\}.
 \label{cut49:found:influence}
\end{align}
Then $I\le0$ on $\R$ and $I=0$ on the full closed support of $P$.
The support is compact. Equivalently, for $\nu=\mathcal L(H-S_{n-1})$,
$\nu([y,\infty))\le\mathfrak q(y)$, with equality on $\supp P$, where
\begin{align}
 n\mathfrak q(y)={}&nF+\tfrac12\{z\phi(z)+C\rho/\sqrt n\}
 -(\sqrt n\phi(z)+3C\beta/\sqrt n)y\notag\\
 &-\tfrac12\{z\phi(z)+3C\rho/\sqrt n\}y^2
 +C|y|^3/\sqrt n.\label{cut49:found:q}
\end{align}
\end{lemma}
\begin{proof}
For fixed $y$, contaminate $P$ by
$P_\varepsilon=(1-\varepsilon)P+\varepsilon\delta_y$, holding $H$ fixed.
Its mean and variance are
\[
 m_\varepsilon=\varepsilon y,\qquad
 v_\varepsilon=1+\varepsilon(y^2-1)-\varepsilon^2y^2.
\]
Write $r_\varepsilon=\E_{P_\varepsilon}|X-m_\varepsilon|^3$.
After standardization, the signed score
\[
 \frac{\sqrt n\,v_\varepsilon^{3/2}}{r_\varepsilon}
 \left[F_{n,P_\varepsilon}(H)-
 \Phi\!\left(\frac{H-nm_\varepsilon}{\sqrt{nv_\varepsilon}}\right)\right]
\]
is at most $C$ and equals $C$ at zero. The raw CDF is a polynomial
in $\varepsilon$, whose derivative there is
$n\{F_{n-1,P}(H-y)-F\}$. Dominated differentiation, using the finite
second moment, gives $r'_0=|y|^3-\rho-3\beta y$.
The derivative of the displayed score is precisely $I(y)$, so $I(y)\le0$.
Integration against $P$ gives $\int I\,dP=0$: in particular
$\int F_{n-1,P}(H-y)\,dP(y)=F$.
The closed CDF is upper semicontinuous. Hence $I$ is upper
semicontinuous, and the closed set $\{I=0\}$ has full $P$ measure;
it contains all of $\supp P$. Rearrangement gives
\eqref{cut49:found:q}. Finally,
$\rho I(y)\le-C|y|^3+O(y^2)$ at this fixed pair, since the CDF term is
bounded and $C>0$. Equality on the support excludes both sufficiently
distant rays. No uniform radius is claimed yet.
\end{proof}

The pointwise equality controls even a support point carrying arbitrarily
little mass. This is why weak closeness to a binary law will not be the
last step of the proof.

\section{A finite base and a first-failure induction}\label{cut49:global:route}
Put $a_n=C_n-C_E$ and $\delta_n=n(C_{n-1}-C_n)_+$.
Different indices may have different extremizers. This section uses
only their scalar suprema.

\subsection{A support-free smoothing estimate}
\begin{lemma}\label{cut49:global:smoothing}
For every integer $n\ge2$,
\[
 C_n\le\sqrt{\frac n{n-1}}\,C_{n-1}
       +\frac{\sqrt n}{15(n-1)^{3/2}}
       +\frac{.08\sqrt n}{(n-1)^2}.
\]
No bounded-support or maximizing-law assumption is required.
\end{lemma}
\begin{proof}
Fix any mean-zero, variance-one iid law with finite
$\rho=\mathbb E|X|^3\ge1$, and let $v=n-1$. Conditional on the last
summand, the defining bound at index $v$ gives
\[
 \left|\mathbb P(S_n\le H)-\mathbb E\Phi((H-X)/\sqrt v)\right|
       \le C_{n-1}\rho/\sqrt v.
\]
Write $F_v(t)=\Phi(t/\sqrt v)$. Taylor's theorem through degree two,
using the actual mean and variance, gives
\[
 \left|\mathbb E F_v(H-X)-F_v(H)-\tfrac12F_v''(H)\right|
       \le\frac{\phi(0)\rho}{6v^{3/2}}<\frac{\rho}{15v^{3/2}}.
\]
Here $\|\phi''\|_\infty=\phi(0)<.4$. Indeed
$\phi''(z)=(z^2-1)\phi(z)$; on $z^2\le1$ the assertion is immediate,
and on $z^2\ge1$ the maximum relative factor is $2e^{-3/2}<1$.
This globally bounded remainder needs only the third absolute moment.

The heat identity $\partial_vF_v=F_v''/2$ gives
$\partial_v^2F_v=F_v^{(4)}/4$. Taylor's theorem in the variance
variable therefore implies
\[
 |F_{v+1}(H)-F_v(H)-\tfrac12F_v''(H)|
 \le\frac{\|\phi'''\|_\infty}{8v^2}<\frac{.08}{v^2}.
\]
For the last bound, Gaussian Fourier inversion gives
$\|\phi'''\|_\infty\le\pi^{-1}\int_0^\infty
t^3e^{-t^2/2}\,dt=2/\pi<.64$.
Combine these estimates, multiply by $\sqrt n/\rho$, and use
$\rho\ge1$. The result is uniform in $H$ and in the law, so taking
both suprema proves the lemma. The bound covers closed events and
their left limits by the same uniform comparison.
\end{proof}

All three terms obtained by substituting
$C_{n-1}<.40974+.7/(n-1)$ decrease with $n$. At $n=46$, rational
square-root enclosures therefore give $C_n<.432$. This is the
current-score cap at a first failure of the $.7/n$ ceiling.

\subsection{The two inputs from the detailed calculation}
The resource enclosure and support estimates prove these five
\emph{supported exclusions}. Each row says that no attained, positively
oriented maximizing pair can have all the listed properties:
\begin{equation}\label{cut49:global:finitecontracts}
\begin{array}{l|c|c|c|c|r}
 &n&|X|\le&C\le&na_n\ge&\text{passes}\\\hline
\mathrm{A40}&40&3.3&.469&1.5&5\\
\mathrm{P41}&41\!:\!49&2.9&.469&1.5&6\\
\mathrm{A45}&45&3.5&.4431&.7&12\\
\mathrm{P46}&46\!:\!98&3.5&.432&.7&13\\
\mathrm{A46}&46&3.4&.42496&.54&4+3
\end{array}
\end{equation}
Every pass preserves the same-law constraints and gives a new upper
bound for the nonnegative defect $D$. Its final bound is negative.
The first two rows use 64 closed cells per pass. A45 and P46 use
18432; A46 uses four signed passes with 18432 cells, followed by
three probability passes with 96 cells. Section~\ref{cut49:bar46:entry}
gives its complete additional calculation.
The upper endpoints 49 and 98 describe those individual certificates;
we only apply them through 45 and 50. These arithmetic exclusions do
not assert an unrestricted ceiling on their own.

The all-index local input is the following precise proposition.
\begin{proposition}[Local exclusion]\label{cut49:global:local}
For any integer $n\ge47$, the conditions
\[
 0\le a_n\le.7/n,\qquad \delta_n\le.19,\qquad
 C_n<.4247,\qquad C_{n-1},C_{n-2}<.4253
\]
cannot hold simultaneously.
\end{proposition}
\begin{proof}[Proof from the estimates established below]
Take the attained positive closed maximum. The all-threshold support
estimate bounds its entire support by 3.84. The resource schedule then
acquires 83 probability panels covering $p\in[.304,.47]$, together with
their moments, signed thresholds, conditional variances and inward-gap
resources. Support and concentration refinements are applied before
the resulting stronger moments are used.

Write the same law as $X=\ell(B+U-p)$ with
$\mathbb E(U\mid B)=0$. If $\lambda=n\mathbb E U^2>0$, the complete
local cover of Section~\ref{cut49:shortloc:section} gives
$\rho_B(C-C_B^+(n,p,k))<-.001\lambda$ at an actual nearest count.
Schulz's all-parameter theorem gives $C_B^+<C_E$, a contradiction.
If $\lambda=0$, then $U=0$ almost surely and the law is itself
Bernoulli; the same published theorem applies directly to its whole
Kolmogorov score. This uses no division by noise or limiting argument.

All entry, support, concentration and comparison estimates are uniform
for $n\ge47$. In the initial arithmetic, the apparent upper endpoint 98
occurs only as $.4095+K/98$, and $K=0$ for this application. Its other
occurrence checks the calculation floor, not the actual index.
Weighted resources such as $nD$ and $n\mathbb E U^2$ are retained
throughout; no unweighted estimate is multiplied by 47 to cover larger
$n$. The excess assumption $.7/n$ is used in the exterior contact
calculation and is therefore included explicitly.
\end{proof}

\subsection{Only six sharp-ceiling indices are needed}
We first prove $a_j<1.5/j$ for $40\le j\le45$. If it failed at 40,
$C_{39}\le.469$ and $C_E>.4095$ would give $\delta_{40}<.88$.
The support row $(n,K,R,q)=(40,1.5,3.3,.88)$ and A40 exclude this.
At a first failure $41\le j\le45$,
$\delta_j<1.5/(j-1)\le1.5/40$. The support row on $[41,49]$
and P41 exclude it.

We now prove
\begin{equation}\label{cut49:global:finiteceiling}
                     a_j<.7/j\qquad(45\le j\le50).
\end{equation}
A failure at 45 would give
\[
 \delta_{45}<45(1.5/44)-.7<.8341,
 \qquad C_{45}<.40974+1.5/45<.4431.
\]
The appropriate support row gives $|X|<3.5$ and A45 excludes the law.
At a first failure $46\le j\le50$,
$\delta_j<.7/(j-1)\le.7/45$, and smoothing gives $C_j<.432$.
The support row on $[46,98]$ and P46 again exclude it.
This proves \eqref{cut49:global:finiteceiling} using no later index.

For $47\le j\le50$, it supplies the local score caps because
\[
 .40974+.7/47<.4247,\qquad .40974+.7/45<.4253.
\]
We strengthen just one auxiliary index:
\begin{equation}\label{cut49:global:ceiling46}
                        a_{46}<.54/46.
\end{equation}
Indeed, if this failed, the preceding finite ceiling gives
\[
 \delta_{46}<46(.7/45)-.54=79/450,
 \qquad C_{46}<.40974+.7/46<.42496.
\]
The support row $(46,46,.54,3.4,79/450)$ gives $|X|<3.4$.
The supported exclusion A46 is now applicable, a contradiction.
This argument uses only the finite ceiling already established.

If $a_{49}\ge0$, Proposition~\ref{cut49:global:local} forces
$a_{48}>a_{49}+.19/49>0$. Repeat at 48 and 47 to obtain
\[
 a_{46}>.19\sum_{j=47}^{49}\frac1j>\frac{.54}{46},\qquad
 .19\sum_{j=47}^{49}\frac1j-\frac{.54}{46}
       =\frac{35419}{254251200}>0.
\]
This contradicts \eqref{cut49:global:ceiling46}. Thus $C_{49}<C_E$.

\subsection{The infinite tail follows from a least failure}
Suppose $n\ge50$ and $C_{n-1}<C_E$. With $s_n=\sqrt{n/(n-1)}$,
smoothing gives
\[
 na_n<C_E n(s_n-1)+s_n^3/15+.08s_n^3/\sqrt{n-1}.
\]
Every term on the right decreases with $n$: in particular
$n(s_n-1)=s_n^2/(s_n+1)$, an increasing function of $s_n$.
Since $s_n<1.011$, $\sqrt{n-1}\ge7$ and $n(s_n-1)<.51$,
\begin{equation}\label{cut49:global:firstfailurecap}
 na_n<.40974(.51)+(1.011)^3(1/15+.08/7)
        =\frac{50691940857}{175000000000}<.29.
\end{equation}

If a failure $C_n\ge C_E$ occurs after 49, choose its least index.
Then $\delta_n=0$, and \eqref{cut49:global:firstfailurecap} gives
$0\le a_n<.29/n<.7/n$ and $C_n<.40974+.29/50<.4247$.
Its predecessor is below $C_E<.4253$. If $n=50$, the second
predecessor is 48 and \eqref{cut49:global:finiteceiling} bounds it by
$.40974+.7/48<.4253$. For $n\ge51$, that predecessor is also
below $C_E$ by minimality. Proposition~\ref{cut49:global:local} now
contradicts the failure. This proves Theorem~\ref{cut49:global:main}.

\section{One support estimate, uniform in the threshold}\label{cut49:support:theorem}
\subsection{The scale correction}
Write $\phi_0=\phi(0)$.
For $n>1$ put $s_n=\sqrt{n/(n-1)}$ and
$e_n=n(s_n-1)-1/2$.
\begin{lemma}\label{cut49:support:scale}
For every integer $n\ge36$, one has $0<e_n<.384/n$.
\end{lemma}
\begin{proof}
Put $x=1/n$ and $c=.384=48/125$. Direct expansion gives
\[
 (1-x)(1+x/2+cx^2)^2-1
 =x^2\{(2c-3/4)-(c+1/4)x+(c^2-c)x^2-c^2x^3\}.
\]
The bracket is decreasing for $x\ge0$, since $0<c<1$.
At $x=1/36$ its exact rational value is $823/4050000>0$.
Squaring positive quantities gives
$s_n<1+x/2+cx^2$ for $0<x\le1/36$. Similarly
$(1-x)(1+x/2)^2-1=-3x^2/4-x^3/4<0$, which gives the lower bound.
\end{proof}

\subsection{A signed Gaussian envelope}
Define, for every real $H,t$ and $n>1$,
\[
 \mathcal H_n(H,t)=n^{3/2}\{\Phi((H-t)/\sqrt{n-1})-
             \Phi(H/\sqrt n)\}
       +nt\phi(H/\sqrt n)+\tfrac H2(t^2-1)\phi(H/\sqrt n).
\]
Write $v=n-1$, $s=\sqrt{n/v}$, $u=|t|$, and
\[
 A=\frac{u^3-3u}{6},\quad B=\frac{u^4-6u^2+3}{24\sqrt v},
 \quad D=\frac{u^5-15u}{120v}.
\]
\begin{lemma}[Fifth-order identity with a uniform finite remainder]
If $u>\sqrt3$, then, with $z=H/\sqrt v$,
\[
 \frac{\mathcal H_n(H,\pm u)}{s^3}
 =\pm A(1-z^2)\phi(z)+B\phi^{(3)}(z)
                  \mp D\phi^{(4)}(z)+R,
\]
where the sign of the remainder need not be fixed and
\[
 |R|\le E^{(5)}_v(u):=
 \frac{51u^6}{14400v^{3/2}}
 +\frac{51(3u^2-2)}{960v^{3/2}}+\frac{u}{8v^2}.
\]
This statement is uniform over every real $H$.
\end{lemma}
\begin{proof}
Let $F_w(x)=\Phi(x/\sqrt w)$, with primes denoting $x$ derivatives.
The heat identity $\partial_w F_w=F_w''/2$ and ordinary Taylor's theorem
give, with unlabelled derivatives at $(v,H)$,
\begin{align*}
 F_{v+1}&=F_v+F_v''/2+F_v^{(4)}/8+R_0,
 &|R_0|&\le\sup_{w\in[v,v+1]}\|F_w^{(6)}\|_\infty/48,\\
 F_{v+1}'&=F_v'+F_v'''/2+F_v^{(5)}/8+R_1,
 &|R_1|&\le\sup_{w\in[v,v+1]}\|F_w^{(7)}\|_\infty/48,\\
 F_{v+1}''&=F_v''+F_v^{(4)}/2+R_2,
 &|R_2|&\le\sup_{w\in[v,v+1]}\|F_w^{(6)}\|_\infty/8.
\end{align*}
Expand $F_v(H-t)$ through degree five, with remainder at most
$u^6\|F_v^{(6)}\|_\infty/720$. In
\[
 \frac{\mathcal H_n}{n^{3/2}}=
 F_v(H-t)-F_{v+1}(H)+tF_{v+1}'(H)
             -\frac{t^2-1}{2}F_{v+1}''(H),
\]
the first three derivative coefficients cancel. The next three are
\[
 \left(\frac t2-\frac{t^3}6\right)F_v'''
 +\frac{t^4-6t^2+3}{24}F_v^{(4)}
 +\left(\frac t8-\frac{t^5}{120}\right)F_v^{(5)}.
\]
The remaining absolute value is bounded by
\[
 \left\{\frac{u^6}{720}+\frac1{48}+\frac{u^2-1}{16}\right\}
       \sup_{w\in[v,v+1]}\|F_w^{(6)}\|_\infty
 +\frac u{48}\sup_{w\in[v,v+1]}\|F_w^{(7)}\|_\infty.
\]
Fourier inversion of the Gaussian, differentiated under its absolutely
integrable Fourier transform, gives
\[
 \|\phi^{(k)}\|_\infty
 \le\frac1\pi\int_0^\infty y^k e^{-y^2/2}\,dy.
\]
Integration by parts yields
$\|\phi^{(5)}\|_\infty\le8/\pi<2.55$ and
$\|\phi^{(6)}\|_\infty\le15\phi_0<6$.
Since $F_w^{(j)}=w^{-j/2}\phi^{(j-1)}$, the displayed remainder
becomes $s^3E^{(5)}_v(u)$ after multiplication by $n^{3/2}$.
The three retained derivative coefficients give the identity.
\end{proof}

\begin{lemma}[A lower tangent bound for the fourth Gaussian derivative]
For every real $z$,
\[
 \phi^{(4)}(z)\ge\phi_0(3-15z^2/2).
\]
\end{lemma}
\begin{proof}
The same Gaussian Fourier integral gives
$\phi^{(4)}(z)=\pi^{-1}\int_0^\infty
y^4e^{-y^2/2}\cos(yz)\,dy$.
Use the globally valid inequality $\cos w\ge1-w^2/2$ inside
this integral. Its fourth and sixth Gaussian moments are $3\phi_0$
and $15\phi_0$, respectively.
\end{proof}

\begin{lemma}[Uniform signed fifth-order envelope]\label{cut49:support:envelope}
Suppose $u>\sqrt3$ and $A,B,D\ge0$. If
$W=A-15D/2>0$, put
\[
 T=A-3D+\frac{9B^2}{4W}.
\]
If $T\ge A/2+2B+3D$, then, for every real $H$ and either sign,
\[
 \mathcal H_n(H,\pm u)\le s^3\{\phi_0 T+E^{(5)}_v(u)\}.
\]
\end{lemma}
\begin{proof}
For positive $t=u$ and $|z|\le1$, divide the retained terms by
$\phi_0$ and use the preceding tangent bound to obtain
\[
 A-3D-(A-15D/2)z^2+3B|z|
 \le A-3D+9B^2/(4W).
\]
Here $(1-z^2)\phi(z)\le(1-z^2)\phi_0$ and
$|\phi^{(3)}(z)|\le3|z|\phi_0$ on this interval.
For $|z|\ge1$ the first term is nonpositive, and the Fourier bounds
$\|\phi^{(3)}\|_\infty<2\phi_0$ and
$\|\phi^{(4)}\|_\infty\le3\phi_0$ give $2B+3D$.
For negative $t=-u$, the first term has global upper bound
$2e^{-3/2}A\phi_0<A\phi_0/2$, and the other terms give
$(2B+3D)\phi_0$. The hypothesis on $T$ covers both exterior cases.
The inequality $2e^{-3/2}<1/2$ follows from four positive terms of
the exponential series at $3/2$. Finally add the uniform remainder.
\end{proof}

\subsection{Finite bands and the exact support polynomial}\label{cut49:support:polynomial}
For integers $36\le N\le L$, use the simultaneous endpoint bounds
\[
 B_+=\frac{u^4-6u^2+3}{24\sqrt{N-1}},\quad
 D_+=\frac{u^5-15u}{120(N-1)},\quad
 D_-=\frac{u^5-15u}{120(L-1)}.
\]
For an infinite index range, set $D_-=0$ instead. Put
\[
 W_*=A-15D_+/2,\qquad
 T_*=A-3D_-+\frac{9B_+^2}{4W_*}.
\]
The proof of Lemma~\ref{cut49:support:envelope} applies uniformly on the whole band
if $W_*>0$ and $T_*\ge A/2+2B_++3D_+$. In particular the favorable
term is frozen at the upper index endpoint, and the adverse terms at
the lower one. No monotonicity of a signed coefficient is assumed.

Let $(P,H)$ be an attained positive closed maximizing pair, standardized
to mean zero and variance one. Write $C=C_n\ge C_E$,
$\rho=\mathbb E|X|^3$, $\beta=\mathbb E[X|X|]$, and
$d=n(C_{n-1}-C_n)_+$. The finite contact/predecessor bound is
\[
 \rho I(t)\le\mathcal H_n(H,t)+
 C\{-|t|^3+3\rho t^2/2+3\beta t+\rho e_n\}+\rho s_nd,
\]
with $I=0$ on the entire closed support. Suppose also, for a finite
band, that $C_n-C_E\ge K/n$, where $K>0$. The published finite cap gives
\[
 c=.4095+K/L<C_n,\qquad
 \rho<r=\min\left\{\frac{.3031(.646)}{c-.3031},
                        \frac{.3322(.429)}{c-.3322}\right\}.
\]
For the unrestricted nonnegative case put $K=0$, $c=.4095$.
The sign-moment inequality
$\rho\ge(1+\beta^2)/\sqrt{1-\beta^2}$ supplies $|\beta|<b$ whenever
$(1+b^2)^2>r^2(1-b^2)$ and $0<b<1$. All these statements concern the
same current law. The lower endpoint cannot replace $L$ in $K/L$.

When a preceding resource enclosure supplies sharper actual bounds
$\rho\le r$ and $|\beta|\le b$, these may instead be used in the
same polynomials. Their validity is the only extra premise.

Use rational $\ell^2\le N-1$ and $S^2>N/(N-1)$. For $u=|t|$ define
\[
 Q(u)=-u^3+3ru^2/2+3bu+.384r/N.
\]
First check $Q<0$, before replacing $C$ by its lower bound $c$.
Replace $1/\sqrt{N-1}$ by $1/\ell$ in the adverse remainder and in
the term $2B_+$ of the domination gate. Write $E_*$ for that rational
remainder. Strict negative contact for every $n$ in the band follows
from the polynomial inequality
\[
 \frac25S^3\left\{(A-3D_-)W_*+
              \frac{(u^4-6u^2+3)^2}{256(N-1)}\right\}
 +S^3E_*W_*+\{cQ+rSq\}W_*<0.
\]
The second polynomial gate is $W_*>0$. The third is
\[
 \{A/2-3D_--3D_+-2B_+\}W_*+
         \frac{(u^4-6u^2+3)^2}{256(N-1)}>0.
\]
This is the cleared-denominator version of the domination hypothesis.
The bound on $T_*$ is positive, so the replacement $\phi_0<.4$ and
$s_n<S$ both have their correct direction.

The exact kernel tests these polynomials on all adjacent closed intervals
of length $.05$ from $R$ to $4.5$. Positive Bernstein coefficients prove
their signs on every point of each cell, including endpoints. On the
rays $u\ge4.5$ use the imported finite bound
$|\mathcal H_n(H,t)|\le u^3/15+u/5+1/2$.
The resulting cubic and its first three derivatives are strictly negative
at $4.5$, as is the coefficient $Q$ before replacing $C$ by $c$.
This excludes the whole rays. Thus, if $d\le q$, actual support is in
$(-R,R)$, including exclusion of arbitrarily small exterior atoms and
exterior accumulation points.

\subsection{The support rows used in this proof}
For the all-index local argument, take $K=0$, $R=3.84$, $q=.19$ on
\[
[45,49],\ [50,59],\ [60,89],\ [90,134],\ [135,199],\ [200,\infty).
\]
Only indices at least 47 are used. The last row uses $D_-=0$;
there is no upper index restriction. The exact checker reconstructs
all closed $u$ cells and the two exterior rays. Missing an index band
or the infinite row is an acceptance failure.

The finite base uses five further rows:
\[
\begin{array}{c|c|c|c}
[N,L]&K&R&q\\\hline
{[40,40]}&1.5&3.3&.88\\
{[41,49]}&1.5&2.9&1.5/40\\
{[45,45]}&.7&3.5&.8341\\
{[46,98]}&.7&3.5&.7/45\\
{[46,46]}&.54&3.4&79/450
\end{array}
\]
The second and fourth are used only on $41\le n\le45$ and
$46\le n\le50$, respectively. Retaining their originally certified
upper endpoints makes their inequalities weaker, so restricting the
application introduces no new arithmetic assertion. These are support
implications, not yet exclusions; the resource calculation supplies
the latter.

\subsection{One bounded-threshold series}\label{cut49:support:series}
Put $c_k=(-1)^k/[2^kk!(2k+1)]$, and let $(x)_j$ denote the rising
factorial. Define the following rational polynomials:
\begin{align*}
 P_j(H,t)={}&\sum_{k=0}^{j+1}
 \frac{c_k(k+1/2)_{j+1-k}}{(j+1-k)!}(H-t)^{2k+1}
       -c_{j+1}H^{2j+3}\\
 &+\frac{(-1)^{j+1}tH^{2j+2}}{2^{j+1}(j+1)!}
       +\frac{(-1)^jH^{2j+1}(t^2-1)}{2^{j+1}j!}.
\end{align*}
In particular,
\[
 P_0=(t^3-3t)/6,\qquad
 P_1=\frac{3t-t^3}{4}H^2+\frac{t^4-6t^2+3}{8}H
                  +\frac{-t^5+10t^3-15t}{40}.
\]
Fix $h>0$, $E(x)=(193/71)^{\lceil x\rceil}$, and
\[
 B_r=\frac{7(h+r)}6E((h+r)^2/6)
       +\left(h+r+\frac{h(r^2+1)}8\right)E(h^2/8).
\]
\begin{lemma}\label{cut49:support:finite-series}
For every $n\ge4$, $|H|\le h$ and $|t|\le r$,
\[
 \left|\mathcal H_n/\phi(0)-\sum_{j=0}^5P_j/n^j\right|
       \le4^7B_r/n^6.
\]
\end{lemma}
\begin{proof}
Put $z=1/n$. Integration of the Gaussian exponential series gives
$\mathcal H_n/\phi(0)=S(z)/z$, where
\begin{align*}
 S(z)={}&\sum_{k\ge0}c_k\{(H-t)^{2k+1}(1-z)^{-k-1/2}
                                  -H^{2k+1}\}z^k\\
 &+te^{-H^2z/2}+\tfrac12zH(t^2-1)e^{-H^2z/2}.
\end{align*}
Its constant coefficient vanishes, and expansion of the binomial factor
gives $[z^{j+1}]S=P_j$. At $z_0=1/4$, the sum obtained by replacing
all coefficients by their absolute-value majorants is at most $B_r$:
use $(1-z_0)^{-1/2}<7/6$, $1/(2k+1)\le1$, and $e<193/71$.
The latter bound follows, for example, by summing the exponential series
through degree six and bounding the rest by a geometric series with
ratio $1/8$. All products and rearrangements are therefore absolutely
convergent. The discarded terms have degrees at least seven in $S$,
so their sum at $0\le z\le z_0$ is at most $(z/z_0)^7B_r$.
Divide by $z$ to get the result. This proves a finite uniform bound,
including every threshold in the stated interval and every $n\ge4$.
\end{proof}

\subsection{A common contact test for radius and strip refinement}
\label{cut49:support:panels}
Suppose $n\ge47$, $0\le C-C_E\le.7/n$, $\delta_n\le.19$ and
$|X|\le R_0$. A previously acquired row gives
\[
 \rho\le r,\quad 0\le b_l\le\beta\le b_u<1,\quad |H|\le h.
\]
Define $Q(t)=-|t|^3+3\rho t^2/2+3\beta t$ and
$L_P(t)=C_EQ(t)+\phi(0)(t^3-3t)/6$.
For $t>0$, upper polynomials for $L_P(t),Q(t)$ are
\[
 (1/15-.4095)t^3+1.5(.40974)rt^2+(3(.40974)b_u-.3988/2)t,
 \quad -t^3+1.5rt^2+3b_ut.
\]
For $t<0$, write $u=|t|$ and use
\[
 (-.3988/6-.4095)u^3+1.5(.40974)ru^2+(.2-3(.4095)b_l)u,
 \quad -u^3+1.5ru^2-3b_lu.
\]
These coefficient choices are valid because $b_l\ge0$ is checked first.
On any closed cell of one sign, let $L_J$ be an upper bound for the
first polynomial and let $Q_J\ge0$ bound the positive part of the second.
For $j=1,2,3$, take $m_j$ to be the positive part of the largest tensor
Bernstein coefficient of $P_j$ on $[-h,h]\times J$. For $j=4,5$, the
absolute monomial bound
$\sum_{a,b}|[H^at^b]P_j|h^aR_J^b$ suffices, where
$R_J=\max_{t\in J}|t|$. The support contact is then bounded above by
\begin{equation}\label{cut49:support:paneltest}
 L_J+\frac{.7Q_J}{47}
 +.4\left(\sum_{j=1}^5\frac{m_j}{47^j}
                +\frac{4^7B_{R_J}}{47^6}\right)
 +r\left(.19S+\frac{.384(.425)}{47}\right),
 \qquad S>\sqrt{47/46}.
\end{equation}
Indeed substitute the finite series in contact, split $C=C_E+(C-C_E)$,
and use $C<.40974+.7/47<.425$. Every error retained in
\eqref{cut49:support:paneltest} is nonnegative and decreases with $n$.
A strictly negative rational result therefore excludes the entire cell
for every integer $n\ge47$, including atoms and accumulation points.

For a coordinate strip, conditional centering gives
\[
 X=\ell(Y-p),\quad
 \ell=\sqrt{(1-\theta)/(p(1-p))},\quad
 \beta=(1-\theta)(1-2p)+\kappa,\quad |\kappa|\le\theta.
\]
If $p\in[p_l,p_u]\subset(0,1/2)$ and $\theta\le\Theta$, put
$\ell_l\le\sqrt{(1-\Theta)/(p_u(1-p_u))}$.
The forbidden strip $Y\notin(-.55,1.55)$ maps into
\[
 X\le-\ell_l(p_l+.55),\qquad X\ge\ell_l(1.55-p_u).
\]
Apply \eqref{cut49:support:paneltest} between these thresholds and $\pm R_0$.
An interval beyond the known radius is empty; an equal endpoint is a
singleton that is still checked. If only an upper $\beta$ bound is given,
use $b_l=(1-\Theta)(1-2p_u)-\Theta$, after checking its nonnegativity.

The selected consequences used in the resource schedule are:
\[
\begin{array}{l|c|c|r}
\text{input rows}&\text{prior radius}&\text{conclusion}&\text{closed cells}\\\hline
10\text{ central preparation rows}&3.84&|X|<2.5&1080\\
45\text{ central rows}&2.5&-.55<Y<1.55&1431\\
38\text{ left rows}&2.7&-.55<Y<1.55&3073
\end{array}
\]
The respective largest contact bounds are below $-.084$, $-.097$ and
$-.01658$. These implications acquire no row by themselves: each is
used only after its input row is proved by the preceding resource step.
The left radii are applications of the all-threshold support polynomial,
with $(N,L,K,q)=(47,\infty,0,.19)$ and the actual moment caps
\[
 (r,b,R)=(1.390,.430,3),\qquad (1.2563,.3926,2.7).
\]
Thus the same heat proof supplies the initial radius, both left
contractions, and the five finite-base rows. The bounded-threshold
series supplies the central contraction and both coordinate strips.

\section{Entry as a resource enclosure}\label{cut49:re:entry}

The entry calculation has one job: turn an attained extremizer into a
short list of simultaneous hypotheses for the local comparison. Its many
successive tables are applications of one enclosure principle. A table
is a cover of possible laws; a refinement intersects that cover with
newly proved inequalities. There is no assertion that an iteration
converges, and no improved bound is used to justify itself.

Throughout, $X$ has mean zero and variance one, $n\ge N$, and $H$ is
the raw sum threshold. Put
\[\begin{gathered}
 \rho=\E|X|^3,\quad \beta=\E[X|X|],\quad\gamma=\E X^3,\quad M_j=\E|X|^j,\\
 C=C_n,\quad e=C-C_E\ge0,\qquad C_E=\frac{3+\sqrt{10}}{6\sqrt{2\pi}}.
\end{gathered}\]
The input is an attained positive closed maximum, together with the
contact inequality stated below, an actual support bound, and an
independently established score cap. The global proof supplies these
inputs. A conditional upper bound on moments is never used to acquire
the support needed to define those moments.

\subsection{The conserved quantity}
Write $\tau(t)=\E[X\mathbf1_{X\ge t}]$. It is nonnegative,
$\int\tau=1$, and $2\int|t|\tau(t)\,dt=\rho$. On each bounded gap
$g=(a,b)$ of the closed support it has the constant value $\tau_g$.
Define
\begin{align*}
 J&=\sum_g\tau_g(b-a)^2,&
 J_1&=\sum_g\tau_g\int_a^b(t-a)t\,dt,\\
 J_2&=\sum_g\tau_g\int_a^b(t-a)t|t|\,dt,&
 D&=(3+\sigma)\rho-\gamma-3J,\qquad\sigma=\sqrt{10},\\
 L&=\sum_g\tau_g\int_{[a,b)}(t-a)\,\nu(dt),&
 \nu&=\mathcal L(H-S_{n-1}).
\end{align*}
The half-open convention includes the lower endpoint with weight zero
and excludes the upper endpoint. In particular it treats atoms without
an implicit midpoint convention. Bounded support makes all sums
absolutely convergent. The gap calculation below proves $D\ge0$;
clearly $L\ge0$.

Put $\varphi=\phi(H/\sqrt n)$ and
\begin{align*}
 K_1&=\E[X|X|^3]-3\beta-\tfrac32\rho\gamma,\\
 K_q&=M_5/4-\rho M_4/2-\tfrac32\beta\gamma+\rho/2,
 &U&=J_1+(M_4-3)/6.
\end{align*}
The exact remainder identity is
\begin{equation}\label{cut49:re:identity}
 ne\rho+nC_E\rho(1-e^{-w})+
 n\phi(0)e^{-w}D/6+n^{3/2}L=\mathcal R,
 \qquad w=H^2/(2n),
\end{equation}
where
\begin{equation}\label{cut49:re:raw}
 \mathcal R=C\{3\beta J/2+3\rho J_1-3J_2-K_q\}
 +H(CK_1+\varphi U)-H^2(C\rho+\gamma\varphi/2)+E_n.
\end{equation}
The complete Fourier estimate supplies a \emph{coherent} pair
$|E_n|\le K_F+s_F|H|$: the constant from one estimate is not combined
with the slope from another unless both were separately proved.
Section~\ref{cut49:re:foundation} derives \eqref{cut49:re:identity} and the whole
Fourier bound from the contact interface.

\begin{lemma}[Removing the Gaussian weight]\label{cut49:re:weight}
Suppose a first pass proves $\mathcal R<r<Nc_0$, where $C\rho\ge c_0$.
Every later valid upper bound $\mathcal R\le r'$ then gives
\begin{equation}\label{cut49:re:unweighted}
 ne\rho+n\phi(0)D/6+n^{3/2}L\le\mathcal R,
 \quad e^{-w}\ge1-\frac{5r'}{2n},
 \quad nD\le\frac{15000r'}{997}.
\end{equation}
The last two bounds use $C_E\rho>2/5$ and $\phi(0)>997/2500$.
\end{lemma}
\begin{proof}
Rearrange the exact identity before changing its Gaussian weight:
$\varphi(\gamma+3J)/6=C\rho+\sqrt nL-\mathcal R/n>0$.
Thus $C_E\rho-\phi(0)D/6>0$. Replacing $e^{-w}$ by one in
\eqref{cut49:re:identity} lowers its left side, proving the first assertion.
The original nonnegative phase term proves
$1-e^{-w}\le r'/(nC_E\rho)<5r'/(2n)$. The defect bound follows from
the first assertion. The initial sign license is proved before any
feedback is used.
\end{proof}

The same identity also gives a useful score update. If this cell forces
$D\ge d_l\ge0$ and a preceding pass gives raw ceiling $r$, then
\[
 C\le .40974+
 \frac{r/N-.3988(1-5r/(2N))d_l/6}{\rho_l}.
\]
A negative numerator excludes the cell because $e\ge0$. Otherwise its
numerator increases with $1/n$, and division by the positive lower
$\rho_l$ is outward. Thus the update is valid for every $n\ge N$ and
uses only the previous raw bound.

For later local comparisons we retain the sharper phase resource.
If previous bounds give $\mathcal R\le r$, $nD\le Nd$, then
\begin{equation}\label{cut49:re:phaseprice}
 ne\rho+\kappa nD+n^{3/2}L
 \le \mathcal B(H),\qquad \kappa=997/15000,
\end{equation}
where an upper bound for $\mathcal B$ is obtained from the raw
quadratic by increasing its curvature by
\begin{equation}\label{cut49:re:extra}
 q_{\rm extra}(\rho)=\frac{1-5r/(2N)}2
                         \left(.4095\rho-\frac d{15}\right)>0.
\end{equation}
Indeed the subtracted exact phase term is
$n(1-e^{-w})(C_E\rho-\phi(0)D/6)$.
Use $1-e^{-w}\ge we^{-w}$, $D\le d$, $\phi(0)<.4$, and the
previous density floor. The bound $Q_{\rm budget}$ includes
\eqref{cut49:re:extra}; it must not be charged a second time.

\subsection{One positive-measure rule}
Suppose zero lies in the gap $(-A,B)$, $A,B>0$. Let
\[\begin{gathered}
 p=\Pp(X>0),\quad q=1-p,\quad m=\E X_+=\E X_-,\quad V=\E[X\mid\mathbf1_{X>0}],\\
 W=X-V,\quad\theta_\pm=\E[W^2;\pm X>0],\quad\theta=\theta_++\theta_-.
\end{gathered}\]
Then $m^2=(1-\theta)pq$, and the conditional means have magnitudes
$\bar A=m/q$, $\bar B=m/p$. Set
$\Delta_A=\bar A-A\ge0$, $\Delta_B=\bar B-B\ge0$.
We also use the unnormalised endpoint displacements
\[
 M_+=\E[(X-B)\mathbf1_{X>0}]=p\Delta_B,\qquad
 M_-=\E[(-X-A)\mathbf1_{X<0}]=q\Delta_A.
\]

\begin{lemma}[Positive-measure pricing]\label{cut49:re:pricing}
For $x=|X|$, $a=|V|$, and a function $h$ constant on each sign class,
\begin{equation}\label{cut49:re:conditional}
 \E[h(|X|^j-|V|^j)]
 =\E\left[hW^2\sum_{k=0}^{j-2}(k+1)a^kx^{j-2-k}\right],\quad j\ge2.
\end{equation}
Also, for $d\mu_0=(|X|-1)^2(|X|+1/2)\,dP$,
\begin{align*}
 \mu_0(\R)&=\rho-1,\\
 M_4-1&=\int\frac{(x+1)^2}{x+1/2}\,d\mu_0,\\
 M_5-1&=\int\frac{x^3+2x^2+3x+3/2}{x+1/2}\,d\mu_0,\\
 K_1+3\rho\beta&=\int
  \frac{\operatorname{sgn}(X)x(x+2-3\rho/2)}{x+1/2}\,d\mu_0.
\end{align*}
For any nonnegative resources $z_i$ satisfying $\sum a_i z_i\le r$,
$a_i>0$, a signed linear cost obeys
\begin{equation}\label{cut49:re:simplex}
 \sum b_i z_i\le r\max\{0,b_1/a_1,\ldots,b_k/a_k\}.
\end{equation}
\end{lemma}
\begin{proof}
Divide $x^j-a^j-ja^{j-1}(x-a)$ by $(x-a)^2$; the linear term has
zero conditional mean. The radial identities follow by multiplying
out and using $\E X=0$, $\E X^2=1$. Finally each $b_i$ is at most
$a_i$ times the displayed maximum, and all $z_i$ are nonnegative.
\end{proof}

This one rule accounts for the moment, signed Taylor, endpoint movement,
outer-gap, and confinement-error estimates. In particular, no two costs
may each spend the full residual and then be added.

\subsection{The simultaneous gap budget}
Let $p_*=(4-\sigma)/2$ and
$T_+=\E X_+^3-mB^2$, $T_-=\E X_-^3-mA^2$.
If $J_o=J-m(A+B)^2$, signwise integration gives
\begin{equation}\label{cut49:re:gapexact}
 D=m(\sqrt{\sigma+1}A-\sqrt{\sigma-1}B)^2
   +(\sigma-1)T_++(\sigma+1)T_-+3(T_++T_--J_o).
\end{equation}
Moreover, writing $R_+=\E[X_+(X_+-B)]$ and similarly $R_-$,
\[
 T_++T_--J_o\ge2BR_++2AR_-\ge0.
\]
If zero belongs to the support, every gap has one sign and $J\le\rho$,
so $D\ge\sigma\rho-\gamma\ge(\sigma-1)\rho$.
Otherwise \eqref{cut49:re:gapexact} proves $D\ge0$.
Consequently $D<(\sigma-1)\rho$ licenses the actual central gap.
We also have $J\le\rho+2m\le\rho+1$; if $D<2$, then
$T_++T_-\le D/2$ and $A+B\ge\sqrt{2-D}$.

To verify the gap estimate, on a positive outer gap $(u,v)$,
$(v-u)^2\le2\int_u^v(t-B)\,dt$.
Summing with weights $\tau_g$ and applying Tonelli gives
$J_o^+\le T_+-2BR_+$; reflect for the other sign.
Expanding the definition of $D$ gives \eqref{cut49:re:gapexact}.
The inequality $AB\le1$ follows by integrating $(X+A)(X-B)\ge0$,
and $m\le1/2$ follows from $\E|X|\le1$.

Conditional expansion about $\bar A,\bar B$ now gives the stronger
single inequality
\begin{equation}\label{cut49:re:joint}
\begin{split}
 D\ge{}&\frac{(6+2\sigma)(1-\theta)^{3/2}}{\sqrt{pq}}(p-p_*)^2
 +6m\{(A+B)(\Delta_A+\Delta_B)+\Delta_A\Delta_B\}\\
 &+3(\sigma+1)B\theta_++3(\sigma+3)A\theta_-.
\end{split}
\end{equation}
For example
$\E X_+^3=m\bar B^2+\E[(X+2\bar B)(X-\bar B)^2;X>0]$
and $R_+=\theta_++m(\bar B-B)$.
The cubic residual is at least $3B\theta_+$.
The deterministic shift left after substituting both signs is
$6m\{\bar A\bar B-AB+B\Delta_B+A\Delta_A\}$, equal to the
second term above. The remaining deterministic expression, using
$m^2=(1-\theta)pq$, is exactly the displayed binary square.

\subsection{A state, its update, and its exported resource}
A state for a closed probability interval $I=[l,u]$ records
\[
 |X|\le R,\quad C\in[c_0,c_1],\quad\mathcal R\le r,
 \quad nD\le Nd,\quad\theta\le\Theta<1,
 \quad A\ge A_l>0,\quad B\ge B_l>0,
\]
together with simultaneous signed-moment and coefficient intervals.
The distinction between $nD\le Nd$ and $D\le d$ is essential.
Let $s_+=\max_Ip(1-p)$, $s_-=\min_Ip(1-p)$ and
\begin{gather*}
 b_I=\frac{(6+2\sigma)(1-\Theta)^{3/2}}{\sqrt{s_+}}
                 \operatorname{dist}(p_*,I)^2,
 \qquad r_I=d-b_I,\\
 K_+=3(\sigma+1)B_l,\quad K_-=3(\sigma+3)A_l,
 \quad K=\min(K_+,K_-),\quad m_l=\sqrt{(1-\Theta)s_-}.
\end{gather*}
If $r_I<0$, exclude the entire closed cell. Otherwise \eqref{cut49:re:joint}
gives, simultaneously,
\begin{equation}\label{cut49:re:residual}
 6m(A_l+B_l)(\Delta_A+\Delta_B)+K_+\theta_++K_-\theta_-\le r_I.
\end{equation}
Thus one may intersect the preceding bounds with
\begin{gather}
 \Theta'\le r_I/K,\qquad
 A_l'\ge\sqrt{(1-\Theta)l/(1-l)}-
                  \frac{r_I}{6m_l(A_l+B_l)},\notag\\
 B_l'\ge\sqrt{(1-\Theta)(1-u)/u}-
                  \frac{r_I}{6m_l(A_l+B_l)}.\label{cut49:re:geometry}
\end{gather}
Each right side is evaluated from established bounds; sequential
improvements may then be used by the next expression. A nonpositive
endpoint estimate fails this update, rather than excluding a law.
Initialization uses the same formulas with the valid lower sum
$A+B\ge\sqrt{2-d}$.

In Bernoulli coordinates set $Y=B_0+U$, $\E[U\mid B_0]=0$,
$B_0\sim\mathrm{Bernoulli}(p)$, $X=\ell(Y-p)$ and
$d_Y=(\Delta_A+\Delta_B)/(\bar A+\bar B)$.
Then $\ell^2=(1-\theta)/(pq)$ and $\lambda=n\E U^2=npq\theta/(1-\theta)$.
The \emph{weighted} versions of \eqref{cut49:re:residual} prove
\begin{equation}\label{cut49:re:exports}
 n\theta\le\frac{Nr_I}{K},\quad
 nd_Y\le\frac{Nr_I}{6(1-\Theta)(A_l+B_l)},\quad
 \lambda\le\frac{Ns_+r_I}{(1-\Theta)K}.
\end{equation}
Indeed $nD\le Nd$ whereas $nb_I\ge Nb_I$.
These exports are derived from the weighted resource; multiplying an
unweighted variance ceiling by $N$ would not prove them.

For $j=3,4,5$ set
\[
 L_3(v)=R+2v,\quad L_4(v)=R^2+2Rv+3v^2,
 \quad L_5(v)=R^3+2R^2v+3Rv^2+4v^3.
\]
Put $v_+=\sqrt{(1-l)/l}$, $v_-=\sqrt{u/(1-u)}$ and
$k_j=[1-(1-\Theta)^{j/2}]/\Theta$, with $k_j=j/2$ at $\Theta=0$.
For $u<1/2$, valid Bernoulli moment ceilings are, with $s=l(1-l)$,
\[
 B_3^*=(1-2s)/\sqrt s,\quad B_4^*=1/s-3,
 \quad B_5^*=(1-4s+2s^2)/s^{3/2}.
\]
They decrease towards $p=1/2$. Convexity places $(1-\theta)^{j/2}$
below its endpoint chord on $[0,\Theta]$; positive-measure pricing gives
\begin{equation}\label{cut49:re:moments}
 M_j\le B_j^*+r_I\max\left\{0,
 \frac{L_j(v_+)-k_jB_j^*}{K_+},
 \frac{L_j(v_-)-k_jB_j^*}{K_-}\right\}.
\end{equation}
The same rule, using
$T_+\le2mv_+\Delta_B+L_3(v_+)\theta_+$ and its reflection, bounds
\begin{equation}\label{cut49:re:outerresource}
 W_o:=(\sigma-1)T_++(\sigma+1)T_-
 \le r_I\max\left\{
 \frac{(\sigma-1)v_+}{3(A_l+B_l)},
 \frac{(\sigma+1)v_-}{3(A_l+B_l)},
 \frac{(\sigma-1)L_3(v_+)}{K_+},
 \frac{(\sigma+1)L_3(v_-)}{K_-}\right\}.
\end{equation}
Intersect every new bound with the previous one. The versions which
merge $\theta_+$ and $\theta_-$ are obtained by replacing both costs by
their maximum and both prices by their minimum; they need no new proof.

\subsection{The enclosure theorem}\label{cut49:re:operator}
On each cell, the formula atlas bounds the four coefficients in
\eqref{cut49:re:raw}. The result, for each actual sign $s=\operatorname{sgn}H$,
is a finite family of quadratics
\begin{equation}\label{cut49:re:quadratics}
 \mathcal R\le\max_j(A_j+F_{j,s}|H|-Q_j^{\rm raw}H^2),
 \qquad Q_j^{\rm raw}>0.
\end{equation}
The maxima retain all score--moment corners. In the initial signed
cover, intervals contain both signs of $\beta,\gamma$; no orientation is
discarded. In the probability cover, the central integral bound is a
maximum of affine functions of $\rho$, hence convex. The other
$\rho$-dependence is affine or the maximum of convex quadratics. For
fixed $H$, the score dependence is convex. Endpoints are therefore taken
\emph{before} completing the square.

\begin{theorem}[Finite resource enclosure]\label{cut49:re:enclosure}
Consider a finite ordered schedule of the following operations:
\begin{enumerate}
\item Split a cell into an explicitly covered family of closed cells.
\item Intersect its enclosures with an inequality proved in the atlas.
\item Exclude a cell by an empty necessary interval or a strictly
negative nonnegative resource.
\item Apply a separately proved support or confinement implication,
checking all its hypotheses against this cell's already acquired data.
\item Export a raw, weighted, or signed phase bound by
\eqref{cut49:re:quadratics}, \eqref{cut49:re:phaseprice}, and \eqref{cut49:re:exports}.
\end{enumerate}
If exact outward arithmetic checks every domain, corner and cover,
every output contains all attaining laws satisfying the input. The
conclusion holds for all integers $n\ge N$ wherever the formula's
monotonicity gates are checked. It includes zero conditional variance,
singleton intersections and arbitrarily small atoms.
\end{theorem}
\begin{proof}
Induct over the schedule. A split preserves the union, intersections use
necessary inequalities, and an excluded cell contains no feasible law.
Each callback consumes only earlier hypotheses for that same law.
For a scalar raw output, use
$fh-qh^2\le(f_+)^2/(4q)$ for $h\ge0$, $q>0$; using $f^2$ is a valid
weaker bound. For a resource output first add \eqref{cut49:re:extra} to each
curvature. A negative maximum excludes the cell; a nonnegative one may
be divided by $\kappa N$ to give its next defect parameter.
The phase bound follows from the positive root of
$A+fh-qh^2\ge0$, replacing $A$ by $A_+$ if necessary.
A strict output is exported only when the certificate retains a strict
upper gap, such as its strict Fourier or target enclosure; closed-domain
coverage does not change strictness. All constants are frozen before evaluating
corners. Thus no iteration
step relies on an unproved fixed point. The stated index uniformity is
that of the displayed finite formulas, not a fixed-law limit.
\end{proof}

For the Fourier cancellation estimate only, reflect $(X,H)$ if needed:
$I_0(-X,-H)=I_0(X,H)$. Both odd coefficients change sign, and
$\int e^{-x^2/2}|zx^2-cx^4|\,dx$ is unchanged by $(z,c)\mapsto(-z,-c)$.
Consequently a lower endpoint used there bounds $|\gamma|$, not necessarily
the original signed $\gamma$. The actual contact, gap and phase data keep
their original orientation. A useful additional signed lower bound is
\begin{equation}\label{cut49:re:gamma}
 \gamma\ge(1-\Theta)^{3/2}\frac{1-2u}{\sqrt{u(1-u)}}
       -\left(R+2\sqrt{\frac{u}{1-u}}\right)\Theta.
\end{equation}
To prove it expand $\E X^3=\E V^3+\E[W^2(X+2V)]$; the positive
residual is nonnegative and the negative residual has the displayed
lower bound. The binary term decreases in both $p<1/2$ and $\theta$.
Thus $\max(0,\text{any signed lower bound})$ is a valid lower bound for
$|\gamma|$ inside this auxiliary Fourier estimate.

\subsection{Support, confinement, and a short acyclic schedule}
Two different support implications are useful. Gaussian contact excludes
a point by a strict negative upper bound on its contact expression;
closed cells cover the finite exterior and a derivative-sign certificate
covers any remaining ray. The shared Gaussian support theorem supplies
this callback, including the bounded-threshold radius and coordinate-strip
contracts; it is a prior theorem, not a test on moments alone. Noise confinement proves $|U|\le a$ on an
already acquired row. It then implies, without a small-noise limit,
\begin{equation}\label{cut49:re:radiusfeedback}
 |X|\le R_*:=
 \frac{\max\{u+a,1-l+a\}}{\sqrt{\min_Ip(1-p)}}.
\end{equation}
The denominator uses $\E U^2\ge0$, so zero noise is included.
The Fourier/moment formulas remain valid for $R\ge1$; the paired Pearson
choice additionally checks $R>2g$. In particular there is no intrinsic
lower radius $2.3$: the radial ratios in Lemma~\ref{cut49:re:pricing} increase
on all $x\ge0$, and the other formulas use their explicit domains.

Here is the selected proof schedule. Each row names \emph{an application}
of Theorem~\ref{cut49:re:enclosure}, not a different analytic theorem.
\begin{center}\small
\begin{tabular}{p{.13\linewidth}p{.32\linewidth}p{.45\linewidth}}
\toprule
Node&Input and operation&Output consumed later\\\midrule
I&$N=47$, $R=3.84$, $C\le.4247$; 13 signed passes plus export&
100 closed probability cells; 35 survive; $\mathcal R<1.885$, $nD<47(.590)$\\
C&Central preparation, contact radius $2.5$, probability refinement&
45 cells $[.38,.47]$, width $.002$; canonical simultaneous input\\
L&Complete complementary cover, contact radius $3$, refinement&
39 cells $[.302,.38]$; all other complementary cells excluded\\
C1&First confinement from C; radius formula; enclosure&
First improved central table\\
C2&Antichain confinement from C; intersect C1; enclosure&
Selected central table: 45 cells, $[.38,.47]$\\
S&Separate sign-resource pricing on C2 and L&
One left cell excluded; 38 left cells remain\\
L2&Left contact radius $2.7$ from S; enclosure&
Selected left table: 38 cells, $[.304,.38]$\\\bottomrule
\end{tabular}
\end{center}
Every shared endpoint belongs to its adjacent closed cells. The
confinement producing C2 consumes C, not C2. The selected central local
comparison uses C2; it does not need the optional further central
refinement in S. The final cover is exactly $[.304,.47]$.

\begin{corollary}[Uniform acquired entry]\label{cut49:re:acquired}
Let an attained positive maximum have $n\ge47$, $C_E\le C\le.4247$,
$0\le n(C-C_E)\le.7$, decrement at most $.19$, and both deleted
score caps at most $.4253$.
Given the shared contact support and strip implications, the exact
schedule above places its law in one of the 38 selected left rows or
45 selected central rows, with every resource in that row holding
simultaneously. This assertion has no upper bound on $n$.
\end{corollary}
\begin{proof}
The shared all-threshold support theorem supplies radius $3.84$.
Apply the finite enclosure theorem in the displayed order, retaining
the weighted bounds at every edge. The initialization's stored upper
index 98 occurs only through $.4095+K/98$; here $K=0$, so it imposes
no upper-index restriction. All subsequent index-dependent inequalities
are checked for every $n\ge47$ by their stated monotonicity gates.
The accepted closed covers exclude the complementary cells and give
the stated 83 rows.
\end{proof}

\paragraph{The local-comparison interface.}
For each final row supply $p\in I$, $\theta\le\Theta$, $\lambda\le\Lambda$,
the actual support/amplitude, $d_Y\le d_0$, both signed phase limits,
and all four coherent quadratics bounding \eqref{cut49:re:phaseprice}.
On a phase/noise subcell, let $T_H$ be the maximum of each concave
quadratic at its clipped vertex, then the maximum over corners and the
minimum with the scalar bound. If a proved moment inequality is
\[
 D\ge\alpha d_Y+b\theta+P,\qquad\alpha>0,
 \quad \ell\ge\ell_-,\quad\lambda\ge L_0,
\]
then the exact strict resource supplied to the comparison is
\begin{equation}\label{cut49:re:localG}
 nd_Y<G:=\frac{T_H/\kappa-NP-b\ell_-^2L_0}{\alpha}.
\end{equation}
This uses $n\theta=\ell^2\lambda$ and $nP\ge NP$. If $G\le0$,
the entire closed subcell is excluded. Formula~\eqref{cut49:re:localG} is
the same resource as \eqref{cut49:re:phaseprice}, not an additional allowance.

\section{The minimal contact--Fourier foundation}\label{cut49:re:foundation}
The only contact input needed here is the following explicit interface.
At the attained positive maximum, with $F=\Pp(S_n\le H)$,
$z=H/\sqrt n$ and $\nu=\mathcal L(H-S_{n-1})$, one has
$\nu([y,\infty))\le\mathfrak q(y)$ for every real $y$, with equality
on the full closed support, where
\begin{align}
 n\mathfrak q(y)={}&nF+\tfrac12\{z\phi(z)+C\rho/\sqrt n\}
 -(\sqrt n\phi(z)+3C\beta/\sqrt n)y\notag\\
 &-\tfrac12\{z\phi(z)+3C\rho/\sqrt n\}y^2
 +C|y|^3/\sqrt n.\label{cut49:re:contact-q}
\end{align}

This contact statement and attainment are supplied by the global section.
Everything that follows in this section is derived from it. We abbreviate
$M_4=\mu_4$, $M_5=\eta_5$ in the coefficient formulas.
\subsection{Compensated smooth inversion}
Write $Z_n=S_n/\sqrt n$, $f(t)=\E e^{itX}$, and let $N$ be standard normal.
The smooth payoff $q_z(x)=(x-z)|x-z|$ gives the exact error
\[
 \epsilon_n(z)=\E q_z(Z_n)-\E q_z(N)-\frac{2\gamma\phi(z)}{3\sqrt n}.
\]
\begin{lemma}\label{cut49:smooth:fourier}\label{cut49:smooth:inversion}
For every real $H$ and $n\ge1$,
\begin{equation}\label{cut49:smooth:smooth-integral}
 \epsilon_n(H/\sqrt n)=-\frac4{\pi n}\int_0^\infty
 \frac{\Im(e^{-itH}V_n(t))}{t^3}\,dt,\qquad
 V_n(t)=f(t)^n-e^{-nt^2/2}(1-in\gamma t^3/6).
\end{equation}
The integral is absolutely convergent. With the split at $t=1/2$,
\begin{equation}\label{cut49:smooth:split}
 I_{\rm low}=\int_0^{1/2}\frac{|\Im(e^{-itH}V_n(t))|}{t^3}\,dt,
 \qquad I_{\rm high}=\int_{1/2}^\infty\frac{|V_n(t)|}{t^3}\,dt,
\end{equation}
one has $|\epsilon_n(z)|\le4(I_{\rm low}+I_{\rm high})/(\pi n)$.
\end{lemma}
\begin{proof}
The scalar identities
\[
 \int_0^\infty\frac{ty-\sin(ty)}{t^3}\,dt=\frac\pi4y|y|,
 \qquad
 \int_0^\infty\frac{|ty-\sin(ty)|}{t^3}\,dt=\frac\pi4y^2
\]
follow by scaling and integration by parts. One elementary evaluation of
the constant uses
$J_a(b)=\int_0^\infty e^{-at}(1-\cos bt)t^{-2}dt$:
$J_a''(b)=a/(a^2+b^2)$, $J_a(0)=J_a'(0)=0$, so
$J_a(1)=\int_0^1\arctan(b/a)db\to\pi/2$.
Monotone convergence removes $a>0$.
For a finite signed measure $\mu$ with zero mass and first moment and
finite absolute second moment, absolute Fubini applied to the
\emph{compensated} scalar integrand therefore gives
\begin{equation}\label{cut49:smooth:signed-inversion}
 \int (s-H)|s-H|\,d\mu(s)
 =-\frac4\pi\int_0^\infty t^{-3}
       \Im(e^{-itH}\widehat\mu(t))\,dt.
\end{equation}
Apply this to the law of $S_n$ minus the pushforward, by $x\mapsto\sqrt n x$,
of the signed density
$\phi(x)[1+\gamma(x^3-3x)/(6\sqrt n)]$.
Gaussian integration by parts gives its transform
$e^{-nt^2/2}(1-in\gamma t^3/6)$, total mass one and first moment zero.
Since $(x^3-3x)\phi=-\phi'''$, integration twice by parts and then across
$z$ gives $-\int q_z\phi'''=4\phi(z)$. This proves the formula and its
scale factor. The high integral is finite from $|f|\le1$ and Gaussian
decay; the low integral is finite by the absolute Fubini argument.
\end{proof}

The payoff is smooth enough for compensated inversion. No discontinuous
CDF is inverted, so no atom midpoint convention is introduced.
\subsection{Derivation of the conserved identity}
\begin{proof}
Compute $\mathcal A=\int\tau(y)[\mathfrak q(y)-\nu([y,\infty))]dy$
in two ways. Put $r=-\mathfrak q'$. Contact at both endpoints of a gap gives
\[
 \nu([a,b))=\int_a^b r(t)dt,
 \quad
 \int_a^b[\mathfrak q(y)-\nu([y,\infty))]dy
 =\int_a^b(t-a)r(t)dt-\int_{[a,b)}(t-a)d\nu(t).
\]
The integrand vanishes on the support and $\tau=0$ on the exterior rays.
Summing, and differentiating \eqref{cut49:re:contact-q}, gives
\begin{equation}\label{cut49:smooth:A-gap}
 \mathcal A=\frac{\phi(z)J}{2\sqrt n}+\frac{G}{n^{3/2}}-L,
 \qquad
 G=\frac32C\beta J+(H\phi(z)+3C\rho)J_1-3CJ_2.
\end{equation}
All sums and integrals are absolutely convergent on the compact convex
hull of the support, since $|J_1|\le R_0J/2$, $|J_2|\le R_0^2J/2$
when $|X|\le R_0$.

For the second computation let $X^*$ have density $\tau$, independent of
$S_{n-1}$. Signwise Fubini gives
$\E g'(X^*)=\E[Xg(X)]$ for continuously differentiable $g$.
In particular,
$\E X^*=\gamma/2$, $\E(X^*)^2=\E X^4/3$ and
$\E|X^*|^3=\E|X|^5/4$.
Applying the identity to each summand shows that
$(S_{n-1}+X^*)/\sqrt n$ has density
$\tau_{Z_n}(t)=\E[Z_n\mathbf1_{Z_n\ge t}]$.
Its CDF at $z$ is
$T_z=\E[(Z_n^2-zZ_n)\mathbf1_{Z_n\le z}]$, by integrating that density;
the boundary integrand is zero. Substituting the three moments into
\eqref{cut49:re:contact-q} therefore gives
\begin{equation}\label{cut49:smooth:zero-bias-average}
 \mathcal A=F-\frac{\gamma\phi(z)}{2\sqrt n}
 +\frac{H\phi(z)(3-\E X^4)/6+CK_q}{n^{3/2}}-T_z.
\end{equation}
Multiply contact equality by $X$ and integrate. By independence and symmetry,
$n\E[X\nu([X,\infty))]$ equals $\E[S_n\mathbf1_{S_n\le H}]$.
Expanding the other side yields, for $M_z=\E[Z_n\mathbf1_{Z_n\le z}]$,
\begin{equation}\label{cut49:smooth:stationarity}
 M_z+\phi(z)=-\frac{\gamma z\phi(z)}{2\sqrt n}+\frac{CK_1}{n}.
\end{equation}
Finally $q_z(x)=(x-z)^2-2(x-z)^2\mathbf1_{x\le z}$, so Gaussian
integration gives
\begin{equation}\label{cut49:smooth:truncated-square}
 T_z-\Phi(z)=-\tfrac12\{\E q_z(Z_n)-\E q_z(N)\}
 +z(M_z+\phi(z))-z^2(F-\Phi(z)).
\end{equation}
Insert \eqref{cut49:smooth:stationarity} and the definition of $\epsilon_n$
into \eqref{cut49:smooth:zero-bias-average}. With
\[
 P_0=H^2(C\rho+\gamma\phi(z)/2)
       -H\{CK_1+\phi(z)(\E X^4-3)/6\}+CK_q,
\]
the result is $\sqrt n\mathcal A=C\rho-\gamma\phi(z)/6+
P_0/n+\sqrt n\epsilon_n(z)/2$. Equating this with
\eqref{cut49:smooth:A-gap} proves
\begin{equation}\label{cut49:smooth:signed-identity}
 C\rho+\sqrt nL=\frac{\phi(z)}6(\gamma+3J)
                 +\frac{G-P_0}{n}-\frac{\sqrt n}2\epsilon_n(z).
\end{equation}
Now substitute $\gamma+3J=(3+\sqrt{10})\rho-D$ and
$\phi(z)=\phio e^{-w}$, and expand $G-P_0$.
\end{proof}

\subsection{One complete low-frequency transfer}
Put $f(t)=\E e^{itX}$, $h(t)=e^{-t^2/2-i\gamma t^3/6}$,
$e_n(t)=e^{-nt^2/2}(1-in\gamma t^3/6)$, and
\[
 I_0(H)=\int_0^{1/2}|\Im(e^{-itH}(f^n-e_n))|\frac{dt}{t^3},\qquad
 I_\infty=\int_{1/2}^\infty|f^n-e_n|\frac{dt}{t^3}.
\]
The inversion bound~\eqref{cut49:smooth:split} charges
$2\sqrt n(I_0+I_\infty)/\pi$; both integrals below are complete.

\begin{lemma}[Signed Gaussian transfer]\label{cut49:fourier:early-low}
\label{cut49:fourier:paired-low}\label{cut49:entryleap:fourier-transfer}
Suppose $n\ge N>2$, $\mu_4\le M\le11.1$, $|\gamma|\le g\le1.85$,
$\mu_6\le M_6$, and the actual coefficients
\[
 a=(3-\mu_4)/24,\quad z=(10\gamma-\E X^5)/120,\quad c=a\gamma/6
\]
satisfy $|a|\le a_*$, $|z|\le Z$. Put
\begin{equation}\label{cut49:fourier:taylor}
 E_6=M_6/720+g^2/72+(1/2+g^2/144)^3/5,\quad
 B_6=a_*+Z/2+E_6/4.
\end{equation}
Let $a_0=(N-1)/(2N)$,
$0<b_0\le(45-M)(1-2/N)/96$, and
$\lambda_N=\sqrt{\pi/2}(N/(N-1))^{3/2}$.
Then, for every real $H$,
\begin{equation}\label{cut49:fourier:low-form}
 I_0(H)\le(L_0+L_H|H|)/\sqrt n+L_2/n,
 \qquad L_2=E_6/(2a_0^2)+B_6^2/(2b_0^3).
\end{equation}
Two admissible choices, each used as a complete triple, are as follows.
The triangle choice is
\begin{equation}\label{cut49:fourier:early}
 L_0=c_0Z+c_1a_*g+251g^3/17280,\quad
 L_H=c_0a_*+251g^2/4800,
\end{equation}
where $c_0\ge\lambda_N$, $c_1\ge\lambda_N/2$.
For the cancellation choice, take any rectangle containing the actual
$(z,c)$, $I_u\ge\sqrt{\pi/2}$ and $F_N\ge(N/(N-1))^{3/2}$. Define
\begin{gather*}
 U(z,c)=\begin{cases}
 I_u(|z|+3|c|),&zc\le0,\\
 I_u\bigl||z|-3|c|\bigr|+6\sqrt{|zc|}e^{-|z|/(2|c|)},&zc>0,
 \end{cases}\\
 J_u=\max_{\rm corners}U(z,c).
\end{gather*}
Then
\begin{align}\label{cut49:fourier:paired}
 L_0&=F_NJ_u+\frac{F_NI_u}{N-1}
 (35a_*g^3/48+35Zg^2/24)+251g^3/17280,\notag\\
 L_H&=F_NI_u\{a_*+5Zg/[2(N-1)]\}+251g^2/4800.
\end{align}
Neither the transfer nor its cancellation choice assumes signs for
$a,z,c$ or $\gamma$.
\end{lemma}
\begin{proof}
Taylor with integral remainder gives
$f-h=-at^4-izt^5+r_6$, $|r_6|\le E_6t^6$ for $0\le t\le1/2$.
Indeed for $w=-t^2/2-i\gamma t^3/6$,
$|w|\le t^2(1/2+g^2/144)<1/6$; the terms through degree five in
$1+w+w^2/2$ are exact, its degree-six term costs $g^2t^6/72$, and
the remaining exponential series costs $|w|^3e^{|w|}/6\le|w|^3/5$.
Add the true sixth-moment remainder.

Four integrations of $1-\cos u\ge0$ give
$\cos u\le1-u^2/2+u^4/24$. Symmetrization, using
$\E(X-X')^4=2\mu_4+6$, therefore gives
$|f(t)|,|h(t)|\le e^{-(45-M)t^2/96}$ on $[0,1/2]$.
Every point on their joining segment has the same bound, so polynomial
Taylor expansion yields
\[
 f^n-h^n=n(f-h)h^{n-1}+R_2,\qquad
 |R_2|\le\tfrac12n^2B_6^2t^8e^{-b_0nt^2}.
\]
The $r_6$ and $R_2$ integrals are at most
\begin{equation}\label{cut49:fourier:common-error}
 \frac1n\{E_6/(2a_0^2)+B_6^2/(2b_0^3)\}=L_2/n.
\end{equation}
For the remaining comparator put $u=n\gamma t^3/6$. Its imaginary part
is bounded by $e^{-nt^2/2}(|u|^3/6+|H|tu^2/2)$, so its divided integral
is at most
\begin{equation}\label{cut49:fourier:comparator-low}
 n^{-1/2}(251g^3/17280+251g^2|H|/4800).
\end{equation}
Here the Gaussian moments follow from integration by parts,
$\int_0^\infty x^{2j}e^{-x^2/2}dx=(2j-1)!!\sqrt{\pi/2}$,
and $\sqrt{2\pi}<251/100$.

It remains to price the leading quartic and fifth terms. Set $b=n-1$,
$v=b\gamma t^3/6$. The triangle estimate uses
$|\sin(Ht+v)|\le|H|t+bg t^3/6$, $|\cos|\le1$; the Gaussian
moments of orders two and four give respectively $\lambda_N$ and
$\lambda_N/2$. This proves~\eqref{cut49:fourier:early}.
For cancellation, after division by $t^3$, rotation by $Ht$ bounds the
leading imaginary part by
\[
 |zt^2-atv|+a_*t|\sin v-v|+Zt^2|\cos v-1|
       +|H|t\{a_*t+Zt^2|\sin v|\}.
\]
Multiply by $ne^{-bt^2/2}$ and use the global sine/cosine remainders.
The first integral equals $nJ(z,c)/b^{3/2}$, where
$J(z,c)=\int_0^\infty e^{-x^2/2}|zx^2-cx^4|dx$.
This norm is convex and $J(-z,-c)=J(z,c)$, so its rectangle maximum
is at a corner. For $zc\le0$ the integral is
$\sqrt{\pi/2}(|z|+3|c|)$. For $z,c>0$, split at $r=\sqrt{z/c}$;
integration by parts gives
\[
 J=\sqrt{\pi/2}(z-3c)+6cr e^{-r^2/2}
       +2(3c-z)\int_r^\infty e^{-x^2/2}dx.
\]
Discard the last term when nonpositive, and otherwise bound its integral
by $\sqrt{\pi/2}$. This proves $J\le U$; simultaneous sign reflection
covers the remaining quadrant. The other four phase terms use moments
of orders two, six, eight and ten and give the coefficients
$1,5/2,35/24,35/48$ in~\eqref{cut49:fourier:paired}. Finally $n/(n-1)$,
$1/(n-1)$, $a_0^{-1}$ and the allowable damping inverses decrease with
$n$, proving the all-index bounds at $N$.
\end{proof}

\begin{lemma}[Retaining the signed Gaussian tail]\label{cut49:entry189:mills}
For $r>0$, put
\[
 \mathcal M(r)=e^{r^2/2}\int_r^\infty e^{-t^2/2}\,dt,
 \qquad
 \mathcal M_k(r)=\cfrac1{r+\cfrac1{r+\cfrac2{\ddots+\cfrac{k}{r}}}}
 \quad(k\ge1),\qquad \mathcal M_0(r)=1/r.
\]
Odd truncations are lower bounds for $\mathcal M$; even truncations
are upper bounds. For $z,c>0$, set $x=z/c$ and
\[
 d_{k+1}=1,\qquad d_j=1+\frac j{xd_{j+1}}\quad(j=k,\ldots,1),
 \qquad m_k(x)=d_1^{-1}.
\]
Then the cancellation norm in Lemma~\ref{cut49:entryleap:fourier-transfer}
satisfies
\begin{equation}\label{cut49:entry189:mills-bound}
 J(z,c)\le \sqrt{\pi/2}(z-3c)+\sqrt{zc}\,e^{-x/2}
 \left(6+\frac{2(3-x)}x m_k(x)\right),
\end{equation}
where $k=11$ if $x\ge3$ and $k=12$ if $x<3$.
The factor in parentheses is positive. One may take the minimum with
the previous bound $U(z,c)$ at each corner of the coefficient rectangle.
\end{lemma}
\begin{proof}
The positive integrals
$I_j(r)=\int_0^\infty t^j e^{-rt-t^2/2}\,dt$ satisfy
$I_1=1-rI_0$ and $I_{j+1}=jI_{j-1}-rI_j$ by integration by parts.
Thus, writing $q_j=I_j/I_{j-1}>0$,
\[
 I_0=\frac1{r+q_1},\qquad q_j=\frac j{r+q_{j+1}}.
\]
Replacing the final positive $q_{k+1}$ by zero reverses the inequality
at each substitution. This proves the alternating finite bounds, since
translation gives $I_0=\mathcal M(r)$. No convergence of an infinite
continued fraction is required. The displayed rational recursion gives
$m_k(x)=r\mathcal M_k(r)$ when $r=\sqrt x$.

Splitting the norm at $r=\sqrt{z/c}$, as in the preceding proof, gives
the exact identity
\[
 J(z,c)=\sqrt{\pi/2}(z-3c)+\sqrt{zc}\,e^{-x/2}
 \left(6+\frac{2(3-x)}x r\mathcal M(r)\right).
\]
The multiplier of $\mathcal M$ is nonpositive when $x\ge3$ and
positive when $x<3$. Substitute the corresponding finite lower or
upper bound. Positivity of the remaining factor is immediate for
$x<3$; for $x\ge3$, use $0<m_{11}(x)<1$ to bound it below by
$6-2(x-3)/x>4$. Finally the exact norm $J$, rather than either
upper-bound formula, is convex, so its containing rectangle is still
exhausted by its four corners.
\end{proof}

At every positive coefficient corner use the displayed finite Mills bound.
A negative Gaussian coefficient uses a lower enclosure for
$\sqrt{\pi/2}$; a nonnegative coefficient uses an upper one. The remaining
factor is positive, so outward upper square-root and exponential enclosures
are valid. A simultaneous reflection of both odd coefficients leaves the
norm unchanged.

\paragraph{The two coefficient choices in the certificates.}
All coefficients come from the actual marginal. For $Q=X^2-\gamma X-1$,
multiplication gives
\begin{align}\label{cut49:fourier:pearson}
 \E Q^2&=\mu_4-1-\gamma^2,&
 \E X^5&=\gamma^3+2\gamma+\E[Q^2(X+2\gamma)],\notag\\
 \mu_6&=1+3\gamma^2+\gamma^4+
       \E[Q^2\{(X+\gamma)^2+2\gamma^2+2\}].
\end{align}
For any radius $R_0\ge|X|$ with $R_0\ge2g$, the two signs of
$\E X^5-10\gamma$ are bounded, at $t=|\gamma|$, by
\begin{align*}
 P_{R_0}(t)&=R_0(M-1)+(10-2M)t-R_0t^2+t^3,\\
 Q_{R_0}(t)&=R_0(M-1)+(2M-10)t-R_0t^2-t^3.
\end{align*}
The gate $R_0-2g\ge0$ permits substituting $M$ for $\mu_4$ in both
signs. The triangle choice uses the valid enclosing radius $R_0=6$, $a_* =\max(2,M-3)/24$,
$M_6=RE_5$, $b_0=(45-M)(1-2/N)/96$ and
\begin{equation}\label{cut49:fourier:z0}
\begin{gathered}
 120Z=\min\{E_5+10g,Z_0\},\\
 Z_0=\begin{cases}
 P_6(g),&M\le5, P'_6(g)\ge0,\\
 6(M-1)+(10-2M)^2/[4(6-g)],&M\le5, P'_6(g)<0,\\
 6(M-1)+(2M-10)^2/24,&M>5.
 \end{cases}
\end{gathered}
\end{equation}
If $M\le5$, $P_6\ge Q_6$ and $P_6''<0$ on $[0,g]$; use the
increasing endpoint or $t^3\le gt^2$ and complete the square.
If $M>5$, $P_6\le6(M-1)$ and the upper quadratic for $Q_6$ gives
the third row. One may also use the actual radius by setting
$r_0=\max(R,2g)>g$ and
\[
 Z_r=\begin{cases}
 r_0(M-1)+(10-2M)t_*-(r_0-g)t_*^2,
 & M\le5,\quad t_*=\min\{g,(10-2M)/[2(r_0-g)]\},\\
 r_0(M-1)+(2M-10)t_*-r_0t_*^2,
 & M>5,\quad t_*=\min\{g,(2M-10)/(2r_0)\}.
 \end{cases}
\]
Indeed for $M\le5$, $P_{r_0}\ge Q_{r_0}$ and $t^3\le gt^2$;
for $M>5$, $P_{r_0}\le r_0(M-1)$ and dropping $-t^3$ bounds
$Q_{r_0}$. Each displayed quadratic is maximized at its clipped vertex.
Thus the enclosure uses
$120Z=\min(E_5+10g,Z_0,Z_r)$.
It uses $c_0\ge\sqrt{\pi/(16a_0^3)}=\lambda_N$ and $c_1=c_0/2$.
The weaker choice $c_1\ge\sqrt{\pi/(256a_0^5)}
=\lambda_N N/[2(N-1)]$ is also permitted by the same transfer lemma.

For the paired choice the already proved cell has
$0\le\gamma_l\le|\gamma|\le g\le1$, $M\le2$, and $R>2g$. After reflecting only this auxiliary Fourier integral if necessary, set
\begin{gather*}
 a_*=(2-\gamma_l^2)/24,\qquad b_0=(43/96)(1-2/N),\\
 M_6=\min\{RE_5,1+3g^2+g^4+(R^2+2Rg+3g^2+2)(M-1-\gamma_l^2)\},\\
 t_*=\min\{g,(10-2M)/[2(R-g)]\},\\
 120Z=\min\{E_5+10g,R(M-1)+(10-2M)t_*-(R-g)t_*^2\},\\
 z_l=\max\{-Z,[8\gamma_l-\gamma_l^3-(R+2\gamma_l)(M-1-\gamma_l^2)]/120\},
 \quad z_u=Z,\\
 c_l=(3-M)\gamma_l/144,\qquad c_u=(2-\gamma_l^2)g/144.
\end{gather*}
Check $M-1-\gamma_l^2\ge0$ and that the rectangle is ordered. Here
$P_R\ge Q_R$, and $P_R(t)\le R(M-1)+(10-2M)t-(R-g)t^2$,
whose maximum is at $t_*$. The numerator of the lower $z$ bound increases
in $\gamma\ge0$ because its derivative is
$10-2M+2R\gamma+3\gamma^2>0$. Pearson's first and third identities
give $a_*$ and $M_6$. If $3g^2\le2$, the increasing function
$2\gamma-\gamma^3$ further licenses $c_u\le(2g-g^3)/144$.
Thus this positive rectangle is one specialization of the signed transfer,
with exactly the coefficients used by the fixed acceptance programs.

\subsection{The complete high-frequency integral}
For $b,\kappa>0$ and $n>0$, put
\[
 T_n=\sqrt n\int_b^\infty e^{-n\kappa t^2}t^{-3}\,dt,
 \qquad x=n\kappa b^2.
\]
\begin{lemma}[HIGH\_LAGUERRE2]
For every $x>0$,
\[
 T_n\le\frac{\sqrt n e^{-x}}{2b^2}
                  \frac{x+2}{x^2+4x+2}.
\]
If $N\kappa b^2\ge1/4$, the right side decreases with $n\ge N$.
\end{lemma}
\begin{proof}
Substitute $t^2=b^2(1+s)$ and use positive Fubini to obtain
\[
 T_n=\frac{\sqrt n e^{-x}}{2b^2}
 \int_0^\infty\frac{e^{-xs}}{(1+s)^2}\,ds
 =\frac{\sqrt n e^{-x}}{2b^2}
 \int_0^\infty\frac{u}{x+u}e^{-u}\,du.
\]
Put $P(u)=u^2-4u+2$ and $D(x)=x^2+4x+2=P(-x)>0$.
Exact division gives
\[
 \frac1{x+u}-\frac{Q_x(u)}{D(x)^2}
       =\frac{P(u)^2}{(x+u)D(x)^2}\ge0,
\]
where
\[
 Q_x(u)=-u^3+(x+8)u^2-(x^2+8x+20)u
                  +x^3+8x^2+20x+16.
\]
Integration uses only $\int_0^\infty u^j e^{-u}du=j!$, proved by
integration by parts. It gives
$\int Q_x(u)e^{-u}du=(x+3)D(x)$ and hence
\[
 \int_0^\infty\frac{u}{x+u}e^{-u}du
 =1-x\int_0^\infty\frac{e^{-u}}{x+u}du
 \le1-\frac{x(x+3)}{D(x)}=\frac{x+2}{D(x)}.
\]
For the all-index statement, differentiate
$\sqrt{x}e^{-x}(x+2)/D(x)$. After multiplication by positive factors,
its derivative has numerator
$4-10x-22x^2-13x^3-2x^4$. This polynomial is negative at $x=1/4$
and strictly decreasing for $x>0$. Thus it is negative on the entire
required half-line.
\end{proof}

Let $EX=0$, $EX^2=1$, $EX^4\le M$, and write
$\gamma=EX^3$, $f(t)=Ee^{itX}$. For independent copies suppose an
actual event $E$ satisfies
\[
 P(E)\ge a>0,\qquad |X-X'|\ge d_0>0\text{ on }E,\qquad
 E[(X-X')^2\mid E]\le M_E.
\]
The event is measured under the product of the same marginal law.
Define the Taylor comparator
$e_n(t)=e^{-nt^2/2}(1-in\gamma t^3/6)$ and
\[
 I_\infty=\int_{1/2}^{\infty}|f(t)^n-e_n(t)|t^{-3}\,dt,
 \qquad L(x)=\frac{x+2}{x^2+4x+2}.
\]
The sign of the imaginary term is immaterial for its absolute-value bound.
Use $b=9/10$, $0<r\le1/4$,
\begin{gather*}
 h=\frac{333}{106}-\frac{(22/7)^3r^2}{6}>0,\quad
 \kappa=a(23/48)^2d_0^2,\\
 S_3=\sum_{j=1}^{20}j^{-3}+\frac1{800},\qquad
 S_5=\sum_{j=1}^{20}(j-r)^{-5}+\frac1{4(20-r)^4}.
\end{gather*}
Set
\[
 \begin{gathered}c=\frac{1-\max_{u\in\{1/4,81/100\}}
              [1-u+(M+3)u^2/12]}2,\\
 c_0=\frac18-\frac{M+3}{384},\quad
 s=\frac{1-(M+3)53/600}{2}.\end{gathered}
\]
The usual damping domain requires $c>0$.

\begin{lemma}[HIGH\_WINDOW50]
For every integer $n\ge5$, the complete upper bound is
\begin{align*}
 \sqrt n I_\infty\le{}& A_n
 +\frac{M_E}{4\pi^{3/2}h\sqrt a}
       \left(S_3+\frac{3S_5}{nah^2}\right)
 +\frac{M_E\sqrt n}{2\pi^2}e^{-nah^2r^2}\\
 &+Z_n+\frac{\sqrt n}{2b^2}e^{-na(23/48)^2}
 +\sqrt n e^{-n/8}\left(2L(n/8)+\frac{\sqrt{M-1}}3\right),
\end{align*}
where always $A_n=2\sqrt n e^{-nc}$ is permitted, and, if $s>0$,
it can be replaced by the smaller of itself and
\[
 2\sqrt n e^{-nc_0}L(ns/4).
\]
The zero-cell term $Z_n$ may be taken as the smaller of
\[
 \frac{e^{-n\kappa b^2}}{2b^4\kappa\sqrt n},\qquad
 \frac{\sqrt n e^{-n\kappa b^2}}{2b^2}L(n\kappa b^2).
\]
Each term may be evaluated at a common floor $N$ for all $n\ge N$
provided
\[
 2Nc'>1\quad(c'\in\{c,ah^2r^2,a(23/48)^2,1/8\}).
\]
For the new first term also require $2Nc_0>1$;
for the second zero-cell option require $N\kappa b^2\ge1/4$.
If either extra condition is unavailable, retain the corresponding old term.
Any fixed collection of permitted windows may be evaluated, and the minimum
of its complete six-term bounds remains valid.
\end{lemma}

\begin{proof}
We give the entire frequency partition to make clear what changes.
On $[1/2,b]$, symmetrization and
$\cos u\le1-u^2/2+u^4/24$ yield
\[
 |f(t)|^2\le1-t^2+(M+3)t^4/12,
 \qquad |f(t)|^n\le e^{-n c(t^2)},
 \quad c(u)=u/2-(M+3)u^2/24.
\]
Bounding $c(t^2)$ below by $c$ and using
$\int_{1/2}^{b}t^{-3}dt<2$ gives the old first term.
Concavity supplies the more precise chord bound
\[
 c(u)\ge c_0+s(u-1/4),\qquad 1/4\le u\le81/100,
\]
because the difference is exactly
$(M+3)(u-1/4)(81/100-u)/24$. If $s>0$, extend only the resulting
positive exponential majorant to the half-line. The preceding rational
Gaussian-tail lemma, with lower endpoint $1/2$, gives
\[
 \sqrt n\int_{1/2}^{b}e^{-nc(t^2)}t^{-3}dt
 \le e^{-n(c_0-s/4)}\sqrt n
       \int_{1/2}^{\infty}e^{-ns t^2}t^{-3}dt
 \le2\sqrt n e^{-nc_0}L(ns/4).
\]

For $t\ge b$, conditioning on the actual event and convex Jensen give
\[
 |f(t)|^n\le e^{-nE\sin^2(t(X-X')/2)}
 \le E[e^{-na\sin^2(td/2)}\mid E],\qquad d=|X-X'|.
\]
Positive Tonelli allows a separate partition for every $d>0$.
Put $\ell=2\pi/d$. The zero cell is $[0,\ell/2]$; the other cells are
$\ell[j-1/2,j+1/2]$, $j\ge1$. In each nonzero cell take the core
$t=\ell(j+v)$, $|v|\le r$. These cores are disjoint even for $r\le1/2$.
For our smaller range $r\le1/4$, the elementary sine bound gives
$|\sin(\pi v)|\ge h|v|$ throughout a core, and
$|\sin(\pi v)|\ge hr$ throughout its wings. The latter uses that
$\sin(\pi|v|)$ increases on $[0,1/2]$.

Since $j-r>0$, Taylor's formula on the two sides gives
\[
 \frac{(j-v)^{-3}+(j+v)^{-3}}2
 \le j^{-3}+6(j-r)^{-5}v^2.
\]
Integrate this majorant against $e^{-nah^2v^2}$ on the full real line.
The result is
$\sqrt\pi[j^{-3}/\sqrt{nah^2}
  +3(j-r)^{-5}/(nah^2)^{3/2}]$.
The Jacobian is
$t^{-3}dt=d^2(4\pi^2)^{-1}(j+v)^{-3}dv$.
Summing and averaging $d^2$ gives the second term. The two series tails
are bounded by their decreasing integrals after $20$, producing exactly
$S_3,S_5$. The total wing weight is at most
$\int_{\ell/2}^{\infty}t^{-3}dt=d^2/(2\pi^2)$, giving the third term.
Extending cores or wings below $b$ only adds a positive integral.

In the zero cell write $u=td/2\le\pi/2$.
For $u\le1/2$, use $\sin u\ge(23/24)u$ and the rational
Gaussian-tail lemma with rate $\kappa$; its two choices give $Z_n$.
For $u\ge1/2$, use $\sin u\ge23/48$ and total weight
$1/(2b^2)$, giving the fifth term. Shared endpoints have zero integration
measure, including when the conditional distance distribution has atoms.

Finally keep the comparator on the entire interval $[1/2,\infty)$.
Its absolute-value bound splits its integral into
\[
 \sqrt n\int_{1/2}^{\infty}e^{-nt^2/2}t^{-3}dt
 +\frac{n^{3/2}|\gamma|}{6}
                 \int_{1/2}^{\infty}e^{-nt^2/2}dt.
\]
The first part is at most $2\sqrt n e^{-n/8}L(n/8)$.
For the second, $t\ge1/2$ implies
\[
 \int_{1/2}^{\infty}e^{-nt^2/2}dt
 \le2\int_{1/2}^{\infty}t e^{-nt^2/2}dt
 =2e^{-n/8}/n.
\]
Pearson's elementary identity
$E[(X^2-1-\gamma X)^2]=EX^4-1-\gamma^2\ge0$
gives $|\gamma|\le\sqrt{M-1}$ and proves the sixth term.
Any separately proved upper bound $|\gamma|\le g$ may replace
$\sqrt{M-1}$ in that term. The enclosure uses
$g=\min(1.85,\sqrt{M-1})$.

For uniformity, $\sqrt n e^{-nc'}$ decreases when $2nc'>1$,
whereas $n^{-1}$ and $n^{-1/2}e^{-nc'}$ always decrease.
The derivative of $L(x)$ is
$-(x^2+4x+6)/(x^2+4x+2)^2<0$.
The new first term is thus decreasing under its stated gate, and the
comparator decreases for every $n\ge5$. The earlier tail lemma proves
the separate $N\kappa b^2\ge1/4$ gate for its new zero-cell option.
These comparisons concern the bounding expressions, so they are uniform
even when the marginal law varies with $n$.
\end{proof}\paragraph{Events, not artificial kernels.}\label{cut49:fourier:events}
The entry uses only the following actual choices, with the direct rate
from the same fourth moment in each case:
\[
\begin{array}{c|c|c|c}
 \text{event}&a&d_0&M_E\\\hline
 (X-X')^2\ge9/25 &(2-9/25)^2/(2M+6)&3/5&2/a\\
 XX'<0, D\le d\le.8 &1/(2r_*^2)&\sqrt{2-d}&4r_*^2\\
 XX'<0, p\in[l,u] &2\min(l(1-l),u(1-u))
 &\max(A_l+B_l,\sqrt{2-d})&2/a
\end{array}
\]
The first row follows from Cauchy using
$\E[(X-X')^2;E]\ge2-d_0^2$ and
$\E(X-X')^4=2\mu_4+6$. The second is Lemma~\ref{cut49:entry201:event}; all its moment and $\alpha<3$ gates are checked at use.
In the third, $D\le d<2$ ensures a genuine central gap,
$P(XX'<0)=2pq$, and restriction of the squared distance gives $M_E=2/a$.
Taking a minimum of \emph{complete} high-frequency bounds is legitimate. No tail is
dropped and the third construction also allows a cell crossing $1/2$.

\section{Coefficient atlas for the enclosure}\label{cut49:re:atlas}
These formulas are inequalities for one actual law. They are the two
coordinate descriptions of the same enclosure: signed moments before a
probability gap is located, and conditional moments afterwards. A bound
is used only with its stated gates, and every refinement is an intersection.
\subsection{Signed moments with a common first moment}
Put $a_\pm=(1\pm\beta)/2$, $b_\pm=(\rho\pm\gamma)/2$.
Cauchy gives $m^2\le a_+a_-$ and $b_\pm m\ge a_\pm^2$.
Consequently
\[
 \frac{(1+\beta)^{3/2}}{\sqrt{1-\beta}}-\rho
 \le\gamma\le
 \rho-\frac{(1-\beta)^{3/2}}{\sqrt{1+\beta}},\qquad
 m\ge\max\left\{\frac{1+\beta^2}{2\rho},
                  \frac{a_+^2}{b_+},\frac{a_-^2}{b_-}\right\}.
\]
These statements also hold with a zero atom. Projection of $X|X|$ onto
the span of $X$ and $|X|$ gives
\[
 \mu_4\ge\rho^2+\frac{(\gamma-\rho\beta)^2}{1-\beta^2}.
\]
For completeness, this follows by minimizing the nonnegative expectation
$E(X|X|-uX-v|X|)^2$ over $u,v$; its Gram matrix has diagonal entries
one and off-diagonal entry $\beta$, and is positive definite.

Two support-localizing moment matrices give the additional upper bounds
as follows. If
$D_\pm=Ra_\pm-b_\pm>0$, then
\[
 2m+\sum_\pm\frac{(Rm-a_\pm)^2}{D_\pm}\le R,
 \qquad
 \mu_4\le R\rho-\sum_\pm\frac{D_\pm^2}{Rm-a_\pm}.
\]
Indeed the first follows by summing the determinant inequalities for
\[E[(R-|X|)(1,|X|)^{\mathsf T}(1,|X|);\pm X>0].\]
The second follows from the corresponding matrices with an additional
factor $|X|$. One more factor $|X|$ yields, sign by sign,
\[
 E[X_\pm^5]\le R E[X_\pm^4]
 -\frac{(Rb_\pm-E[X_\pm^4])^2}{D_\pm}.
\]
Zero denominators are treated by dropping their nonnegative correction;
they never imply exclusion without a separate contradictory numerator.
The common $m$ is retained throughout.

A second useful upper bound for $m$ follows by integrating
$(x-a)^2(R-x)\ge0$ for $x=|X|\in[0,R]$. For every fixed $a>0$,
\[
 2m\le\frac{R-\rho+2a+Ra^2}{a(2R+a)}.
\]
The cover uses only the eleven rational choices $a=2/10,\ldots,12/10$.

Write $\delta=\gamma-\rho\beta$. Cauchy and a square completion on
each sign give the signed constant bound
\begin{align*}
 \frac\rho2\mu_4-\frac14\eta_5
 &\le\frac{\rho^3}{4}
 -\frac{(\rho+\gamma)(\delta_+)^2}{8(1+\beta)^2}
 -\frac{(\rho-\gamma)((-\delta)_+)^2}{8(1-\beta)^2}.
\end{align*}
Indeed on either sign, write $a=\E X_s^2$, $b=\E X_s^3$,
$c=\E X_s^4$, $f=\E X_s^5$. Cauchy gives $c\ge b^2/a$,
$f\ge c^2/b$, and
$\rho c/2-f/4\le\rho^2b/4-(c-\rho b)^2/(4b)$.
Distance to the half-line $c\ge b^2/a$ proves the stated penalty
on summing both signs. Independently,
$\mu_4^2\le\rho\eta_5$ and the projection bound above give
\[
 \frac\rho2\mu_4-\frac14\eta_5
 \le\frac{\rho^3}{4}
 -\frac{(\gamma-\rho\beta)^4}{4\rho(1-\beta^2)^2}.
\]
The larger of the two nonnegative deductions is retained.

\subsection{An actual central gap without assuming $D<2$}
If zero belongs to the support, the signwise gap calculation gives $J\le\rho$.
Thus $D\ge\sqrt{10}\rho-\gamma\ge(\sqrt{10}-1)\rho$.
Whenever $D<(\sqrt{10}-1)\rho$, there is therefore a central gap
$(-A,B)$, $A,B>0$, with no zero atom. Put $p=P(X>0)$, $q=1-p$,
$T_+=b_+-mB^2$ and $T_-=b_--mA^2$. The exact gap decomposition is
\[
 D=mQ+(\sigma-1)T_++(\sigma+1)T_-+3(T-J_o),\qquad
 Q=(\sqrt{\sigma+1}A-\sqrt{\sigma-1}B)^2,
 \quad\sigma=\sqrt{10}.
\]
All terms are nonnegative. Hence $T\le D/(\sigma-1)$ and
\[
 (A+B)^2\ge A^2+B^2\ge\frac{\rho-D/(\sigma-1)}m.
\]
Moreover $A\le a_-/m$, $B\le a_+/m$ and $J\le\rho+2m$.
The last bound gives the compulsory defect
\[
 D\ge \sigma\rho-\gamma-6m.
\]
The inequalities $T-J_o\ge2BR_++2AR_-$ give the useful additional
endpoint constraint
\[
 D\ge(\sigma-1)b_++6a_+B-(\sigma+5)mB^2.
\]
If $D<(\sigma-1)b_+$, this downward quadratic forces $B$ to be at
least its positive root. The reflected constraint has $a_-,b_-,A$
and the coefficients $\sigma+1,\sigma+7$. In a closed moment box,
positive coefficients are rounded down, negative coefficients up, and
the lower square-root endpoint is used in the resulting root bound.

The actual opposite-sign event has probability $2pq$.
The first two moment constraints give lower bounds for both $p$ and $q$.
For independent copies one has exactly
\[
 E[(X-X')^2\mid XX'<0]
 =\frac{1+\beta(1-2p)+4m^2}{2pq}
 =\frac{1+4m^2+\beta}{2p}
  +\frac{1+4m^2-\beta}{2q}.
\]
For an upper bound, use the upper $m$ endpoint and both endpoints of
$p$ and $\beta$. The expression is convex in $p$ and affine in $\beta$.
This improves the former $2/P(XX'<0)$ bound inside the complete HIGH
estimate. Both copies still have the same actual marginal.

\subsection{Preserving central and outer integral costs}
Let
\[
 P_1(A,B)=B^3/3+AB^2/2-A^3/6,\quad
 P_2(A,B)=B^4/4+AB^3/3-A^4/12,\quad G=\rho P_1-P_2.
\]
Then $J_{1,c}=mP_1$, $J_{2,c}=mP_2$. Since
$\partial_BG=B(A+B)(\rho-B)$,
\[
 G(A,B;\rho)\le G(A,\rho;\rho)
 =\rho^4/12+\rho^3A/6-\rho A^3/6+A^4/12.
\]
The last expression increases with $A$ because its derivative is
$(\rho-A)^2(\rho+2A)/6$, and is convex in $\rho$ because its second
$\rho$ derivative is $\rho^2+\rho A$. Its closed-rectangle maximum is
therefore bounded using the upper $A$ endpoint and both $\rho$ endpoints.

Positive outer gaps satisfy
$(t-a)(\rho-t)\le(\rho-B)_+^2/4$. On negative outer gaps the positive
part is bounded by $(R-\rho)^2/4$ after extracting $|t|$.
Integration gives
\[
 3(\rho J_1-J_2)_o
 \le\frac38\{(\rho-B)_+^2T_++(R-\rho)^2T_-\}.
\]
The available weighted outer budget is $D-mQ$, not the entire $D$.
The cover intersects this budget with the individual bounds
$T_+\le b_{+,u}-m_lB_l^2$ and $T_-\le b_{-,u}-m_lA_l^2$.

For the signed phase term $U=J_1+(\mu_4-3)/6$, retain
\[
 -RT_-/2\le J_{1,o}\le RT_+/2.
\]
The central minimum of $P_1$ occurs at $B_l$ and an $A$ endpoint;
its maximum occurs at $B_u$ and $A$ clipped to the interval at $A=B_u$.
A nonnegative polynomial bound is multiplied by $m_u$ for an upper
bound and $m_l$ for a lower bound; a negative one uses the opposite
endpoint. For example, the binary endpoints $A=3/2$, $B=2/3$ make this
maximum negative; the sign switch is necessary even when $D<2$.
The fourth-moment sign intervals also sharpen
$K_1=E[X|X|^3]-3\beta-3\rho\gamma/2$.
The complete $K_1,U,\phi$ rectangles are used in the signed expression
$|CK_1+\phi U|$, without replacing it by two separate absolute values.

\paragraph{A complete coarse bound before a central gap is licensed.}
Always $-\mu_4/6\le J_1\le\mu_4/2$ and
$-3J_2-\eta_5/4\le0$. On a gap of length $h=b-a$, these follow from
\[
 \int_a^b(t-a)t\,dt+\tfrac12\int_a^bt^2dt=hb^2/2,\quad
 \tfrac32\int_a^bt^2dt-\int_a^b(t-a)t\,dt=h(2a+b)^2/6,
\]
and, on a negative gap $(-v,-u)$, from
$\tfrac13\int_u^v x^3dx-\int_u^v(v-x)x^2dx=u^3(v-u)/3$.
For $1\le r\le\rho$ and actual support radius $R\le6$, set
\[
 \Gamma(R,r)\ge\max_{r\le v\le R}
 \left\{-\frac{v^2}4+\frac{(R+r)v}3-\frac{Rr}2+
                      \frac{r^3(2R-r)}{12v^2}\right\}_+.
\]
Then
\[
 3(\rho J_1-J_2)\le\frac32\Gamma(R,r)(\rho-\gamma)
                              +\frac38\rho^2J.
\]
On negative gaps extract $x=|t|$ and discard the nonpositive part
$x<\rho$. Replacing $\rho$ by $r$ and the outer endpoint by $R$
gives $\int_r^v(R-x)x(x-r)\,dx$. Dividing this integral by $v^2$
gives exactly the displayed expression. Signwise Tonelli now charges
$\Gamma\E X_-^3=\Gamma(\rho-\gamma)/2$.
On positive gaps $(t-a)(\rho-t)\le\rho^2/4$, giving $\rho^2J/8$.
We use 64 closed $v$-intervals: on each, maximize the
concave quadratic at its clipped vertex $2(R+r)/3$ and the positive
reciprocal term at the left endpoint. This proves its interval bound.
The alternative bound $\Gamma\le5/2$ follows at $r=1,R=6$ by
$3v^4-28v^3+66v^2-11>0$ on $[1,6]$: the endpoint values are $30,205$,
and the interior critical values are $(85\mp35\sqrt5)/2>0$.
These monotone enclosing parameters cover all $r\ge1,R\le6$.
\subsection{Moment and gap coefficients for a whole cell}
\label{cut49:mom:conditional-boxes}

The following coefficient rules are collected here to make the finite
calculation transparent. Every rule is an inequality for an actual law;
there is no extremal-law assumption inside a cell.
Write $v=1-\theta$. For $p\in[l,u]\subset(0,1)$ set
\begin{gather*}
 s_0=\min\{l(1-l),u(1-u)\},\quad
 v_* =\max\{\sqrt{u/(1-u)},\sqrt{(1-l)/l}\},\\
 r_B=(1-2s_0)/\sqrt{s_0},\quad
 g_B=\max_{t=l,u}|1-2t|/\sqrt{t(1-t)},\quad b_B=\max_{t=l,u}|1-2t|,\\
 L_3=R+2v_*,\quad L_4=R^2+2Rv_*+3v_*^2,\\
 L_5=R^3+2R^2v_*+3Rv_*^2+4v_*^3,\quad
 e_B=r_B\frac{1-4s_0+2s_0^2}{s_0(1-2s_0)}.
\end{gather*}
The polynomials $T_3=x+2a$, $T_4=x^2+2ax+3a^2$ and
$T_5=x^3+2ax^2+3a^2x+4a^3$ are nonnegative. Their maxima at
$x=R$, $a=v_*$, in Lemma~\ref{cut49:re:pricing}, give, if $\theta\le\Theta<1$,
\begin{align}
 \rho&\le r_B+(L_3-r_B)\Theta,&
 |\gamma|&\le g_B+(L_3-g_B)\Theta,\notag\\
 |\beta|&\le b_B+(1-b_B)\Theta,&
 \mu_4&\le(1-\Theta)^2(1/s_0-3)+L_4\Theta,\notag\\
 \eta_5&\le(1-\Theta)^2e_B+L_5\Theta.&&\label{cut49:mom:boxes}
\end{align}
Replacing $\theta$ by $\Theta$ uses the explicitly checked inequalities
$L_3>r_B,g_B$, $L_4>2(1/s_0-3)$ and $L_5>2e_B$.
The exact binary factors are $v^{3/2},v,v^2,v^{5/2}$; the upper bounds
use $v^{3/2}\le v$ and $v^{5/2}\le v^2$. Thus cells crossing $1/2$ and
zero conditional variances are included. The same cubic expansion, in the
later domain $p\in[.37,.47]$, $A\ge.75$, $B\ge1.05$, gives
\begin{equation}\label{cut49:mom:rho-lower}
 \rho\ge v^{3/2}\rho_B(p)+3\min(A,B)\theta
       \ge\rho_B(p)+.45\theta,\qquad
 \rho_B(p)=\frac{1-2pq}{\sqrt{pq}}.
\end{equation}
Indeed $\rho_B<1.11$ there and $v^{3/2}\ge1-3\theta/2$.
Also $\beta\ge1-2p-2(1-p)\theta$.

For the signed coefficient
$K_1=\E[X|X|^3]-3\beta-3\rho\gamma/2$, define
\[
 K(p,v)=(1-2p)\{(3v^3/2-v^2)/(pq)-3v^3+v^2+3v\}.
\]
Direct substitution gives $K_1(V)=-K(p,v)$, with the coefficient $3$
unchanged. For $v\ge2/3$, $K$ decreases in $p$: both coefficients in
braces are nonnegative and $(1-2p)/(pq)$ decreases. Its derivative in $v$
has the sign of $1-2p$, since the derivative of the braces is at least
$9v^2-6v+3>0$. Hence both variance endpoints are retained. The error is
\begin{equation}\label{cut49:mom:K-error}
 |K_1(X)-K_1(V)|\le
 \Theta\{L_4+3+\tfrac32L_3(R_\rho+g_B)\}
\end{equation}
when $\rho\le R_\rho$. This follows from the conditional expansions and
$\rho\gamma-\rho_V\gamma_V=
\rho(\gamma-\gamma_V)+\gamma_V(\rho-\rho_V)$.

Here are the remaining signed charges. They also explain the two
alternative bounds retained in the initial, coarse moment cover:
\begin{align}
 -\mu_4/6&\le J_1\le\mu_4/2,&
 -3J_2-\eta_5/4&\le0,                                  \label{cut49:mom:gap-signs}\\
 \rho J_1-J_2&\le\tfrac54(\rho-\gamma)+\tfrac18\rho^2(\rho+1),
                                                              \label{cut49:mom:oldgap}\\
 -75D/104&\le J_{1,o}\le3D/2.                           \label{cut49:mom:outer}
\end{align}
To prove the first line, a gap of length $h=b-a$ satisfies
\[
 \int_a^b(t-a)t\,dt+\tfrac12\int_a^bt^2dt=hb^2/2,\quad
 \tfrac32\int_a^bt^2dt-\int_a^b(t-a)t\,dt=h(2a+b)^2/6.
\]
Sum with weight $\tau_g$. On a negative piece $(-v,-u)$,
$\frac13\int_u^v x^3dx-\int_u^v(v-x)x^2dx=u^3(v-u)/3\ge0$;
Tonelli proves the $J_2$ assertion. For~\eqref{cut49:mom:oldgap}, the positive
part is at most $\rho^2J/8$. The negative part is at most
$\E[V\int_1^V(6-x)x(x-1)dx;V\ge1]$, $V=X_-$, hence at most
$(5/2)\E X_-^3$. The last comparison is equivalent to
$3V^4-28V^3+66V^2-11>0$ on $[1,6]$: the endpoint values are $30,205$,
and at the possible interior critical points $(7\pm\sqrt5)/2$ its value
is $(85\mp35\sqrt5)/2>0$.
For~\eqref{cut49:mom:outer}, use $J_{1,o}\ge-3T_-$,
$J_{1,o}\le3J_o\le3T$ and~\eqref{cut49:re:gapexact}.

The central integrals are $J_{j,c}=mP_j(A,B)$, with
\[
 P_1=B^3/3+AB^2/2-A^3/6,\qquad
 P_2=B^4/4+AB^3/3-A^4/12.
\]
They increase in $B$; their $A$ derivative changes from positive to negative
at $A=B$. A rectangle minimum is therefore at $B_l$ and an $A$ endpoint.
Multiply a negative minimum by an upper bound for $m$, and a nonnegative
minimum by a lower bound. For the upper central charge, maximizing at
$(\bar a,\bar b)$ if $p\le1/2$ gives
$J_{1,c}\le v^2(2-3p)/(6pq)$, a decreasing function of $p$ on that interval.
If $p\ge1/2$, maximize first at $A=B$ to get
$J_{1,c}\le(2/3)v^2q^2/p\le1/3$. Thus every cell may use
\begin{equation}\label{cut49:mom:central-upper}
 J_{1,c}\le\max\{1/3,(2-3l)/(6l(1-l))\}.
\end{equation}

For $|X|\le R$, $R\ge1$, define
\[
 k_4(R)=\frac{(R+1)^2}{R+1/2},\qquad
 k_5(R)=\frac{R^3+2R^2+3R+3/2}{R+1/2}.
\]
Then
\begin{equation}\label{cut49:mom:universal-moments}
 \mu_4\le\min\{R\rho,1+k_4(R)(\rho-1)\},\quad
 \eta_5\le\min\{R\mu_4,1+k_5(R)(\rho-1)\},\quad
 \mu_6\le R\eta_5.
\end{equation}
Writing $u=|X|$ in~\eqref{cut49:re:pricing}, the first two constants
are the endpoint values of the ratios
$u+3/2+1/[4(u+1/2)]$ and
$u^2+3u/2+9/4+3/[8(u+1/2)]$, whose derivatives are nonnegative.
The other bounds are pointwise. The following coarse cancellation estimates
continue to use the valid radius-six majorants:
The last identity in~\eqref{cut49:re:pricing} gives
\begin{equation}\label{cut49:mom:K-cancellation}
 |K_1+3\rho\beta|\le(96-18\rho)(\rho-1)/13.
\end{equation}
Indeed $u(u+2-3r/2)/(u+1/2)$ is convex on $[0,6]$, its maximum is
$(96-18r)/13$, and its negative part is at most $31/40$, smaller than that
maximum for $1\le r\le37/20$.
The other bounds needed below are
\begin{align}
 |K_1|&\le6\rho-3\rho^2/2+3|\beta|,\notag\\
 K_1&\le M+3|\beta|-3\rho g_-/2,&
 -K_1&\le3\rho+3|\beta|+(3\rho/2-3)g_-,\notag\\
 -K_q&\le(\rho/2-3/8)M-\rho/2+1/8+3|\beta|g/2,&
 -K_q&\le8\rho^3/27+(3|\beta|/2-1/2)\rho,               \label{cut49:mom:K-bounds}
\end{align}
where $\mu_4\le M$, $|\gamma|\le g$, $\gamma\ge g_-$, and
$K_q=\eta_5/4-\rho\mu_4/2-3\beta\gamma/2+\rho/2$.
For the first two lines use
$|\E[X|X|^3]-3\gamma|\le3\rho$, and note the negative coefficients
of $\gamma$. For the last line use
$\eta_5\ge\mu_4^{3/2}\ge3\mu_4/2-1/2$, or maximize
$\rho y^2/2-y^3/4$ over $y\ge0$.

\paragraph{Compatible weaker resource constants.}
One may also use the weaker binary charge
$b_I=(616/25)(1-\Theta)^2\operatorname{dist}(p_*,I)^2$.
It follows from $2(6+2\sigma)>616/25$ and
$(1-\theta)^{3/2}\ge(1-\theta)^2$.
With $K=\min(K_+,K_-)$ and $v_* =\max(v_+,v_-)$, the merged charge
\eqref{cut49:re:residual} gives
\[
 T_++T_-\le T_*=(d-b_I)\max\{v_*/[3(A_l+B_l)],(R+2v_*)/K\},
 \qquad d_o=\min\{d,(\sigma+1)T_*\}.
\]
Thus $W_o\le d_o$ and, for $R\le6$,
\[
 -75d_o/104\le J_{1,o}\le3d_o/2,\qquad
 3(\rho J_{1,o}-J_{2,o})\le\frac38
 \max\{\rho^2/(\sigma-1),(R-\rho)^2/(\sigma+1)\}d_o.
\]
These are precisely the merged-cost versions of the sign-specific
resource rule; they do not introduce another available budget.

\subsubsection{Conditional lower moments improve two signed charges}

Let $s_u=\max_{p\in[l,u]}pq$ and $\theta\le\Theta$.
The unsigned cases of~\eqref{cut49:re:conditional} give
\begin{align}
 \mu_4&\ge M_l:=\max\{1,(1-\Theta)^2(1/s_u-3)\},\label{cut49:entry201:m4lower}\\
 \eta_5&\ge E_l:=
 (1-\Theta)^{5/2}\frac{1-4s_u+2s_u^2}{s_u^{3/2}}.\label{cut49:entry201:e5lower}
\end{align}
Both binary moment functions decrease in $s\in(0,1/4]$.
For the fifth this follows because the derivative has the sign of
$s^2+2s-3/2<0$; the fourth is immediate.
The fourth-moment lower bound sharpens the lower endpoint of
$J_1+(\mu_4-3)/6$; its explicit use appears in the cell map below.

The constant signed charge is
\[
 -K_q=\frac{\rho\mu_4}{2}-\frac{\eta_5}{4}
       +\frac32\beta\gamma-\frac{\rho}{2}.
\]
The old coherent estimate used $\eta_5\ge3\mu_4/2-1/2$ and
$\mu_4\le M$. A second estimate uses \eqref{cut49:entry201:e5lower}.
Their difference is independent of $\rho$. Consequently the improved
bound is exactly the old one minus
\begin{equation}\label{cut49:entry201:eta-improvement}
 \frac14\max\{0,E_l-3M/2+1/2\}.
\end{equation}
This is a constant shift on the cell and does not spoil convexity.
It does not assert the generally invalid implication
$\eta_5\ge3M/2-1/2$ from an upper bound on $\mu_4$.

\subsubsection{Signed conditional Taylor cancellation}

Let $K_1=\E[X|X|^3]-3\beta-\tfrac32\rho\gamma$.
On a sign class put $x=|X|$, $a=|V|$, $W=X-V$ and
$\varepsilon=1$ on the positive class, $-1$ on the negative class.
In the actual central-gap domain there is no atom at zero.
With $w_3=x+2a$, $w_4=x^2+2ax+3a^2$ and
$\gamma_V=\E V^3$, exact conditional expansion gives
\begin{equation}\label{cut49:entry201:Kexact}
 K_1(X)-K_1(V)
 =\E\!\left[W^2
 \left\{\varepsilon(w_4-3-\tfrac32\rho w_3)
                   -\tfrac32\gamma_Vw_3\right\}\right].
\end{equation}
Apply~\eqref{cut49:re:conditional} to $j=2,3,4$, with $h=1$
or $h=\varepsilon$, and use
$\rho\gamma-\rho_V\gamma_V
=\rho(\gamma-\gamma_V)+\gamma_V(\rho-\rho_V)$.

The positive coefficient in \eqref{cut49:entry201:Kexact} is
$F(x,a;t)=x^2+2ax+3a^2-3-\tfrac32t(x+2a)$
with $t=\rho+\gamma_V$; the negative coefficient is $-F$
with $t=\rho-\gamma_V$. Their ranges are bounded on complete
rectangles of actual conditional means and support magnitudes.
The positive-support magnitude is at most $R$.
The same formula permits separate actual bounds on the two signed support magnitudes.
On the late domain $p<1/2$, bounds for $\gamma_V$ are
\[
 \sqrt{(1-\Theta)^3(1-2u)^2/[u(1-u)]}
 \le\gamma_V\le(1-2l)/\sqrt{l(1-l)}.
\]
The positive mean lies in
$[\sqrt{(1-\Theta)(1-u)/u},\sqrt{(1-l)/l}]$, and the negative
magnitude mean lies in
$[\sqrt{(1-\Theta)l/(1-l)},\sqrt{u/(1-u)}]$.
Use the lower bound $\rho\ge\rho_B(u)$ only after checking
$3\min(A_l,B_l)\ge(3/2)\sup_I\rho_B$. Indeed conditional expansion
and $(1-\theta)^{3/2}\ge1-3\theta/2$ then imply
$\rho\ge(1-\theta)^{3/2}\rho_B(p)+3\min(A,B)\theta\ge\rho_B(p)$.
Otherwise retain the universal lower bound one.

Since $x+2a>0$, $F$ decreases with $t$.
For fixed $t$ its Hessian is
$\left(\begin{smallmatrix}2&2\\2&6\end{smallmatrix}\right)>0$.
Its rectangle maximum is at a corner. Its minimum is among the four
corners, the minimizer on each edge, and the stationary point
$(x,a)=(3t/8,3t/8)$ if that lies in the rectangle.
On fixed-$x$ edges use $a=t/2-x/3$, clipped to the edge; on
fixed-$a$ edges use $x=3t/4-a$, also clipped.
This is a complete finite extremum list.
Multiplication by $0\le\theta\le\Theta$ uses
$\Theta\min(0,L)$ and $\Theta\max(0,U)$ for the lower and upper
coefficient bounds. It intersects this new interval with the retained
absolute-error interval; neither need be discarded.

\subsection{Generic event and complete central rectangle}
\begin{lemma}[Generic small-defect event]\label{cut49:entry201:event}\label{cut49:mom:third}
Assume $|X|\le R\le6$ and $\rho\le37/20$. Let $0\le d<2$ and $D\le d$. Define
\[
 r_*=\min\{37/20,(3+d)/(\sigma-1)\},\quad
 W=\sqrt{37d/10},\quad c=\sqrt{\sigma+1},\quad e=\sqrt{\sigma-1},\quad C_0=c+e,
\]
\[
 h_*=\frac{C_0}{6}\sqrt{12+W^2}+\frac{W}{3C_0},\qquad
 L_*=\max\left\{h_*,\frac{(2c-e)h_*+3W}{C_0}\right\}.
\]
Then $\rho\le r_*$, $A+B\le h_*$ and $|2B-A|\le L_*$.
If $\alpha=L_*/6\le3$, then
\[
 |J_1|\le\alpha(\rho+1)+(3-\alpha)d/2,\quad
 \mu_4\le h_*\rho+Rd/2,\quad
 \eta_5\le h_*^2\rho+R^2d/2.
\]
For independent copies, the actual event $E=\{XX'<0\}$ obeys
\[
 P(E)\ge a:=1/(2r_*^2),\quad
 |X-X'|\ge d_0:=\sqrt{2-d}\text{ on }E,\quad
 \E[(X-X')^2\mid E]\le2/a.
\]
\end{lemma}
\begin{proof}
The inequality $J\le\rho+1$ gives $D\ge(\sigma-1)\rho-3$.
Equation \eqref{cut49:re:gapexact} gives $w^2\le2\rho D\le37d/10$ for
$w=cA-eB$. Since $ce=3$,
\[
 A+B=\frac{C_0}{6}\sqrt{w^2+12AB}-\frac{w}{3C_0},
 \qquad 2B-A=\frac{(2c-e)(A+B)-3w}{C_0}.
\]
Use $AB\le1$, $|w|\le W$; also $2B-A\ge-(A+B)$.
Now $J_{1,c}=J_c(2B-A)/6$, $|J_{1,o}|\le3J_o$ and
$0\le J_o\le T\le d/2$ imply
$|J_1|\le\alpha J+(3-\alpha)J_o$.

For either inner endpoint $b_*$ and $b_*\le x\le R$,
\[
 x^4-h_*x^3\le R x(x^2-b_*^2),\qquad
 x^5-h_*^2x^3\le R^2x(x^2-b_*^2).
\]
Indeed $h_*\ge b_*$, and after factoring $x-b_*$ the remaining
comparisons follow from $x\le R$. Integrate and use $T\le d/2$.
Finally $m^2\le pq$ and $m\ge1/(2\rho)$ give $2pq\ge a$;
the conditional squared distance is bounded by its unconditional
expectation $2$, divided by the event probability.
\end{proof}

The central part of $\rho J_1-J_2$ is exactly
$mG(A,B;\rho)$, where $G=\rho P_1-P_2$.
Let a proved probability cell give
$A\in[A_l,A_u]$, $B\in[B_l,B_u]$, $m\in[m_l,m_u]$.
Here
\[
 A_u=\sqrt{u/(1-u)},\quad B_u=\sqrt{(1-l)/l},\quad
 m_l=\sqrt{(1-\Theta)s_0},\quad
 m_u=\sqrt{\max_{p\in[l,u]}pq}.
\]

\paragraph{The closed-box Bernstein identity.}\label{cut49:re:bernstein}
For a polynomial $P(x)=\sum_{k=0}^dt_kx^k$,
\[
 P(x)=\sum_{i=0}^d b_i\binom di x^i(1-x)^{d-i},\qquad
 b_i=\sum_{k\le i}t_k\binom ik/\binom dk.
\]
Indeed $\binom di\binom ik=\binom dk\binom{d-k}{i-k}$ and the
binomial theorem recover each monomial $x^k$. The basis terms are
nonnegative and sum to one on the closed unit interval. Tensor products
prove the two-variable version below, including all boundary points.

\begin{lemma}[Complete central rectangle]\label{cut49:entry201:central}
Write $G(A_l+(A_u-A_l)x,B_l+(B_u-B_l)y;\rho)
=\sum t_{kh}x^ky^h$. Define the 25 rational coefficients
\[
 b_{ij}(\rho)=\sum_{\substack{k\le i\\h\le j}}
 t_{kh}\frac{\binom{i}{k}}{\binom4k}
          \frac{\binom{j}{h}}{\binom4h},\qquad 0\le i,j\le4,
\]
and
\[
 \mathcal G(\rho)=\max_{i,j}
       \{m_l b_{ij}(\rho),\,m_u b_{ij}(\rho)\}.
\]
Then $\rho J_{1,c}-J_{2,c}\le\mathcal G(\rho)$.
Moreover $\mathcal G$ is convex in $\rho$.
\end{lemma}
\begin{proof}
Apply the Bernstein identity~\eqref{cut49:re:bernstein} separately to the
two coordinates after the affine substitution. The Bernstein basis is nonnegative and sums to one on the closed square,
so $G$ is a convex combination of its 25 coefficients.
For each coefficient retain both $m$ endpoints, or equivalently use
$m_u$ when it is nonnegative and $m_l$ otherwise. This gives a uniform
bound on the entire three-dimensional rectangle, including boundaries.
Each endpoint-multiplied coefficient is affine in $\rho$; their maximum
is convex. No interior stationary point is omitted.
\end{proof}

\subsection{One interval-feedback theorem}
\label{cut49:entry:feedback}\label{cut49:fourier:uniform}

The signed identity from the preceding section is
\begin{equation}\label{cut49:entry:remainder}
 \cR=C\{3\beta J/2+3\rho J_1-3J_2-K_q\}
 +H(CK_1+\varphi U)-H^2(C\rho+\gamma\varphi/2)+E_n,
\end{equation}
where $U=J_1+(\mu_4-3)/6$, $\varphi=\phi(H/\sqrt n)$.
For any one coherent low-frequency triple and complete high-frequency value $K_H$ at floor $N$,
\begin{equation}\label{cut49:entry:fourier-budget}
 |E_n|\le K_F+s_F|H|,\quad
 K_F=c_\pi(L_0+K_H+L_2/\sqrt N),\quad s_F=c_\pi L_H,\quad c_\pi\ge2/\pi.
\end{equation}
We use $16/25$ before the signed phase and $212/333$
afterward; both are valid since $\pi>333/106$.

\begin{lemma}[Continuous cells and feedback]\label{cut49:entry:cell-theorem}\label{cut49:entry:corners}
Suppose bounds for one actual cell give
\[
 r_0\le\rho\le r_1,\quad0<c_0\le C\le c_1,\quad
 |\beta|\le b,\quad |\gamma|\le g,\quad \mu_4\le M,\quad D\le d,
\]
$K_1\in[k_l,k_u]$, $U\in[u_l,u_u]$, and
$0\le\varphi_-\le\varphi\le2/5$. Freeze the shared outer charge $d_o$,
lower fifth moment $E_l$, upper $s_u=\max_{p\in[l,u]}p(1-p)$, and
central envelope $\mathcal G$ from the preceding bounds.
Define
\begin{align*}
 A(r)&=\tfrac32 b(r+2\sqrt{s_u})+3\mathcal G(r)
 +\tfrac38\max\{r^2/(\sigma-1),(R-r)^2/(\sigma+1)\}d_o\\
 &\quad +(r/2-3/8)M-r/2+1/8+\tfrac32bg
 -\tfrac14\max\{0,E_l-3M/2+1/2\},\\
 B(c)&=\max_{k\in\{k_l,k_u\},\,u\in\{u_l,u_u\},\,
                      \psi\in\{\varphi_-,2/5\}}|ck+\psi u|,\\
 Q(c,r)&=cr+\tfrac12\varphi_-(3.16r-3-d)
       -\tfrac12(2/5-\varphi_-)(d-.16)_+.
\end{align*}
If all four corner values of $Q$ are positive, then, for every real $H$,
\begin{align}
 \cR&\le\max_{(c,r)\in\{c_0,c_1\}\times\{r_0,r_1\}}
       \{K_F+cA(r)+(B(c)+s_F)|H|-Q(c,r)H^2\}\notag\\
 &\le\max_{c,r\text{ corners}}
       \{K_F+cA(r)+(B(c)+s_F)^2/(4Q(c,r))\}.           \label{cut49:entry:four-corners}
\end{align}
If $T$ is a strict upper bound for the final maximum in
\eqref{cut49:entry:four-corners}, the weight-removal step in
Lemma~\ref{cut49:re:weight} gives
\begin{equation}\label{cut49:entry:feedback-values}
 D<\frac{15000T}{997N},\quad
 \varphi>\frac{997}{2500}\left(1-\frac{5T}{2N}\right),\quad
 C<\CE+V/N\quad\text{if also }\cR/\rho<V.
\end{equation}
Only positive density floors are used. These statements remain valid at
$H=0$, on closed cell edges, and for zero conditional variance.
\end{lemma}
\begin{proof}
The constant coefficient follows from the shared outer budget,
Lemma~\ref{cut49:entry201:central}, and~\eqref{cut49:mom:K-bounds}, using
\[
 J\le\rho+2ABm\le\rho+2(1-\theta)^{3/2}\sqrt{pq}
       \le\rho+2\sqrt{s_u}.
\]
This follows from $A\le\bar a$, $B\le\bar b$, $m^2=(1-\theta)pq$;
it is used only in the charge $3|\beta|J/2$. Since $\gamma\ge3.16\rho-3-d$, its negative
part is at most $(d-.16)_+$. Using the upper density only for that negative
part gives $C\rho+\gamma\varphi/2\ge Q$. Absolute value is separately
convex in each coordinate of $ck+\psi u$, proving the formula for $B$.
For fixed $r$ and $|H|$, the resulting upper bound is convex in $c$.
Take those endpoints \emph{before} completing the square. At either
endpoint $c>0$, it is convex in $r$, since $\mathcal G$ and the maximum of the two outer quadratics are
convex, every other term of $A$ is affine, and $Q$ is affine in $r$. This proves the first line of~\eqref{cut49:entry:four-corners};
$fh-qh^2\le f^2/(4q)$ proves the second. Bilinearity of $Q$ proves
its positivity throughout the rectangle from its corners.

\end{proof}

\subsection{Signed phase curvature in the resource schedule}
\label{cut49:re:signed-phase-repair}
The auxiliary Fourier orientation does not change the signed third moment
in the actual transfer identity. The paired phase estimate must therefore
retain its sign. On a fixed resource cell put
\[
 L=\frac{997}{2500}\left(1-\frac{5r}{2N}\right),\qquad
 g_0(\rho)=3.16\rho-3-d,\qquad U=.4,
\]
where $r,d$ are the preceding raw and defect upper bounds, and the
checker requires $0\le L\le U$. The actual density lies in $[L,U]$
and the actual signed moment satisfies $\gamma\ge g_0(\rho)$. Consequently
\begin{equation}\label{cut49:re:sign-aware-price}
 C\rho+\tfrac12\gamma\phi
 \ge q_{\rm s}(C,\rho):=
 C\rho+\tfrac12\min\{L g_0(\rho),U g_0(\rho)\}.
\end{equation}
This follows by first multiplying the lower bound on $\gamma$ by the
nonnegative actual density, and then minimizing over its two endpoints.
The simpler expression $C\rho+Lg_0(\rho)/2$ is licensed if $g_0\ge0$,
or if nonnegativity of the actual $\gamma$ has been separately proved.
It is not licensed merely by reflecting the auxiliary Fourier integral.

The corrected price preserves the four-corner reduction. For fixed $C$,
$q_{\rm s}$ is concave in $\rho$, being affine plus the minimum of two
affine functions. For fixed $\rho$ it is affine in $C$. In the phase
envelope $A+f|H|-q_{\rm s}H^2$, the remaining terms are separately
convex in $C,\rho$, by the moment and coefficient bounds already proved.
Thus the entire envelope is separately convex, and two successive
one-dimensional endpoint bounds place its maximum at the four corners.
The same argument applies after adding the affine conserved-resource
curvature $e_\rho$ and subtracting a barrier price $K\rho$.
If the four corrected divisors are positive, concavity makes them
positive throughout the rectangle, and square completion gives the
valid raw and resource bounds
\[
 \max_{C,\rho\ {\mathrm{corners}}}\left(A+\frac{f^2}{4q_{\rm s}}\right),
 \qquad
 \max_{C,\rho\ {\mathrm{corners}}}
       \left(A+\frac{f^2}{4(q_{\rm s}+e_\rho)}-K\rho\right).
\]
Any threshold-root export likewise uses the corrected corner curvature.

For completeness, nonnegativity in the later acquired domains follows
directly from their constraints. Write $B=\mathbf1_{\{X>0\}}$,
$V=\mathbb E(X\mid B)$,
$W=X-V$, $\theta=\mathbb E W^2$, and $p=\mathbb P(X>0)$. Conditional
centering and $W\ge-R-V$ imply
\[
 \gamma=\mathbb EV^3+3\mathbb E(VW^2)+\mathbb EW^3
 \ge\mathbb EV^3+2\mathbb E(VW^2)-R\theta.
\]
For $p\in[l,u]\subset(0,1/2)$ and $\theta\le\Theta<1$ this gives
\begin{equation}\label{cut49:re:actual-gamma-license}
 \gamma\ge
 \frac{(1-\Theta)^{3/2}(1-2u)}{\sqrt{u(1-u)}}
 -\left(R+2\sqrt{\frac{u}{1-u}}\right)\Theta.
\end{equation}
Here the binary cubic decreases in $p<1/2$, and the right side decreases
as its variance bound increases. These are bounds on the same actual law
throughout every later call. The fixed acquired input domains license
\eqref{cut49:re:actual-gamma-license} before those calls are made.
The exact domain check covers 302 input rows: 45 central-input rows,
45 left-input rows, 45 first-confinement feedback rows, 45 Sperner feedback
rows, 84 signed-resource rows, and 38 final left-radius rows. With the
common already acquired radius $R=3.84$, each right side is greater than
$.096$. Thus all later paired calls have an actual nonnegative signed
third moment, independently of any auxiliary Fourier orientation.

There is also a convenient cell-constant correction. Since $\rho\ge1$,
$(-g_0(\rho))_+\le(d-.16)_+$, so the weaker price
\[
 C\rho+\tfrac12L g_0(\rho)
       -\tfrac12(U-L)(d-.16)_+
\]
is valid. This is the correction selected for the additional index--46
calculation. Both forms keep the raw remainder separate from the
barrier-subtracted resource budget.

The preliminary and complementary resource bounds use
\eqref{cut49:re:sign-aware-price}. Every exported resource, including
threshold curvature, must satisfy these corrected inequalities. Equality
of two corner maxima alone is insufficient, since the threshold bounds
also depend on the phase coefficients. Later uses of the simpler formula require their own explicit
nonnegative signed-moment or small-defect license.

\subsection{The unsigned first-stage cover is the same enclosure}
\label{cut49:re:unsigned}
The first finite barrier uses a cover of $\rho$ alone before introducing
the signed cover. Here are its exact coefficient choices, so that the
first barrier needs no separate inherited entry theorem. Fix a closed
cell $\rho\in[l,u]\subset[1,1.85]$, actual radius $1\le R\le6$, and
write $k_j(R)$ for the radial ratios of Lemma~\ref{cut49:re:pricing}. Set
\begin{gather*}
 M=\min\{11.1,Ru,1+k_4(R)(u-1)\},\quad
 g=\min\{1.85,u,\sqrt{M-1}\},\\
 E_5=\min\{RM,1+k_5(R)(u-1)\},\qquad
 b=\sqrt{[u\sqrt{u^2+8}-u^2-2]/2}.
\end{gather*}
The last expression solves $\rho\ge(1+\beta^2)/\sqrt{1-\beta^2}$,
proved by $\rho m\ge a_+^2+a_-^2$ and $m^2\le a_+a_-$.
If a previous pass licenses $D\le d<2$, the generic-event lemma also
permits the ceilings $h_*u+Rd/2$ and $h_*^2u+R^2d/2$ for $M,E_5$;
all dependent quantities are recomputed after taking these minima.
Use the triangle LOW triple and the minimum of complete permitted HIGH
values, with the same moment cap. Put
\[
 K_F=\tfrac{16}{25}(L_0+K_H+L_2/\sqrt N),\qquad s_F=\tfrac{16}{25}L_H.
\]
The radial identity bounds $|K_1|$ by
\[
 k_{\rm rad}=3ub+
     \frac{R(R+2-3l/2)}{R+1/2}(u-1),
\]
provided the displayed radial coefficient is at least $31/40$.
Indeed its positive maximum is at $R$ and its negative part is at most
$31/40$ on $[0,R]$ for $1\le\rho\le1.85$; convexity of
$x(x+2-3\rho/2)/(x+1/2)$ proves this directly.
Two additional valid bounds are
\[
 k_{\rm old}=6u-3u^2/2+3b,\qquad
 k_R=Ru/2+\max_{t=l,u}|R/2-3t/2|g+3b.
\]
For $k_R$, bound the positive and negative fourth moments separately by
$R(\rho\pm\gamma)/2$. Before defect feedback use
$k=\min(k_{\rm rad},k_{\rm old},k_R)$.
After feedback one may replace $k_{\rm old}$ by $\max(k_+,k_-)$, where
\begin{gather*}
 g_l=3.16l-3-d,\quad g_u=3.16u-3-d,\\
 k_+=M+3b-\tfrac32\min_{r=l,u;,v=g_l,g_u}rv,\\
 k_-=3u+3b+\max_{r=l,u;,v=g_l,g_u}(3r/2-3)v.
\end{gather*}
These substitutions bound the lower affine graph for the actual $\gamma$,
not an independently chosen skewness.

Put $\Delta=2u$ initially and $\Delta=(1-3.16)l+3+d$ after feedback.
The constant coefficient of the raw identity is bounded by
\begin{align*}
 A_* =\min\bigl\{&\tfrac32b(2u+1)+\tfrac32\Gamma(R,l)\Delta
       +\tfrac38u^2(u+1)+\tfrac8{27}u^3-l/2,\\
      &\tfrac32b(2u+1)+2uM-l/2\bigr\}.
\end{align*}
This is the coarse gap bound above together with the two bounds for
$-K_q$. The density coefficient obeys
\[
 v_H=\min\{.4j+.4\max(2,M-3)/6,
                     .4\max(1/2,2M/3-1/2)\},
\]
where $j=3(u+1)$ initially; the generic-gap choice is
$j=\alpha(u+1)+(3-\alpha)d/2$, with its checked $\alpha<3$.
For a score endpoint enclosure $C\in[c_-,c_+]$, the resulting raw
bound is
\begin{equation}\label{cut49:re:unsignedbound}
 K_F+c_+A_*+\frac{(c_+k+v_H+s_F)^2}{4q_*},
 \qquad q_*>0,\quad k\ge0,\quad A_*>0,
\end{equation}
where initially $q_*=c_-l-g/5$. After feedback use the larger of that
quantity and
\[
 c_-l+\varphi_-(3.16l-3-d)/2
                  -(2/5-\varphi_-)(d-.16)_+/2.
\]
The score endpoints are the lower score floor and the upper allowed
score cap; if the latter depends on $\rho$, use its values at $u,l$
for $c_-,c_+$, respectively. Convexity in the actual score precedes the
completion of the square, and these outward endpoint intervals cover it.
For the phase resource add \eqref{cut49:re:extra} before maximizing. A barrier
$e\ge K/n$ subtracts $K l$ from the nonnegative resource. The first raw
target must pass the license of Lemma~\ref{cut49:re:weight} first.

\paragraph{Optional signed improvement inside an unsigned cell.}
If $g_0=\max(-g,\sigma_l l-3-d)>g$, exclude the cell. When $g_0\ge0$,
the sign-moment constraint forces
$(1-\beta)^3\le(1+\beta)(u-g_0)^2$.
This inequality is monotone on $(-1,1)$, so exact bisection yields a
lower $b_0\le\beta$; if it fails even at $b$, the cell is empty.
For $x\in[g_0,g]$, set
\begin{align*}
 f_-(x)&=(l+x)^2/[2(1+b)]-R(u-x)/2-3b-3ux/2,\\
 f_+(x)&=R(u+x)/2-(l-x)_+^2/[2(1-b_0)]-3b_0-3lx/2.
\end{align*}
The signwise fourth-moment Cauchy bounds prove
$\min f_-\le K_1\le\max f_+$. The first minimum is at its clipped
vertex $3u(1+b)/2-l-R(1+b)/2$. The second maximum is among its endpoints,
the clipped kink $l$, and the clipped vertex
$l+(R/2-3l/2)(1-b_0)$. Intersect with $[-k,k]$ and replace the linear
term in \eqref{cut49:re:unsignedbound} by the smaller full rectangle bound
\[
 \max_{c=c_-,c_+;,k=k_l,k_u;,\psi=\varphi_-,.4;,
             v=-1/2,\,2M/3-1/2}|ck+\psi v|+s_F.
\]
All sixteen corners are retained. If $b_0>0$, the opposite-sign event
also has probability at least $(1+b_0^2)^2/(2u^2)$, from
$m\ge(1+\beta^2)/(2\rho)$ and $m^2\le pq$. It may improve the complete
HIGH value. These are intersections with necessary conditions; the cover
still examines every $\rho$-cell.

\paragraph{Moving from a small defect to probability coordinates.}
For $0<d\le.7$, start with
\[
 A_l=\frac{1.47\sqrt{2-d}-\sqrt{3.7d}}{3.52},\qquad
 B_l=\frac{2.04\sqrt{2-d}-\sqrt{3.7d}}{3.52}.
\]
They follow from $A+B\ge\sqrt{2-d}$ and
$|\sqrt{\sigma+1}A-\sqrt{\sigma-1}B|\le\sqrt{3.7d}$;
their positivity is checked, rather than assumed at every $d$.
The pointwise inequality
$x(x+b)\ge(29/5)b(x-b)$ for $x>b>0$, applied to the two residual
third moments, gives
$D\ge(29/5)\min((\sigma+1)A,(\sigma-1)B)\theta$.
Start $\Theta=1$ and perform the declared finite number of updates:
\begin{align*}
 \Theta&\leftarrow\min\{\Theta,d/[(29/5)
                 \min((104/25)A_l,(54/25)B_l)]\},\\
 \epsilon&=\sqrt{d/[(616/25)(1-\Theta)^2]},\qquad
 [l,u]=[(4-\sigma_u)/2-\epsilon,(4-\sigma_l)/2+\epsilon],\\
 A_l&\leftarrow\max\{A_l,\sqrt{(1-\Theta)l/(1-l)}
                           -d/[2(1-u)(179/25)A_l^2]\},\\
 B_l&\leftarrow\max\{B_l,\sqrt{(1-\Theta)(1-u)/u}
                           -d/[2l(129/25)B_l^2]\}.
\end{align*}
The last two bounds use
$D\ge2(\sigma+2)B^2M_++2(\sigma+4)A^2M_-$ and
$\Delta_A=M_-/q$, $\Delta_B=M_+/p$.
Every division, interval endpoint and variance domain is gated.
Subdivide this entire probability interval, then use the common
cell update and coefficient atlas. Thus the unsigned first barrier,
the signed initialization and the final probability refinements are
three data choices for the same finite-enclosure theorem.

\section{One confinement lemma with two concentration choices}\label{cut49:conf:re:confinement}
A proved amplitude is used as a support update, never as a premise of
its own proof. The following derivation applies to every acquired row,
with the same formulas for the central and left probability ranges.
\subsection{Statement and hypotheses}
Let $n\ge N=47$, $Y=B+U$, $\Pp(B=1)=p<1/2$ and $\E(U\mid B)=0$.
Put $q=1-p$, $s=\E U^2$, $\lambda=ns$ and $X=\ell(Y-p)$, where
$\E X^2=1$. Thus $\ell^2=(1-\theta)/(pq)$ and $\theta=\ell^2s$.
The conditional label-zero support is at most $a\ge0$ and the label-one
support is at least $b\le1$, where $a,b$ are support points and
$b-a=1-d$. In particular $0\le a\le d$, $1-d\le b\le1$.
Assume the initial strip $\operatorname{supp}Y\subset(-.55,1.55)$.
Write $\rho=\E|X|^3$, $\beta=\E[X|X|]$ and $\epsilon=\E|U|$.

Select an acquired closed probability panel $[l,h]$.
Its fields give upper bounds $\Theta,d_0,\rho_0,\beta_+,T,\Lambda$,
a lower bound $\beta_-$, and separate bounds $H_+,H_-$ for the two signs of
$H$. Assume
\[
 \theta\le\Theta,\quad d\le d_0,\quad \rho\le\rho_0,
 \quad \beta_-\le\beta\le\beta_+,\quad 0\le\lambda\le\Lambda,
 \quad .4<C\le C_0=.4253.
\]
For the positive sign $0\le H\le H_+$; for the negative sign
$-H_-\le H\le0$. Let $x=np+H/\ell$ and
$\nu=\mathcal L(x-\sum_{i<n}Y_i)$. Assume a contact majorant $Q$ for
$g(y)=\nu([y,\infty))$, with equality at each support point of $Y$,
and derivative
\begin{equation}\label{cut49:conf:derivative}
 \sqrt n(-Q'(y))=\ell\phi(H/\sqrt n)
 +\frac\ell n\{3C\beta+(H\phi(H/\sqrt n)+3C\rho)t-3Ct|t|\},
 \quad t=\ell(y-p).
\end{equation}
Both actual deleted convolutions satisfy, separately,
\begin{equation}\label{cut49:conf:concentration}
 Q_j(w)\le\frac{\ell\phi(0)w+2C_0\rho}{\sqrt j},
 \qquad j=n-1,n-2,
\end{equation}
for every interval of width $w\ge0$.

Here is the precise common resource assumption. Put
$\kappa=997/15000<\phi(0)/6$, $J=\int_{[a,b)}(y-a)\nu(dy)$,
and note that the central gap contributes $m\ell J=(1-\theta)J$ to $L$.
Thus $L\ge(1-\theta)J$; retain $e\ge0$ and
\begin{equation}\label{cut49:conf:budget}
 n e\rho+n\phi(0)D/6+n^{3/2}L<\mathcal B(H),\qquad
 \mathcal B(H)=\min\{T,\max_j(A_j+F_{j,\pm}|H|-Q_jH^2)\}.
\end{equation}
The four quadratics are the table's quadratic rows, using
$Q_{\rm budget}$, not $Q_{\rm raw}$. The Gaussian phase contribution
has already been included in $Q_{\rm budget}$.

For each panel let $E,F_\ell,m_0$ be respectively strict lower, upper,
and lower rational bounds for
\[
 \sqrt{(1-\Theta)/(h(1-h))},\qquad
 1/\sqrt{l(1-l)},\qquad \sqrt{(1-\Theta)l(1-l)}.
\]
The bounds use the grid $10^{-6}$, with an extra unit outward. Define
\begin{align*}
 c_+&=2(5.16)E^2(1-h-d_0)^2,&
 c_-&=2(7.16)E^2(l-d_0)^2,\\
 b_*&=\min\{3(4.16)E(1-h-d_0),3(6.16)E(l-d_0)\},\\
 P&=m_0E^2\pos{3.5\operatorname{dist}([l,h],[.41886,.41887])-2.041d_0}^2.
\end{align*}
The simultaneous moment inequality has the following usable form:
\begin{equation}\label{cut49:conf:moment}
 D\ge c_+M_++c_-M_-+b_*\theta+P,\quad
 d=(M_+/p+M_-/q)/\ell,\quad
 \epsilon\le2(M_++M_-)/\ell,\quad M_\pm\ge0.
\end{equation}
To derive it, the endpoint expansion of \eqref{cut49:re:gapexact} gives
$D\ge m(\sqrt{\sigma+1}A-\sqrt{\sigma-1}B)^2
+2(\sigma+2)B^2M_++2(\sigma+4)A^2M_-
+3(\sigma+1)B\theta_++3(\sigma+3)A\theta_-$.
Insert $A\ge E(l-d_0)$, $B\ge E(1-h-d_0)$ and the rational lower
coefficients displayed above. For the square, write
$A=\ell(p-a)$, $B=\ell(1-p-(1-b))$, with
$a+(1-b)=d\le d_0$. Its unshifted coefficient is
$(\sqrt{\sigma+1}+\sqrt{\sigma-1})(p-p_*)$;
$\sqrt{\sigma+1}+\sqrt{\sigma-1}>3.5$ and
$\sqrt{\sigma+1}<2.041$ prove the stated lower square $P$.
For completeness, on the positive class put $Z=X-B\ge0$.
Then $T_+=2B^2M_++3B\E[Z^2;X>0]+\E[Z^3;X>0]$ and
$R_+=BM_++\E[Z^2;X>0]$. Consequently
$(\sigma-1)T_++6BR_+\ge2(\sigma+2)B^2M_+
+3(\sigma+1)B\theta_+$, since $\E[Z^2;X>0]\ge\theta_+$.
The reflected expansion gives the negative-class bound.
The centered absolute deviation on each class is at most twice its
mean distance from the inner endpoint, proving the bound on $\epsilon$.
All quantities refer to the same marginal.

\subsection{Concentration from an antichain}
\begin{lemma}[Strict-gap concentration]
Suppose $Y=U_0+B(1+U_1-U_0)$, where $B$ is Bernoulli$(p)$ independent
of the pair $(U_0,U_1)$, and $1+U_1-U_0\ge c>0$ almost surely.
For independent copies $Y_1,\ldots,Y_j$ and any interval $I$ of diameter
strictly less than $c$,
\[
 \Pp\left(\sum_{i=1}^jY_i\in I\right)
 \le \max_{0\le k\le j}\binom jk p^k(1-p)^{j-k}.
\]
The same bound holds after adding any independent random variable.
\end{lemma}
\begin{proof}
Condition on all noise pairs. The sum is a fixed constant plus
$\sum B_iw_i$, with every $w_i\ge c$. The family $\mathcal F$ of subsets
whose weighted sums lie in $I$ is an antichain: distinct comparable
subsets have sums differing by at least $c$, larger than the diameter.
A uniformly random permutation of $\{1,\ldots,j\}$ defines the maximal
chain of its initial subsets. The chain contains a fixed $k$-subset with
probability $1/\binom jk$, and meets an antichain at most once. Hence
\[
 \sum_{S\in\mathcal F}\binom j{|S|}^{-1}\le1.
\]
Multiplying each summand by its binomial level mass proves the bound
conditionally. Integration removes the conditioning. Conditioning on an
additional independent variable proves the last assertion. This argument
includes all atom collisions and any open or closed interval endpoints.
\end{proof}
For the actual marginal we may couple the two conditional noises
independently and then sample an independent label. Since $U_0\le a$ and
$U_1\ge b-1$, the step is at least $b-a=1-d$.
Strict diameter is necessary: a Bernoulli law on $\{0,c\}$ assigns mass
one to $[0,c]$, exceeding its largest atom when $0<p<1$.

\paragraph{A rational cap for every remaining index.}
Put $J_0=N-2=45$ and define
\[
 r_L=l-(1-l)/J_0,\qquad r_R=h+h/J_0,\qquad
 v_* =\min\{r_L(1-r_L),r_R(1-r_R)\}.
\]
The checker verifies $0<r_L\le r_R<1$. Adjacent binomial-mass ratios
show that any mode $k$ at index $j\ge J_0$ satisfies
$(j+1)p-1\le k\le(j+1)p$, hence $r=k/j\in[r_L,r_R]$.
In particular $1\le k\le j-1$. Finite Stirling bounds give
\[
 \sqrt j\binom jk p^k(1-p)^{j-k}
 \le\frac{\exp(1/(12j)-jD(r\Vert p))}{\sqrt{2\pi r(1-r)}}
 <\frac{.4}{\sqrt{v_*}(1-1/(12J_0))}.
\]
Here $D(r\Vert p)\ge0$ follows from Jensen's inequality for the logarithm;
$\exp x\le(1-x)^{-1}$ for $0\le x<1$ follows term by term from the
power series. The finite Stirling inequalities used are
$\sqrt{2\pi m}(m/e)^m<m!<\sqrt{2\pi m}(m/e)^m\exp(1/(12m))$.
Choose the strict $10^{-6}$ grid upper square root $M$ of
$.16/[v_*(1-1/(12J_0))^2]$. Thus, whenever $w<1-d$,
\begin{equation}\label{cut49:conf:antichainconcentration}
 Q_j(w)<M/\sqrt j,\qquad j\ge n-2\ge J_0.
\end{equation}
This covers the $n-1$ full summands of $\nu$ and the $n-2$ full summands
remaining after one mark or deletion. An independent remaining conditional
noise preserves the cap. We do not replace the latter index by $n-1$.

Either the score-derived cap or the strict-gap cap may be selected on
an interval where its hypotheses hold. The first confinement uses the
score cap; later central and left schedules use the strict-gap cap.
For the left panels replace the conditional cubic coefficient 24 below
by 28; in general any $c_\rho$ satisfying
$c_\rho^2[l(1-l)]^3>[3(1-l)+.55]^2$ is valid.
\subsection{Signed deletion, including atoms}
Let $m=n-1$, $Z=x-\sum_{i=1}^mY_i$, $\mu=x-mp$ and let $K=\sum B_i$.
For a fixed removed label put
$T_b(I)=\Pp(x-\sum_{i<m}Y_i-b-U_b\in I)$.
Exchangeability and $mp-K=Z-\mu+\sum U_i$ give the exact identities
\[
 \nu(I)-T_1(I)=\frac{\E[(Z-\mu+\sum U_i)1_I(Z)]}{mp},\qquad
 T_0(I)-\nu(I)=\frac{\E[(Z-\mu+\sum U_i)1_I(Z)]}{mq}.
\]
There is no absolute value around the sum of noises.
For a full removed pair $(Y_i,U_i)$, condition on the pair. The remaining
sum gives a function $f(Y_i)\in[0,Q_{m-1}(w)]$. Since $\E U_i=0$,
\[
 \E[U_i f(Y_i)]\le Q_{m-1}(w)\E(U_i)_+
 =\tfrac12\epsilon Q_{m-1}(w).
\]
This also works when one conditional law is deterministic.

Let independent conditional noises satisfy $\Delta=U_1-U_0\ge-d$.
For $\tau>0$ put $I^+=I+1+[-d,\tau]$. On $\Delta\le\tau$, the
point for $T_1$ in $I$ maps to the point for $T_0$ in $I^+$. On the
complement, conditioning on both noises gives error at most
$Q_{m-1}(w)\Pp(\Delta>\tau)$. Moreover
\begin{equation}\label{cut49:conf:bad}
 \Pp(\Delta>\tau)\le
 \min\left\{\frac{s}{s+p\tau^2},\frac d{d+\tau}\right\}.
\end{equation}
Indeed $\E\Delta=0$, $\operatorname{Var}\Delta\le s/p$, and
$\Delta+d\ge0$ has mean $d$. Cantelli and Markov prove the two bounds.
The exceptional event is strict, so $\Delta=\tau$ belongs to the good event.

If $I,I^+$ are disjoint, combine the two signed noise terms before
bounding them. For the same remaining-sum law their conditional
probabilities $f_I,f_J$ satisfy
\[
 0\le f_I\le Q_I,\quad0\le f_J\le Q_J,\quad f_I+f_J\le Q_{\rm hull}.
\]
Since $p<q$,
\[
 \frac{f_I}{p}+\frac{f_J}{q}
 \le\min\left\{\frac{Q_I}p+\frac{Q_J}q,
 \frac{Q_{\rm hull}}q+\left(\frac1p-\frac1q\right)Q_I\right\}.
\]
Multiplication by $\E U_+=\epsilon/2$ proves the combined bound.
Closed, open and half-open intervals are all covered by containment in
the corresponding closed interval. No density or atom separation is used.

\subsection{Slope bounds and the two exterior intervals}
Split $|H|$ into 16 closed equal cells $[h_l,h_h]$ and $\lambda$ into
24 closed equal cells $[s_l,s_h]$. On each box put
\[
 \rho_* =\min\{\rho_0,\rho_B(l)+24s_h/N\},\quad
 b_- =\max\{\beta_-,1-2h-2(1-h)\Theta\},\quad
 b_+ =\min\{\beta_+,1-2l+\Theta\}.
\]
Here $\rho_B(p)=(1-2p(1-p))/\sqrt{p(1-p)}$. Write $s_i=\E U_i^2$. Conditional centering and the two sign supports give
\[
 \rho=\ell^3\{pq(1-2pq)+3pq(s_0+s_1)
                  +p\E U_1^3-q\E U_0^3\}.
\]
The initial strip and $d_0<.55$ imply $|U_i|<.55$.
Thus $\rho\le\rho_B+\ell^3(3q+.55)s\le\rho_B+24s$; the last
coefficient inequality is checked as
$24^2[l(1-l)]^3>[3(1-l)+.55]^2$.
Similarly $\beta=(1-\theta)(1-2p)+\ell^2(ps_1-qs_0)$, and
$|\ell^2(ps_1-qs_0)|\le\theta$. These identities prove the stated
outward bounds on $\rho$ and $\beta$.
Set $f=.4F_\ell$, $\varphi_-=.39886(1-h_h^2/(2N))>0$.
The exact upper bound for $H\phi(H/\sqrt n)$ used below is
\[
 v_+=\begin{cases}.4h_h&H\ge0,\\-\varphi_-h_l&H\le0.\end{cases}
\]
For all $t\le0$, completing the square in \eqref{cut49:conf:derivative} gives
\[
 \mathcal M_-=F_\ell\pos{(v_++3C_0\rho_*)^2/(12(.4))-3(.4)b_-},
 \qquad L_-=E\varphi_- -\mathcal M_-/N.
\]
A positive bound $B_*>\sqrt{(1-l)/l}$ gives valid central-gap lower
and positive-half-line upper slope bounds
\begin{align*}
 \mathcal M_+&=F_\ell\pos{.4h_hB_*+3C_0B_*(B_*-1)-1.2b_-},\\
 L_g&=E\varphi_--\max(\mathcal M_-,\mathcal M_+)/N,\\
 U_+&=f+\frac{F_\ell}{N}
 \left\{3C_0b_++\frac{(.4h_h+3C_0\rho_*)^2}{12(.4)}\right\}.
\end{align*}
They follow by completing a square on the positive half-line and using
$t(t-1)\le B_*(B_*-1)$ for $1\le t\le B_*$; on $0\le t\le1$ that
term has the favorable sign. Negative values are covered by $\mathcal M_-$.
All divisions and positivity requirements are explicit gates.

\paragraph{The sharper lower-ray slope.}
Take $\tau_-=5r_0/12$, $\eta=10^{-5}$ and
$I=[a-r,a-d-\tau_--\eta]$, where $r\ge r_0$.
Since $a\le d$, the entire interval has
$t\le t_*=-E(l+\tau_-+\eta)<0$.
Write $B(t)=b_-+\rho_*t+t^2$ and
$B_*^{\min}=B(\min(t_*,-\rho_*/2))$. For the actual $C$ and $t\le t_*$,
\[
 3C(\beta+\rho t+t^2)+H\phi t
 \ge1.2B(t)+v_+t+3(C_0-.4)\min(0,B_*^{\min}).
\]
The inequality uses $t<0$ and separates the possible negative part of
$B(t)$ before varying $C$. Minimize the right side at the clipped vertex
$\min(t_*,-(v_++1.2\rho_*)/2.4)$ and call its value $q_*$. Then
\[
 L=E\varphi_--F_\ell\max(0,-q_*)/N
\]
is a lower slope on the whole interval, for every $n\ge N$.
Positive corrections are discarded. Therefore contact gives
\[
 \sqrt n\nu(I)\ge L(r-d-\tau_--\eta).
\]
This one slope replaces the previous separate bounds on the long and
short parts of the interval.

\paragraph{Lower interval geometry.}
Its translate is $I^+=[a-r+1-d,b-\eta]$, of width $r-\eta$.
The initial strip implies $I^+\subset[a+g,b)$,
$g=1-.55-2d_0>0$, hence $\nu(I^+)\le J/g$.
The intervals are disjoint throughout $r\le.55+d_0$, since
$\inf I^+-\sup I=1-r+\tau_-+\eta>0$.
At $r=r_0$, concentration and directed distance bounds are
\begin{align*}
 Q_I&=f(r_0-\tau_--\eta)+2C_0\rho_*,&
 Q_J&=f(r_0-\eta)+2C_0\rho_*,\\
 Q_{\rm hull}&=f(1+r_0-\eta)+2C_0\rho_*,\\
 D_I&=\pos{v_H-l-\tau_--\eta},&
 D_J&=\pos{v_H+1-l-\eta},
\end{align*}
where $v_H=-h_l/F_\ell$ for $H\ge0$ and $v_H=h_h/E$ for $H\le0$.
These signs follow from $\mu=p+H/\ell$.

\paragraph{Upper interval geometry.}
For a support point $y\ge1+r_0$, take $\tau_+=7r_0/15$,
$I=[a,a+\tau_+]$ and $I^+=[b,a+1+2\tau_+]$.
The strict gates $d_0+2\tau_+<r_0$ and $1-d_0-\tau_+>0$ give
$I^+\subset[b,y)$ and disjointness. Contact, including the atom at $b$,
and the leakage definition give
\[
 \sqrt n\nu(I)\ge L_g(1-d)-\sqrt nJ/\tau_+,\qquad
 \sqrt n\nu(I^+)\le U_+(.55+d).
\]
Use
\begin{align*}
 Q_I&=f\tau_++2C_0\rho_*,&Q_J&=f(2\tau_++d_0)+2C_0\rho_*,\\
 Q_{\rm hull}&=f(1+2\tau_+)+2C_0\rho_*,\\
 D_I&=\pos{d_0+\tau_++h_h/E-l},&
 D_J&=\pos{1+d_0+2\tau_++h_h/E-l}.
\end{align*}
The upper side deliberately keeps the absolute phase bound.

\paragraph{The strict-gap alternative.}
If $1-.55-2d_0>0$ and the displayed lower intervals have width below
$1-d$ throughout $r\in[r_0,.55+d_0]$, take $Q_I=Q_J=M$.
For the upper intervals the corresponding gate is
$1-2d_0-2\tau_+>0$; then again $Q_I=Q_J=M$.
Use the same directed distances. Their concentration caps are now
independent of $r$, and use
$a_e=F_dM\{1/l+1/(1-h)\}/2$ in the next formulas.
The lower-ray gate reduces to $L>0$: its contact side increases with
$r$, while the entire error is constant on the acquired finite strip.
For the score-derived cap retain the explicit affine slope gate below.
No assertion outside that strip is needed.

\subsection{One budget, two error bounds}
Let $F_d$ be the strict grid upper bound for
$\sqrt{\max\{(47/46)^3,47/45\}}$. It simultaneously covers the two
concentration factors in \eqref{cut49:conf:concentration} after multiplication by
$\sqrt n$, uniformly for $n\ge47$. For either pair of intervals set
\begin{align*}
 a_e&=\frac{F_d}{2}\min\left\{\frac{Q_I}l+\frac{Q_J}{1-h},
 \frac{Q_{\rm hull}}{1-h}+\left(\frac1l-\frac1{1-l}\right)Q_I\right\},\\
 A&=F_d\left\{Q_ID_I/l+Q_JD_J/(1-h)\right\},\qquad B=F_dQ_I.
\end{align*}
The probability denominators are outward because $1/p-1/(1-p)$ is
decreasing on $(0,1/2)$.
Use $a_d=L$, $a_J=1/g$ on the lower side, and
$a_d=L_g+U_+$, $a_J=1/\tau_+$ on the upper side.
For the phase cell, maximize each concave budget quadratic at its
vertex clipped to $[h_l,h_h]$, take their maximum, then the minimum
with $T$; call this upper bound $T_*$. Put
\[
 R=T_*/\kappa-b_*E^2s_l,
 \quad\alpha=\min(c_+lE,c_-(1-h)E),\quad G=(R-NP)/\alpha.
\]
If $G\le0$, the strict common budget excludes the entire closed box.
Otherwise $nd<G$ and
\[
 B_\tau=\min\{s_h/(s_h+Nl\tau^2),G/(N\tau+G)\}
\]
bounds \eqref{cut49:conf:bad}. With
\[
 z=\max\left\{\frac{a_d/l+2a_e}{Ec_+},
 \frac{a_d/(1-h)+2a_e}{Ec_-},\frac{\kappa a_J}{1-\Theta}\right\},
\]
one valid total error bound is
\begin{equation}\label{cut49:conf:oldbound}
 \mathcal E_0=(zR+A)/N-zP+B B_\tau.
\end{equation}
This follows by multiplying the single budget by $z/\kappa$.
The inverse-index coefficient $zR+A$ is checked positive.

A stronger bound spends the same resource on the pair event. Let
$R_0=R-NP>0$ and define the three pairs
\[
 (v_+,w_+)=\left(\frac{a_d/l+2a_e}{Ec_+}\frac{R_0}{N},
                    \frac{R_0}{Ec_+l}\right),
\]
\[
 (v_-,w_-)=\left(\frac{a_d/(1-h)+2a_e}{Ec_-}\frac{R_0}{N},
                    \frac{R_0}{Ec_-(1-h)}\right),\qquad
 (v_J,w_J)=\left(\frac{\kappa a_J}{1-\Theta}\frac{R_0}{N},0\right).
\]
They are upper cost/displacement coordinates at the vertices of the
simplex with resources $nc_+M_+$, $nc_-M_-$ and $n^{3/2}L/\kappa$.
For any $t\ge0$,
\[
 \frac{D}{N\tau+D}\le\frac{N\tau}{(N\tau+t)^2}D+
 \frac{t^2}{(N\tau+t)^2}\qquad(D\ge0).
\]
The difference is $N\tau(D-t)^2/[(N\tau+D)(N\tau+t)^2]$.
Consequently a further valid error bound is
\[
 \mathcal E_t=\frac AN+
 \max_{i\in\{+,-,J\}}\left(v_i+B\frac{N\tau}{(N\tau+t)^2}w_i\right)
 +B\frac{t^2}{(N\tau+t)^2}.
\]
The checker takes the minimum of \eqref{cut49:conf:oldbound} and these eleven
bounds with $t=jG/10$, $j=0,\ldots,10$. Each is individually valid;
no numerical minimization or derivative of their minimum is assumed.

The lower contradiction is the strict gate
$Lr_0-L(\tau_-+\eta)>\mathcal E$.
To extend it to every $r\ge r_0$, bound increments directly in the original
affine-in-width signed deletion bound for the same law. The budget gives
$n\epsilon<E_*=2T/(\kappa E\min(c_+,c_-))$. The sufficient slope gate is
\[
 L>\frac{F_df}{N}\left\{D_I/l+D_J/(1-h)
 +\frac{E_*}{2}\left(1/l+1/(1-h)\right)\right\}+F_dfB_\tau.
\]
Using the uncoupled increment is safe even when the hull bound was used
at $r_0$. The upper contradiction is $L_g-.55U_+>\mathcal E$.
Thus no support point lies below $a-r_0$ or above $1+r_0$.
Since $a\ge0$, $b\ge1-d_0$ and $d_0<r_0$, the label-zero and label-one
supports both satisfy $|U|<r_0$.

\paragraph{The selected schedules.}
The ordinary, central-antichain, and left-antichain certificates substitute
their recorded $r_0$ into these same gates. Noise-dependent amplitudes
are the same theorem with a positive lower noise endpoint in
$R=T_*/\kappa-b_*E^2s_l$. In the central rows the additional amplitudes
are $.35,.30,.275$ on $[.05,.08],[.08,.10],[.10,\Lambda]$;
on the left they are $.375,.325,.30$. Intersect each interval with its
own row's $\Lambda$ and each amplitude with that row's global amplitude.
The fixed acceptance checks both signs of $H$, all closed phase/noise
cells, clipped endpoints, and strict budget exclusions. No proof of
LOW, MID, or a binomial score comparison is used by this confinement lemma.

\section{A smaller excess ceiling at index 46}\label{cut49:bar46:entry}
This section adds one finite calculation to the resource argument. Write
$C_E=(3+\sqrt{10})/(6\sqrt{2\pi})$ and $a_j=C_j-C_E$.
The external published inputs and the attained-contact reduction are
exactly those of the accompanying manuscript. All new entry estimates
refer to the same attaining marginal and threshold.

\begin{theorem}[Supported barrier at 46]\label{cut49:bar46:supported}
There is no attained positive maximum with $n=46$,
\[
 a_{46}\ge\frac{.54}{46},\qquad C_{46}\le.42496,
 \qquad \operatorname{supp}X\subset[-3.4,3.4].
\]
The assertion includes either sign of the raw threshold and arbitrary
atoms. It does not assume positive conditional variance.
\end{theorem}

\paragraph{The support edge.}
The existing finite ceiling $a_j<.7/j$ for $j=45,46$ gives
\[
 C_{46}<.40974+.7/46<.42496,\qquad
 46(a_{45}-a_{46})_+<46(.7)/45-.54=79/450.
\]
The independently checked all-threshold contact support implication
with $(N,L,K,R,q)=(46,46,.54,3.4,79/450)$ therefore supplies the
support premise of Theorem~\ref{cut49:bar46:supported}.
This implication is checked on 22 closed radial cells through $4.5$
and the complete remaining ray, for both signs and every threshold.
No entry output is used to prove this initial support edge.
Consequently the supported theorem implies $a_{46}<.54/46$.

\paragraph{Unchanged quantities and formulas.}
Let $\mathcal R$ and $D\ge0$ be the raw remainder and gap defect,
respectively. The conserved identity \eqref{cut49:re:identity}, weight removal
\eqref{cut49:re:unweighted}, and coherent phase price \eqref{cut49:re:extra}
have already been proved. In particular, after a prior strict raw bound
$\mathcal R<r<.4n$ has licensed weight removal,
\[
 na_n\rho+\frac{997}{15000}nD
 \le\mathcal B(H),\qquad
 \phi(H/\sqrt n)\ge\phi_-:=\frac{997}{2500}
                \left(1-\frac{5r}{2n}\right).
\]
Here $\mathcal B$ is the raw quadratic with the already proved extra
curvature
\[
 e_\rho=\frac{1-5r/(2n)}2\left(.4095\rho-\frac d{15}\right)>0,
 \qquad nD<nd.
\]
In this single-index calculation $n=46$ throughout. The extra curvature
uses only the preceding bounds $r,d$. On the barrier branch
$na_n\ge K=.54$, it follows that
\begin{equation}\label{cut49:bar46:resource}
 D<\frac{15000}{997n}\max_H\{\mathcal B(H)-K\rho\}.
\end{equation}
The subtraction is made only in this resource bound, never in the raw
remainder passed to the next stage.

\paragraph{Four signed-moment covers.}
The first four stages use the signed-moment and common-first-moment
formulas of Sections~\ref{cut49:re:foundation}--\ref{cut49:re:atlas}, with fixed parameters
\[
 N=L=46,\ R=3.4,\ K=.54,\ C\le.42496,
 \qquad C\ge.4095+.54/46.
\]
The published moment caps determine the outer $\rho$ interval.
Each pass uses its complete $24\times24\times8\times4$ closed cover
in $(\rho,\beta,\gamma,m)$, including both signs of $\beta,\gamma$.
The LOW estimate and all six HIGH pieces are the complete estimates
of Section~\ref{cut49:re:foundation}; the signed core and its boundary rules are unchanged.
The first raw output is below $.4(46)=18.4$, before any defect output
is used. Every later pass consumes only its accepted predecessor state.

\paragraph{Three probability covers.}
After the fourth pass $D<.2371<.7$. The general gap geometry and the
24 ordered conditional updates in Section~\ref{cut49:re:unsigned} provide a closed
probability interval $[p_-,p_+]$. Split it into 96 equal closed cells.
The simultaneous resource \eqref{cut49:re:joint} proves each cell's moment and inner-gap
bounds. A cell may be discarded only when its binary defect charge
already exceeds the available defect. Otherwise apply the general
four-corner enclosure \eqref{cut49:entry:four-corners}, using the score floor $.4095+.54/46$.
All closed endpoints remain covered.

Where its explicit moment and rectangle gates hold, also use paired
Fourier cancellation \eqref{cut49:fourier:paired}. In particular the checker verifies
\[
 0<l<u<1/2,\quad 1\le M\le2,\quad 0\le g\le1,\quad R>2g,
 \quad 3\min(A_l,B_l)\ge\tfrac32\max_{p\in[l,u]}\rho_B(p),
\]
and the nonnegative paired absolute-skewness lower endpoint and every
coefficient ordering gate. The last displayed inequality licenses the
actual lower bound $\rho\ge\rho_B(u)$. Reflect only the auxiliary
Fourier integral when necessary, as in Theorem~\ref{cut49:re:enclosure}; the actual signed
contact and phase data retain their original orientation. If a paired
premise fails, retain the valid general enclosure for that cell.

\paragraph{The signed phase correction.}
The actual coefficient of $H^2$ is $C\rho+\gamma\phi/2$.
The available signed lower bound is $\gamma\ge3.16\rho-3-d$.
It can be negative during the first probability pass. Define
\[
 \epsilon_\gamma=\frac{.4-\phi_-}{2}(d-.16)_+.
\]
Since $\rho\ge1$, its negative part is at most $(d-.16)_+$, so
\begin{equation}\label{cut49:bar46:phase}
 C\rho+\gamma\phi/2
 \ge C\rho+\phi_-(3.16\rho-3-d)/2-\epsilon_\gamma.
\end{equation}
Thus every paired corner $(A,f,q)$ returned before this correction
is used with $\widetilde q=q-\epsilon_\gamma>0$. Both the raw and the
resource bounds are recomputed:
\[
 R_{\rm pair}=\max_{C,\rho\ {\rm corners}}
       \left(A+\frac{f^2}{4\widetilde q}\right),\qquad
 B_{\rm pair}=\max_{C,\rho\ {\rm corners}}
       \left(A+\frac{f^2}{4(\widetilde q+e_\rho)}-K\rho\right).
\]
The four score--moment corners use $C\in\{.4095,.42496\}$ and
$\rho\in\{\rho_B(u),\rho_u\}$. The lower score corner is a valid
relaxation of the stronger barrier score floor. The checker requires
all corrected divisors to be positive.

For clarity, the endpoint argument precedes square completion. At fixed
$H$, the constant and linear-phase envelopes are convex in $C$ and in
$\rho$, as proved for \eqref{cut49:entry:four-corners}; the curvature is affine in each variable.
The new correction is constant on the cell, $e_\rho$ is affine in
$\rho$, and $-K\rho$ is affine. Hence taking those same endpoints
still gives an upper bound. Only then use
$f|H|-qH^2\le f^2/(4q)$. The maximum for the general estimate and the
maximum for the paired estimate are separate complete bounds; their
minimum is therefore valid. No coefficients from different estimates
are combined.

\paragraph{Fixed exact schedule.}
All decimals in the table denote exact rational target bounds. Each
row gives strict simultaneous upper bounds for $\mathcal R,D$.
\begin{center}\small
\begin{tabular}{c|c|r|r|r|r|r}
Pass&Cover&$\mathcal R$&$D$&Active&Excluded&Paired\\\hline
1&signed&5.8415&1.6324&17752&680&0\\
2&signed&2.8322&.6195&12871&5561&0\\
3&signed&1.8719&.3550&5515&12917&0\\
4&signed&1.4410&.2371&3127&15305&0\\
5&probability&.9953&.1077&92&4&83\\
6&probability&.7416&.0379&95&1&95\\
7&probability&.5840&-.0046&96&0&96
\end{tabular}
\end{center}
There are $4(18432)+3(96)=74016$ closed cells in the recorded covers.
The final row contradicts $D\ge0$, proving
Theorem~\ref{cut49:bar46:supported}.

\section{One local comparison theorem}\label{cut49:shortloc:section}

A law concentrated near two points has a natural integer count and a small
centered error. We compare it with its own Bernoulli count. The same exact
identity handles every noise regime. There are two useful estimates of its
right side: discard all but the nearest outgoing count, or retain all counts
through a random cutoff. The estimates share their normalization, moments,
and tail bounds; only the final scalar witness changes.

\subsection{The identity and the common numerical test}

Let $0<p\le1/2$, $q=1-p$, $v=pq$, and let
\[
 Y=B+U,\qquad \Pr(B=1)=p,\qquad \mathbb E(U\mid B)=0.
\]
The two actual conditional laws $U_0,U_1$ obey
\begin{equation}\label{cut49:shortloc:support}
 |U_j|\le A,\quad U_0\le a_0,\quad U_1\ge-a_1,\quad
 a_0,a_1\ge0,\quad d=a_0+a_1\le d_0<\min(p,q),\quad d_0\le A.
\end{equation}
Write $s_j=\mathbb E U_j^2$, $s=qs_0+ps_1$, $\lambda=ns$,
$\ell=(v+s)^{-1/2}$, $\theta=s/(v+s)$ and
$\rho=\mathbb E|\ell(Y-p)|^3$. At the actual threshold choose
\[
 x=np+H/\ell=k+u,\qquad k\in\{1,\ldots,n-1\},\quad |u|\le\tfrac12,
 \qquad \delta=k-np=H/\ell-u.
\]
Either nearest integer is allowed at a tie. Put
\[
 \pi_j=\binom njp^jq^{n-j},\quad F_B(j)=\Pr(\operatorname{Bin}(n,p)\le j),
 \qquad \rho_B=(1-2v)/\sqrt v,
\]
with $\pi_j=0$ outside $[0,n]$. Let $W_j$ be an independent sum of $j$
copies of $U_1$ and $n-j$ copies of $U_0$. The signed scores are
\[
 C=\frac{\sqrt n}{\rho}\{\Pr(\textstyle\sum Y_i\le x)-\Phi(H/\sqrt n)\},
 \qquad C_B=\frac{\sqrt n}{\rho_B}\{F_B(k)-\Phi(\delta/\sqrt{nv})\}.
\]
Throughout this section $C\ge c_*>0$; the selected application uses $c_*=.4$.
Define
\begin{align}
 f(u)&=\frac{\Phi((\delta+u)/\sqrt{n(v+s)})-
                    \Phi(\delta/\sqrt{n(v+s)})}{\pi_k},\label{cut49:shortloc:f}\\
 Q&=\Pr(\textstyle\sum Y_i\le k+u)-F_B(k),\nonumber\\
 \mathcal R&=\sqrt n\{\Phi(\delta/\sqrt{nv})-
                 \Phi(\delta/\sqrt{n(v+s)})\}+C(\rho_B-\rho).\nonumber
\end{align}
Direct subtraction gives the exact identity
\begin{equation}\label{cut49:shortloc:master}
 \rho_B(C-C_B)=\sqrt n\pi_k\{Q/\pi_k-f(u)\}+\mathcal R.
\end{equation}
In particular, the phase $u$, the sign of $H$, and both conditional laws
refer to the same marginal throughout.

\begin{theorem}[Parameterized local comparison]\label{cut49:shortloc:theorem}
Fix $n\ge N$, $\lambda>0$ and a closed parameter cell on which
\eqref{cut49:shortloc:support} holds. Suppose $\sqrt n\pi_k\ge m>0$.
Either of the following witnesses proves $\rho_B(C-C_B)<-\varepsilon\lambda$.
\begin{enumerate}
\item \emph{Nearest-count witness.} Valid bounds $O,T,R_0,\kappa$ satisfy
\[
 \Pr(W_k>u)+f(u)\ge O,\quad
 \sum_{i\ge1}\Pr(W_{k+i}\le u-i)\le T,\quad
 \pi_{k+i}/\pi_k\le R_0,\quad \mathcal R\le\kappa\lambda/n.
\]
Only legal counts occur in the sum; adding positive upper bounds for
other counts is allowed. If $D=O-R_0T>0$ and
\begin{equation}\label{cut49:shortloc:directtest}
             mD-\kappa\lambda/n>\varepsilon\lambda,
\end{equation}
the conclusion follows.
\item \emph{Whole-count witness.} Bounds $B(\lambda)<0$ and $I(\lambda)\ge0$
satisfy
\[
 \rho_B(C-C_B)\le\sqrt n\pi_k\lambda B(\lambda)
                                      +\lambda I(\lambda)/n.
\]
It is enough that
\begin{equation}\label{cut49:shortloc:whole-test}
                  mB(\lambda)+I(\lambda)/N<-\varepsilon.
\end{equation}
\end{enumerate}
The witnesses may differ from cell to cell. A finite closed cover and a
separate treatment of $\lambda=0$ therefore suffice for a complete local
comparison.
\end{theorem}
\begin{proof}
Conditioning on the binomial count gives, with all atom conventions intact,
\begin{equation}\label{cut49:shortloc:count}
 Q=-\sum_{i=0}^{k}\pi_{k-i}\Pr(W_{k-i}>u+i)
       +\sum_{i=1}^{n-k}\pi_{k+i}\Pr(W_{k+i}\le u-i).
\end{equation}
Dropping the outgoing terms with $i>0$ gives $Q/\pi_k-f(u)\le-D$.
Use \eqref{cut49:shortloc:master} and $D>0$. For the second witness, use
$B<0$ before replacing $\sqrt n\pi_k$ by $m$, and $I\ge0$ before replacing
$n$ by $N$. These sign requirements are part of the test, not conventions.
\end{proof}

\subsection{Normalization and the repair, once for every witness}

Here is an explicit finite recipe for the constants required above.
Let $p\in[p_-,p_+]\subset(0,1/2)$, $n\ge N$, $\lambda\le\Lambda$,
$\theta\le\Theta$, and $|\delta|\le\Delta$, where $\Delta\ge1$ also includes
every binomial mode. Choose $0<r_-<p_- -\Delta/N$ and
$p_++\Delta/N<r_+<1/2$. Let $v_*>0$ be a lower bound for $p(1-p)$.
Choose positive $E,L,R$ bounding, on $[r_-,r_+]$,
\[
 |D'''(r\Vert p)|=\frac{1-2r}{r^2(1-r)^2},\qquad
 |(\log(r(1-r)))'|=\frac{1-2r}{r(1-r)},\qquad
 \frac1{12r(1-r)},
\]
respectively, where $D(r\Vert p)=r\log(r/p)+(1-r)\log((1-r)/q)$.
Then, with $Z(t)=Lt/2+R+Et^3/(6N)$,
\begin{equation}\label{cut49:shortloc:mass}
 \sqrt n\pi_k=\frac{\phi(0)}{\sqrt v}
       e^{-\delta^2/(2nv)+\zeta},\qquad n|\zeta|<Z(|\delta|).
\end{equation}
For completeness, the finite Stirling formula used here is
\[
 \log(j!)=(j+\tfrac12)\log j-j+\tfrac12\log(2\pi)+r_j,
 \qquad 0<r_j<1/(12j).
\]
If $a_j=\log(j!)-(j+1/2)\log j+j$ and $z=(2j+1)^{-1}$, the integrated
geometric series gives
$0<a_j-a_{j+1}=\sum_{i\ge1}z^{2i}/(2i+1)<1/[12j(j+1)]$.
Summing gives the remainder relative to $\lim a_j$. The integration-by-parts
recursion for $\int_0^{\pi/2}\sin^j t\,dt$ and the inequalities between three
successive integrals give Wallis's product; its finite product equals
$e^{4a_j-2a_{2j}}j/[2(2j+1)]$ and tends to $\pi/2$. Thus
$\lim a_j=\tfrac12\log(2\pi)$. Applying this formula to $n!,k!,(n-k)!$
and Taylor-expanding $D$ and the logarithmic prefactor proves
\eqref{cut49:shortloc:mass}; the factorial error is at most $1/[12nr(1-r)]$.

Consequently a valid mass lower bound is any positive
\begin{equation}\label{cut49:shortloc:m}
 m<\frac{.39886}{\sqrt{p_+(1-p_+)}}
          \left[1-\frac{\Delta^2/(2v_*)+Z(\Delta)}N\right].
\end{equation}
The same formula at a mode proves $\sqrt n\max_j\pi_j<1$ whenever
$v_*[1-Z(\Delta)/N]^2>.4^2$ and $Z(\Delta)<N$.
The rational bounds $.39886<\phi(0)<.4$ follow, for example, by the
alternating arctangent series in Machin's identity.

Differentiating \eqref{cut49:shortloc:f} gives
\begin{equation}\label{cut49:shortloc:logder}
 \log f'=-\tfrac12\log(1+s/v)+\delta^2/(2nv)
                   -(\delta+u)^2/(2n(v+s))-\zeta.
\end{equation}
Thus $n|\log f'|\le B_0\lambda+P(D,a)$ on $|\delta|\le D\le\Delta$,
$|u|\le a\le1/2$, where
\[
 B_0>\frac1{2v_*}+\frac{\Delta^2}{2Nv_*^2},\qquad
 P(D,a)=\frac{2Da+a^2}{2v_*}+Z(D).
\]
Choose $0<b<1$ above $[B_0\Lambda+P(\Delta,1/2)]/N$ and put
$F=(1-b)^{-1}$. The elementary estimates $e^z\le(1-z)^{-1}$ and
$1-e^{-z}\le z$ imply
\begin{equation}\label{cut49:shortloc:derivative}
 |f'-1|\le(B_1\lambda+P_1)/n,\qquad B_1=FB_0,\quad P_1=FP(D,a).
\end{equation}
This supplies both a derivative interval $l_f\le f'\le h_f$ and the
finer phase-error bound $|u-f(u)|\le a(B_1\lambda+P_1)/n$.

The exact count variance is
\begin{equation}\label{cut49:shortloc:variance}
 t_j=\mathbb E W_j^2=\lambda+(j-np)(s_1-s_0),\qquad
 |s_1-s_0|\le s/p\le c_v s,\quad c_v\ge1/p_-.
\end{equation}
It gives $\omega_-\lambda\le t_k\le\omega_+\lambda$ with
$\omega_\pm=1\pm c_vD/N$. All uses below require the checked condition
$\omega_->0$. Also
$t_{k+i}\le iV_*\lambda$ for $V_*=1+c_v(\Delta+1)/N$.
The binomial recursion gives, for every legal $i\ge1$,
\begin{equation}\label{cut49:shortloc:ratio}
 \frac{\pi_{k+i}}{\pi_k}\le R_0,
 \qquad R_0>1+\frac{p_++\Delta-1}
 {Nv_*[1-(\Delta-1)/(Np_-)]},
\end{equation}
provided the denominator is positive and $p_++\Delta<2$.
Indeed the first ratio has this bound and all subsequent rising factors
are at most one. No count outside $[0,n]$ is used in this recursion.

Finally \eqref{cut49:shortloc:support} fixes the signs inside the absolute cubic.
Expansion and centering give
\[
 \mathbb E|Y-p|^3=v(1-2v)+3pq(s_0+s_1)
                   +p\mathbb E U_1^3-q\mathbb E U_0^3
 \ge v(1-2v)+(3p-d)s.
\]
The inequality $(1+z)^{3/2}\le1+3z/2+3z^2/8$ for $z\ge0$ yields
\begin{equation}\label{cut49:shortloc:cubic}
 (v+s)^{3/2}(\rho-\rho_B)\ge s
     \left(6p-3p^2-\tfrac32-d-\frac{3(1-2v)s}{8v}\right).
\end{equation}
Let $Q_*>0$ be a lower bound for the bracket and
$m_c\le c_*/[p_+(1-p_+)+\Lambda/N]^{3/2}$.
The normal variance correction is nonpositive for $\delta\le0$;
for $\delta>0$ its derivative in $s$ is at most $.4\delta/(2v^{3/2})$.
Therefore
\begin{equation}\label{cut49:shortloc:repair}
 \mathcal R\le(P_*D_+-m_cQ_*)\lambda/n,
 \qquad P_*\ge\frac{.4}{2v_*^{3/2}},\quad D_+\ge\delta_+.
\end{equation}
One may take $D_+=D$, but retaining the positive part is often decisive.
All bounds concern $v+s$, which can exceed $1/4$.

\subsection{Two elementary tools for probabilities}

\begin{lemma}[A shifted quartic]\label{cut49:shortloc:quartic}
Suppose $W$ is a sum of independent centered variables bounded in absolute
value by $A$, and $t=\mathbb E W^2$. Let $K_0\ge(9+6\sqrt3)/4$ and $a>0$.
For $r\ge0$ and $\sigma\in\{-1,1\}$,
\begin{equation}\label{cut49:shortloc:quarticformula}
 \Pr(W+\sigma r>0)\ge
 \frac{(3a^2-A^2)t-3t^2-4Art-6r^2t-r^4
                      +3a^2r^2+2\sigma a^3r}{K_0a^4}.
\end{equation}
\end{lemma}
\begin{proof}
The polynomial $h_a(z)=(-z^4+3a^2z^2+2a^3z)/(K_0a^4)$ is below
$\mathbf1_{\{z>0\}}$: it is nonpositive for $z\le0$, and its maximum for
$z>0$ is at $z/a=(1+\sqrt3)/2$ and is at most one. Independence gives
$|\mathbb E W^3|\le At$ and $\mathbb E W^4\le A^2t+3t^2$.
Expand $\mathbb Eh_a(W+\sigma r)$ and apply these two inequalities.
The value at the threshold is zero, so atoms require no limiting argument.
\end{proof}

\begin{lemma}[One common tail bound]\label{cut49:shortloc:tails}
If independent centered $X_i\le A$ have total variance at most $V$, then,
for $b>0$, their sum has upper tail bounded by the minimum of one and
\begin{equation}\label{cut49:shortloc:tailmin}
 \begin{gathered}
 \frac{V}{V+b^2},\qquad
 \mathcal B(A,V,b)=e^{b/A}\left(\frac{V}{V+Ab}\right)^{b/A+V/A^2},\qquad
 \frac{\mathbb E(P_\mu-j)_+^3}{(b/A+\mu-j)^3},\\
 \mu=V/A^2,\qquad 0\le j<b/A+\mu.
 \end{gathered}
\end{equation}
where $P_\mu$ is Poisson. The same holds for a lower tail when $|X_i|\le A$.
All bounds apply to closed events, hence also to strict events.
\end{lemma}
\begin{proof}
Cantelli's bound follows by applying Markov to $(\sum X_i+V/b)^2$.
Taylor's integral remainder gives
$\mathbb E e^{zX_i}\le\exp\{\mathbb E X_i^2(e^{zA}-1-zA)/A^2\}$
for $z\ge0$. Multiply and use $z=A^{-1}\log(1+Ab/V)$ to get $\mathcal B$.
The zero-variance case is immediate.

For the third bound, smooth $f(x)=(x-t)_+^3$ by a nonnegative smooth
compactly supported density. Then $f'',f''',f''''\ge0$. For
$\mathcal Af(x)=f(x+A)-f(x)-Af'(x)$, the Taylor quotient
\[
 \frac{f(x+y)-f(x)-yf'(x)}{y^2}
       =\int_0^1(1-r)f''(x+ry)\,dr
\]
is nondecreasing in $y$. A centered $X\le A$ of variance $v$ consequently
satisfies $\mathbb E f(x+X)\le f(x)+(v/A^2)\mathcal Af(x)$.
For $\mathcal P_zf(x)=\mathbb E f(x+A(P_z-z))$,
$\partial_z\mathcal P_zf=\mathcal P_z\mathcal Af$ and
\[
 \partial_z^2\mathcal P_zf=\mathcal P_z\mathcal A^2f\ge0,\qquad
 \mathcal A^2f(x)=\int_0^A\!\int_0^A(A-r)(A-w)f''''(x+r+w)\,dr\,dw.
\]
Thus $\mathcal P_{v/A^2}f\ge f+(v/A^2)\mathcal Af$. Replace the independent
variables successively; translations and expectations preserve the derivative
signs and Poisson rates add. Increasing the total variance to $V$ is permitted
by Jensen after adding independent centered Poisson noise. Removing the
smoothing is justified by bounded actual upper supports and finite Poisson
third moments. Markov at $t=A(j-\mu)$ proves the displayed bound.
\end{proof}

These bounds have finite rational evaluations. For example,
\begin{align}\label{cut49:shortloc:poissonfinite}
 \mathbb E(P_\mu-j)_+^3={}&\mu^3+3(1-j)\mu^2+(1-3j+3j^2)\mu-j^3\\
 &+e^{-\mu}\sum_{h=0}^{\lfloor j\rfloor}\frac{\mu^h(j-h)^3}{h!}.\nonumber
\end{align}
Replacing $e^{-\mu}$ by $(\sum_{h=0}^{24}\mu^h/h!)^{-1}$ enlarges the
nonnegative correction. Alternatively sum the positive Poisson series
through 32; the remaining tail is at most its first term divided by
$1-\mu(34-j)^3/[34(33-j)^3]$, provided $j<33$ and that ratio is below one.
For Bennett, put $x=Ab/V$ and
$L=[(b/A+\mu)\log(1+x)-b/A]$. The positive series
$\log(1+x)=2\sum_{h\ge0}z^{2h+1}/(2h+1)$, $z=x/(2+x)$, gives lower
bounds for $L$; any positive lower exponential series gives an upper bound
for $e^{-L}$. Rounding a positive lower $L$ down before exponentiation is
also valid. These are choices of rational witnesses, not optimization claims.

\subsection{The nearest-count witness, including vanishing noise}

For $u=-r\le0$, add $-h_fr$ to \eqref{cut49:shortloc:quarticformula} with
$\sigma=1$; for $u=r\ge0$, add $l_fr$ with $\sigma=-1$.
Both bound $\Pr(W_k>u)+f(u)$ from below. For a cell
$t\in[t_-,t_+]$, $r\in[r_-,r_+]$, this expression is concave in $t$
(second derivative $-6/(K_0a^4)$), so only $t_-,t_+$ are needed.
At each endpoint, bound a monomial with positive coefficient at $r_-$
and one with negative coefficient at $r_+$. Any rational $a>0$ is a valid
witness. Other valid choices are
\begin{equation}\label{cut49:shortloc:ownother}
 \frac{r_-^2}{t_++r_-^2}-h_fr_+\quad(u\le0),\qquad l_fr_-\quad(u\ge0).
\end{equation}
They follow from Cantelli and nonnegativity, respectively. Their maximum
or one selected witness is valid without an optimality proof.

For incoming counts use $t_{k+i}\le iV_*\lambda$. On $u\in[-r_+,-r_-]$
the first distance is $1+r_-$ and every later distance is at least $i$.
On $u\in[r_-,r_+]$ they are $1-r_+$ and $i(1-r_+/2)$ for $i\ge2$.
Let $b,b_o$ be these first and later distances. On $\lambda\in[L,U]$,
let $T_1$ be any bound from Lemma~\ref{cut49:shortloc:tails} at $(A,V_*U,b)$.
Since $\mathcal B(A,iV,ib)=\mathcal B(A,V,b)^i$, all later counts total at most
$B^2/(1-B)$ for $B=\mathcal B(A,V_*U,b_o)<1$. Thus
\begin{equation}\label{cut49:shortloc:midtest}
 D=O-R_0\{T_1+B^2/(1-B)\}>0,\qquad
 mD-\kappa U/N>\varepsilon U
\end{equation}
proves the entire cell when $\kappa\ge0$. The selected MID normalizers
explicitly verify this sign. This is the common central and left MID proof.

A cell touching $\lambda=0$ requires estimates proportional to $\lambda$;
a positive lower endpoint may not be substituted. The following rescaling
provides them. Put $r=zt$, $t_*=\omega_+U$ and $d_a=K_0a^4$.
The shifted quartic plus $br$, divided by $t$, becomes
\begin{equation}\label{cut49:shortloc:scaledown}
 F_{\sigma,b}(z,t)=\frac{3a^2-A^2-3t}{d_a}
 +\left(\frac{2\sigma}{K_0a}+b-\frac{4At}{d_a}\right)z
 +\frac{3a^2t-6t^2}{d_a}z^2-\frac{t^3z^4}{d_a}.
\end{equation}
It extends to $t=0$ and is concave in $t$. On $0\le z\le z_0$, replacing
$z^4$ by $z_0^2z^2$ gives a lower quadratic at $t=t_*$; at $t=0$ it is
linear. Endpoints and a clipped quadratic vertex therefore give an exact
lower bound $J(\sigma,b,z_0)$ for the whole rectangle.

For negative phase choose $z_-,c_->0$ and verify
\begin{gather}\label{cut49:shortloc:tinynegative}
 J(1,-h_f,z_-)>c_-,\quad
 \frac{z_-^2}{1+z_-^2t_*}-h_fz_->c_-,\quad
 \frac{2z_-}{(1+z_-^2t_*)^2}>h_f,\\
 \frac1{1+4t_*}-\frac{h_f}{2}>c_-t_*,\qquad z_-t_*<1/2.\nonumber
\end{gather}
The quartic handles $r\le z_-t$. Beyond that point use
$r^2/(t+r^2)-h_fr$. Its derivative at $r=z_-t$ is positive by the third
gate; its second derivative has just one sign change, so any later critical
point is a maximum. Its minimum is at $z_-t$ or $1/2$, covered by the other
gates. Hence the own count exceeds $c_-t$.
For positive phase choose $z_+,c_+>0$ with
\begin{equation}\label{cut49:shortloc:tinypositive}
 J(-1,l_f,z_+)>c_+,\qquad l_fz_+>c_+,\qquad z_+t_*<1/2.
\end{equation}
The quartic handles $r\le z_+t$ and the normal term handles the rest.
The own count exceeds $c_+t$ and also exceeds $l_fr$.

For incoming tails proportional to $\lambda$, fix a Poisson anchor $j$.
A positive series term of index $h$ divided by
$\lambda(b/A+\mu-j)^3$ has logarithmic derivative, multiplied by $\lambda$,
\[
 h-1-\mu-\frac{3\mu}{b/A+\mu-j},\qquad \mu=V_*\lambda/A^2.
\]
If $b/A-j>0$ and this is nonnegative at $\lambda=U$ for the first positive
index, it is nonnegative for every term and all $0<\lambda\le U$.
Therefore the Poisson bound divided by $\lambda$ is largest at $U$.
Likewise $B^2/[\lambda(1-B)]$ is increasing whenever
$\lambda B'/B\ge1/2$. Differentiating the Bennett formula gives
\[
 \lambda B'/B=b/A-\mu\log(1+b/(A\mu));
\]
The function $g(\mu)=b/A-\mu\log(1+b/(A\mu))$ decreases because
$\log(1+x)>x/(1+x)$ for $x>0$. Moreover the logarithmic derivative of
$B^2/[\lambda(1-B)]$, multiplied by $\lambda$, is
$2g-1+gB/(1-B)\ge0$ when $g\ge1/2$. Thus checking the gate at $U$
proves the whole interval. A rational upper bound for the logarithm
verifies that gate. On cells using the fixed positive bound $l_fr_-$ instead,
unscaled endpoint tails suffice: after a positive bracket is checked, divide
it by $U$ to obtain a lower bound per $\lambda$. These arguments, followed by
\eqref{cut49:shortloc:directtest}, prove TINY uniformly down to arbitrary positive
noise, including a deterministic conditional component.

\subsection{The whole-count witness}

Assume $c=1-d_0>1/2$. Couple all $2n$ conditional variables independently and
set $W_j=\sum_{i\le j}U_{1,i}+\sum_{i>j}U_{0,i}$, extending it constantly
outside $[0,n]$. Every increment of $j+W_j$ is at least $c$. Define
\[
 K=\max\{j\in\mathbb Z:j+W_j\le k+u\},\qquad R=K-k.
\]
An independent binomial selector has the correct actual count laws, so
\begin{equation}\label{cut49:shortloc:cutoff}
                  \Pr(\textstyle\sum Y_i\le k+u)=\mathbb EF_B(k+R).
\end{equation}
Let $\chi=(n+1)p-k<2$ and $m_B=\max_j\pi_j$. For every integer $r$,
\begin{equation}\label{cut49:shortloc:curvature}
 F_B(k+r)-F_B(k)-r\pi_k\le\frac{m_B}{(n+1)v}
 \left[(\chi-1)_+r_++\chi_+\binom{r_-}{2}+\binom{r_-}{3}\right].
\end{equation}
Indeed, the positive part of $\pi_{k+i}-\pi_{k+i-1}$ is at most
$m_B(\chi-i)_+/((n+1)v)$, by
$\pi_j-\pi_{j-1}=\pi_j((n+1)p-j)/(p(n-j+1))$.
Double summation gives \eqref{cut49:shortloc:curvature}: for $r>0$ only $i=1$
can contribute; for $r=-h<0$ use $(\chi+i)_+\le\chi_++i$ and sum the
multiplicities $h-1-i$. Zero extension makes the same calculation valid
outside the legal count range.

Write $t=\mathbb EW_k^2$. Crossing monotonicity and Cauchy--Schwarz give
\begin{equation}\label{cut49:shortloc:crossingmoments}
 \mathbb ER_+\le\frac{t}{(c-u)^2},\quad
 \mathbb E\binom{R_-}{2}\le\frac{t}{[c+\min(u,c/3)]^2},\quad
 \mathbb E\binom{R_-}{3}\le\frac{t\sqrt{A^2+3t}}{(2c-1/2)^3}.
\end{equation}
For the first, $R=i\ge1$ implies $-W_k\ge ci-u$ and
$(ci-u)^2\ge i(c-u)^2$. For $R=-h\le-2$, the next count is outside the
closed event, so $W_k>c(h-1)+u$. Put $j=h-1$.
The function $j(j+1)/(2(cj+u)^2)$ decreases for $j\ge1$ if $u\le c/3$;
for larger $u$ reduce to that endpoint. Also
$(j^3-j)/(6(cj-1/2)^3)$ decreases for $j\ge2$, since its derivative has
the sign of $2cj-(3j^2-1)/2<0$. These give the last two bounds using
$\mathbb E(W_k)_+^3\le\sqrt{\mathbb EW_k^2\mathbb EW_k^4}$.

\paragraph{The marked crossing.}
Let $D_i=U_{1,i}-U_{0,i}$, $\vartheta=\mathbb E(D_i)_+=\mathbb E(D_i)_-$,
$0<\eta\le1$, and $R_f=\mathbb E\{u+\eta-W_k\}$. The two deleted laws are
\[
 A_0\overset{d}=kU_1+(n-k-1)U_0,\qquad
 A_1\overset{d}=(k-1)U_1+(n-k)U_0.
\]
Then
\begin{align}
 \mathbb ER&\le u+\eta-R_f+J,\label{cut49:shortloc:mark}\\
 J&=\frac{\vartheta}{\eta}\sum_{i\ge1}
 \{\Pr(A_0\le u+d_0-ci)+\Pr(A_1>u+\eta+i-A)\}.\nonumber
\end{align}
Here are the pathwise details. Put $z=u-W_k$ and mark
\[
 Z_f=\sum_{i=1}^{n-k}(D_{k+i})_-\mathbf1_{\{W_k\le u-ci\}},\quad
 Z_b=\sum_{i=1}^{k}(D_{k-i+1})_+\mathbf1_{\{W_k>u+\eta+i\}}.
\]
If $R>z+\eta$, the crossing activates the marks through $R$: for $R>0$
use $z\ge cR$, while for $R<0$ use $z<R-\eta$. In either case
$R\le z+Z_f+Z_b$. Hence if $h=R-\lfloor z+\eta\rfloor>0$,
$Z_f+Z_b>h-1+\eta\ge\eta h$. Thus
$R\le\lfloor z+\eta\rfloor+(Z_f+Z_b)/\eta$ on every path.
On an active forward mark $U_0>U_1\ge-d_0$; deleting its pair implies
$A_0\le u+d_0-ci$. On a backward mark delete $U_1\le A$ to obtain the
strict event for $A_1$. The remaining sum is independent of the entire
pair, so the expectations factor. Extending to all $i\ge1$ adds only
nonnegative bounds. Finally $|R|\le1+(|W_k|+|u|)/c$ proves integrability
and permits expectations. No coupling has replaced either deleted marginal.

\paragraph{A finite bound for every mark.}
Fix $\lambda\in[L,U]$, $L>0$, and $t\le\omega_+\lambda$; put
$T_L=\omega_+L$, $T_U=\omega_+U$. On a phase subcell $u\in[v_l,v_h]$,
for $i=1,2$ let $b_1=ci-d_0-v_h>0$ and $b_2=i-A+v_l>0$.
Lemma~\ref{cut49:shortloc:tails} gives a bound $\omega_+\lambda K_i^{\rm ind}$,
where
\[
 K_i^{\rm ind}=\sum_{j=1}^2\min\left\{\frac1{b_j^2+T_L},
                          \frac{\overline P(A,T_U,b_j)}{T_L}\right\}.
\]
The denominator $T_L$ is a lower bound for the chosen upper envelope
$\omega_+\lambda$, not for any actual variance.
Alternatively couple $A_0=Z+U_1$, $A_1=Z+U_0$ with
$Z\overset d=(k-1)U_1+(n-k-1)U_0$. The two events project into
$Z\le-\alpha$ and $Z>\beta$, where
$\alpha=ci-2d_0-v_h>0$, $\beta=i-A-d_0+v_l>0$.
For any $z\in(-\alpha,\beta)$ their sum is at most
$(\omega_+\lambda+z^2)/\min(z+\alpha,\beta-z)^2$.
Consequently a second kernel is
\[
 K_i^{\rm com}=\frac{1+z^2/T_L}{\min(z+\alpha,\beta-z)^2},\qquad
 z=\max\{-T_L/\beta,\min((\beta-\alpha)/2,T_L/\alpha)\}.
\]
This center is inside the interval; differentiating the two quadratic
ratios on either side of the midpoint proves its optimality, although
any interior rational center suffices. Take the smaller kernel on each
phase subcell and the maximum over their closed cover. Denote it $K_i$.
For $i\ge3$ use fourth moments and the convex midpoint integral for
$(cx-b)^{-4}$. On $u\in[u_l,u_h]$ this gives
\begin{equation}\label{cut49:shortloc:K}
 K=K_1+K_2+\frac{A^2+3T_U}{3c}
  \left\{\frac1{((5/2)c-A-u_h)^3}
                 +\frac1{((5/2)c-A+u_l)^3}\right\},\qquad
 J\le\frac{\vartheta}{\eta}\omega_+\lambda K.
\end{equation}
All denominators must be positive. Indeed, the tail fourth moment is at
most $A^2t+3t^2$, and integration from $5/2$ to infinity covers every
remaining integer, including the harmless added illegal counts.

If the same-law resource supplies $nd<G$ and
$G^2<c_vNL$, then
\begin{equation}\label{cut49:shortloc:theta}
 \vartheta\le\frac{c_vG\lambda}{c_vn\lambda+G^2}.
\end{equation}
To prove this, for centered $X\ge-b$, $\mathbb EX^2\le V$ and $b^2\le V$,
the quadratic $b(X-V/b)^2/(b+V/b)^2$ majorizes $X_-$ on $[-b,\infty)$.
Its expectation is at most $bV/(V+b^2)$. Apply this to $D_i\ge-d$,
using $\operatorname{Var}(D_i)=s_0+s_1\le c_v\lambda/n$, and monotonicity
in $b\le\sqrt V$. The case $d=0$ is immediate. With $\eta=z\lambda$,
$z>0$, $zU\le1$, \eqref{cut49:shortloc:K} becomes
\begin{equation}\label{cut49:shortloc:Jfinal}
 \frac J\lambda\le\frac{\omega_+Kc_vG}{z(c_vN\lambda+G^2)}.
\end{equation}

\paragraph{The fractional-part gain.}
If each centered summand has support span at most $b<1/2$ and the sum has
variance $t$, then, for any real $w$ and
$0<r_*\le2\sin^2(\pi b)/b^2$,
\begin{equation}\label{cut49:shortloc:harmonic}
             \mathbb E\{w-W\}\ge\frac5{23}(1-e^{-r_*t}).
\end{equation}
Indeed $1-\cos(2\pi x)\le(23/5)x$ for $0\le x\le1$.
A finite proof of this sole scalar inequality is recorded below.
For an independent copy $Z'$ of one summand, concavity of sine on
$[0,\pi b]$ gives
\[
 |\mathbb Ee^{2\pi iZ}|^2
 =1-2\mathbb E\sin^2\pi(Z-Z')
 \le e^{-4\sin^2(\pi b)\operatorname{Var}(Z)/b^2}.
\]
Multiply square roots over the independent summands and use periodicity.
This argument includes zero variances and integer collisions.
Here the conditional spans are at most $A+d_0$, so this is the $b$ to use.

For a reproducible scalar proof, $x\ge10/23$ follows from $1-\cos\le2$.
On $[0,10/23]$ use the degree-ten alternating upper polynomial for
$1-\cos(2\pi x)$, with directed coefficients:
\[
 Q(x)=\sum_{j=1}^5\frac{(-1)^{j+1}(2\pi_jx)^{2j}}{(2j)!},\qquad
 \pi_j=\begin{cases}3.141592653589794,&j\text{ odd},\\
                       3.141592653589793,&j\text{ even}.\end{cases}
\]
Put
$P(x)=23/5-Q(x)/x$. On the eight intervals with left endpoint $a_j=5j/92$
and width $w=5/92$, its degree-nine Bernstein coefficients are
\[
 b_{jk}=\sum_{i=0}^k\frac{\binom ki}{\binom9i}
           \sum_{h=i}^9[x^h]P(x)\binom hi a_j^{h-i}w^i,
 \qquad 0\le j<8,\quad0\le k\le9.
\]
The 80 rational coefficients are positive; the retained scalar acceptor
reconstructs them. Expansion in powers of the unit-interval coordinate
proves the formula; Bernstein weights are nonnegative and sum to one.
The alternating remainder and the stated rational bounds on $\pi$ follow
from the same elementary series above. This completely specifies the
finite scalar obligation. A rational rate follows by replacing
$\sin(\pi b)$ with $x-x^3/6+x^5/120-x^7/5040$, $x=3.14159b$, on the
checked alternating-series domain.

Set $r=r_*\omega_-$, require $rU<3$, and define
\begin{equation}\label{cut49:shortloc:gain}
 g=\frac5{23}\left(\frac{r^3L^2}{6}-\frac{r^4L^3}{24}
                     +\frac{r^5L^4}{120}-\frac{r^6L^5}{720}\right).
\end{equation}
The sixth-order Taylor lower bound for $1-e^{-r\lambda}$ gives
$R_f/\lambda\ge5r/23-5r^2\lambda/46+g$. The discarded higher-order
part is increasing as a function of $\lambda$ for $r\lambda<3$:
its derivative has the sign of $48-18x+(24/5)x^2-x^3>0$ on $[0,3]$.
For instance this polynomial has positive degree-three Bernstein
coefficients $48,30,132/5,51/5$ there. Thus freezing it at $L$ is valid.

\paragraph{One final expression.}
On a signed phase cell find an interval $[x_l,x_h]$ containing $H/\ell$
by its four endpoint quotients, using positive bounds for $\ell$.
Let
\[
 D=\max(|x_l-u_h|,|x_h-u_l|),\quad
 D_+=\max(0,x_h-u_l),\quad a=\max(|u_l|,|u_h|),\quad b_\chi=p_+-x_l.
\]
The same formula allows coarser upper bounds $D,D_+,b_\chi$, as used by
the central certificate. Check $D\le\Delta$ and $b_\chi+u_h<2$.
With $\omega_\pm=1\pm c_vD/N$, the curvature coefficient is
\begin{equation}\label{cut49:shortloc:V}
 V_0=\frac{\omega_+}{v_*}\left[
 \sup_{u_l\le u\le u_h}\left\{
 \frac{(b_\chi+u-1)_+}{(c-u)^2}
 +\frac{(b_\chi+u)_+}{[c+\min(u,c/3)]^2}\right\}
 +\frac{\sqrt{A^2+3T_U}}{(2c-1/2)^3}\right].
\end{equation}
On a rational phase subcell, both numerators are bounded at the right
endpoint; the first denominator is smallest there, and the second is
smallest at the left endpoint. This proves the supremum bound, including
both kinks. Equations \eqref{cut49:shortloc:cutoff}--\eqref{cut49:shortloc:crossingmoments},
$\sqrt n m_B\le1$, and \eqref{cut49:shortloc:master} now give the whole-count
witness with
\begin{align}\label{cut49:shortloc:BI}
 B(\lambda)&=z+\frac{\omega_+Kc_vG}{z(c_vN\lambda+G^2)}
                      -\frac{5r}{23}+\frac{5r^2\lambda}{46}-g,\\
 I(\lambda)&=V_0+aB_1+\frac{aP_1}{\lambda}+P_*D_+-m_cQ_*.
 \nonumber
\end{align}
We used $|u-f(u)|\le a(B_1\lambda+P_1)/n$ and $\sqrt n\pi_k\le1$.
Every constant is fixed before $\lambda$ varies. Both $B$ and
$mB+I/N$ are convex on $\lambda>0$; $I$ is decreasing. Thus their two
endpoint gates $B<0$, $I>0$, $mB+I/N<-.001$ prove the entire closed noise
cell. This proves all three LOW variants by exactly the same argument.

\subsection{The selected finite cover and the zero-noise endpoint}

The retained tables contain 45 central probability intervals covering
$[.38,.47]$ and 38 left intervals covering $[.304,.38]$. Each row supplies
its own $\Theta$, signed $H$ bounds, $d_0$, $\Lambda$, and global amplitude;
the confinement profiles may improve the amplitude with $\lambda$.
All these are hypotheses of the present section, acquired earlier without
using its conclusion. The shared resource has the form
\[
 nD_{\rm mom}<\frac{15000}{997}T_H,\qquad
 D_{\rm mom}\ge\alpha d+\beta\theta+P.
\]
Consequently a noise cell $[L,U]$ has
\[
 nd<G=\{(15000/997)T_H-NP-\beta\ell_-^2L\}/\alpha.
\]
Here $T_H$ is the maximum of the acquired concave quadratic budget on the
closed signed $H$ cell, further bounded by its scalar ceiling. It already
contains the phase charge. If $G\le0$, the entire cell is impossible and
receives no score test; the exclusion is strict even at $G=0$.

For positive noise the complete cover is the following. Each interval is
intersected with the row ceiling $\Lambda$; shared endpoints are included.
\begin{center}\small
\begin{tabular}{@{}llll@{}}\toprule
Noise & Central phase & Left phase & Witness\\\midrule
$(0,.05]$ & all & all & nearest count (scaled)\\
$[.05,.06]$ & $|u|\ge.09$ & $|u|\ge.09$ & nearest count\\
$[.06,.08]$ & $|u|\ge.13$ & $|u|\ge.14$ & nearest count\\
$[.08,.10]$ & $|u|\ge.18$ & $|u|\ge.19$ & nearest count\\
$[.10,.12]$ & $|u|\ge.29$ & $u\le-.29$ & nearest count\\
$[.10,.12]$ & --- & $u\in[.29,.5]$ & whole count\\
$[.05,\Lambda]$ & remaining phase & remaining phase & whole count\\\bottomrule
\end{tabular}
\end{center}
Beyond $.12$ the last row covers the full phase interval. The three LOW
objects are one central cover, one left cover with phase slices, and one
extra positive left cover. Their distinction is only the closed domain
and the chosen outward constants. Signed $H$ cells meet at zero; nearest
count ties meet at $u=\pm1/2$. The fixed coverage checker also checks every
row-ceiling singleton and every entirely excluded LOW row.

The common TINY parameters are
\[
 N=47,\quad p\in[.304,.47],\quad |H|\le1.393,\quad
 \Theta=.0154,\quad d_0=.0181,\quad A=.425,\quad\Lambda=.05.
\]
For its scaled own witness use $a=.4784$, $K_0=9/(4\cdot.4641)$,
$(z_-,c_-)=(1.66,.514)$ and $(z_+,c_+)=(1.2,1.036)$.
Its six phase cells are $[-.5,0],[0,.1],\ldots,[.4,.5]$.
The first two use scaled Poisson anchors $1.25,1$, respectively; the other
four use the fixed positive normal bound. The exact normalizers and gates
specified in the accompanying certificate map verify these constants.
The central MID amplitudes are $.35,.35,.30,.275$ in its four rows.
The left ones are $.375,.375,.325,.30$, each further reduced by the row's
global amplitude. No interpolation of the tables is used.

All six accepted comparisons give the same strict conclusion
\[
                        \rho_B(C-C_B)<-.001\lambda.
\]
Their proofs have now reduced to Theorem~\ref{cut49:shortloc:theorem}, the shared
normalization, and the two explicitly derived witnesses. The arithmetic
objects still have their original domains and targets; no certificate
branch has been silently deleted or relabeled.

If $\lambda=0$, conditional centering and $p,q>0$ imply $U_0=U_1=0$
almost surely. The law is exactly Bernoulli. For an attained positive closed
maximum its threshold is a binomial atom: lowering a noninteger threshold
to the preceding atom leaves the CDF unchanged and strictly improves the
normal term, while thresholds below the first atom have nonpositive score.
Thus $u=0$ at such a maximum and $C=C_B$. Schulz's all-probability
Bernoulli theorem~\cite[Theorem~1]{Schulz2016} gives $C_B<C_E$ for every
$0<p<1$ and in particular for all the indices $n\ge47$ used here. The same
external theorem closes the positive-noise conclusion $C<C_B<C_E$; no
finite Bernoulli reserve is needed. No estimate
in this section divides by zero or passes from fixed-law convergence to
uniformity in varying laws.

\section{Exact arithmetic and the proof boundary}\label{cut49:short:arithmetic}

The preceding sections reduce the theorem to finitely many rational
inequalities on closed parameter cells. The executable supplement is
part of this computer-assisted proof. It fixes the row data, rational
witnesses and subdivisions; acceptance reconstructs their inequalities.
The numerical tables are not evidence of a sampled inequality, and
no optimizer's success or infeasibility report is used as acceptance.

\subsection{A small arithmetic language}
Integers and reduced fractions are exact. A square root of a rational
$x\ge0$ is enclosed using an integer $k$ with
$k^2\le m^2x<(k+1)^2$; the endpoints are $k/m,(k+1)/m$.
Every denominator's sign is checked before division. Exponentials,
logarithms, sines and cosines occur only through the finite rational
series and remainder bounds proved in the text. Those bounds include
the endpoint $1/n=0$ where an infinite index band is used.

For clarity, the polynomial test used throughout can be reconstructed
without a library. If $P(x)=\sum_{j=0}^d a_jx^j$ and
$x=l+(u-l)y$, then its coefficients in powers of $y$ are
\[
 c_i=\sum_{j=i}^d a_j\binom ji l^{j-i}(u-l)^i.
\]
Its Bernstein coefficients are
\[
 b_k=\sum_{i=0}^k\frac{\binom ki}{\binom di}c_i,
 \qquad
 P(l+(u-l)y)=\sum_{k=0}^d b_k\binom dk y^k(1-y)^{d-k}.
\]
The basis is nonnegative and sums to one on $0\le y\le1$.
Thus $\min b_k\le P\le\max b_k$ on the entire closed interval.
Applying the transform one coordinate at a time proves the tensor
version. Strict positive coefficients suffice for positivity; failure
to obtain them is a rejection, not a proof of negativity. On an infinite
ray, strict signs of a polynomial and all its derivatives at the endpoint
supply the separate ray certificate. The independent transform checks
expand the Bernstein expression back into monomials.

The resource iteration is a finite induction over rows and passes.
Each accepted state encloses the same unknown law; subsequent states
intersect it with proved inequalities. Counts of feasible and excluded
cells must sum to the prescribed complete cover. An excluded cell
requires a strict contradiction, including at boundary values. The
local score tests check sign premises before substituting a mass lower
bound or the index floor. Convexity in the noise parameter then reduces
each closed noise cell to its two endpoints.

\subsection{The dependency graph and the supplement}
The new logical graph is
\[
\begin{gathered}
 \text{finite caps, attainment, contact}
 \longrightarrow \text{support and resource formulas},\\
 \text{five supported exclusions}
 \longrightarrow \text{finite ceilings through }50\text{, sharpened at }46,\\
 \text{acquired parameter cover and confinement}
 \longrightarrow \text{uniform local comparison},\\
 \text{sharper ceiling + three losses}\longrightarrow C_{49}<C_E,\\
 \text{smoothing + least later failure}\longrightarrow C_n<C_E\ (n\ge49).
 \\
 \text{fourth-index contact + common defect measure + finite cover}
 \longrightarrow C_4\le4097/10000<C_E.
\end{gathered}
\]
The all-parameter Bernoulli theorem closes the local comparison.
No edge enters from an earlier numerical cutoff. The complete Fourier,
gap, moment and count arguments are shared tools proved once in this
manuscript; retaining those tools does not retain their former global
applications.

The electronic supplement supplies the rational row data, witnesses and
closed subdivisions used in these finite verifications. Every accepted
cell satisfies the inequalities stated above; every discarded cell has a
strict contradiction with its necessary constraints. The covers include
all shared endpoints and the infinite index band. Together with the
analytic implications proved here, these finite inequalities complete
the computer-assisted proof.

\section{Two summands: a rational tensor certificate}\label{n2:section}
\subsection{Result and proof boundary}
For a standardized law $P$, write
\[
 \E X=0,\qquad \E X^2=1,\qquad \rho(P)=\E|X|^3<\infty,
\]
\[
 \Delta_n(P)=\sup_x\left|P^{*n}(( -\infty,x\sqrt n])-\Phi(x)\right|,
 \qquad C_n(P)=\frac{\sqrt n\,\Delta_n(P)}{\rho(P)},\qquad C_n=\sup_P C_n(P).
\]
Let $\phi=\Phi'$ and $\phi_0=(2\pi)^{-1/2}$. Set
\[
 K=\frac{2\sqrt2}{7}=0.404061017820884\ldots,
 \qquad \CE=\frac{3+\sqrt{10}}{6\sqrt{2\pi}}=0.409732183702396\ldots.
\]
\begin{theorem}[Two-summand bound]\label{n2:thm:main}
For every standardized law $P$ with finite third absolute moment,
\[
 \boxed{\Delta_2(P)\le\frac27\rho(P).}
\]
Equivalently, $C_2\le K<\CE$.
\end{theorem}
The analytic reduction is proved below. The finite positivity assertions
are proved by the accompanying rational certificate and checkers. The finite certificate proves the stated upper bound, not an identification of the exact value of $C_2$.

The sole external numerical probability estimate is Shevtsova's published
structural inequality \cite{Shevtsova2013}:
\begin{equation}\label{n2:eq:struct}
 \frac{\sqrt n\,\Delta_n(P)}{\rho(P)}
 \le a\left(1+\frac b{\rho(P)}\right),
 \qquad a=\frac{1661}{5000},\quad b=\frac{429}{1000}.
\end{equation}
It implies
\begin{equation}\label{n2:eq:large}
 \rho\ge2\quad\Longrightarrow\quad
 C_n(P)\le a(1+b/2)<K.
\end{equation}
The last comparison is rational after squaring, since $K^2=8/49$.

\subsection{Reduction of a putative maximizer to \texorpdfstring{$[-6,6]$}{[-6,6]}}
\begin{lemma}[Attainment under the contrary hypothesis]\label{n2:lem:attain}
If $C_2>K$, the supremum is attained by a standardized law $P$ with $\rho(P)<2$.
\end{lemma}
\begin{proof}
By \eqref{n2:eq:large}, the supremum may be restricted to standardized laws with
$\rho\le2$. This class is weakly compact: the second moment gives tightness,
and the third-moment bound gives uniform integrability of the first two
moments, so their equalities survive weak limits. The third-moment bound is
closed by lower semicontinuity. Convolution preserves weak convergence.
The Kolmogorov distance to the continuous cdf $\Phi$ is continuous under weak
convergence, as follows from the L\'evy inequalities and uniform continuity
of $\Phi$. Since $\rho\ge1$ is lower semicontinuous, the quotient is upper
semicontinuous and attains its maximum. A maximizer with $\rho=2$ would
contradict \eqref{n2:eq:large}.
\end{proof}

Reflect a maximizing law if needed. Upper semicontinuity of the two signed
cdf discrepancies and their vanishing at infinity give a finite positive
maximizing point. Write
\[
 \begin{gathered}
 t=x\sqrt2,\qquad F=P^{*2}(( -\infty,t]),\qquad D=F-\Phi(x)>0,\\
 C=\frac{\sqrt2D}{\rho}=C_2,\qquad
 \lambda=\frac D\rho,\qquad m=\E[X|X|].
 \end{gathered}
\]
Here $|m|\le1$. The negative branch uses the left limit of the cdf; reflection
converts it to this positive right-continuous branch.

\begin{lemma}[Explicit contact inequality]\label{n2:lem:contact}
Define
\begin{align}
 Q(y)={}&\lambda|y|^3-\frac{x\phi(x)+3D}{2}y^2
       -(\sqrt2\phi(x)+3\lambda m)y\notag\\
       &\hspace{20mm}+2F+\frac{x\phi(x)+D}{2}.\label{n2:eq:Q}
\end{align}
Then
\begin{equation}\label{n2:eq:contact}
 2P(( -\infty,t-y])\le Q(y)\quad(y\in\R),
 \qquad Q(y)=2P(( -\infty,t-y])\quad(y\in\supp P).
\end{equation}
\end{lemma}
\begin{proof}
Contaminate the raw law, $R_\epsilon=(1-\epsilon)P+\epsilon\delta_y$, and then
standardize. Its mean and variance are
\[
 \mu_\epsilon=\epsilon y,\qquad
 \sigma_\epsilon^2=1+\epsilon(y^2-1)-\epsilon^2y^2.
\]
Evaluate the standardized sum at
$x_\epsilon=(t-2\mu_\epsilon)/(\sqrt2\sigma_\epsilon)$.
The event in the raw variables is the fixed event $Y_1+Y_2\le t$, so its
probability is a polynomial in $\epsilon$; no differentiation across a
moving cdf jump is used. At zero,
\[
 x'_0=-\sqrt2y-\frac x2(y^2-1),\qquad
 \rho'_0=|y|^3-3my+\frac\rho2-\frac{3\rho}2y^2,
\]
\[
 \left.\frac d{d\epsilon}R_\epsilon^{*2}(( -\infty,t])\right|_0
 =2\{P(( -\infty,t-y])-F\}.
\]
The derivative of the selected normalized score is nonpositive by global
maximality. Thus
$2\{P(( -\infty,t-y])-F\}-\phi(x)x'_0-\lambda\rho'_0\le0$,
which is the inequality in \eqref{n2:eq:contact}. Its two sides have the same
$P$-integral $2F$. The nonnegative difference vanishes $P$-almost surely.
It is lower semicontinuous because the cdf expression on the left is upper
semicontinuous. Its zero set therefore contains $\supp P$.
\end{proof}

\begin{lemma}[Six-unit support bound]\label{n2:lem:support}
Under $C_2>K$, every maximizing witness in Lemma~\ref{n2:lem:attain} satisfies
$\supp P\subset(-6,6)$.
\end{lemma}
\begin{proof}
Equation \eqref{n2:eq:struct} gives $C<1/2$ and $C\rho<a(2+b)<1$.
Hence $\lambda>2/7$, $D<1/\sqrt2<5/7$, and $\lambda<5/14$.
Using $\phi(x)<2/5$, $|x\phi(x)|<1/4$, $\sqrt2<10/7$,
$|m|\le1$, and $F,D\ge0$, we obtain
\[
 Q(y)\ge P_6(|y|),\qquad
 P_6(r)=\frac27r^3-\frac{67}{56}r^2-\frac{23}{14}r-\frac18.
\]
The exact identity
\[
 P_6(6+u)-2
 =\frac27u^3+\frac{221}{56}u^2+\frac{104}{7}u+\frac{373}{56}>0
 \quad(u\ge0)
\]
contradicts $Q\le2$ on the support whenever $|y|\ge6$.
\end{proof}
No finite-support assumption appears in this reduction.

\subsection{The two moment parameters}
Let $\gamma=\E X^3$. Every standardized law satisfies
\begin{equation}\label{n2:eq:momentregion}
 1\le\rho,\qquad |\gamma|\le\rho,\qquad
 \rho^2(\gamma^2+4)-(\gamma^2+2)^2\ge0.
\end{equation}
For completeness, the last inequality follows from the elementary defect
factorization of Theorem~\ref{orig:thm:defect}. Put
\[
 s=\sqrt{\gamma^2+4},\quad a_0=(s-\gamma)/2,\quad b_0=(s+\gamma)/2,
 \quad c_\gamma=\gamma(\gamma^2+6)/s^3.
\]
The function
\[
 w_\gamma(y)=|y|^3-c_\gamma y^3-12y^2/s^3+6\gamma y/s^3+4/s^3
\]
factors on the two half-lines as
\[
 w_\gamma(y)=\frac{2a_0^2}{s^3}(y-b_0)^2\{2+(a_0+3b_0)y\}
 \quad(y\ge0),
\]
\[
 w_\gamma(-z)=\frac{2b_0^2}{s^3}(z-a_0)^2\{2+(3a_0+b_0)z\}
 \quad(z\ge0).
\]
Thus $w_\gamma\ge0$ and
$\E w_\gamma(X)=\rho-(\gamma^2+2)/\sqrt{\gamma^2+4}\ge0$.
Squaring gives \eqref{n2:eq:momentregion}. In the remaining problem $\rho\le2$,
so it suffices to consider the compact domain
\begin{equation}\label{n2:eq:D}
 \mathcal D=\{(\gamma,\rho)\in[-2,2]\times[1,2]:
 \rho^2(\gamma^2+4)-(\gamma^2+2)^2\ge0\}.
\end{equation}
Only necessity of these moment constraints is used; no assertion that they
characterize all laws on $[-6,6]$ is needed.

\subsection{The tensor majorants}
Define
\[
 v_0(y)=1,\quad v_1(y)=y,\quad v_2(y)=y^2,
 \quad v_3(y)=y^3,\quad v_4(y)=|y|^3,
 \qquad \mu=(1,0,1,\gamma,\rho).
\]
For rational coefficients $c=(c_{ij})_{0\le i\le j\le4}$ set
\begin{align}
 M_c(u,v)&=\sum_i c_{ii}v_i(u)v_i(v)
 +\sum_{i<j}c_{ij}\{v_i(u)v_j(v)+v_j(u)v_i(v)\},\label{n2:eq:tensor}\\
 q_c(\gamma,\rho)&=\sum_i c_{ii}\mu_i^2
                 +2\sum_{i<j}c_{ij}\mu_i\mu_j.\label{n2:eq:expect}
\end{align}
Independence and the common distribution give the exact identity
\begin{equation}\label{n2:eq:factor}
                      \E M_c(X_1,X_2)=q_c(\gamma,\rho).
\end{equation}
On each quadrant $M_c$ is a polynomial of degree at most three in each
coordinate. Its expectation is a polynomial of total degree at most two
in $(\gamma,\rho)$.

For $j=0,\ldots,69$, put
\[
 a_j=-\frac52+\frac j{20},\qquad b_j=a_j+\frac1{20}.
\]
The certificate supplies one vector $c^{(j)}$ for each interval. It proves
\begin{align}
 M_{c^{(j)}}(u,v)&\ge\mathbf1_{\{u+v\le b_j\}}
                 &&(u,v\in[-6,6]),\label{n2:eq:C1}\\
 q_{c^{(j)}}(\gamma,\rho)&\le\Phi(a_j/\sqrt2)+\frac27\rho
                 &&((\gamma,\rho)\in\mathcal D).\label{n2:eq:C2}
\end{align}
The rational certificate contains exactly 70 threshold rows and 1050 coefficients;
coefficients are ordered lexicographically by $(i,j)$ with $i\le j$.

\begin{lemma}[Certificate implication]\label{n2:lem:certimp}
Equations \eqref{n2:eq:C1}--\eqref{n2:eq:C2} imply
$\Delta_2(P)\le2\rho(P)/7$ for every standardized law on $[-6,6]$ with
$\rho\le2$.
\end{lemma}
\begin{proof}
If $t\in[a_j,b_j]$, monotonicity, \eqref{n2:eq:factor}, and the certificate give
\begin{align*}
 P(X_1+X_2\le t)-\Phi(t/\sqrt2)
 &\le\E M_{c^{(j)}}(X_1,X_2)-\Phi(a_j/\sqrt2)\\
 &\le\frac27\rho.
\end{align*}
For $t\le-5/2$, Cantelli's inequality gives
$P(X_1+X_2\le t)\le2/(2+t^2)\le8/33<2/7\le2\rho/7$.
For $t\ge1$, $1-\Phi(t/\sqrt2)\le1-\Phi(1/\sqrt2)<2/7$.
The last strict inequality is checked by the rational Gaussian enclosure below.
Thus the positive branch is bounded at every real threshold.
Applying the same result to $-X$ bounds
$\Phi(t/\sqrt2)-P(X_1+X_2<t)$, and hence also the negative right-cdf error.
Atoms and exact sum collisions are included, not removed by genericity.
\end{proof}

\subsection{Exact checking of the continuous domains}
\subsubsection{Piecewise-polynomial positivity}
Each of the four quadrants of $[-6,6]^2$ is intersected with the two closed
half-planes $u+v\le b_j$ and $u+v\ge b_j$. Each resulting rational polygon is
triangulated. On a triangle, the checker proves nonnegativity of
$M_{c^{(j)}}-1$ on the first side or $M_{c^{(j)}}$ on the second.
The shared boundary belongs to both closed regions.

Write a degree-six polynomial in barycentric Bernstein form on a triangle:
\[
 p=\sum_{\alpha_0+\alpha_1+\alpha_2=6}
 B_\alpha\frac{6!}{\alpha_0!\alpha_1!\alpha_2!}
 \lambda_0^{\alpha_0}\lambda_1^{\alpha_1}\lambda_2^{\alpha_2}.
\]
If every $B_\alpha\ge0$, then $p\ge0$ on the whole triangle. Otherwise the
longest edge is bisected and both children are checked. A resource limit or
negative vertex value is a rejection, never an acceptance.

The first checker constructs the coefficients from a monomial expansion and
subdivides them using rational de Casteljau identities. The second checker
independently reconstructs the homogeneous polynomial on \emph{every} child.
If $D_c$ clears polynomial denominators and $D_v$ clears vertex denominators,
write $U=D_vu$, $V=D_vv$, and $S=\lambda_0+\lambda_1+\lambda_2$.
For $p(u,v)=\sum a_{rs}u^rv^s$ the second checker expands
\[
 \sum_{r,s}(D_ca_{rs})D_v^{6-r-s}U^rV^sS^{6-r-s}.
\]
All its coefficients are integers. They are nonnegative exactly when the
corresponding Bernstein coefficients are nonnegative, because the omitted
denominators and multinomial factors are positive. This computation does not
use the first checker's subdivision formula.

\subsubsection{Moment-parameter positivity}
The scalar expression
$\Phi(a_j/\sqrt2)+2\rho/7-q_{c^{(j)}}(\gamma,\rho)$ is enclosed from below
on rational rectangles using degree-two Bernstein coefficients in each
variable. A rectangle outside $\mathcal D$ is discarded. This test is exact:
if $g_*$ is the minimum of $\gamma^2$ on its $\gamma$ interval and $r_+$ is
its upper $\rho$ endpoint, the rectangle is outside whenever
\[
 r_+^2(g_*+4)-(g_*+2)^2<0.
\]
Indeed, $r^2(w+4)-(w+2)^2$ increases in $r\ge0$ and decreases in
$w\ge0$ for $r\le2$. All other rectangles are certified or subdivided.

\subsubsection{The normal cdf}
For $|t|\le5/2$,
\begin{equation}\label{n2:eq:normal}
 \Phi(t/\sqrt2)=\frac12+\frac1{2\sqrt\pi}
 \sum_{k=0}^\infty\frac{(-1)^k t^{2k+1}}{4^k k!(2k+1)}.
\end{equation}
For $t\ge0$ the term magnitudes decrease: their successive ratio is at most
$(25/16)(2k+1)/((k+1)(2k+3))<1$. Alternating partial sums therefore give
rational enclosures after bounding $1/(2\sqrt\pi)$. Machin's arctangent
identity and integer square roots supply those bounds. The second checker
uses different rational prefactor bounds and different truncation orders.
Negative arguments are handled by symmetry. No platform implementation of
$\Phi$, exponential, sine, or cosine is used for acceptance.

\subsubsection{The complete finite certificate}
The supplied certificate has exactly seventy consecutive threshold
rows and \(1050\) rational coefficients. Complete exact verification
establishes both the pointwise majorant and moment-parameter inequalities
on their stated closed domains. The two independent positivity
calculations described above verify the same finite obligations.

\subsection{Completion of the two-summand proof}
\begin{proof}[Proof of Theorem~\ref{n2:thm:main}]
Suppose $C_2>K$. Lemma~\ref{n2:lem:attain} supplies a global maximizing law
with $\rho<2$. Lemma~\ref{n2:lem:support} confines it to $(-6,6)$.
The complete rational verification establishes \eqref{n2:eq:C1}--\eqref{n2:eq:C2};
Lemma~\ref{n2:lem:certimp} then gives $C_2(P)\le K$, a contradiction.
Thus $C_2\le K$. Finally, exact rational enclosures give $K<0.405<0.409<\CE$
using rational bounds and squaring.
\end{proof}
The bounded-support theorem in Lemma~\ref{n2:lem:certimp} treats arbitrary
probability measures, not merely finitely supported measures or an assumed
number of atoms. The passage to unrestricted laws uses global extremality,
not a truncation approximation.


\section{Three summands: independence-null certificates}\label{n3:section}
\subsection{Result and external input}
For a standardized law $P$, let
\[
 \E X=0,\quad \E X^2=1,\quad \rho(P)=\E|X|^3<\infty,
\]
\[
 \Delta_n(P)=\sup_x\left|P^{*n}(( -\infty,x\sqrt n])-\Phi(x)\right|,
 \qquad C_n=\sup_P\frac{\sqrt n\,\Delta_n(P)}{\rho(P)}.
\]
Write $\phi=\Phi'$ and set
\[
 \Lambda=\Lam,\qquad K=\sqrt3\Lambda
 =0.408763990586255\ldots,
 \qquad \CE=\frac{3+\sqrt{10}}{6\sqrt{2\pi}}.
\]
\begin{theorem}[Three summands]\label{n3:thm:main}
Every standardized law with finite third absolute moment satisfies
\[
 \boxed{\Delta_3(P)\le\Lambda\rho(P).}
\]
Consequently $C_3\le59\sqrt3/250<\CE$.
\end{theorem}
The analytic reduction and the finite-certificate implication are proved
below. The polynomial inequalities are verified by the accompanying exact
checker and rational data. This upper bound does not identify the exact value of $C_3$.

The sole external numerical probability bound is Shevtsova's structural
estimate \cite{Shevtsova2013}:
\begin{equation}\label{n3:eq:struct}
 C_n(P)\le A\left(1+\frac B{\rho(P)}\right),
 \qquad A=\frac{1661}{5000},\quad B=\frac{429}{1000}.
\end{equation}
It implies
\begin{equation}\label{n3:eq:large}
 \rho(P)\ge\rhoMax\quad\Longrightarrow\quad
 C_n(P)\le\frac{159456}{390625}<K.
\end{equation}
The last inequality is rational after squaring. No result about the optimal
two-point constant, no near-lattice transfer, and no effective eventual
threshold is used.

\subsection{A bounded witness would attain any violation}
\begin{lemma}[Attainment under the contrary hypothesis]\label{n3:lem:attain}
If $C_3>K$, a standardized law $P$ with $\rho(P)<15/8$ attains $C_3$.
\end{lemma}
\begin{proof}
The laws with $\E X=0$, $\E X^2=1$ and $\E|X|^3\le15/8$ form a weakly
compact set. Tightness follows from the moment bound. For $j=1,2$, the
uniform tail estimate
\[
 \E\bigl[|X|^j\one_{\{|X|>R\}}\bigr]\le\frac{15}{8R^{3-j}}
\]
preserves the first two moment equalities under weak limits; lower
semicontinuity preserves the third-moment bound. Convolution preserves weak
convergence. The Kolmogorov distance to the fixed continuous cdf $\Phi$ is
continuous under weak convergence: the L\'evy inequalities change the cdf
argument by at most an arbitrarily small amount, whose cost is bounded by
the uniform modulus of continuity of $\Phi$. Since $\rho\ge1$ is lower
semicontinuous, $\Delta_3/\rho$ is upper semicontinuous on this compact set.
It attains its supremum there. Bound \eqref{n3:eq:large} excludes all remaining
laws and excludes $\rho=15/8$ at a violating maximizer.
\end{proof}
After reflecting a maximizing law if necessary, its maximum is attained
at a positive right-cdf discrepancy. Indeed, the functions
$F(x)-\Phi(x)$ and $\Phi(x)-F(x-)$ are upper semicontinuous and vanish at
infinity; reflection interchanges these two signed discrepancies. Put
\[
 \begin{gathered}
 t=x\sqrt3,\quad F=P^{*3}(( -\infty,t]),\quad D=F-\Phi(x)>0,\\
 C=C_3=\frac{\sqrt3D}{\rho},\quad \lambda=\frac D\rho,
 \quad m=\E[X|X|],\quad |m|\le1.
 \end{gathered}
\]

\begin{lemma}[Explicit contamination contact inequality]\label{n3:lem:contact}
Define
\begin{align}
 Q(y)={}&\lambda|y|^3-\frac{x\phi(x)+3D}{2}y^2
 -(\sqrt3\phi(x)+3\lambda m)y \notag\\
 &\hspace{22mm}+3F+\frac{x\phi(x)+D}{2}.
 \label{n3:eq:Q}
\end{align}
Then
\begin{equation}\label{n3:eq:contact}
 3P^{*2}(( -\infty,t-y])\le Q(y)\qquad(y\in\R),
\end{equation}
with equality throughout $\supp P$.
\end{lemma}
\begin{proof}
For fixed $y$, contaminate the raw law by
$R_\epsilon=(1-\epsilon)P+\epsilon\delta_y$. Its mean and variance are
\[
 \mu_\epsilon=\epsilon y,\qquad
 \sigma_\epsilon^2=1+\epsilon(y^2-1)-\epsilon^2y^2.
\]
Standardize $R_\epsilon$ and evaluate at
\[
 x_\epsilon=\frac{t-3\mu_\epsilon}{\sqrt3\,\sigma_\epsilon}.
\]
This choice keeps the event in the raw variables exactly
$Y_1+Y_2+Y_3\le t$, even at cdf jumps. Its probability is a polynomial
in $\epsilon$. Direct differentiation gives
\[
 x'_0=-\sqrt3 y-\frac x2(y^2-1),
\]
\[
 \rho'_0=|y|^3-3my+\frac\rho2-\frac{3\rho}{2}y^2,
\qquad
 F'_0=3\{P^{*2}(( -\infty,t-y])-F\}.
\]
Differentiation of the third absolute moment is justified by the integrable
quadratic majorant for the derivative of $|z-\epsilon y|^3$.
Global maximality makes the right derivative of the selected normalized
score nonpositive. Substitution and rearrangement give \eqref{n3:eq:contact}.
Both sides have integral $3F$ under $P$: for the polynomial this follows
from $\lambda\rho=D$, $\E X=0$ and $\E X^2=1$.
The nonnegative gap is therefore zero $P$-almost surely. The function
$y\mapsto P^{*2}(( -\infty,t-y])$ is upper semicontinuous, so the gap is
lower semicontinuous. It must vanish everywhere on $\supp P$.
\end{proof}

\begin{lemma}[Uniform support bound]\label{n3:lem:support}
A violating maximizer in Lemma~\ref{n3:lem:attain} satisfies
$\supp P\subset(-6,6)$.
\end{lemma}
\begin{proof}
The structural bound and $1\le\rho<15/8$ give
\[
 \frac25<C<\frac{19}{40},\qquad C\rho<\frac{77}{100}.
\]
Using $12/7<\sqrt3<7/4$ and $\sqrt3>17/10$, we obtain
\[
 D<\frac9{20},\qquad \frac8{35}<\lambda<\frac27.
\]
Also $\phi(x)<2/5$ and $|x\phi(x)|<1/4$. The first follows from
$\pi>25/8$, and the second from the maximum
$(2\pi e)^{-1/2}<1/4$. Thus \eqref{n3:eq:Q} implies, with $r=|y|$,
\[
 Q(y)\ge\frac8{35}r^3-\frac45r^2-\frac85r-\frac18.
\]
For $r=6+u$, the difference of the right side and $3$ equals
\[
 \frac8{35}u^3+\frac{116}{35}u^2+\frac{472}{35}u+
 \frac{2197}{280}>0\qquad(u\ge0).
\]
But \eqref{n3:eq:contact} is an equality on the support, and its left side is
at most $3$. This excludes every support point with $|y|\ge6$.
\end{proof}
\begin{remark}
This is an extremizer reduction, not a truncation of arbitrary laws.
The original theorem assumes only a finite third absolute moment.
Higher moments are used below only after the hypothetical violating
maximizer has been proved bounded.
\end{remark}

\subsection{The admissible two-dimensional moment region}
Write $\gamma=\E X^3$. The defect identity in
Theorem~\ref{orig:thm:defect} gives an exact necessary condition. For completeness,
the elementary proof is included.
Set
\[
 s=\sqrt{\gamma^2+4},\quad a=\frac{s-\gamma}{2},\quad
 b=\frac{s+\gamma}{2},\quad c=\frac{\gamma(\gamma^2+6)}{s^3},
\]
so $ab=1$. Define
\[
 w_\gamma(y)=|y|^3-cy^3-\frac{12y^2}{s^3}
             +\frac{6\gamma y}{s^3}+\frac4{s^3}.
\]
For $y,z\ge0$ direct multiplication gives
\[
 w_\gamma(y)=\frac{2a^2}{s^3}(y-b)^2\{2+(a+3b)y\},
\]
\[
 w_\gamma(-z)=\frac{2b^2}{s^3}(z-a)^2\{2+(3a+b)z\}.
\]
Therefore
\[
 0\le\E w_\gamma(X)=\rho-\frac{\gamma^2+2}{\sqrt{\gamma^2+4}}.
\]
Together with $|\gamma|\le\rho$ and $\rho\ge1$, every relevant law belongs to
\begin{equation}\label{n3:eq:momdomain}
 \mathcal D=\left\{(\gamma,\rho):
 |\gamma|\le\frac{15}{8},\quad1\le\rho\le\frac{15}{8},\quad
 \rho^2(\gamma^2+4)-(\gamma^2+2)^2\ge0\right\}.
\end{equation}
This is an outer domain: every point need not be realizable by a law.
The certificate checks all of it, so non-realizable points cannot cause
an invalid acceptance.

\subsection{Independence-null controls}
For a single coordinate define the features
\begin{equation}\label{n3:eq:features}
 (v_0,v_1,v_2,v_3,v_4,v_5,v_6)(y)
 =(1,y,y^2-1,y^3,|y|^3,y^4,y^5).
\end{equation}
Let $\mathcal I$ consist of the nondecreasing triples
$\alpha=(i,j,k)$ with $0\le i\le j\le k\le6$ such that either
$k<5$, or at least one entry is $1$ or $2$. There are $59$ such triples,
ordered lexicographically in the certificate files.
For each $\alpha$, let
\[
 V_\alpha(x_1,x_2,x_3)
 =\frac1{|\operatorname{Perm}(\alpha)|}
   \sum_{(i,j,k)\in\operatorname{Perm}(\alpha)}
        v_i(x_1)v_j(x_2)v_k(x_3),
\]
where the sum is over the \emph{distinct} permutations of the index triple.
For $c\in\mathbb Q^{59}$ put
\begin{equation}\label{n3:eq:majorant}
 M_c=\sum_{\alpha\in\mathcal I}c_\alpha V_\alpha.
\end{equation}

\begin{lemma}[Exact expectation without higher-moment parameters]\label{n3:lem:expect}
For independent standardized identically distributed bounded variables,
\begin{equation}\label{n3:eq:expect}
 \E M_c(X_1,X_2,X_3)=q_c(\gamma,\rho)
 :=\sum_{\substack{\alpha\in\mathcal I\\
                    \text{entries only }0,3,4}}
        c_\alpha\gamma^{\#3(\alpha)}\rho^{\#4(\alpha)}.
\end{equation}
In particular, $q_c$ has total degree at most three and has no unknown
fourth- or fifth-moment parameters.
\end{lemma}
\begin{proof}
The expectations of $v_1$ and $v_2$ are zero. Independence makes every
product containing either factor have expectation zero. By the definition
of $\mathcal I$, every product involving $v_5$ or $v_6$ contains such a
zero-mean factor in another coordinate. The only remaining products use
$v_0,v_3,v_4$, whose means are $1,\gamma,\rho$.
\end{proof}
This adds $24$ coefficients to the $35$ low-feature tensor coefficients.
In the resulting family, $49$ basis elements have zero expectation and
only $10$ contribute monomials in $(\gamma,\rho)$.
For example, all permutations of
\[
 X_1^4X_2,
 \qquad X_1^5(X_2^2-1),
 \qquad X_1^5X_2|X_3|^3
\]
have mean zero. These are exact independence identities, not assumptions
that the fourth or fifth moment has a prescribed value.

The general principle is useful beyond three summands. For any integrable
function of the other coordinates,
\[
 \E\{X_i A(X_1,\ldots,\widehat X_i,\ldots,X_n)\}=0,
\]
\[
 \E\{(X_i^2-1)B(X_1,\ldots,\widehat X_i,\ldots,X_n)\}=0.
\]
Such terms may reshape a pointwise majorant without changing its i.i.d.
expectation. For the present proof all functions are polynomials on a
bounded cube, so integrability is automatic.

\subsection{Finite-certificate implication}
\begin{proposition}[Continuous-domain certificate]\label{n3:prop:certificate}
For each row $j$ of the supplied rational certificate, let $[a_j,b_j]$
be its threshold interval and $c^{(j)}\in\mathbb Q^{59}$ its coefficients.
The rows form an exact consecutive partition of $[-7/2,3/2]$ and satisfy
\begin{align}
 M_{c^{(j)}}(u,v,w)&\ge\one_{\{u+v+w\le b_j\}}
       &&((u,v,w)\in[-7,7]^3),\label{n3:eq:pointwise}\\
 q_{c^{(j)}}(\gamma,\rho)&\le
      \Phi(a_j/\sqrt3)+\Lambda\rho
       &&((\gamma,\rho)\in\mathcal D).\label{n3:eq:scalar}
\end{align}
These assertions are verified by the exact algorithm in Section~\ref{n3:sec:check}.
\end{proposition}

\begin{proof}[Proof of Theorem~\ref{n3:thm:main}]
Suppose $C_3>K$. Lemmas~\ref{n3:lem:attain}--\ref{n3:lem:support} give a globally
maximizing violating law on $(-6,6)$, with $\rho<15/8$, and a positive
maximizing raw threshold $t$.
If $t\in[a_j,b_j]$, then independence and Proposition~\ref{n3:prop:certificate}
give
\begin{align*}
 P(X_1+X_2+X_3\le t)-\Phi(t/\sqrt3)
 &\le\E M_{c^{(j)}}(X_1,X_2,X_3)-\Phi(a_j/\sqrt3)\\
 &\le\Lambda\rho.
\end{align*}
The larger certificate cube $[-7,7]^3$ contains the entire witness support.
For $t\le-7/2$, Cantelli's inequality for the variance-three sum yields
\[
 P(X_1+X_2+X_3\le t)-\Phi(t/\sqrt3)
 \le\frac3{3+t^2}\le\frac{12}{61}<\Lambda\le\Lambda\rho.
\]
For $t\ge3/2$, the discrepancy is at most
\[
 1-\Phi(\sqrt3/2)<\Lambda\le\Lambda\rho.
\]
All these inequalities contradict the choice of the maximizing witness.
The reduction by reflection already handles the other signed discrepancy;
no atom or boundary event is discarded.
\end{proof}

\subsection{Exact verification algorithm}\label{n3:sec:check}
Each coefficient is rational. Verification reconstructs the polynomial from its defining features and requires complete threshold coverage. Every unresolved positivity test is rejected.

\subsubsection{Piecewise polynomial representation and symmetry}
On each sign orthant, $|y|^3$ is $y^3$ or $-y^3$. Thus $M_c$ is a polynomial
of degree at most five in each variable. The checker expands this polynomial
from the feature definition, not from a supplied power expansion. Symmetry
permits restriction to $u\le v\le w$. The four initial sign boxes are
\[
 [-7,0]^3,\quad[-7,0]^2\times[0,7],\quad
 [-7,0]\times[0,7]^2,\quad[0,7]^3.
\]
A dyadic box $[L_1,H_1]\times[L_2,H_2]\times[L_3,H_3]$ may be discarded
only if $L_i>H_j$ for some $i<j$. Such a box has no ordered point. All other
accepted inequalities are proved on the whole box, a safe enlargement.

\subsubsection{An affine half-space Bernstein test}
For a box, let $B_\alpha$ be the tensor Bernstein basis of degree five in
three variables, so $B_\alpha\ge0$ and $\sum B_\alpha=1$. Write
\[
 M_c=\sum_\alpha b_\alpha B_\alpha,\qquad
 u+v+w-b_j=\sum_\alpha s_\alpha B_\alpha,
\]
where
\[
 s_\alpha=\sum_{i=1}^3\left(L_i+\frac{\alpha_i}{5}(H_i-L_i)\right)-b_j.
\]
\begin{lemma}[Exact affine test]\label{n3:lem:affine}
If a real $\ell\le0$ satisfies $b_\alpha\ge1+\ell s_\alpha$ for all
$\alpha$, then $M_c\ge1$ on the portion $u+v+w\le b_j$.
If $\ell\ge0$ and $b_\alpha\ge\ell s_\alpha$ for all $\alpha$, then
$M_c\ge0$ on the portion $u+v+w\ge b_j$.
\end{lemma}
\begin{proof}
Average the coefficient inequalities with the nonnegative Bernstein weights.
In each case the sign restriction makes $\ell(u+v+w-b_j)$ nonnegative
on the specified half-space.
\end{proof}
Existence of $\ell$ is checked by intersecting rational half-lines. A
coefficient with $s_\alpha=0$ requires $b_\alpha\ge1$ or $b_\alpha\ge0$.
The remaining coefficients give upper or lower rational bounds on $\ell$.
Only signed integer products are needed to compare these bounds.

A box wholly inside the event is checked by $\min b_\alpha\ge1$; a box
strictly outside by $\min b_\alpha\ge0$. A crossing box uses the affine
tests. An unresolved box is bisected along a widest coordinate. Dyadic
subdivision uses the exact degree-five de Casteljau identities
\[
 b_i^{\rm left}=2^{-i}\sum_{j=0}^i\binom ij b_j,
 \qquad
 b_i^{\rm right}=2^{-(5-i)}\sum_{j=i}^5\binom{5-i}{j-i}b_j.
\]
After multiplying the common denominator by $32$, all new numerators are
integers. Ties at the event plane remain on the included side, and that
plane is required to have value at least one.

\subsubsection{Moment region and Gaussian enclosure}
The scalar polynomial
\[
 \underline\Phi(a_j/\sqrt3)+\Lambda\rho-q_{c^{(j)}}(\gamma,\rho)
\]
is checked by rational tensor Bernstein coefficients of degree three on
moment rectangles. Here $\underline\Phi$ is a rigorous rational lower bound.
For $s=\gamma^2$ the domain constraint is
\[
 h(s,\rho)=\rho^2(s+4)-(s+2)^2\ge0.
\]
When $\rho\le15/8<2$, this expression decreases in $s\ge0$ and increases
in $\rho\ge0$. If its maximum at the smallest possible $\gamma^2$ and
largest possible $\rho$ on a rectangle is negative, that rectangle is
outside the domain and is discarded. Otherwise Bernstein positivity proves
the scalar inequality or the rectangle is bisected.

For rational $t$ in the required range,
\[
 \Phi(t/\sqrt3)=\frac12+\frac1{\sqrt{6\pi}}
 \sum_{k=0}^{\infty}\frac{(-1)^k t^{2k+1}}{6^k k!(2k+1)}.
\]
The checker uses consecutive alternating partial sums through indices
$34,35$. For $|t|\le4$ the magnitudes decrease, so these enclose the sum.
Rational Machin bounds enclose $\pi$, and integer square roots enclose
$\sqrt{6\pi}$ to rational precision. No floating-point normal cdf enters
acceptance.

\subsubsection{The complete finite certificate}
The certificate contains \CertificateCells\ consecutive threshold
intervals and \CertificateCoefficients\ rational coefficients.
Exact integer calculations reconstruct the polynomial expansions and
verify the continuous cube and moment inequalities. Undecided cells
are rejected.

The initial expansions are also checked on a complete \(6^3\)
tensor interpolation grid in each sign box. Since the degree is at
most five in each coordinate, agreement on this grid proves the
polynomial identity. This is an independent identity check, not a
second verification of all positivity inequalities.


\section{Four summands: an unrestricted bound}\label{n4all:section}

\begin{theorem}\label{n4all:thm:main}
For every standardized real law with finite third absolute moment,
\[
 R_4(P)\le\frac{4097}{10000}<\CE.
\]
In particular, $C_4\le4097/10000$. No restriction on the number of
support points is imposed.
\end{theorem}

Put $T=4097/10000$ and $c_0=256/625$. The proof has two steps.
A hypothetical unrestricted maximizer above $T$ lies in a compact
moment class. A finite rational certificate bounds every law in that
class by $T$. The three-point law introduced below is a reference law,
not a reduction of the admissible laws to three points.

\subsection{The canonical reference and its excess}
For a standardized $X$, write
\[
 \begin{gathered}
 \vartheta=\E[X|X|],\qquad u=\E X_+=\E X_->0,\qquad
 v_\pm=\frac{1\pm\vartheta}{2},\\
 a=\frac{v_-}{u},\qquad b=\frac{v_+}{u},\qquad
 w=\frac{4u^2}{1-\vartheta^2}.
 \end{gathered}
\]
Both $v_\pm$ are positive. Cauchy's inequality on the two half-lines
gives $0<w\le1$. Define
\begin{equation}\label{n4all:eq:reference}
 Q=\frac{u}{a}\delta_{-a}+(1-w)\delta_0+\frac{u}{b}\delta_b,
 \qquad
 r=\rho(Q)=\frac{1+\vartheta^2}{\sqrt{w(1-\vartheta^2)}}.
\end{equation}
This is a probability law of mean zero and variance one. Its contacts
are
\[
 a=\sqrt{\frac{1-\vartheta}{w(1+\vartheta)}},\qquad
 b=\sqrt{\frac{1+\vartheta}{w(1-\vartheta)}}.
\]

\begin{lemma}\label{n4all:lem:excess}
Set
\[
 W(X)=X_+(X_+-b)^2+X_-(X_--a)^2,\qquad
 \eta=\E W=\rho(P)-r.
\]
Then $\eta\ge0$, and $\eta=0$ if and only if $P=Q$.
Moreover,
\begin{align}
 \E X^3&=b-a+\E[\sgn(X)W(X)],\notag\\
 \E X^4&=a^2-ab+b^2+
   \E[(|X|+2c(X))W(X)],\label{n4all:eq:physical-moments}
\end{align}
whenever the fourth moment is finite, where $c(X)=b$ on $X>0$ and
$c(X)=a$ on $X<0$.
\end{lemma}
\begin{proof}
On a tagged half-line put $y=|X|$ and $c=a$ or $b$. The identities
\[
 y^3=y(y-c)^2+2cy^2-c^2y,
 \qquad
 y^4=(y+2c)y(y-c)^2+3c^2y^2-2c^3y
\]
integrate using $\E[y;\text{tag}]=u$ and
$\E[y^2;\text{tag}]=uc$. They give all three formulas.
If $\eta=0$, the support is contained in $\{-a,0,b\}$; the first
two moments fix its masses as in \eqref{n4all:eq:reference}.
The converse is immediate.
\end{proof}

The excess $\eta$ is not the two-point defect $d$ of Part~I.
Its zero set permits an atom at zero; the reference is two-point
precisely when $w=1$.

\subsection{Localization of an unrestricted extremizer}
Write $K_3=59\sqrt3/250$. Theorem~\ref{n3:thm:main} gives
$C_3\le K_3$, and
\[
 c_0^2-K_3^2=\frac{1069}{1562500}>0,\qquad c_0<T.
\]

\begin{lemma}\label{n4all:lem:owner}
If $C_4>c_0$, an unrestricted maximizing pair $(P,H)$ can be chosen
as a positive closed discrepancy. It satisfies
\begin{equation}\label{n4all:eq:owner-class}
 \supp P\subset(-101/25,101/25),\quad
 \rho(P)<46/25,\quad
 \frac{625}{2116}<w\le1,\quad |\vartheta|<\frac{321}{500}.
\end{equation}
\end{lemma}
\begin{proof}
Attainment is Lemma~\ref{cut49:global:attainment}.
Shevtsova's estimate \eqref{cut49:global:finitecaps} gives, at this pair,
\[
 \rho<\frac{(3031/10000)(323/500)}{c_0-3031/10000}
       =\frac{979013}{532500}<\frac{46}{25}.
\]
This bound is obtained before any support or fourth-moment assumption.
Put $C=C_4$, $z=H/2$, $F=F_{4,P}(H)$ and $G=F_{3,P}$.
Lemma~\ref{cut49:global:contact} gives
\begin{align}
 \rho I(y)={}&8\{G(H-y)-F\}+4y\phi(z)+z\phi(z)(y^2-1)\notag\\
 &+\tfrac32C\rho(y^2-1)
   -C\{|y|^3-\rho-3\vartheta y\}.\label{n4all:eq:contact}
\end{align}
Here $I\le0$ everywhere and $I=0$ at every point of the closed
support, including accumulation points. Using the all-law predecessor
bound and $F=\Phi(z)+C\rho/2$, with $s=\sqrt{4/3}$, yields
\begin{align}
 \rho I(y)\le{}&\mathcal H_4(z,y)+4\rho(K_3-C)
                  +4\rho(s-1)K_3\notag\\
 &+C\{-|y|^3+\tfrac32\rho y^2+3\vartheta y-\rho/2\},
 \label{n4all:eq:predecessor}
\end{align}
where
\[
 \mathcal H_4(z,y)=
 8\left[\Phi\!\left(\frac{z-y/2}{\sqrt{3/4}}\right)-\Phi(z)\right]
 +4y\phi(z)+z\phi(z)(y^2-1).
\]

We use the following elementary uniform estimate:
\begin{equation}\label{n4all:eq:gaussian}
 |\mathcal H_4(z,y)|<\frac{\phi(0)}6|y|^3+\frac{|y|}{5}+\frac1{15}.
\end{equation}
Indeed, $\phi(0)<2/5$, $\sup|\phi''|=\phi(0)$ and
$\sup|(x\phi(x))'|=\phi(0)$. Also
\[
 \sup_x|x(3-x^2)\phi(x)|<3/5.
\]
For the last inequality, set $t=x^2$. On $0\le t\le3$,
\[
 \frac94(1+t+t^2/2+t^3/6)-t(3-t)^2
 \ge\frac34\left[7(t-9/14)^2+3/28\right]>0.
\]
Compare with $e^t$ and use $\phi(0)^2<4/25$.
On $t\ge3$, use $|x|^3\phi(x)<1/2$, which follows from $e>8/3$.
For $f(\lambda)=\Phi((z-y/2)/\sqrt{1-\lambda})$,
\[
 f'(0)=\tfrac12(z-y/2)\phi(z-y/2),\qquad
 |f''(\lambda)|<4/15\quad(0\le\lambda\le1/4).
\]
Thus $8\{f(1/4)-f(0)\}=(z-y/2)\phi(z-y/2)+E$, with
$|E|<1/15$. A second-order translation of $\Phi$ gives the cubic
term in \eqref{n4all:eq:gaussian}; the remaining difference
$(z-y/2)\phi(z-y/2)-z\phi(z)$ is bounded by $|y|/5$.

By \eqref{n4all:eq:reference}, $r\le\rho<46/25$ and $w\le1$.
At $|\vartheta|=321/500$,
\[
 \frac{(1+\vartheta^2)^2}{1-\vartheta^2}
   =\frac{124637947681}{36739750000}>(46/25)^2.
\]
Monotonicity in $\vartheta^2$ proves $|\vartheta|<321/500$.
Drop the negative term $4\rho(K_3-C)$ in
\eqref{n4all:eq:predecessor}. Since $4(s-1)<5/8$, its remaining
$C$-bracket is bounded, for $r_y=|y|\ge4$, by
\[
 B(r_y)=-r_y^3+\frac{69}{25}r_y^2
                 +\frac{963}{500}r_y+\frac{23}{100}.
\]
The values of $B,B',B'',B'''$ at 4 are
$-5953/500,-11997/500,-462/25,-6$; hence $B<0$ on that ray.
Only now substitute $C>c_0$ in the negative bracket. Equation
\eqref{n4all:eq:gaussian} gives
\[
 \rho I(y)<P(r_y):=-\frac{643}{1875}r_y^3
 +\frac{17664}{15625}r_y^2+\frac{77257}{78125}r_y
 +\frac{7541}{46875}.
\]
At $101/25$ the values of $P,P',P'',P'''$ are respectively
\[
 -\frac{51557}{9765625},\quad -\frac{520966}{78125},\quad
 -\frac{94558}{15625},\quad -\frac{1286}{625}.
\]
They are negative, so $P<0$ thereafter. Contact excludes every
support point with $|y|\ge101/25$. Finally $r^2\ge1/w$ gives
$w>625/2116$. This proves \eqref{n4all:eq:owner-class}.
\end{proof}

\subsection{One defect measure for all frequencies}
For $\eta>0$, let $\nu=W\,dP/\eta$, viewed as a probability on
tagged points $(\sigma,y)$, $\sigma\in\{-1,1\}$, $y>0$.
Write $c_+=b$, $c_-=a$, $p_0=1-w$, and
\[
 k_0(\sigma,y)=\frac1{c_\sigma^2y},\qquad
 k_2(\sigma,y)=\frac1{(y-c_\sigma)^2},\qquad
 k_1(\sigma,y)=\frac1{y-c_\sigma}.
\]
The contacts have zero $\nu$ mass. Singular expressions are integrated
only off the contacts; the following identities prove integrability.

\begin{lemma}\label{n4all:lem:common-measure}
With $g(q)=\varphi_Q(q)$ and
$r_q(\sigma,y)=\exp(\ii\sigma q\,\cdot)[0,c_\sigma,c_\sigma,y]$,
\begin{align}
 \varphi_P(q)-g(q)&=\eta\int r_q\,d\nu,\notag\\
 \eta\int k_0\,d\nu&=p_0-P\{0\}\le p_0,\notag\\
 \eta\int_{\sigma}k_2\,d\nu&=u-c_\sigma P\{\sigma c_\sigma\}\le u,
 &\int_{\sigma}k_1\,d\nu&=0.\label{n4all:eq:budgets}
\end{align}
The same $\nu$ occurs at every frequency, and $|r_q|\le |q|^3/6$.
For $\eta=0$, $P=Q$ and every characteristic-function difference is
zero, without defining $\nu$.
\end{lemma}
\begin{proof}
Hermite interpolation at $0,c,c$ has remainder
$y(y-c)^2\exp(\ii\sigma q\,\cdot)[0,c,c,y]$.
On each tag, the linear and quadratic terms have the same integrals
against $P$ and $Q$, by their matching first two tagged moments.
Both interpolating polynomials have constant term 1. Their constant
contributions agree only after summing the tags and adding the atom
at zero, since both total masses are 1.
Hermite--Genocchi gives the response bound. Multiplication of the
three costs by $W=y(y-c)^2$ gives
$(y-c)^2/c^2$, $y$, and $y(y-c)$, respectively.
Their integrals give \eqref{n4all:eq:budgets}, including the omitted
contact masses. Absolute integrability of $k_1$ follows from Cauchy
and the $k_2$ bound. Lemma~\ref{n4all:lem:excess} handles $\eta=0$.
\end{proof}

Fix a threshold $H$ and finite coefficients $\kappa_j(H)$, defined
below. For the bounded class $|X|\le R_*=17/4$, all suprema below
are over $\sigma\in\{-1,1\}$ and $0<y\le R_*$, off the contact.
The linear response is
\[
 A_H(\sigma,y)=\Re\sum_{j=0}^{N-1}
        4\kappa_j(H)g(q_j)^3r_{q_j}(\sigma,y).
\]
For any fixed $\lambda_0,\lambda_+,\lambda_-\ge0$ and
$\theta_+,\theta_-\in\R$, a bound
\begin{equation}\label{n4all:eq:priced-supremum}
 S_H\ge\sup_{\sigma,y}
 \{A_H-\lambda_0k_0-\lambda_\sigma k_2-\theta_\sigma k_1\}
\end{equation}
implies
\begin{equation}\label{n4all:eq:gamma}
 \eta\int A_H\,d\nu\le
 \Gamma_H(\eta):=p_0\lambda_0+u(\lambda_++\lambda_-)+S_H\eta.
\end{equation}
Upper bounds for $p_0,u$ may replace them. The two equality prices
retain their signs. In particular, $S_H$ and $\Gamma_H$ need not be
nonnegative. The prices are fixed before taking the supremum over
the whole tagged domain. This preserves the cancellation between
frequencies; separate absolute estimates for each $r_q$ do not.

\subsection{The fourth-power remainder}
For complex $\psi=g+\delta$,
\begin{equation}\label{n4all:eq:quartic}
 \psi^4-g^4=4g^3\delta+
 \delta^2(\psi^2+2\psi g+3g^2).
\end{equation}
Suppose $|g|\le A\le1$, $|\psi|\le1$ and $|\delta|\le E$.
Two complete bounds for the remainder are
\begin{equation}\label{n4all:eq:quartic-basic}
 E^2(p^2+2pA+3A^2),\quad p=\min(1,A+E),\qquad
 v(4A^2+2-v),\quad v=\min(E^2,2A^2+1).
\end{equation}
The first follows directly from \eqref{n4all:eq:quartic}. For the
second, use
\[
 \psi^4-g^4-4g^3\delta
   =12\delta^2\int_0^1(1-t)(g+t\delta)^2\,dt
\]
and
$|g+t\delta|^2=(1-t)|g|^2+t|\psi|^2-t(1-t)|\delta|^2$.

A sharper two-disk bound is also finite algebra. Rotate $g=a\ge0$
and put $\delta=x+\ii y$, $s=|\delta|^2$. Then
\[
 |\psi^4-g^4-4g^3\delta|^2
 =s^2\{36a^4+4a^2s+s^2+8ax(6a^2+s)+24a^2x^2\}.
\]
For fixed $a,s$, the expression in braces is convex in $x$.
Its maximum on the two disks occurs at a collinear endpoint or on
$|\psi|=1$. Put $r_+=\min(E,1)$, $m=\min(A,1-r_+)$,
$r_-=\min(E,1+A)$. The collinear bounds are
\[
 r_+^2(6m^2+4mr_++r_+^2),\qquad
 r_-^2(6A^2-4Ar_-+r_-^2),\qquad
 \begin{cases}r_-^4,&r_-\le1,\\
 r_-^2(3r_-^2-8r_-+6),&r_-\ge1.
 \end{cases}
\]
On the unit circle it remains to maximize
\begin{equation}\label{n4all:eq:lens-polynomial}
 Q_A(s)=s^2\{18A^4+12A^2+6-(12A^2+8)s+3s^2\}
\end{equation}
on $[(1-A)^2,\min(E^2,(1+A)^2)]$. Its endpoints and the sole
interior local maximum
\[
 s_-=(3A^2+2-A\sqrt{4-3A^2})/2
\]
suffice. For $A=0$ the interval is degenerate or empty.
Taking the maximum of the collinear values and the square root of
this last maximum gives the lens bound. To see that no $a<A$ is
lost, note that the unit-circle expression is convex in $a^2$;
its lower feasible endpoint is already collinear. The collinear
expressions increase with their admissible radius after maximization
over $a$. Thus these are whole-disk bounds, not sampled values.
Their minimum with \eqref{n4all:eq:quartic-basic} is valid.

On a closed excess interval $[L,U]$, an affine radius
$e_j(\eta)=A_j+B_j\eta$ is used only when it is nonnegative at both
ends and bounds $|\delta_j|$ throughout. Freeze
$p_j=\min(1,a_j+\max\{e_j(L),e_j(U)\})$ in the first bound of
\eqref{n4all:eq:quartic-basic}. This gives a quadratic in $\eta$
with nonnegative leading coefficient. Alternatively, freezing $E$
on the whole cell gives a constant disk or lens bound. A choice
between complete bounds is allowed; a coefficientwise minimum is not.

\subsection{A physical Prawitz comparison}
We use the closed-CDF Prawitz majorant~\cite{Prawitz1972}, in the
normalization of Lemma~\ref{orig:lem:BPcardinal}. To distinguish the
factor from the kernel $K$ used there, put
\[
 \widehat K(u)=2K(u)=(1-u)+\ii R(u),\qquad
 R(u)=(1-u)\cot(\pi u)+1/\pi.
\]
For $Z=(X_1+\cdots+X_4)/2$,
\begin{equation}\label{n4all:eq:prawitz}
 F_Z(x)\le J_P(x):=\frac12+
 \int_0^1\Re\{\widehat K(u)e^{-6\ii ux}\varphi_P(3u)^4\}\,du.
\end{equation}
The real integrand is regular at zero; no convergence of the complex
integral by itself is asserted. With $D=x-Z$ and $S(u)=uR(u)$ it is
\[
 F_x(u)=(1-u)\E\cos(6Du)+S(u)\E[6D\operatorname{sinc}(6Du)].
\]
The cotangent partial fractions, extended to the endpoints, give
\[
 0\le S\le1/3,\qquad |S'|<2,\qquad |S''|<5.
\]
Here is a direct derivative check. Separate the first pole:
\[
 S=\frac1\pi-\frac2\pi(A+MG),\quad
 A=\frac{u^2}{1+u},\quad M=u^2(1-u),\quad
 G=\sum_{j\ge2}\frac1{j^2-u^2}.
\]
On $[0,1]$, $|A'|\le3/4$, $|A''|\le2$, $|M|\le1/6$,
$|M'|\le1$ and $|M''|\le4$. The derivative series converge
uniformly, and
\[
 G\le3/4,\qquad |G'|\le10/27<1/2,\qquad
 |G''|\le106/135<7/6.
\]
The first sum telescopes. For the other two use
$j^2-u^2\ge3j^2/4$ and
$\sum_{j\ge2}j^{-k}\le2^{-k}+\int_2^\infty t^{-k}\,dt$,
with $k=4,6$. The product rule gives
$|S'|<19/18<2$ and $|S''|<223/54<5$.
For $u>1/2$, put $v=1-u$ and use
$v\cot(\pi v)\le1/\pi$; for $u\le1/2$ the cotangent is
nonnegative. This proves $R\ge0$. The positive partial-fraction
terms then give $0\le S\le1/\pi<1/3$.
Since $\operatorname{sinc}(v)=\int_0^1\cos(tv)\,dt$,
\begin{equation}\label{n4all:eq:integrand-curvature}
 |F_x''(u)|\le42\E|D|+108\E D^2+24\E|D|^3.
\end{equation}
These differentiations are justified by the displayed moments.

Now assume $|X|\le R_*=17/4$ and $\rho<46/25$.
The fourth moment follows from these hypotheses; it is not an
additional input. Write $\kappa=\E X^3$, $M_4=\E X^4$.
Lemma~\ref{n4all:lem:excess} gives the useful cell bounds
\begin{align}
 |\kappa|&\le\min\{\rho,|b-a|+\eta\},\notag\\
 M_4&\le\min\{R_*\rho, a^2-ab+b^2+
                         (R_*+2\max(a,b))\eta\}.
 \label{n4all:eq:local-moments}
\end{align}
For $x=H/2$, choose whole-cell bounds $\kappa_*$ and $M_{4,*}$
from these formulas and put
\[
 D_2=x^2+1,\qquad
 D_4=x^4+6x^2+2|x|\kappa_*+(M_{4,*}+9)/4.
\]
Then $\E|D|\le\sqrt{D_2}$ and
$\E|D|^3\le\sqrt{D_2D_4}$. Uniformly on $|x|\le2$, the simpler
bounds $\E D^2\le5$, $\E D^4<54$, $\E|D|<9/4$,
$\E|D|^3<33/2$ give $|F_x''|<2061/2$.

Take $N=256$, $u_j=(j+1/2)/N$, $q_j=3u_j$ and
\[
 \kappa_j(H)=\widehat K(u_j)e^{-\ii q_jH}/N,\qquad
 C_H(Q)=1+2\Re\sum_j\kappa_j(H)g(q_j)^4-2\Phi(H/2).
\]
Composite midpoint quadrature gives one raw numerator reserve
\begin{equation}\label{n4all:eq:midpoint}
 Q_H=\min\left\{\frac{687}{8N^2},
 \frac{42\sqrt{D_2}+108D_2+24\sqrt{D_2D_4}}{12N^2}\right\}.
\end{equation}
It controls the integral of the original law $P$ once. No separate
reference, linear-transfer or remainder quadrature error is added.

The raw thresholds are $H_\ell=-4+\ell/8$, $0\le\ell\le57$.
To fill the intervals between them, bound the characteristic function
on entire frequency cells, not just at the midpoints. If
$a_j\ge|g(q_j)|$ and $e_j$ bounds $|\varphi_P(q_j)-g(q_j)|$ on the
whole excess cell, then
\[
 P_j=\min\{1,a_j+e_j+3/(2N)\}
\]
bounds $|\varphi_P(3u)|$ throughout $j/N\le u\le(j+1)/N$,
since $\E|X|\le1$. Two complete kernel bounds are
\[
 u^2|\widehat K(u)|\le u^2(1-u)+u/3,\qquad
 u^2|\widehat K(u)|\le u^2(1-u)+u(1-u)/2.
\]
The second follows by retaining the first term of the cotangent
partial fractions:
$S(u)\le(1-u)(1+2u)/(\pi(1+u))<(1-u)/2$.
Let $d_j$ be the smaller of the two exact integrals over the $j$th
frequency cell. Their separate totals are $1/4$ and $1/6$.
Thus, for the same whole geometry and excess cell,
\[
 |J_P''(x)|\le B=36\sum_jd_jP_j^4.
\]
Since $\sup|\Phi''|<1/4$, the smooth physical function
$H\mapsto2\{J_P(H/2)-\Phi(H/2)\}$ needs the single interpolation
reserve
\begin{equation}\label{n4all:eq:H-reserve}
 E_H=\frac{h^2}{16}(B+1/4),\qquad h=1/8.
\end{equation}
It is not an interpolation of the discontinuous CDF or of a minimum
of computed bounds. The elementary cube, mass and unit radii
\[
 \eta q_j^3/6,\qquad \max(a,b)^2p_0q_j^2,\qquad 1+a_j
\]
suffice for this curvature calculation. The mass bound follows from
$|r_q|\le q^2/y$ and the first budget in
\eqref{n4all:eq:budgets}.

For $H\le-4$, Cantelli gives $F_Z(H/2)\le1/5$.
For $H\ge25/8$, $\Phi(H/2)>4/5$. The latter inequality follows,
for example, by integrating the alternating Gaussian Taylor
polynomial through degree six on $[0,25/16]$ and using
$\phi(0)>19947/50000$. Consequently both threshold tails have
positive normalized score at most $2/5$.

\subsection{The low-excess certificate}
On a geometry rectangle and a closed excess cell $[L,U]$, let
$M_H$ bound $C_H(Q)-Tr$ for the entire rectangle. Let
$\Gamma_H$ be obtained from \eqref{n4all:eq:gamma}, and let
\[
 R_2(\eta)=r_0+r_1\eta+r_2\eta^2,\qquad r_2\ge0,
\]
bound the sum of the fourth-power remainders weighted by
$|\kappa_j(H)|$. Combining \eqref{n4all:eq:quartic} at the same
$H$ gives
\begin{equation}\label{n4all:eq:low-polynomial}
 2\{F_{4,P}(H)-\Phi(H/2)\}-T\rho
 \le P_H(\eta):=M_H+Q_H+E_H+2\Gamma_H(\eta)+2R_2(\eta)-T\eta.
\end{equation}
When used at grid points, $E_H$ reserves interpolation to adjacent
thresholds. The leading coefficient is nonnegative. Endpoint
nonpositivity therefore proves the bound on the entire closed excess
cell. Equivalently its degree-two Bernstein coefficients are
\[
 P_H(L),\quad
 \frac{P_H(L)+P_H(U)}2-\frac{p_2(U-L)^2}{2},\quad P_H(U),
\]
where $p_2$ is the quadratic coefficient. For a quantitative signed
score, with $r_L\le r$ and $\rho<46/25$, use
\[
 T+\max_{\eta=L,U}
 \left\{\frac{P_H(\eta)}{46/25},
        \frac{P_H(\eta)}{r_L+\eta}\right\}.
\]
Both denominator faces are needed when the numerator is negative.
The cap branch is convex; for the other branch the derivative
numerator has derivative $P_H''(\eta)(r_L+\eta)\ge0$, so an
interior critical point can only be a minimum.

Here are the finite enclosures that produce these polynomials.
The center is enclosed by signed Taylor expansion in $(w,\vartheta)$
or in the contacts $(a,b)$ along probability-preserving paths.
For either pair of coordinates $i,j$,
\[
 \partial_i C_H=8\Re\sum_j\kappa_jg^3\partial_i g,\qquad
 |\partial_{ik}C_H|\le8\sum_j|\kappa_j|
       (3D_iD_k+D_{ik}),
\]
where $D_i,D_{ik}$ bound the derivatives of the explicit
three-atom characteristic function. The first-order terms retain
their signs; the second-order error bounds the whole rectangle.
The function $r(w,\vartheta)$ is jointly convex. Indeed, writing
$r=w^{-1/2}f(\vartheta)$, $f=(1+\vartheta^2)/\sqrt{1-\vartheta^2}$,
the Hessian determinant has the positive factor
\[
 3ff''-(f')^2=
 \frac{9+9\vartheta^2+15\vartheta^4-\vartheta^6}{(1-\vartheta^2)^3}.
\]
Subtracting its supporting plane therefore encloses $C_H-Tr$
without an absolute Hessian charge for $r$.

For the signed response, put $R_*=17/4$, $y_k=k/64$,
$0\le k\le272$, and use their nearest-node closed cells,
clipped to $[0,R_*]$. Hermite--Genocchi gives
\begin{align}
 r_q(\sigma,y)&=\sum_{m=0}^{60}(\ii\sigma)^{m+3}b_m h_m+E_q,
 &|E_q|&\le\frac{q^3}{6}\frac{(qR_*)^{61}}{61!},\notag\\
 b_0&=q^3/6,&b_m&=b_{m-1}qR_* /(m+3),\notag\\
 h_0&=1,&h_m&=(y/R_*)h_{m-1}+(m+1)(c/R_*)^m.
 \label{n4all:eq:response-series}
\end{align}
The bound $|\partial_y r_q|\le q^4/24$ adds the half-cell reserve
$q^4/(48\cdot64)$. The signed frequency coefficients are summed
before bounding their polynomial on a $y$ cell. If $g_0,c_0'$ are
nominal values, $e_q\ge|g-g_0|$ is the geometry error,
$p_q\ge|g|$ and $a_{0,q}\ge|g_0|$, then
\[
 |g^3-g_0^3|\le e_q(p_q^2+p_qa_{0,q}+a_{0,q}^2),
 \qquad |\partial_c r_q|\le q^4/12.
\]
These give a whole-rectangle correction for
$g^3r_q-g_0^3r_{q,0}$, with no second copy of the series or mesh
reserve. Thus the priced supremum bounds actual references throughout
the rectangle, not only its nominal center.

The prices used in the finite certificate are chosen from
\begin{align*}
 \lambda_0&\in\{0,1/1000,1/400,1/200,1/100,1/50,1/25,
                         2/25,4/25,8/25,16/25\},\\
 \lambda_\pm&\in\{0,1/2000,1/1000,1/400,1/200,1/100,1/50\},\\
 \theta_\pm&\in\{0,\pm1/200,\pm1/100,\pm1/50\}.
\end{align*}
For each fixed choice, all 273 cells on each of the two tags enter
the supremum.
When $\lambda_\sigma=0$, only $\theta_\sigma=0$ is used;
otherwise the reciprocal contact rays would give an infinite supremum.
On a cell $[y_L,y_U]$, with $c\in[c_L,c_U]$,
$k_0\ge1/(c_U^2y_U)$. The remaining cost is the concave quadratic
$-\lambda_\sigma z^2-\theta_\sigma z$ in $z=1/(y-c)$.
Its endpoints and admissible vertex suffice; when the cell crosses
the contact both reciprocal rays are included. No division at the
contact or at zero is used for the response series.

The same budgets yield whole-cell radius models by pricing
$|r_q|$ or its signed coordinate projections. They are combined
with the three elementary radii above and the complete two-disk
remainder. An independent absolute branch uses
\[
 |\psi^4-g^4|\le E(p^3+p^2A+pA^2+A^3),\qquad p=\min(1,A+E),
\]
or the complete linear-plus-remainder bound. It is evaluated as a
whole alternative. In every branch, the actual reference modulus,
the actual normalization $\rho=r+\eta$, one midpoint reserve and
one physical threshold reserve are retained. The final maximum over
$H$ is taken only after assembling each complete row.

\subsection{Closing the large excess}
Let $p=P(X\ge0)$ and $q=1-p$; zero belongs to the nonnegative class.
Condition on exactly $j$ nonnegative observations among four.
For $Z=(X_1+\cdots+X_4)/2$, the conditional moments are
\[
 \mu_j=\frac{u(j-4p)}{2pq},\qquad
 V_j=\frac14\left[j\left(\frac{v_+}{p}-\frac{u^2}{p^2}\right)
 +(4-j)\left(\frac{v_-}{q}-\frac{u^2}{q^2}\right)\right].
\]
Closed-event Cantelli gives, for every $x$,
\begin{equation}\label{n4all:eq:count-cantelli}
 F_Z(x)\le\sum_{j=0}^4\binom4j p^jq^{4-j}A_j(x),\qquad
 A_j(x)=\begin{cases}1,&x\ge\mu_j,\\
 V_j/(V_j+(\mu_j-x)^2),&x<\mu_j.
 \end{cases}
\end{equation}
This includes variance zero and atoms at the threshold. The necessary
sign-mass interval is
\[
 w(1-\vartheta)/2\le p\le1-w(1+\vartheta)/2.
\]
It is covered by closed cells. On a threshold cell $[x_L,x_U]$,
bound $A_j(x_U)$ and subtract a lower bound for $\Phi(x_L)$.
The five degree-four Bernstein coefficients bound the binomial
mixture jointly; its masses are not maximized independently.
Each retained sign-mass cell supplies a nonnegative raw numerator
bound $d_i$ and a correlated lower bound $r_i\le r$, so its score
is at most $d_i/(r_i+\eta)$.

If $r_B$ is the whole-box lower bound for $r$, put
\[
 c_C^0=\max\left\{0,\max_i(d_i/T-r_i),\frac{2}{5T}-r_B\right\},
 \qquad
 c_C=10^{-4}\lceil10^4c_C^0\rceil+10^{-4}.
\]
This closes every $\eta\ge c_C$ by a strict bound below $T$,
including the threshold tails. Independently, Shevtsova closes
every $\rho\ge979013/533000$:
\[
 \frac{3031}{10000}\left(1+\frac{323}{500\rho}\right)\le T.
\]
Hence $c_S=\max\{0,979013/533000-r_B\}$ is another closed sufficient
cutoff. Take $c_B=\min(c_C,c_S)$, selecting one complete route.
The low-excess interval $[0,c_B]$ and the large-excess range
$[c_B,\infty)$ include their common endpoint.

\subsection{The finite cover and conclusion}
\begin{proposition}[Complete bounded-law certificate]
\label{n4all:prop:certificate}
For every standardized law $P$ with
\[
 |X|\le17/4,\quad \rho(P)<46/25,\quad
 625/2116\le w\le1,\quad |\vartheta|\le13/20,
\]
and every real $H$,
\[
 \frac{2\{F_{4,P}(H)-\Phi(H/2)\}}{\rho(P)}\le T.
\]
The same assertion holds for the genuine reflected law $-P$.
\end{proposition}
\begin{proof}
Divide $[625/2116,1]$ into 20 equal closed bands and $[0,13/20]$
into 13 closed bands of width $1/20$. There are 260 rectangles.
The minimum of
$r^2=(1+\vartheta^2)^2/[w(1-\vartheta^2)]$ is attained at the
largest $w$ and smallest $|\vartheta|$. A rectangle with
$\min r^2\ge(46/25)^2$ is empty for the strict physical cap.
For every other rectangle construct its own large-excess cutoff.
If $c_B=0$, the large-excess bound covers all $\eta\ge0$ directly.
If $c_B>0$, divide $[0,c_B]$ into 16 equal closed intervals. At each interval
check all 58 thresholds, their 57 joins, the two tails, and both
actual reflected orientations, using the preceding inequalities.
A parent is accepted by its complete bound or by complete bounds
for all four children obtained by bisecting both coordinates.
Each child has its own cutoff; no part of another child's domain
or numerical conclusion is used in its place.

The exact certificate gives the following complete cover:
\begin{center}
\begin{tabular}{lr}
\toprule
Closed parent rectangles & 260\\
Empty by the strict moment cap & 54\\
Nonempty, bounded directly & 129\\
Bounded by all four children, absolute route & 43\\
Bounded by all four children, full comparison portfolios & 34\\
Uncovered rectangles & 0\\
Uniform upper bound & $4097/10000$\\
\bottomrule
\end{tabular}
\end{center}
The electronic supplement accompanying this revision supplies the
four-summand certificate: the finite rational inequalities for all
260 rectangles, the complete child portfolios and their reconstruction.
All elementary operations
are outward rational enclosures; square roots are checked by integer
squaring. Trigonometric, Gaussian and cotangent values are enclosed
by finite series with analytic remainders, using the arithmetic
language of Section~\ref{cut49:short:arithmetic}. The response uses
\eqref{n4all:eq:response-series}, not divided quotients at contacts.
The complete arithmetic reconstruction reproduces the whole cover
and its upper bounds; no sampled value or optimizer's status supplies
an inequality.

For completeness, acceptance has a direct mathematical interpretation.
Every actual canonical point belongs to a closed parent.
In a nonempty parent or child, the large-excess inequalities cover
$\eta\ge c_B$, while the low-excess inequalities cover all of
$[0,c_B]$, including $\eta=0$. The same unknown law supplies the
single $\nu$ in every priced frequency response. Quadrature and
physical interpolation extend the finite rows to every central
threshold; the tails cover the rest. Thus the accepted rational
inequalities prove the proposition for all actual laws, including
bounded infinite supports and zero atoms.
\end{proof}

\begin{proof}[Proof of Theorem~\ref{n4all:thm:main}]
Suppose $C_4>T$. Since $T>c_0$, Lemma~\ref{n4all:lem:owner} gives
an unrestricted positive closed maximizing pair in the class of
Proposition~\ref{n4all:prop:certificate}. The proposition bounds its
score by $T$, a contradiction. Both discrepancy signs are covered:
\[
 \Phi(x)-F_Z(x)\le\Phi(x)-F_Z(x-)
                     =F_{-Z}(-x)-\Phi(-x).
\]
This is the closed-CDF reflection identity, with the atom retained.
Finally, elementary Machin bounds give $\pi<3927/1250$ and hence
$\phi(0)>19947/50000$; also $\sqrt{10}>15811/5000$. Therefore
\[
 \CE>\frac{204862339}{500000000}
       =\frac{4097}{10000}+\frac{12339}{500000000}>T.
\]
No equality case $C_4=T$ is excluded.
\end{proof}

\section{A sharper fourth-index bound for three-point laws}\label{n4:section}

We prove the fourth-summand bound for every nondegenerate law supported
on at most three points, with arbitrary masses and spacing. The support
bound used by the finite verification is derived for a hypothetical
restricted maximizer. No reduction of unrestricted maximizers to three
points is assumed.

\begin{theorem}\label{n4:thm:main}
For centered variance-one probability laws supported on at most three
points, put $\rho(P)=\E_P|X|^3$ and
\[
 C_{4,\le3}=\sup_P\frac{2}{\rho(P)}
   \sup_x\left|\Pp_P\{(X_1+\cdots+X_4)/2\le x\}-\Phi(x)\right|.
\]
Then
\[
                    C_{4,\le3}<\frac{409}{1000}<\CE.
\]
Both inequalities are strict, including the first supremum over laws.
Nonzero affine transformations give the same conclusion for every
nondegenerate law on at most three points.
\end{theorem}

\subsection{The external input and critical-level attainment}
For all centered variance-one iid laws with finite $\rho$, use
\begin{equation}\label{n4:eq:external}
 R_n(P):=\frac{\sqrt n\,\Delta_n(P)}{\rho}
 \le\frac{3031}{10000}\left(1+\frac{323}{500\rho}\right).
\end{equation}
This is Shevtsova's structural estimate in
\href{https://www.mathnet.ru/eng/ia252}{\emph{On the absolute constants in
the Berry--Esseen inequality and its structural and nonuniform improvements},
Informatics and its Applications 7(1) (2013), 124--125}.
In the notation of that structural estimate, substituting
$X_i/\sqrt n$ gives $\ell_n=\rho/\sqrt n$ and $\tau_n=1/\sqrt n$.
Only this finite estimate is used. In particular, neither $C_3\le\CE$, a sharp Bernoulli theorem,
nor the project's large-index bound is needed for this proof.

Jensen gives $\rho\ge1$, and hence
\[
 R_n(P)\le\frac{2494513}{5000000}<\frac12.
\]
At $\rho=37/20$, the right side of \eqref{n4:eq:external} is strictly
below $409/1000$. More specifically, if
$R_n(P)>40899/100000$, then
\[
             \rho<\frac{979013}{529450}<\frac{37}{20}.
\tag{2}
\]

Suppose $C_{4,\le3}\ge409/1000$. Take an optimizing sequence with ratios
above $40899/100000$, and reflect each law if necessary so that the
chosen discrepancy is positive and upper closed. Reflection is legitimate
because a negative open-CDF edge at $x$ becomes a positive closed edge
at $-x$; continuity of $\Phi$ equates the relevant suprema.

The laws are tight, and (2) gives uniform integrability of their first
two moments: $\E[X^2;|X|>M]\le(37/20)/M$. A weakly converging subsequence
therefore has mean zero and variance one. The class of laws supported on
at most three points is weakly closed: four disjoint open neighborhoods
of positive limiting mass would, by Portmanteau, force four distinct
support points in all sufficiently late approximants. Lower
semicontinuity gives $\rho(P)\le\liminf\rho(P_j)$; uniform integrability
of third moments is not assumed.

For the raw sum threshold $H$, Cantelli gives
$\Pp(S_4\le H)\le4/(4+H^2)\le1/5$ when $H\le-4$.
When $H\ge2$, $1-\Phi(H/2)\le1-\Phi(1)<1/5$.
Thus positive scores greater than $2/5$ have owners in $[-4,2]$.
The elementary last normal bound follows, for instance, from
$\phi(0)>39/100$ and
$\int_0^1 e^{-t^2/2}\,dt\ge5/6$.
Each fixed-law positive discrepancy attains its positive supremum:
a CDF minus a continuous function is upper semicontinuous and tends
to zero in both tails.

Pass to a subsequence with $H_j\to H$ and $\rho(P_j)\to r$.
Weak convergence of the actual fourfold products and moving-threshold
Portmanteau give
\[
 \limsup_j\Pp_{P_j}(S_4\le H_j)\le\Pp_P(S_4\le H).
\]
The limiting positive discrepancy is therefore at least
$C_{4,\le3}r/2$. Since $1\le\rho(P)\le r$, its normalized score is
at least $C_{4,\le3}$. The definition of the supremum forces equality,
and also forces $r=\rho(P)$. We have obtained an actual restricted
maximizing pair $(P,H)$ with
\[
 c=C_{4,\le3}\ge409/1000,\quad \rho<37/20,\quad c<1/2.
\tag{3}
\]
This is attainment only in the stated closed support-cardinality class.

\subsection{Existing-atom variations give the needed compact support}
Write $z=H/2$, $F=\Pp_P(S_4\le H)$, $\beta=\E_P[X|X|]$, and
$G(y)=\Pp_P(S_3\le H-y)$, with the same iid marginal throughout.
For each existing atom $y$, contaminate by
$P_\epsilon=(1-\epsilon)P+\epsilon\delta_y$, then center and scale.
This variation remains supported on at most three points. Hold the
\emph{raw} event $S_4\le H$ fixed. Its probability is a polynomial in
$\epsilon$; no derivative of an atomic CDF in its threshold is taken.
The required derivatives at zero are
\[
 F'=4(G(y)-F),\qquad z'=-2y-z(y^2-1)/2,
\]
\[
 \rho'=|y|^3-3\beta y+\rho/2-3\rho y^2/2.
\tag{4}
\]
Nonpositive score derivative implies $Q(y)\ge4G(y)$ at every existing
atom, where
\begin{align}
 Q(y)={}&\frac c2|y|^3-
 \left(\frac{z\phi(z)}2+\frac{3c\rho}4\right)y^2
 -\left(2\phi(z)+\frac{3c\beta}2\right)y\notag\\
 &\quad+4F+\frac{z\phi(z)+c\rho/2}{2}.\tag{5}
\end{align}
Direct integration gives $\E Q(X)=4F=4\E G(X)$, so equality holds
at every atom. No contact inequality at a newly introduced fourth atom
is asserted or used. In particular, $Q(y)\le4$ on the actual support.

We have $|\beta|\le1$, $\phi(z)<2/5$ and $|z\phi(z)|<1/4$ on the
whole real line. The latter follows by maximizing $x\phi(x)$ at $x=1$,
using $e^{1/2}>13/8$ and $\phi(0)<2/5$.
By (3), for $r=|y|$,
\[
 Q(y)-4\ge L(r):=\frac{409}{2000}r^3-\frac{131}{160}r^2
                         -\frac{31}{20}r-\frac{33}{8}.
\tag{6}
\]
The complete exterior is covered by the exact identity
\[
 L(6+s)=\frac{159}{125}+\frac{10711}{1000}s
         +\frac{11449}{4000}s^2+\frac{409}{2000}s^3>0
          \qquad(s\ge0).\tag{7}
\]
Consequently the restricted maximizing law has support in $(-6,6)$.
This reduction is independent of any predecessor-row bound.

\subsection{The complete bounded three-point parameter space}
It remains to prove the following finite-dimensional statement.
\begin{lemma}\label{n4:lem:bounded}
If $P$ has at most three support points, mean zero, variance one,
support in $[-6,6]$ and $\rho\le37/20$, then $R_4(P)<409/1000$.
\end{lemma}
Affine invariance permits raw support $\{0,v,1\}$, with masses
\[
 p_0=a,\quad p_1=(1-a)b,\quad p_2=(1-a)(1-b),\quad
                         0\le a,b\le1.\tag{8}
\]
The endpoints 0 and 1 are actual support points. Three-point laws have
$0<a,b<1$ and $0<v<1$; two-point laws are represented by $b=0$ and
$v=1/2$. A degenerate standardized law is impossible. Reflection and
checking both discrepancy signs reduce $v$ to $[0,1/2]$.

Put $\mu=p_1v+p_2$, $V=\operatorname{Var}Y$, and $M=\E|Y-\mu|^3$.
The exact positive formulas are
\begin{align}
 V&=p_0p_1v^2+p_0p_2+p_1p_2(1-v)^2,\tag{9}\\
 M&=p_0\mu^3+p_1|v-\mu|^3+p_2(1-\mu)^3,\quad
 \rho=M/V^{3/2}.\tag{10}
\end{align}
Since both endpoints are present, the support assumption implies
\[
                  V\ge\max\{\mu,1-\mu\}^2/36.\tag{11}
\]
We also retain $0<V\le1/4$ and $\rho\ge1$.

The fifteen fourfold count vectors are
$\mathcal A=\{(i,j,k)\in\mathbb N_0^3:i+j+k=4\}$.
Their sum locations and actual product probabilities are
\[
 h_{ijk}=jv+k,\qquad
 w_{ijk}=\frac{24}{i!j!k!}p_0^ip_1^jp_2^k.\tag{12}
\]
All order changes in $0<v<1/2$ occur at $v=1/4,1/3$; $v=1/2$
is another possible collision boundary. This follows by solving
$(j-j')v=k'-k$ for all pairs in $\mathcal A$, retaining precisely
the ratios in the stated range. The acceptance code independently
orders counts at the rational midpoints $1/8,7/24,5/12$.

For each of the three closed chambers, let $A_1,\ldots,A_{15}$ be that
order. Every upper-closed positive discrepancy and lower-open negative
discrepancy is covered by the thirty expressions
\[
 \frac2\rho\left[\sum_{r\le t}w_{A_r}-
     \Phi\left(\frac{h_{A_t}-4\mu}{2\sqrt V}\right)\right],\qquad
 \frac2\rho\left[\Phi\left(\frac{h_{A_t}-4\mu}{2\sqrt V}\right)
                       -\sum_{r<t}w_{A_r}\right].\tag{13}
\]
At a collision boundary, use the last member of a tied block for the
positive closed value and the first for the negative left value. Both
are in (13). Thus no mass is lost at a coincident sum or zero-probability
occupancy. Between successive sum atoms, the relevant extrema lie at
these closed or left edges, and the two infinite tails tend to zero.

\subsection{Leaf inequalities and exact domain coverage}
A certificate cell is a closed box in $(a,b,t)\in[0,1]^3$.
The chamber coordinate is respectively
\[
 v=t/4,\qquad v=1/4+t/12,\qquad v=1/3+t/6.\tag{14}
\]
Its dyadic endpoints in $a,b,t$ are exact. They are mapped to outward
enclosures in $v$. Subscripts $l,h$ below denote valid lower and upper
enclosures, not independently adjustable law moments.

The mean decreases with $a,b$ and increases with $v$:
\[
 \mu=(1-a)(1-b+bv),\quad
 \mu_l=(1-a_h)(1-b_h+b_hv_l),\quad
 \mu_h=(1-a_l)(1-b_l+b_lv_h).\tag{15}
\]
Nonnegative products in (9) give $V_l,V_h$. With
$d=\max(0,\mu_l-v_h,v_l-\mu_h)$, (10) gives
\[
 M_l=(p_0)_l\mu_l^3+(p_1)_ld^3+(p_2)_l(1-\mu_h)^3,\quad
 \rho_l=\max\{1,M_l/V_h^{3/2}\}.\tag{16}
\]
A leaf may be rejected as infeasible if $\rho_l>37/20$ or if
$V_h<\max\{\mu_l,1-\mu_h\}^2/36$.

For each fixed prefix in (13), its probability increases with $a$ and
with $b$. Increasing $a$ replaces a draw in $\{v,1\}$ by 0; increasing
$b$ replaces a draw at 1 by $v$. An independent coordinatewise coupling
therefore decreases every sum. Thus the upper prefix value is bounded
by its value at $(a_h,b_h)$, and the lower prefix by its value at
$(a_l,b_l)$. The same argument applies to strict prefixes, since the
ordering is fixed within the chamber. At chamber boundaries the fixed
prefix polynomial is the one-sided limit; monotonicity persists by
continuity in the masses. This is a common-marginal multinomial
calculation, not a coordinatewise optimization with unrelated laws.

The numerator $h_{A_t}-4\mu$ is enclosed directly using (15).
Its sign determines whether the lower or upper endpoint of $2\sqrt V$
minimizes or maximizes its ratio. If $V_l=0$, the corresponding negative
lower phase or positive upper phase is conservatively infinite. A
finite Gaussian table bounds $\Phi$ at both enclosed phases.
After retaining the positive part of either discrepancy, division by
$\rho_l\ge1$ gives a valid upper bound in (13). A score leaf is
accepted only if every such bound is strictly below $409/1000$.

The certificate consists of three finite binary subdivision trees,
one for each chamber. An internal node splits one dyadic coordinate
at its exact midpoint, and both closed children must be verified.
Each leaf satisfies either a strict constraint exclusion or the score
inequality above. The three roots cover the entire three chambers.
Induction on these trees therefore proves complete closed-domain
coverage, including all collisions and shared boundaries.

\subsection{Gaussian enclosures and outward interval arithmetic}
For the Gaussian table, Machin's identity
$\pi=16\arctan(1/5)-4\arctan(1/239)$ and alternating rational sums
certify
\[
 \frac{3989422804014326}{10^{16}}<\phi(0)
                   <\frac{3989422804014328}{10^{16}}.
\]
At $x=i/1024$, $0\le i\le4096$, integrate the degree-80 Taylor polynomial
of $e^{-s^2/2}$. Its remainder lies between zero and the first omitted
term, with the known negative sign. Multiplication by the displayed
positive bounds encloses $\Phi(x)$; the table stores outward multiples
of $2^{-50}$. Its 4097 entries are rechecked from their defining rational
expressions. Monotonicity brackets arbitrary phases by adjacent entries,
reflection handles negative phases, and outside $[-4,4]$ the valid
outer bounds 0 and 1 are retained.

The interval calculation assumes IEEE--754 binary64 basic arithmetic
with rounding to nearest. Each operation is enlarged outward to the
adjacent representable endpoints, without contraction or reassociation.
A square-root approximation is only an initial guess: its endpoints
must also satisfy outward squared-endpoint tests. No Gaussian special
function or library exponential is used to decide acceptance.
These enclosure conditions and the complete closed forest certify all
leaf inequalities.

Lemma~\ref{n4:lem:bounded} follows. It contradicts the actual maximizing
pair (3), whose support was proved bounded by (7), proving
Theorem~\ref{n4:thm:main}. Finally
\[
 \CE^2>\frac{6643}{39600}>\left(\frac{409}{1000}\right)^2,
\]
using $\sqrt{10}>79/25$ and $\pi<22/7$.
The bound in this section is sharper than the unrestricted estimate
of Theorem~\ref{n4all:thm:main}, but applies only on at most three
points. Its restricted first variations do not reduce the unrestricted
optimization to this class; that optimization was treated separately
in Section~\ref{n4all:section}.

\clearpage
\appendix
\section{Auxiliary Bernstein coefficients and binomial localization}\subsection{The exact Bernstein data}\label{orig:app:bernstein}

For a polynomial $f(x)=\sum_{j=0}^9a_jx^j$, the Bernstein coefficients of
degree nine are
\[
 b_k=\sum_{j=0}^k a_j\frac{\binom{k}{j}}{\binom{9}{j}}.
\]
After affine substitution to the two half-intervals, the coefficients used
in Lemma~\ref{orig:lem:J2} are listed below:
\begin{center}\scriptsize
\begin{tabular}{c|c}
$j$ & $(b_{0j},b_{1j})$\\ \hline
0 & $(243/1250,\ 2079582664341761/1197443957149440000)$\\
1 & $(5623/45000,\ 1709192148941903/1347124451793120000)$\\
2 & $(1334869/17010000,\ 450741032531969/206258289269760000)$\\
3 & $(13004779/272160000,\ 10866565157710711/2958390952957440000)$\\
4 & $(146380277/5239080000,\ 198519592136938909/37719484650207360000)$\\
5 & $(28525339727/1838917080000,\ 125666379478652267/18859742325103680000)$\\
6 & $(487000027957/60071291280000,\ 48022819751421589/6286580775034560000)$\\
7 & $(37879182239513/9336794987520000,\ 10619877339829403/1347124451793120000)$\\
8 & $(11879475383308237/5388497807172480000,\ 1147322830325941/168390556474140000)$\\
9 & $(2079582664341761/1197443957149440000,\ 29384282680087/9355030915230000)$
\end{tabular}
\end{center}
Every entry is positive.  The Bernstein convex-hull property therefore
proves positivity of the polynomial on both halves without numerical
evaluation.

\subsection{Binomial localization}\label{orig:app:binloc}

Only the following previously established result is used in passing from
the local beta integral to the global supremum over lattice indices.

\begin{proposition}[Schulz's binomial localization, normalized]\label{orig:prop:schulzlocal}
Let $n\in\mathbb N$, $p\in[1/6,5/6]$, $q=1-p$,
$D=p^2+q^2$, and $\rho(p)=D/\sqrt{pq}$.  For an integer $k$ put
\[
 z_k={k-np\over\sqrt{npq}},
\]
\[
 R_{n,p}(k):={\sqrt n\over\rho(p)}
       \left|B_{n,p}(k)-\Phi(z_k)\right|,\qquad
 L_{n,p}(k):={\sqrt n\over\rho(p)}
       \left|B_{n,p}(k-1)-\Phi(z_k)\right|,              \tag{A.1}
\]
where $B_{n,p}(-1)=0$.  Define
\[
 \widehat G(p):={\phi_0\{3+|p-q|\}\over6D},\qquad
 0<c_0:={2\over\sqrt{2\pi}\,\CE}-1
      =12\sqrt {10}-37<{237\over250}=0.948.             \tag{A.2}
\]
Then, with $k_-:=\lfloor np\rfloor$ and $k_+:=\lceil np\rceil$,
\[
 {\sqrt n\,\Delta_n(Q_p)\over\rho(p)}
 \le\max\{R_{n,p}(k_-),L_{n,p}(k_-),
            R_{n,p}(k_+),L_{n,p}(k_+),c_0\widehat G(p)\}. \tag{A.3}
\]
Moreover, for every $n\in\mathbb N$ and $p\in[1/3,2/3]$,
\[
             {\sqrt n\,\Delta_n(Q_p)\over\rho(p)}
             <\widehat G(p).                              \tag{A.4}
\]
\end{proposition}

\begin{proof}[Source and normalization]
The unnormalized statement (A.3) is Proposition~6.1 of \cite{Schulz2016},
p.~31 (the proposition is restated at the beginning of Chapter~7); (A.4)
is the sharper assertion in Theorem~1 there, p.~1.  Schulz writes
$\sigma=\sqrt{npq}$ and bounds the noncentral alternative in Proposition~6.1
by
\[
 c_0\,{3+|p-q|\over6\sqrt{2\pi}\,\sigma}.
\]
The rational bound in (A.2) is exact: after moving $37$ and dividing by
$12$ it reduces to $\sqrt {10}<9487/3000$, whose square is
$10<90003169/9000000$.
Multiplication by $\sqrt n/\rho(p)$ gives exactly
$c_0\widehat G(p)$ because
\[
 \rho(p)={D\over\sqrt{pq}},\qquad
 {\sqrt n\over\rho(p)\sigma}={1\over D}.
\]
The same calculation without $c_0$ converts the sharper part of
Schulz's Theorem~1 into (A.4).  For $p\le1/2$,
\[
 \widehat G(p)={\phi_0(2-p)\over3D}=G(p),                \tag{A.5}
\]
while reflection $p\leftrightarrow q$, $k\leftrightarrow n-k$ exchanges
the two edge orientations.  Thus (A.1)--(A.4) are exactly the cited
claims, not an asymptotic reformulation of them.
\end{proof}

This citation is used solely for lattice-index localization.  The
$n^{-1}$ coefficient and its uniform $o(n^{-1})$ remainder are proved in
Lemma~\ref{orig:lem:binsaddle}; no unquantified Edgeworth expansion is used.

\section{Exact coefficient ledger}\label{bin:app:ledger}

With the scaled notation in the proof of Lemma~\ref{bin:lem:fifth}, exact
Gaussian integration gives
\begin{align*}
 \widetilde A_1={}&\eta+\frac{2-p}{3},\\[1mm]
 \widetilde\delta_2={}&
 \frac{6\eta^2+6\eta+p^2-p+1}{12},\\[1mm]
 \widetilde A_3={}&\frac1{135}
 \bigl(45\eta^3+45p\eta^2+45\eta^2+45p\eta
 -4p^3+6p^2+6p-4\bigr),\\[1mm]
 \widetilde\delta_4={}&\frac1{288}
 \bigl(36\eta^4+96p\eta^3+24\eta^3+12p^2\eta^2
 +132p\eta^2-24\eta^2\\
 &\hspace{25mm}{}+12p^2\eta+36p\eta-12\eta
 +p^4-2p^3+3p^2-2p+1\bigr),\\[1mm]
 \widetilde A_5={}&\frac1{2835}
 \bigl(189\eta^5+945p\eta^4+630p^2\eta^3+1260p\eta^3
 -315\eta^3\\
 &\hspace{13mm}{}-42p^3\eta^2+1008p^2\eta^2+63p\eta^2-42\eta^2
 -42p^3\eta+378p^2\eta\\
 &\hspace{13mm}{}-252p\eta+84\eta
 +8p^5-20p^4+8p^3+8p^2-20p+8\bigr).
\end{align*}
Substitution into \eqref{bin:eq:Hledger} yields \eqref{bin:eq:P}.  This ledger is
short enough to be checked independently without reproducing the full
polynomials $E_4$ and $E_5$.

\section{Arithmetic supplement to the elementary bound \(C\le9/20\)}\label{app:elementary045}

This early bound is retained as a short historical supplement.
Its mechanism is classical smoothing, comparison with a Rademacher law,
positive elementary estimates, and a finite interval proof.
It does not require the asymptotic or extremizer arguments of the main text.

Let \(X,X_1,X_2,\ldots\) be iid, with \(\E X=0\), \(\E X^2=1\), and
\(\rho=\E|X|^3<\infty\). Put
\[
 \D=\sup_{x\in\R}
 \left|\mathbb P\!\left(\frac{X_1+\cdots+X_n}{\sqrt n}\le x\right)
       -\Phi(x)\right|,
 \qquad C=\sup_{n,X}\frac{\sqrt n\,\D}{\rho}.
\]
\begin{theorem}\label{app:045:theorem}
The interval certificate described below gives \(C\le9/20=0.45\).
\end{theorem}

\begin{proof}
\emph{The finite reduction.}
The one-summand estimate of Bentkus--Kir\v{s}a
\cite{BentkusKirsa1989} and two estimates of Shevtsova
\cite{Shevtsova2013,Shevtsova2012} give
\[
 \Delta_1\le0.3704\rho,\qquad
 \D\le\frac{0.3322(\rho+0.429)}{\sqrt n},\qquad
 \D\le \CE\frac{\rho}{\sqrt n}+\frac{3\rho^2}{n},
 \quad \CE=\frac{3+\sqrt{10}}{6\sqrt{2\pi}}.
\]
All terminating decimals here are exact rationals. Since \(\rho\ge1\), set
\[
 \rho_0=\frac{0.3322\cdot0.429}{0.45-0.3322}
       =\frac{712569}{589000}.
\]
The second estimate settles \(\rho\ge\rho_0\). The third settles
\(n\ge8124\) when \(1\le\rho\le\rho_0\), because
\(\CE+3\rho_0/\sqrt{8124}<9/20\). For example, rational enclosures give
\[
 \CE<\frac{409732184}{10^9},\qquad
 \sqrt{8124}>\frac{90133234}{10^6},
\]
and
\[
 \frac9{20}-\frac{409732184}{10^9}
 -\frac{3(712569/589000)}{90133234/10^6}
 =\frac{3121000875001}{3318029676625000000}>0.
\]
The square-root enclosure follows by squaring; the enclosure of \(\pi\)
needed for the first comparison follows from Machin's formula and the
alternating arctangent series. It remains to prove the bound on
\[
                 2\le n\le8123,\qquad 1\le\rho\le\rho_0.       \tag{C.1}
\]

\emph{A distribution-free smoothing bound.}
Write
\[
 \varphi(u)=\E e^{\ii uX},\quad
 f_n(t)=\varphi(t/\sqrt n)^n.
\]
\[
 K(s)=\frac{1-s}{2}+
       \frac{\ii}{2}\left((1-s)\cot(\pi s)+\frac1\pi\right),
 \quad L(s)=K(s)-\frac{\ii}{2\pi s}.
\]
For \(T>0\) and \(0<t_0<1\), Prawitz's inequality
\cite{Prawitz1972} bounds \(\D\) by
\begin{align}
 U_n(\rho;T,t_0)={}&
 \int_0^{t_0}2|K(s)|R_{n,\rho}(Ts)\dd s
 +\int_{t_0}^1 2|K(s)|M_{n,\rho}(Ts)\dd s \notag\\
 &+\int_0^{t_0}2|L(s)|e^{-T^2s^2/2}\dd s
 +\frac1{2\pi}E_1(T^2t_0^2/2),                         \tag{C.2}
\end{align}
where \(E_1(x)=\int_x^\infty e^{-y}\dd y/y\),
\(R_{n,\rho}(t)\ge|f_n(t)-e^{-t^2/2}|\), and
\(M_{n,\rho}(t)\ge|f_n(t)|\).

These majorants are explicit. The Mattner--Shevtsova comparison
\cite{MattnerShevtsova2019} is
\[
 |\varphi(u)-\cos u|\le\frac{B(\rho)|u|^3}{6},
 \qquad
 B(\rho)=
 \sqrt{\frac{\rho^2+\rho\sqrt{\rho^2+8}}2-2}.
\]
Put \(u=t/\sqrt n\), \(c=|\cos u|\),
\(\delta=B(\rho)|u|^3/6\), and \(a=\min(1,c+\delta)\). The identity for
a difference of powers gives
\[
 D(n,\rho,t)=\min\left(2,n\delta,\,
       \delta\sum_{j=0}^{n-1}a^jc^{n-1-j}\right),\qquad
 R_{n,\rho}(t)=\min\left(2,\,
             |\cos^n u-e^{-t^2/2}|+D(n,\rho,t)\right).
                                                               \tag{C.3}
\]
For the modulus we take the smaller of
\(\min(1,c^n+D(n,\rho,t))\) and Shevtsova's characteristic-function
majorant. With \(T=q\sqrt n/(\rho+1)\), the latter is
\[
 M_P(n,\rho,Ts)=
       \left(1-\frac{2\psi(qs)}{(\rho+1)^2}\right)^{n/2}.
\]
Here
\[
 \psi(z)=
 \begin{cases}
 z^2/2-\kappa z^3,&0\le z\le\vartheta_0,\\
 1-\cos z,&\vartheta_0<z\le2\pi,\\
 0,&z>2\pi,
 \end{cases}
 \qquad
 \kappa=\frac{\cos\vartheta_0-1+\vartheta_0^2/2}{\vartheta_0^3},
\]
and \(\vartheta_0\) is the unique root in \((\pi,2\pi)\) of
\(z^2+2z\sin z+6(\cos z-1)=0\).
Elementary differentiation gives \(0\le\psi\le2\);
the base of \(M_P\) therefore lies in \([0,1]\).
The resulting \(R\) and \(M\) bound every admissible law with moment \(\rho\),
not merely the Rademacher law.

\emph{Stable bounds at the endpoints.}
The following identities remove the apparent singularities:
\[
 2B(\rho)^2=(\rho^2-1)
   \left(1+\frac{\rho^2+9}{\rho\sqrt{\rho^2+8}+3}\right),
 \qquad
 \ell(s):=\frac1{\pi s}-\cot(\pi s)
       =\frac{2s}{\pi}\sum_{k=1}^{\infty}\frac1{k^2-s^2}.
                                                               \tag{C.4}
\]
The series is positive and increasing. Its omitted tail after \(m\) terms
is at most
\[
             \frac{2s}{\pi m(1-s^2/(m+1)^2)}.
\]
Also \(2|L(s)|=(1-s)\sqrt{1+\ell(s)^2}\), while the exact identity
\[
 2|K(1-h)|=h\sqrt{1+\ell(h)^2}
\]
gives, for \(0<H\le1/2\),
\[
 \int_{1-H}^1 2|K(s)|M_{n,\rho}(Ts)\dd s
          \le\frac{H^2}{2}\sqrt{1+\ell(H)^2}.             \tag{C.5}
\]
At the other endpoint, if \(0<\eta,\lambda\le\min(t_0,1/2)\) and
\(T\eta/\sqrt n\le1\), elementary Taylor bounds give
\[
 \int_0^\eta 2|K|R_{n,\rho}(Ts)\dd s
 \le\frac{B(\rho)T^3\eta^3}{9\sqrt n}
       +\frac{T^4\eta^4}{12n},\qquad
 \int_0^\lambda2|L|e^{-T^2s^2/2}\dd s
 \le\lambda\sqrt{1+4/\pi^2}.                             \tag{C.6}
\]
For \(0\le u<\pi/2\), the Gaussian comparison in (C.3) can be evaluated
as the nonnegative product
\[
 e^{-nu^2/2}-\cos^n u
 =(e^{-u^2/2}-\cos u)
   \sum_{j=0}^{n-1}e^{-ju^2/2}\cos^{n-1-j}u .
\]
Outside this range the signed power \(\cos^n u\) is enclosed directly.
Thus odd sample sizes require no change of sign convention.

\emph{Why the finite arithmetic proves the theorem.}
For a closed rational cell \(I=[\rho_-,\rho_+]\), choose rational
\(q,t_0,\eta,\lambda\) and use \(T(\rho)=q\sqrt n/(\rho+1)\).
Require
\[
 q>0,\quad 0<t_0<1,\quad
 0<\eta,\lambda\le\min(t_0,1/2),\quad
 \frac{q\eta}{\rho_-+1}\le1.
\]
The last inequality makes (C.6) valid uniformly on \(I\).
Each remaining proper integral in (C.2) is bounded by an upper sum
\[
 \sum_j(s_{j+1}-s_j)b_j,\qquad
 b_j\ge\sup_{\rho\in I,\ s\in[s_j,s_{j+1}]}G(\rho,s)
\]
for its nonnegative integrand \(G\). The \(b_j\) are obtained with outward
rounding: the interval assigned to every operation contains all its exact
values on the whole rectangle. Equations (C.4)--(C.6) bound the omitted
endpoint pieces. A cell crossing a branch point of \(\psi\) may safely use
the lower bound \(\psi\ge0\).

For the first endpoint use \(B(\rho_+)\) and
\(T_+=q\sqrt n/(\rho_-+1)\). For the exponential-integral tail use
\(T_-=q\sqrt n/(\rho_++1)\), since \(E_1\) is decreasing.
Let \(\widehat U_{n,I}\) be the sum of these upper bounds.
Then, throughout the entire closed cell,
\[
 \frac{\sqrt n\,\D}{\rho}
 \le\frac{\sqrt n}{\rho_-}\widehat U_{n,I}.                \tag{C.7}
\]
The certificate supplies a finite rational cover of \([1,\rho_0]\)
for every integer \(n=2,\ldots,8123\), including all shared endpoints.
Its \(886\,998\) cells satisfy
\[
 \max_{n,I}\frac{\sqrt n}{\rho_-}\widehat U_{n,I}
 \le\frac{449999}{10^6}<\frac9{20}.
\]
This is the finite computational step of the proof.
It controls continuous intervals of \(\rho\) and of the integration
variable, not sampled points. Equation (C.7) proves (C.1), and the initial
reductions prove \(C\le9/20\) for every sample size.
\end{proof}

\paragraph{Historical note.}
The original certificate records completion on 6 August 2026;
the public source archive begins on 25 September 2026.
The full argument and interval certificate are available in
\href{https://github.com/rodziewiczk/berry-esseen-045-certificate}
{Rodziewicz's certificate repository}.
Stronger universal bounds, including \(0.4395\) by Xiao--Li and \(0.423\)
by Lin, are now reported in their
\href{https://github.com/haonan-xiao/iid-berry-esseen}{respective}
\href{https://github.com/CoolRmal/BerryEsseen}{repositories}.
The present appendix is retained for its elementary mathematical mechanism,
not as the current best bound.

\section*{AI assistance declaration}
Artificial-intelligence systems were used in developing the mathematical
arguments and proofs, constructing and checking computational certificates,
and drafting and revising this manuscript. This included assistance from
ChatGPT and Codex using the GPT-5.6 Sol, GPT-6.1 Sol and GPT-6 Astra models.
Kacper Rodziewicz and Bartosz
Ko{\l}odziejek remain responsible for the mathematical claims and the
final text.

\clearpage
\begin{thebibliography}{99}
\bibitem{BentkusKirsa1989}
V.~Bentkus and K.~Kir\v{s}a, Estimates of the closeness of a
distribution to the normal law, \emph{Lithuanian Mathematical Journal}
\textbf{29} (1989), 321--332.
\bibitem{MattnerShevtsova2019}
L.~Mattner and I.~Shevtsova, An optimal Berry--Esseen type theorem
for integrals of smooth functions, \emph{ALEA} \textbf{16} (2019),
487--530. \url{https://doi.org/10.30757/ALEA.v16-19}.
\bibitem{Shevtsova2014}
I.~Shevtsova, On the accuracy of the approximation of the complex
exponent by the first terms of its Taylor expansion with applications,
\emph{Journal of Mathematical Analysis and Applications}
\textbf{418}(1) (2014), 185--210.
\url{https://doi.org/10.1016/j.jmaa.2014.03.075}.
\bibitem{Shevtsova2012}
I.~Shevtsova, Moment-type estimates with asymptotically optimal
structure for the accuracy of the normal approximation,
\emph{Annales Mathematicae et Informaticae} \textbf{39} (2012),
241--307.
\url{https://ami.uni-eszterhazy.hu/uploads/papers/finalpdf/AMI_39_from241to307.pdf}.
\bibitem{Prawitz1972}
H.~Prawitz, Limits for a distribution, if the characteristic function
is given in a finite domain, \emph{Skandinavisk Aktuarietidskrift}
(1972), 138--154. \url{https://doi.org/10.1080/03461238.1972.10404645}.
\bibitem{Schulz2016}
J.~Schulz, \emph{The optimal Berry--Esseen constant in the binomial
case}, doctoral dissertation, Universit\"at Trier, 2016.
\url{https://d-nb.info/1197702695/34}.
\bibitem{Shevtsova2013}
I.~G.~Shevtsova, On the absolute constants in the Berry--Esseen
inequality and its structural and nonuniform improvements,
\emph{Informatika i Ee Primeneniya} \textbf{7}(1) (2013), 124--125.
\url{https://www.mathnet.ru/eng/ia252}.
\bibitem{HC}
H.~He and G.~Cheng, \emph{The Berry--Esseen Bound is Sharp for All
Sufficiently Large Sample Sizes}, arXiv:2609.06358v1 (2026).
\url{https://arxiv.org/abs/2609.06358v1}.
\bibitem{HC2026v3}
H.~He and G.~Cheng, \emph{The Berry--Esseen Constant Conjecture is
Eventually True}, arXiv:2609.06358v3 (2026), revised 30 September.
\url{https://arxiv.org/abs/2609.06358v3}.
\bibitem{EffectiveSupplement}
K.~Rodziewicz and B.~Ko{\l}odziejek, \emph{Effective finite-index
continuation and certificates for the iid Berry--Esseen problem},
supplementary proof material (2026).
\url{https://github.com/rodziewiczk/berry-esseen-045-certificate/tree/main/moment-defects}.
\bibitem{ZNC}
A.~Zolotukhin, S.~Nagaev and V.~Chebotarev, On a bound of the
absolute constant in the Berry--Esseen inequality for i.i.d.
Bernoulli random variables, \emph{Modern Stochastics: Theory and
Applications} \textbf{5} (2018), 385--410.
\bibitem{Certificate045}
K.~Rodziewicz, \emph{An elementary interval certificate for the
Berry--Esseen bound \(0.45\)}, certificate archive (2026).
\url{https://github.com/rodziewiczk/berry-esseen-045-certificate}.
\end{thebibliography}
\end{document}

